\chapter{A Gauduchon variational volume}
\label{chap:gauduchon-variational-volume}

\begingroup
\predisplaypenalty=10000

For a closed K\"ahler background, integration determines the scalar in the normalized Monge--Amp\`ere equation. For a nonclosed background, the total mass can vary with the potential. The smooth equation nevertheless selects one scalar for each positive probability measure, and Gauduchon duality expresses its $n$th root as an infimum of linear pairings. This formula extends the scalar to the pseudoeffective cone. On the positive cone, its first and second variations recover the solution form and define a metric modulo smooth $\ddc$ terms.

The dual problem also separates three questions: whether a current maximizes the determinant integral, whether a smooth Gauduchon metric attains the quotient, and whether the subsolution threshold is an exact equation scalar. Weighted tail estimates give current attainment. Smoothness or controlled degeneracy of the maximizing test supplies the additional rigidity needed for the equation.

Varying the probability measure recovers the upper and lower Monge--Amp\`ere masses. A threefold with zero lower mass shows that positivity for each smooth density need not be uniform over densities. At the nef boundary, the distinction from the closed K\"ahler case is stronger: intersection polynomials give way to fractional, logarithmic or oscillating decay, and a single blowup can change its order.

Throughout the chapter, $X$ is a connected compact complex manifold of dimension $n\geq1$, and $\mu$ is a smooth strictly positive probability form, unless an example specifies otherwise.

\section{The quotient and the smooth equation}\label{sec:gaud-volume-definition}
\label{sec:gaud-volume-positive-cone}

Gauduchon Stokes makes the numerator of the quotient invariant under a smooth $\ddc$ change of background. Its denominator is chosen so that the quotient has the same homogeneity as the root of the equation scalar. For a Gauduchon metric $\gamma$, put
\[
 D_\mu(\gamma)=\int_X
           \left(\frac{\mu}{\gamma^n}\right)^{1/n}\gamma^n.
\]
This number is positive. It scales by $r^{n-1}$ when $\gamma$ is multiplied by a positive constant $r$.

\Needspace{12\baselineskip}
\begin{definition}\label{def:gaud-variational-volume}\index{Gauduchon variational volume!definition}
For a pseudoeffective background $\theta$, define
\[
 \lambda_\mu(\theta)=
 \inf_{\gamma\text{ Gauduchon}}
 \frac{\displaystyle\int_X\theta\wedge\gamma^{n-1}}
      {D_\mu(\gamma)},
 \qquad
 \mathcal V_\mu(\theta)=\lambda_\mu(\theta)^n.
\]
In dimension one every Hermitian metric is permitted in the infimum; its quotient is simply $\int_X\theta$.
\end{definition}

\begin{proposition}\label{prop:gaud-volume-elementary}
The number $\lambda_\mu(\theta)$ is finite and nonnegative. It is unchanged by a smooth $\ddc$ change of $\theta$, positively homogeneous of degree one, and concave on the pseudoeffective cone. If $\beta$ is a smooth real $(1,1)$-form and $\beta-\theta$ is pseudoeffective, then
\[
 \lambda_\mu(\beta)\geq\lambda_\mu(\theta).
\]
The function is upper semicontinuous under smooth convergence within that cone.
\end{proposition}
\begin{proof}
Choose $u\in\PSH(X,\theta)$. Distributional Stokes and the Gauduchon condition give
\[
 \int_X\theta\wedge\gamma^{n-1}
 =\int_X(\theta+\ddc u)\wedge\gamma^{n-1}\geq0.
\]
Thus every quotient is nonnegative. Testing one fixed Gauduchon metric gives a finite upper bound. The same identity proves gauge invariance and monotonicity for a pseudoeffective difference.

Each quotient is a linear function of the background. The infimum of these linear functions is concave and positively homogeneous. Each individual quotient is continuous in the smooth topology. Given a smoothly converging sequence, bound the infimum by any one quotient, take the upper limit, and finally take the infimum over the test metric. This proves upper semicontinuity.
\end{proof}

The smooth equation identifies both the value of the quotient and every minimizing metric. Let $\mathscr P$ consist of backgrounds $\theta$ admitting a smooth real $u$ with $\theta+\ddc u>0$. For $\theta\in\mathscr P$, the smooth equation theorem gives a unique positive form
\[
 \Omega_\theta=\theta+\ddc u_\theta,
 \qquad \Omega_\theta^n=c_\theta(\mu)\mu,
 \qquad c_\theta(\mu)>0.
\]
The potential is unique after any fixed additive normalization. To apply \cref{thm:tw-smooth}, first choose a Hermitian representative of the background and express $\mu$ as a smooth positive density relative to its volume form.

\Needspace{16\baselineskip}
\begin{theorem}\label{thm:gaud-volume-dual-formula}\index{Gauduchon variational volume!smooth equation}
For every $\theta\in\mathscr P$,
\[
 \lambda_\mu(\theta)=c_\theta(\mu)^{1/n},
 \qquad \mathcal V_\mu(\theta)=c_\theta(\mu).
\]
The infimum in \cref{def:gaud-variational-volume} is attained. When $n\geq2$, its minimizing Gauduchon metrics are precisely the constant multiples of the Gauduchon metric conformal to $\Omega_\theta$.
\end{theorem}
\begin{proof}
For any Hermitian metric $\gamma$, the arithmetic--geometric mean inequality for the eigenvalues of $\Omega_\theta$ relative to $\gamma$ gives the pointwise inequality
\[
 \Omega_\theta\wedge\gamma^{n-1}
 \geq\left(\frac{\Omega_\theta^n}{\gamma^n}\right)^{1/n}\gamma^n
 =c_\theta(\mu)^{1/n}
       \left(\frac{\mu}{\gamma^n}\right)^{1/n}\gamma^n.
\]
If $\gamma$ is Gauduchon, integration replaces $\Omega_\theta$ in the numerator by $\theta$. Every test quotient is therefore at least $c_\theta(\mu)^{1/n}$.

For $n\geq2$, take the Gauduchon representative $\gamma=\mathrm{e}^g\Omega_\theta$ given by \cref{thm:gauduchon}. The eigenvalues are equal at every point, so equality holds pointwise and then in the quotient. This proves the formula and attainment. Conversely, if a metric attains the infimum, the nonnegative smooth difference in the displayed inequality has zero integral. It vanishes everywhere, and equality in the arithmetic--geometric mean inequality makes $\gamma$ conformal to $\Omega_\theta$. Uniqueness of the Gauduchon metric in a conformal class, up to constant scale, proves the last assertion.

For $n=1$, integrate the equation and use $\int_X\ddc u_\theta=0$. It gives $c_\theta(\mu)=\int_X\theta$, equal to every quotient in the definition.
\end{proof}

For nonzero nonnegative $L^1$ densities in a Hermitian background, the analogous quotient computes the subsolution threshold by \cref{thm:integrable-gauduchon-formula}. Attainment and regularity for those data will be studied below.

\begin{corollary}\label{cor:gaud-volume-big-positive}
If $\theta$ is big, then $\lambda_\mu(\theta)>0$.
\end{corollary}
\begin{proof}
Choose a Hermitian form $h$ and $\delta>0$ such that $\theta-\delta h$ is pseudoeffective. Monotonicity, homogeneity and the theorem give
$\lambda_\mu(\theta)\geq\delta\lambda_\mu(h)>0$.
\end{proof}

Whether positive $\mathcal V_\mu(\theta)$ forces bigness remains open here, even for a general nef background.

\Needspace{17\baselineskip}
\section{Concavity, equality and variation}

A minimizing Gauduchon metric supports the concave function $\lambda_\mu$. Equality in concavity forces the same metric to minimize at both endpoints; differentiating the smooth equation identifies that supporting functional with the derivative. For $\theta\in\mathscr P$, choose a minimizing metric $\gamma_\theta$. The variation directions are smooth real $(1,1)$-forms.

\begin{theorem}\label{thm:gaud-volume-concavity-equality}\index{Gauduchon variational volume!concavity and equality}
Let $\theta_0,\theta_1\in\mathscr P$ and $0<s<1$. Then
\[
 \lambda_\mu((1-s)\theta_0+s\theta_1)
 \geq(1-s)\lambda_\mu(\theta_0)+s\lambda_\mu(\theta_1).
\]
Equality holds if and only if, for some $r>0$ and some smooth real $f$,
\[
 \theta_1=r\theta_0+\ddc f.
\]
\end{theorem}
\begin{proof}
Concavity was proved in \cref{prop:gaud-volume-elementary}. Suppose first that $n\geq2$ and equality holds. Take a minimizing Gauduchon metric for the middle background. Its quotient is the indicated convex combination of its two endpoint quotients. Each endpoint quotient is at least the corresponding infimum. Equality therefore forces this same metric to minimize both endpoint quotients.

By \cref{thm:gaud-volume-dual-formula}, the solution forms $\Omega_i=\Omega_{\theta_i}$, with scalars $c_i=c_{\theta_i}(\mu)$ for $i=0,1$, are conformal to this metric. Their volume forms are $c_0\mu$ and $c_1\mu$, so their pointwise proportionality factor is the constant $r=(c_1/c_0)^{1/n}$. Hence $\Omega_1=r\Omega_0$, which gives the stated smooth $\ddc$ relation. Conversely that relation, gauge invariance and homogeneity prove equality.

When $n=1$, the function is $\int_X\theta$. Take $r=\int_X\theta_1/\int_X\theta_0$. The difference $\theta_1-r\theta_0$ has integral zero, so the Green operator on the compact curve writes it as $\ddc f$ with $f$ smooth. Thus the same equivalence holds.
\end{proof}

\begin{proposition}\label{prop:gaud-volume-first-variation}
The function $\lambda_\mu$ is smooth along smooth finite-dimensional families in $\mathscr P$. For a smooth real $(1,1)$-form $\chi$ and $n\geq2$, its derivative at $\theta$ is
\[
 (\mathrm{d}\lambda_\mu)_\theta(\chi)
 =\frac{\displaystyle\int_X\chi\wedge\gamma_\theta^{n-1}}
        {D_\mu(\gamma_\theta)},
\]
where $\gamma_\theta$ is any minimizing Gauduchon metric. In dimension one the derivative is $\int_X\chi$.
\end{proposition}
\begin{proof}
Choose a fixed smooth potential $u_*$ making $\theta+\ddc u_*>0$. Nearby backgrounds have the positive representatives $\theta_t+\ddc u_*$. The implicit-function argument in the proof of \cref{prop:smooth-background-potential-stability} applies to this family and gives smooth dependence of the solution and its scalar. Adding $u_*$ to these solution potentials and fixing the mean gives the normalized solutions for $\theta_t$.

\Needspace{9\baselineskip}
Write $\lambda=\lambda_\mu(\theta)$, $\Omega=\Omega_\theta$, and $\gamma_\theta=\mathrm{e}^g\Omega$. Along $\theta_t=\theta+t\chi$ for small real $t$, put $\Omega_t=\theta_t+\ddc u_t$ and $\lambda_t=\lambda_\mu(\theta_t)$. Dots denote derivatives at $t=0$. Differentiating $\Omega_t^n=\lambda_t^n\mu$ gives
\[
 (\chi+\ddc\dot u)\wedge\Omega^{n-1}
                   =\lambda^{n-1}\dot\lambda\,\mu.
\]
Multiply by $\mathrm{e}^{(n-1)g}$ and integrate. The term with $\ddc\dot u$ vanishes by the Gauduchon condition. Moreover
\[
 D_\mu(\gamma_\theta)
       =\lambda^{n-1}\int_X \mathrm{e}^{(n-1)g}\mu.
\]
Substitution proves the derivative formula. The dimension-one formula follows directly from the definition.
\end{proof}

\section{Second variation on background classes}
\label{sec:gaud-volume-hessian}

Differentiating the smooth equation twice along an affine family of backgrounds computes the second derivative of the root volume. After integration against the adjoint probability of the solution form, the second derivative of the potential disappears, and what remains is the integrated squared norm of the linearized solution form. Consequently the Hessian of $-\log\lambda_\mu$ is positive definite once smooth $\ddc$ directions are factored out. On the cone of closed K\"ahler classes, where $\lambda_\mu(\theta)=(\int_X\theta^n)^{1/n}$, it is $1/n$ times the Hessian of $-\log\int_X\theta^n$.

Fix $\theta\in\mathscr P$ and write $\Omega=\Omega_\theta$, $u=u_\theta$ and $\lambda=\lambda_\mu(\theta)$. Let $\nu_\Omega$ be the adjoint probability of $\Omega$ from \cref{prop:adjoint-gauduchon}, that is, the smooth positive probability measure, given by the Gauduchon multiplier of $\Omega$, which satisfies $\int_X\operatorname{tr}_\Omega(\ddc v)\,\mathrm{d}\nu_\Omega=0$ for every smooth real $v$. In dimension one this probability is $\Omega/\int_X\Omega=\mu$.

Let $\theta_t=\theta+t\chi$, where $t$ is real and $\chi$ is a smooth real $(1,1)$-form on $X$. For sufficiently small $|t|$, the background $\theta_t$ stays in $\mathscr P$. We normalize the solution potentials $u_t$ by one fixed condition on their mean, and dots denote derivatives of $u_t$ and of $\lambda_\mu(\theta_t)$ at $t=0$. Define the Hermitian endomorphism
\[
 A_\chi=\Omega^{-1}(\chi+\ddc\dot u)
\]
of $T^{1,0}X$, where $\Omega^{-1}$ means raising an index with the positive form $\Omega$; traces of powers of $A_\chi$ are independent of coordinates. We write $I$ for the identity endomorphism of $T^{1,0}X$ and $\|\cdot\|_\Omega$ for the Hilbert--Schmidt norm on its endomorphisms induced by $\Omega$.

\Needspace{16\baselineskip}
\begin{proposition}\label{prop:gaud-volume-hessian}\index{Gauduchon variational volume!Hessian}
Along this affine variation,
\begin{align}
 \frac{\ddot\lambda}{\lambda}
   &=-\frac1n\int_X
      \left\|A_\chi-\frac{\dot\lambda}{\lambda}I\right\|_\Omega^2
                                      \,\mathrm{d}\nu_\Omega,\label{eq:gaud-root-hessian}\\
 -\frac{\mathrm{d}^2}{\mathrm{d}t^2}\log\lambda_\mu(\theta_t)
   &=\frac1n\int_X\operatorname{tr}(A_\chi^2)\,\mathrm{d}\nu_\Omega.
                                      \label{eq:gaud-log-hessian}
\end{align}
The first expression vanishes exactly when
$\chi=r\theta+\ddc f$ for a real constant $r$ and a smooth real $f$.
The second vanishes exactly when $\chi=\ddc f$ for a smooth real $f$.
\end{proposition}
\begin{proof}
Smooth dependence on $t$ was proved in \cref{prop:gaud-volume-first-variation}. We differentiate the logarithmic form $\log(\Omega_t^n/\mu)=n\log\lambda_\mu(\theta_t)$ of the equation. Since $\mu$ is fixed,
\[
 \operatorname{tr}A_\chi=n\frac{\dot\lambda}{\lambda}.
\]
In particular this trace is constant on $X$. Differentiating once more along the affine path $\theta_t$ gives
\[
 \operatorname{tr}_\Omega(\ddc\ddot u)
       -\operatorname{tr}(A_\chi^2)
   =n\left(\frac{\ddot\lambda}{\lambda}
                  -\frac{\dot\lambda^2}{\lambda^2}\right).
\]
Integration against the probability $\nu_\Omega$, which is fixed since everything is evaluated at $t=0$, removes the first term. This proves \eqref{eq:gaud-log-hessian}.

Put $r=\dot\lambda/\lambda$. Since $A_\chi$ is Hermitian and $\operatorname{tr}A_\chi=nr$,
\[
 \|A_\chi-rI\|_\Omega^2
      =\operatorname{tr}(A_\chi^2)-nr^2.
\]
Substituting this into \eqref{eq:gaud-log-hessian} gives \eqref{eq:gaud-root-hessian}. Since the adjoint measure has a strictly positive smooth density, the right side of \eqref{eq:gaud-root-hessian} vanishes exactly when $A_\chi=rI$ everywhere, that is, when
\[
 \chi+\ddc\dot u=r(\theta+\ddc u).
\]
This is the stated relation with $f=ru-\dot u$. Conversely, if $\chi=r\theta+\ddc f$, then $\theta_t$ differs from $(1+rt)\theta$ by a smooth $\ddc$ term for small $t$. Homogeneity gives $\lambda_\mu(\theta_t)=(1+rt)\lambda$, whose second derivative is zero.

The integrand in \eqref{eq:gaud-log-hessian} is the squared norm of the Hermitian endomorphism $A_\chi$. Its integral therefore vanishes exactly when $A_\chi=0$, that is, when $\chi=-\ddc\dot u$. Conversely, $\lambda_\mu(\theta+t\ddc f)=\lambda$ for all $t$. This proves both equality statements.
\end{proof}

The linearized solution and the derivative of the constant depend linearly on $\chi$. For smooth real $(1,1)$-forms $\chi_1,\chi_2$ on $X$, polarization therefore gives a symmetric bilinear form
\begin{equation}\label{eq:gaud-background-quadratic-form}
 \mathfrak g_\theta(\chi_1,\chi_2)
    =-D^2\log\lambda_\mu|_\theta(\chi_1,\chi_2)
    =\frac1n\int_X\operatorname{tr}(A_{\chi_1}A_{\chi_2})\,\mathrm{d}\nu_\Omega.
\end{equation}
We call the classes of smooth real $(1,1)$-forms modulo $\ddc C^\infty(X,\mathbb R)$ \emph{background classes}. By \cref{prop:gaud-volume-hessian}, the form $\mathfrak g_\theta$ descends to background classes of directions and is positive definite there. A change of representative of $\theta$ leaves its solution form $\Omega$ unchanged, and a change of representative of a direction is absorbed by the derivative of the solution potential. Thus $\mathfrak g$ is intrinsically defined on the background classes in $\mathscr P$. On every finite-dimensional affine family whose directions are linearly independent modulo smooth $\ddc$ terms, it is a smooth Riemannian metric.

These variations change the background class. Directions obtained by varying a potential inside one fixed class are exactly the directions annihilated by \eqref{eq:gaud-background-quadratic-form}. Thus this construction concerns a different space from the potential-space geometry of \cref{chap:energy}.

\Needspace{13\baselineskip}
\section{Recovering a positive class from its Gauduchon derivative}
\label{sec:gaud-volume-gradient-recovery}

By \cref{prop:gaud-volume-first-variation}, the derivative of $\lambda_\mu$ at $\theta\in\mathscr P$ is represented by a minimizing Gauduchon metric, which is determined up to a constant factor. The solution form is conformal to this metric and has a prescribed volume form, so the value and the derivative of $\lambda_\mu$ together reconstruct the solution form, and hence the background class. Fix $\theta\in\mathscr P$ and a minimizing metric $\gamma_\theta$.

For $n\geq2$ define the smooth positive $(n-1,n-1)$-form
\[
 \Xi_\theta=\frac{\gamma_\theta^{n-1}}{D_\mu(\gamma_\theta)}.
\]
Since numerator and denominator are both homogeneous of degree $n-1$ in $\gamma_\theta$, this form does not depend on the choice of the minimizing metric. It is $\ddc$-closed, and by \cref{prop:gaud-volume-first-variation} it represents the derivative in every smooth real $(1,1)$-direction $\chi$:
\[
 (\mathrm{d}\lambda_\mu)_\theta(\chi)=\int_X\chi\wedge\Xi_\theta.
\]

\begin{proposition}\label{prop:gaud-derivative-recovery}
Let $n\geq2$ and $\theta\in\mathscr P$. The value $\lambda_\mu(\theta)$ and the derivative of $\lambda_\mu$ at $\theta$ on all smooth real $(1,1)$-forms together determine the background class of $\theta$, and the derivative alone determines this class up to multiplication by a positive constant. More precisely, if $\widehat\gamma$ is the positive $(1,1)$-form with $\widehat\gamma^{n-1}=\Xi_\theta$, then
\[
 D_\mu(\widehat\gamma)=1,
 \qquad
 \Omega_\theta=
   \lambda_\mu(\theta)
       \left(\frac{\mu}{\widehat\gamma^n}\right)^{1/n}
       \widehat\gamma.
\]
Moreover, assigning to a background class in $\mathscr P$ with $\lambda_\mu=1$ the ray of its minimizing Gauduchon metrics is a bijection onto the set of rays of Gauduchon metrics.
\end{proposition}
\begin{proof}
The values of the derivative on all smooth $(1,1)$-forms determine the smooth form $\Xi_\theta$. Its positive root exists and is unique by the positive-root construction preceding the proof of \cref{thm:lamari}. If $D=D_\mu(\gamma_\theta)$, this root is $\widehat\gamma=D^{-1/(n-1)}\gamma_\theta$, and the homogeneity of $D_\mu$ gives $D_\mu(\widehat\gamma)=1$.

By \cref{thm:gaud-volume-dual-formula}, the solution form $\Omega_\theta$ is conformal to $\widehat\gamma$, and its volume form is $\lambda_\mu(\theta)^n\mu$. These two requirements determine the conformal factor pointwise and give the formula for $\Omega_\theta$. Thus the value and the derivative recover $\Omega_\theta$, a representative of the class of $\theta$. Without the value, the same formula recovers $\Omega_\theta$ up to a positive constant factor. Conversely, multiplying the class by a positive constant multiplies $\Omega_\theta$ by the same constant and leaves its minimizing Gauduchon ray, and hence $\Xi_\theta$, unchanged; so the derivative cannot determine more than the ray of the class.

For the surjectivity in the last assertion, take any Gauduchon metric $\gamma$ and set
\[
 \Omega_\gamma=\left(\frac\mu{\gamma^n}\right)^{1/n}\gamma.
\]
This Hermitian form is unchanged by constant rescaling of $\gamma$ and satisfies $\Omega_\gamma^n=\mu$. The zero potential therefore solves the equation for the background $\Omega_\gamma$ with constant one, so its class has $\lambda_\mu=1$, and $\gamma$ is a minimizing Gauduchon metric by \cref{thm:gaud-volume-dual-formula}. Injectivity follows from the recovery formula: two classes with $\lambda_\mu=1$ and the same Gauduchon ray have the same solution form. This proves the bijection.
\end{proof}

In dimension one, a background class is determined by its integral, by the Green operator of the curve. The derivative is then the fixed functional $\chi\mapsto\int_X\chi$, and the value alone determines the class.


\section{The Hessian metric and its completion}
\label{sec:gauduchon-metric-geometry}

We now study the metric $\mathfrak g$ of \eqref{eq:gaud-background-quadratic-form} on finite-dimensional cones of backgrounds, where it defines an ordinary Riemannian path distance. Homogeneity of $\lambda_\mu$ separates the scale of a positive class from its position on a level set of $\lambda_\mu$, and the logarithm of this scale turns out to be a Euclidean coordinate orthogonal to the level sets. The metric is therefore a product, and its completeness reduces to that of the unit-volume hypersurface.

Let $V$ be a finite-dimensional real vector space of smooth real $(1,1)$-forms such that
\[
 V\cap\ddc C^\infty(X,\mathbb R)=\{0\},
 \qquad \mathscr C=V\cap\mathscr P\ne\varnothing.
\]
Then $\mathscr C$ is an open convex cone in $V$. Write $\lambda=\lambda_\mu|_{\mathscr C}$, and let $g=-D^2\log\lambda$ be the restriction to $\mathscr C$ of the metric \eqref{eq:gaud-background-quadratic-form}; it is a smooth Riemannian metric on $\mathscr C$ because $V$ meets $\ddc C^\infty(X,\mathbb R)$ only in zero. Put
\[
 \mathscr S=\{\theta\in\mathscr C:\lambda(\theta)=1\},
\]
and equip this smooth hypersurface with the induced metric $g_{\mathscr S}$.

\begin{proposition}\label{prop:gaud-volume-metric-product}\index{Gauduchon variational volume!metric geometry}
The map
\[
 \Phi:\mathbb R\times\mathscr S\longrightarrow\mathscr C,
 \qquad (r,q)\longmapsto \mathrm{e}^r q
\]
is an isometry from the product metric $\mathrm{d}r^2+g_{\mathscr S}$ onto $g$. In particular, if $d_g$ and $d_{\mathscr S}$ are the intrinsic distances, then
\[
 d_g(\mathrm{e}^{r_0}q_0,\mathrm{e}^{r_1}q_1)^2
    =|r_1-r_0|^2+d_{\mathscr S}(q_0,q_1)^2.
\]
The metric on $\mathscr C$ is complete if and only if the induced metric on $\mathscr S$ is complete. Every path for which $\lambda$ tends to zero or to infinity has infinite length.
\end{proposition}
\begin{proof}
Homogeneity gives $\lambda(s\theta)=s\lambda(\theta)$ for $s>0$. Hence $r=\log\lambda(\theta)$ and $q=\theta/\lambda(\theta)$ define the smooth inverse of $\Phi$. The level set is smooth because the Euler identity gives $D\lambda_\theta(\theta)=\lambda(\theta)>0$.

Differentiating the Euler identity $D\log\lambda_\theta(\theta)=1$ in a direction $\chi\in V$ gives
\[
 g_\theta(\theta,\chi)=D\log\lambda_\theta(\chi),
 \qquad g_\theta(\theta,\theta)=1.
\]
For directions $\chi,\eta\in V$, differentiating homogeneity twice also gives
\[
 g_{s\theta}(s\chi,s\eta)=g_\theta(\chi,\eta).
\]
Let $q\in\mathscr S$, let $\xi,\eta\in T_q\mathscr S$, and let $a,b,r$ be real numbers. Since $\xi$ and $\eta$ are tangent to a level set of $\lambda$, the first of these identities shows that both are $g_q$-orthogonal to the radial vector $q$. Using also the second identity and the scaling rule, we obtain
\[
 g_{\mathrm{e}^rq}\bigl(\mathrm{e}^r(aq+\xi),\mathrm{e}^r(bq+\eta)\bigr)
       =ab+g_q(\xi,\eta),
\]
which proves the metric identity.

It remains to prove the distance formula and its consequences. A path $(r(t),q(t))$ in the product has length
$\int\sqrt{|r'|^2+|q'|_{g_{\mathscr S}}^2}\,\mathrm{d}t$.
The triangle inequality in $\mathbb R^2$ (applied to the integral of the vector $(|r'|,|q'|_{g_{\mathscr S}})$) bounds this below by the square root of the sum of the squared endpoint difference in $r$ and the squared length of the projected path in $\mathscr S$. For the opposite bound, choose a path in $\mathscr S$ of length within $\epsilon$ of the distance of its endpoints, parametrize it with constant speed, and let $r$ vary linearly; then let $\epsilon$ decrease to zero. This proves the distance formula. It identifies Cauchy sequences in $\mathscr C$ with pairs of Cauchy sequences in the two factors; since $\mathbb R$ is complete, the completeness assertion follows. Finally, the length of any path is at least the total change of $\log\lambda$ along it, and this change is infinite when $\lambda$ tends to zero or to infinity.
\end{proof}

\subsection{Product and blowup examples}
\label{sec:gaud-volume-profile-example}

The following two examples of closed classes show that the unit-volume hypersurface may be complete or may have a boundary at finite distance. The first also shows why the derivative is needed in \cref{prop:gaud-derivative-recovery}: knowing the volume for every smooth probability measure need not identify the class.

For the first example take $X=\mathbb P^1\times\mathbb P^1$, and let $\alpha,\beta$ be the pulled-back Fubini--Study forms of integral one on their respective factors. They satisfy
\[
 \alpha^2=\beta^2=0,\qquad\int_X\alpha\wedge\beta=1.
\]
For $a,b>0$ the form $\theta_{a,b}=a\alpha+b\beta$ is K\"ahler. Since it is closed, integrating the smooth equation gives, for every smooth positive probability measure $\mu$,
\[
 \mathcal V_\mu(\theta_{a,b})=2ab,
 \qquad \lambda_\mu(\theta_{a,b})=\sqrt{2ab}.
\]
Consequently, for real $t$, all the forms $\theta_{\mathrm{e}^t,\mathrm{e}^{-t}}$ have volume $2$ for every $\mu$. Their background classes are nevertheless distinct: for any fixed point $p$ of the second factor, integration over $\mathbb P^1\times\{p\}$ gives $\mathrm{e}^t$, and this integral is unchanged by the addition of $\ddc f$.

\Needspace{6\baselineskip}
By direct differentiation, the restriction of \eqref{eq:gaud-background-quadratic-form} to this two-dimensional cone is
\[
 \mathfrak g=\frac12\left(\frac{\mathrm{d}a^2}{a^2}+\frac{\mathrm{d}b^2}{b^2}\right).
\]
The coordinates $((\log a)/\sqrt2,(\log b)/\sqrt2)$ identify this metric with the Euclidean metric on $\mathbb R^2$, so the metric on this family is complete. The curve $a=\mathrm{e}^t$, $b=\mathrm{e}^{-t}$ has unit speed. All distances here are intrinsic to this two-dimensional family.

In the second example, the blowup of $\mathbb P^2$ at a point, the cone is incomplete, and the boundary added by the completion consists of classes of positive volume.

\begin{example}\label{ex:gaud-metric-blowup}
Let $\pi:X\to\mathbb P^2$ be the blowup at one point. Write $H$ for the pullback of the hyperplane class and $E$ for the exceptional divisor class. In the two-dimensional space they span, the cone of classes with a positive smooth representative is
\[
 \mathscr C=\{aH-bE:a>b>0\}.
\]
These are closed classes, and we represent them by fixed smooth forms in order to apply \cref{prop:gaud-volume-metric-product}. Their intersection numbers are
$H^2=1$, $H\cdot E=0$, $E^2=-1$.
Consequently, for every smooth positive probability measure $\mu$,
\[
 \lambda_\mu(aH-bE)=\sqrt{a^2-b^2}.
\]
In the coordinates
\[
 a=\mathrm{e}^r\cosh s,\qquad b=\mathrm{e}^r\sinh s,
 \qquad r\in\mathbb R,\quad s>0,
\]
the Hessian metric is $\mathrm{d}r^2+\mathrm{d}s^2$. Its completion is therefore the Euclidean closed half-plane. The added boundary $s=0$ consists of the nef classes $\mathrm{e}^rH$, each of positive volume $\mathrm{e}^{2r}$.
\end{example}
\begin{proof}
The projection from the blown-up point extends to a morphism $f:X\to\mathbb P^1$ with class $L=H-E$. The incidence map $(\pi,f):X\to\mathbb P^2\times\mathbb P^1$ is an embedding. If $h$ and $\ell$ are pullbacks of Fubini--Study forms of integral one, then
\[
 (a-b)h+b\ell
\]
is positive for $a>b>0$, by the product embedding. Conversely, a positive representative has positive integrals over the exceptional curve and the strict transform of a line through the blown-up point. These integrals are respectively $b$ and $a-b$, proving the cone description. The same curves show that no nonzero linear combination in this two-dimensional space is smoothly $\ddc$-exact.

\Needspace{10\baselineskip}
The intersection numbers follow from the hyperplane pullback and the normal bundle $\mathcal O_{\mathbb P^1}(-1)$ of the exceptional curve; equivalently $H^2=H\cdot L=1$ and $L^2=0$. Since the classes are closed, integrating the smooth equation gives $\mathcal V_\mu=a^2-b^2$, independently of $\mu$. Direct differentiation of $-\log\sqrt{a^2-b^2}$ gives
\[
 g=\frac{(a^2+b^2)(\mathrm{d}a^2+\mathrm{d}b^2)-4ab\,\mathrm{d}a\,\mathrm{d}b}
            {(a^2-b^2)^2}.
\]
Substituting the stated coordinates gives $\mathrm{d}r^2+\mathrm{d}s^2$. The coordinate domain is exactly $\mathbb R\times(0,\infty)$, so the intrinsic distance and completion are those of that Euclidean half-plane. Finally, $H$ has a semipositive pullback representative and is a limit of positive classes, hence is nef. Its square is one, but its integral over $E$ is zero, so it has no positive representative. This proves all assertions about the added boundary.
\end{proof}

\subsection{Invariant backgrounds on compact Lie-group quotients}
\label{sec:parallelizable-gauduchon-cones}

A normalized invariant volume reduces the Gauduchon quotient to a matrix determinant. The resulting metric controls the eigenvalues of the normalized minimizing form. This gives complete cones of nonclosed backgrounds, with an explicit distinction between the coefficient boundary and the unit-volume completion.

Let $G$ be a connected complex Lie group of dimension $n\geq2$, and let $X=\Gamma\backslash G$ be compact, where the discrete subgroup $\Gamma$ acts on the left. Choose a left-invariant holomorphic coframe $\alpha_1,\ldots,\alpha_n$, which descends to $X$. For a Hermitian matrix $A$ put
\[
 \begin{gathered}
 \theta_A=\frac{\mathrm i}{2\pi}\sum_{j,k}A_{jk}\alpha_j\wedge\bar\alpha_k,
 \qquad h_0=\theta_I,\\
 H=\int_Xh_0^n,\qquad \mu=\frac{h_0^n}{H}.
 \end{gathered}
\]
Fix a real linear subspace $W\subset\operatorname{Herm}(n)$ containing a positive matrix, and write $V_W=\{\theta_A:A\in W\}$. Here the probability volume is the normalized invariant volume just specified.

\Needspace{12\baselineskip}
\begin{proposition}\label{prop:parallelizable-gaud-volume}\index{Gauduchon variational volume!invariant backgrounds}
The backgrounds in $V_W$ admitting a smooth positive representative after a smooth $\ddc$ change are precisely $\theta_A$ with $A>0$. Moreover,
\[
 V_W\cap\ddc C^\infty(X,\R)=\{0\}.
\]
Every positive $\theta_A$ is balanced, and its unique minimizing Gauduchon metric normalized by $D_\mu=1$ is
\begin{equation}\label{eq:parallelizable-gaud-volume}
 \lambda_\mu(\theta_A)=(H\det A)^{1/n},\qquad
 \gamma_A=\frac{\theta_A}{\lambda_\mu(\theta_A)},\qquad
 \gamma_A^n=\mu.
\end{equation}
On this coefficient cone the Hessian metric is
\begin{equation}\label{eq:parallelizable-gaud-hessian}
 g_A(U,Z)=\frac1n\operatorname{tr}(A^{-1}UA^{-1}Z),
 \qquad U,Z\in W.
\end{equation}

For $A\in W$, a positive-current representative $\theta_A+\ddc u$, with a real distributional potential $u$, exists exactly when $A\geq0$. The coefficient closure is $W\cap\{A\geq0\}$, and its canonical representative is the smooth semipositive form $\theta_A$. The zero background has only the zero positive-current representative. If $G$ is nonabelian, every positive background in this cone is nonclosed and has no closed representative obtained by adding a smooth $\ddc$ potential.
\end{proposition}

\begin{proof}
\emph{Invariant averages and balancedness.}
Let $Z_j$ be dual to $\alpha_j$, and set $\Upsilon=\alpha_1\wedge\cdots\wedge\alpha_n$. This holomorphic top form and its conjugate are closed. For an invariant holomorphic vector field $Z$, evaluation on the invariant frame gives
\[
 \mathcal L_Z\Upsilon=-\operatorname{tr}(\operatorname{ad}Z)\Upsilon,
 \qquad \mathcal L_Z\bar\Upsilon=0.
\]
The second identity also follows from Cartan's formula. Since $h_0^n$ is a constant multiple of $\Upsilon\wedge\bar\Upsilon$, with the positive orientation factor, compact Stokes gives
\[
 0=\int_X\mathcal L_Zh_0^n=-H\operatorname{tr}(\operatorname{ad}Z).
\]
Thus the fields and their conjugates have zero divergence. In particular,
\begin{equation}\label{eq:parallelizable-zero-mean}
 \int_XZ_jf\,\mu=\int_X\bar Z_jf\,\mu=0
 \qquad(f\in C^\infty(X,\mathbb C)).
\end{equation}
Cartan's formula also gives $\mathrm d(\iota_{Z_j}\Upsilon)=0$. These contractions span the invariant $(n-1,0)$-forms. The form $\theta_A^{n-1}$ is a constant linear combination of their wedges with conjugates, so $\mathrm d\theta_A^{n-1}=0$.

The holomorphic fields commute with the antiholomorphic fields, and $\mathrm d\alpha_j$ has type $(2,0)$. The coefficient matrix of $\ddc u$ is therefore $(Z_j\bar Z_k u)$. For smooth $u$ its entries have mean zero by \eqref{eq:parallelizable-zero-mean}. The mean matrix of $\theta_A+\ddc u$ is $A$. Smooth positivity implies positivity of the mean quadratic form on every nonzero constant vector, hence $A>0$. The converse follows by taking $u=0$. The same average proves the absence of nonzero smooth $\ddc$ directions in $V_W$.

For distributions the adjoints with respect to $\mu$ annihilate constants, so the mean of each $Z_j\bar Z_k u$ is still zero. The coefficients of a positive current in this global frame are a positive matrix of measures. Their mean is semipositive, giving $A\geq0$. Conversely $\theta_A$ itself is positive when $A\geq0$. If $A=0$, the positive trace measure has total mass zero; polarization then makes the entire current zero. Finally, for a fixed positive $A_+\in W$, the matrices $A+\varepsilon A_+$ approximate every semipositive $A\in W$ from the positive cone. This proves the closure assertion.

\emph{The quotient and its equality case.}
Put $Q=H\det A$. The determinant identity is $\theta_A^n=Q\mu$. For every smooth positive Hermitian form $\gamma$, the matrix arithmetic--geometric mean inequality gives
\[
 \frac{\theta_A\wedge\gamma^{n-1}}{\gamma^n}
 =\frac1n\operatorname{tr}_\gamma\theta_A
 \geq\left(\frac{\theta_A^n}{\gamma^n}\right)^{1/n}.
\]
Integration yields $\int_X\theta_A\wedge\gamma^{n-1}\geq Q^{1/n}D_\mu(\gamma)$. This bounds every Gauduchon test in \cref{def:gaud-variational-volume}. The balanced test $\gamma=\theta_A$ has $D_\mu(\theta_A)=Q^{(n-1)/n}$ and attains equality, proving the root-volume formula.

At a minimizing smooth Gauduchon test the continuous nonnegative deficit has integral zero. Equality pointwise forces $\gamma=f\theta_A$ for a smooth positive $f$. Set $q=f^{n-1}$ and $T=\theta_A^{n-1}$. Since $\mathrm dT=0$, the Gauduchon equation becomes $\ddc q\wedge T=0$, and
\[
 0=\int_Xq\ddc q\wedge T=-\int_X\mathrm dq\wedge\dc q\wedge T.
\]
The last integrand is nonnegative and vanishes exactly where $\mathrm dq=0$. Thus $f$ is constant. As $n\geq2$ and
\[
 D_\mu(c\theta_A)=c^{n-1}\lambda_\mu(\theta_A)^{n-1},
\]
normalization determines $c=\lambda_\mu(\theta_A)^{-1}$. This proves the unique minimizer and its volume identity.

Differentiating $\log\lambda_\mu(\theta_A)=(\log H+\log\det A)/n$ gives \eqref{eq:parallelizable-gaud-hessian}. Its square norm is $\|A^{-1/2}UA^{-1/2}\|_{\mathrm F}^2/n$, which is positive for $U\ne0$. Here $\|B\|_{\mathrm F}^2=\operatorname{tr}(B^*B)$ is the Frobenius norm.

Finally,
\[
 \partial\theta_A=\frac{\mathrm i}{2\pi}\sum_{j,k}A_{jk}\,\mathrm d\alpha_j\wedge\bar\alpha_k.
\]
For invertible $A$, vanishing forces every $\mathrm d\alpha_j=0$, by independence of the conjugate coframe. The Lie algebra would then be abelian. Addition of $\ddc u$ does not change $\mathrm d\theta_A$, proving the last assertion.
\end{proof}

\Needspace{12\baselineskip}
\begin{lemma}\label{lem:positive-matrix-hessian-distance}\index{Gauduchon variational volume!matrix distance}
The full positive Hermitian matrix cone, with the metric in \eqref{eq:parallelizable-gaud-hessian}, has distance
\begin{equation}\label{eq:positive-matrix-hessian-distance}
 d_{\mathrm{full}}(A,B)=\frac1{\sqrt n}
       \|\log(A^{-1/2}BA^{-1/2})\|_{\mathrm F}.
\end{equation}
For endpoints with equal determinant, this is also the intrinsic distance in their full fixed-determinant hypersurface.
\end{lemma}

\begin{proof}
Fixed congruence is an isometry by the trace formula. Reduce to endpoints $I,C>0$. The curve $\exp(t\log C)$ has constant speed $\|\log C\|_{\mathrm F}/\sqrt n$.

For the reverse inequality, diagonalization gives the convergent resolvent formula and its derivative
\begin{align*}
 \log M&=\int_0^\infty\bigl((1+s)^{-1}I-(M+sI)^{-1}\bigr)\,\mathrm ds,\\
 D\log_M[U]&=\int_0^\infty(M+sI)^{-1}U(M+sI)^{-1}\,\mathrm ds.
\end{align*}
On a neighborhood with a positive spectral lower bound, differentiation of the inverse is dominated by an integrable function. In a unitary basis diagonalizing $M$, with eigenvalues $d_i>0$, the derivative has entries $c_{ij}U_{ij}$, where
\[
 c_{ij}=\frac{\log d_i-\log d_j}{d_i-d_j}\quad(d_i\ne d_j),
 \qquad c_{ij}=\frac1{d_i}\quad(d_i=d_j).
\]
Writing $d_i/d_j=\mathrm e^{2s}$ gives $\sqrt{d_id_j}|c_{ij}|=|s/\sinh s|\leq1$, including the limiting value at $s=0$. Hence
\[
 \|D\log_M[U]\|_{\mathrm F}
 \leq\|M^{-1/2}UM^{-1/2}\|_{\mathrm F}.
\]
For any piecewise $C^1$ positive matrix path from $I$ to $C$, integration bounds its length below by $\|\log C\|_{\mathrm F}/\sqrt n$. This proves the distance formula. Restoring the endpoints, the attaining curve is
\begin{equation}\label{eq:positive-matrix-exponential-path}
 M(t)=A^{1/2}(A^{-1/2}BA^{-1/2})^tA^{1/2}.
\end{equation}
Its determinant is $(\det A)^{1-t}(\det B)^t$. Thus for equal-determinant endpoints it stays in the indicated hypersurface.
\end{proof}

Write $\lambda(A)=\lambda_\mu(\theta_A)$ and
\[
 S_W=\{A\in W:A>0,\ H\det A=1\}.
\]
The radial derivative of $\lambda$ is positive, so $S_W$ is a smooth hypersurface of the coefficient cone. All distances $d_W$ below are intrinsic path distances on this fixed hypersurface with its induced Hessian metric.

\Needspace{12\baselineskip}
\begin{theorem}\label{thm:parallelizable-unit-volume-completion}\index{Gauduchon variational volume!complete invariant cone}
The space $(S_W,d_W)$ is complete and proper. A sequence is Cauchy exactly when its coefficients converge to a positive matrix in $S_W$. For $A,B\in S_W$,
\begin{align}
 d_W(A,B)&\geq\frac1{\sqrt n}
       \|\log(A^{-1/2}BA^{-1/2})\|_{\mathrm F},\label{eq:parallelizable-intrinsic-lower}\\
 d_W(A,B)&\leq\frac{\|A-B\|_{\mathrm F}}
       {\sqrt n\min\{\lambda_{\min}(A),\lambda_{\min}(B)\}}.
       \label{eq:parallelizable-intrinsic-upper}
\end{align}
The actual normalized minimizing forms satisfy pointwise on $X$
\begin{equation}\label{eq:parallelizable-minimizer-order}
 \mathrm e^{-\sqrt n d_W(A,B)}\gamma_A
 \leq\gamma_B\leq\mathrm e^{\sqrt n d_W(A,B)}\gamma_A.
\end{equation}
Thus every completion point has its canonical smooth positive current representative, and the completion adds no semipositive coefficient boundary point.

If $A_j>0$ in $W$ converges to a nonzero singular semipositive $A_0$, then $\widehat A_j=A_j/\lambda(A_j)$ lies in $S_W$ and, for fixed $A_\star\in S_W$,
\begin{equation}\label{eq:parallelizable-singular-boundary-rate}
 d_W(A_\star,\widehat A_j)
 \geq\frac{|\log\det A_j|}{n\sqrt n}+O(1)
 \longrightarrow\infty.
\end{equation}
Normalization of a sequence tending to zero can instead converge to an existing positive point of $S_W$.
\end{theorem}

\begin{proof}
\emph{Paths in the specified subspace.}
Every path in $S_W$ is a path in the full fixed-determinant hypersurface, so \cref{lem:positive-matrix-hessian-distance} gives \eqref{eq:parallelizable-intrinsic-lower}. To obtain an upper bound within $W$, put
\[
 U=B-A,\qquad C(t)=A+tU,\qquad N(t)=\frac{C(t)}{\lambda(C(t))}.
\]
The positive cone is convex and $W$ is linear, so $N(t)$ stays in $S_W$ and joins $A$ to $B$. If $K=C^{-1/2}UC^{-1/2}$, differentiation gives
\[
 \begin{gathered}
 g_N(N',N')=\frac1n\operatorname{tr}(K^2)
                 -\frac1{n^2}(\operatorname{tr}K)^2
 \leq\frac{\|U\|_{\mathrm F}^2}{nm^2},\\
 m=\min\{\lambda_{\min}(A),\lambda_{\min}(B)\}.
 \end{gathered}
\]
Indeed $C(t)\geq mI$ and $(\log\lambda(C))'=\operatorname{tr}(C^{-1}U)/n$. Integrating proves \eqref{eq:parallelizable-intrinsic-upper}, and also connectedness of $S_W$.

Put $d=d_W(A,B)$. The logarithmic spectral lower bound gives
\begin{equation}\label{eq:parallelizable-coefficient-order}
 \mathrm e^{-\sqrt n d}A\leq B\leq\mathrm e^{\sqrt n d}A.
\end{equation}
Since $\gamma_A=\theta_A$ on $S_W$, this proves the pointwise minimizing-form estimate.

\emph{Cauchy sequences and bounded balls.}
For an intrinsic Cauchy sequence, comparison with one fixed tail element gives $mI\leq A_j\leq MI$ with fixed positive $m,M$. Comparing two tail elements then gives
\[
 \|A_j-A_k\|_{\mathrm{op}}
 \leq M\bigl(\mathrm e^{\sqrt n d_W(A_j,A_k)}-1\bigr).
\]
Its coefficients are Cauchy. Their limit lies in the closed linear space $W$, remains positive, and has $H\det A=1$. The intrinsic upper bound proves convergence to that point. Conversely coefficient convergence to a positive point gives intrinsic convergence by the same bound. This proves completeness and the exact Cauchy criterion. The minimizing forms converge in $C^\infty$, because their coefficients multiply a fixed smooth coframe.

The two topologies agree by these estimates. A closed ball of radius $R$ around $A_\star$ is contained in the compact matrix interval
\[
 \begin{gathered}
 \mathrm e^{-\sqrt nR}A_\star\leq A\leq\mathrm e^{\sqrt nR}A_\star,\\
 A\in W,\qquad H\det A=1.
 \end{gathered}
\]
It is closed in that interval, hence compact. This proves properness. Distances from a fixed reference are bounded exactly when the relative eigenvalues are bounded above and away from zero.

\emph{The coefficient boundary.}
For $A_j\to A_0\geq0$ singular and nonzero, $\lambda(A_j)\to0$ while $\lambda_{\max}(A_j)$ has a positive limit. The largest eigenvalue of $\widehat A_j$ tends to infinity. Applying \eqref{eq:parallelizable-coefficient-order} to $A_\star,\widehat A_j$ gives
\[
 \begin{aligned}
 d_W(A_\star,\widehat A_j)
 &\geq\frac1{\sqrt n}\log
       \frac{\lambda_{\max}(A_j)}{\lambda(A_j)\lambda_{\max}(A_\star)}\\
 &=\frac{|\log\det A_j|}{n\sqrt n}+O(1).
 \end{aligned}
\]
The error is bounded because $A_0\ne0$. This proves \eqref{eq:parallelizable-singular-boundary-rate}. For comparison, $A_j=t_jA_\star$ with $t_j\downarrow0$ normalizes identically to $A_\star$. Its raw zero-current limit supplies no new point of the unit-volume completion.
\end{proof}

By \cref{prop:gaud-volume-metric-product}, the full prescribed cone is complete as well. The invariant probability and the fixed linear subspace are the hypotheses used here; other geometries can have the finite-distance boundary of \cref{ex:gaud-metric-blowup}.

The exponential path need not stay in a proper $W$. For $n=3$, take
\[
 W=\operatorname{span}_{\R}\{I,\operatorname{diag}(0,1,2)\},
 \qquad a=H^{-1/3}.
\]
Its unit-volume endpoints $aI$ and $a6^{-1/3}\operatorname{diag}(1,2,3)$ have full-matrix midpoint $a6^{-1/6}\operatorname{diag}(1,\sqrt2,\sqrt3)$. The diagonal entries fail the relation $x_1-2x_2+x_3=0$, since $1-2\sqrt2+\sqrt3\ne0$. Thus this path leaves $W$, whereas the normalized chord in the proof remains in it.

\Needspace{12\baselineskip}
\begin{example}\label{ex:iwasawa-complete-gaud-cone}\index{Iwasawa manifold!Gauduchon background cone}
On $G=\mathbb C^3$ take the multiplication
\[
 z\cdot w=(z_1+w_1,z_2+w_2,z_3+w_3+z_1w_2),
\]
the Gaussian lattice $\Gamma=\mathbb Z[\mathrm i]^3$, and $X=\Gamma\backslash G$. This is a compact non-K\"ahler manifold of dimension three, with invariant coframe
\[
 \alpha_1=\mathrm dz_1,\qquad \alpha_2=\mathrm dz_2,\qquad
 \alpha_3=\mathrm dz_3-z_1\mathrm dz_2.
\]
For its full nine-dimensional cone $W=\operatorname{Herm}(3)$ and the normalized invariant probability,
\[
 H=\frac6{\pi^3},\qquad
 \mathcal V_\mu(\theta_A)=\frac6{\pi^3}\det A,\qquad
 S=\{A>0:\det A=\pi^3/6\}.
\]
Every positive background is nonclosed. The eight-dimensional unit-volume hypersurface is complete and proper, with exact intrinsic distance
\begin{equation}\label{eq:iwasawa-unit-volume-distance}
 d_S(A,B)=\frac1{\sqrt3}\|\log(A^{-1/2}BA^{-1/2})\|_{\mathrm F}.
\end{equation}
\end{example}

\begin{proof}
The bilinear cross term makes the multiplication associative, with identity zero and inverse $(-z_1,-z_2,-z_3+z_1z_2)$. These operations preserve the Gaussian lattice. Its action is free and properly discontinuous: elements with $\gamma K\cap K\ne\varnothing$ belong to the compact set $KK^{-1}$, containing only finitely many lattice elements.

Choose Gaussian integers $\gamma_1,\gamma_2$ to reduce $z_1,z_2$ to their half-open unit squares, then $\gamma_3$ to reduce $z_3+\gamma_1z_2$. This gives a representative in the real six-dimensional unit cube. Two representatives in the half-open cube force $\gamma_1=\gamma_2=0$, then $\gamma_3=0$. The cube is therefore a measurable fundamental domain, and the image of its closure proves compactness.

The displayed coframe is left-invariant, descends to $X$, and has
\[
 \mathrm d\alpha_3=-\alpha_1\wedge\alpha_2,\qquad
 \alpha_1\wedge\alpha_2\wedge\alpha_3
       =\mathrm dz_1\wedge\mathrm dz_2\wedge\mathrm dz_3.
\]
The density of $h_0^3$ is $6/\pi^3$ times real Lebesgue volume on the fundamental cube, giving the stated $H$. The coefficient of $\alpha_1\wedge\alpha_2\wedge\bar\alpha_3$ in $\partial\theta_A$ is $-\mathrm iA_{33}/(2\pi)$, so positivity makes every such background nonclosed.

To exclude an arbitrary K\"ahler metric, suppose a positive closed $\kappa$ existed and average its coefficient functions in this coframe against $\mu$. Their mean is a positive matrix $A$. In the coefficients of $\partial\kappa$, derivative terms have mean zero by \eqref{eq:parallelizable-zero-mean}; the remaining means are exactly the coefficients of $\partial\theta_A$. Closedness would force them to vanish, contradicting $A_{33}>0$. Thus $X$ is non-K\"ahler.

All volume, minimizer and completeness assertions follow from \cref{prop:parallelizable-gaud-volume,thm:parallelizable-unit-volume-completion}. As $W$ is the full matrix space, the equal-determinant exponential path stays in $S$, and \cref{lem:positive-matrix-hessian-distance} gives \eqref{eq:iwasawa-unit-volume-distance}. For $A_\star=aI$, $a=(\pi^3/6)^{1/3}$, this becomes
\[
 d_S(A_\star,A)^2=\frac13\sum_{j=1}^3
       \left(\log\frac{\lambda_j(A)}a\right)^2.
\]
The three logarithms sum to zero.

For a concrete degeneration take $A_\varepsilon=\operatorname{diag}(1,1,\varepsilon)$, $0<\varepsilon\leq1$. Its coefficient limit gives the nonzero smooth semipositive current $\theta_{\operatorname{diag}(1,1,0)}$. The normalized minimizing matrices and their distances are instead
\begin{equation}\label{eq:iwasawa-diagonal-degeneration}
 \begin{gathered}
 \widehat A_\varepsilon
   =a\operatorname{diag}(\varepsilon^{-1/3},\varepsilon^{-1/3},\varepsilon^{2/3}),\\
 d_S(A_\star,\widehat A_\varepsilon)=\frac{\sqrt2}{3}|\log\varepsilon|.
 \end{gathered}
\end{equation}
The coefficient boundary current therefore lies at infinite normalized parameter distance.
\end{proof}


\subsection{Warped product cones and flat completion}
\label{sec:warped-gauduchon-completion}

The invariant cones above have complete unit-volume hypersurfaces. On the product of a compact K\"ahler manifold with a curve, a nonconstant coefficient in the fiber direction gives cones of nonclosed backgrounds whose unit-volume hypersurfaces are flat and may instead have a Euclidean strip or half-plane as their metric completion. The volume is the reciprocal of an integral of the reciprocal fiber coefficient, and the variance of that reciprocal coefficient measures the distance to the boundary. The probability measure in this variance is a trace of a minimizing balanced metric.

Let $(B,\kappa)$ be a connected compact K\"ahler manifold of dimension $m\geq1$, with $\int_B\kappa^m=1$. Let $C$ be a connected compact complex curve with a smooth positive $(1,1)$-form $w$ of integral one. Fix a smooth strictly positive probability form $\sigma$ on $B$ and a smooth nonconstant real function $a$ on $B$ with $\min a=0$ and $\max a=M>0$. On $X=B\times C$, of dimension $n=m+1$, put $\mu=\sigma\wedge w$; we omit pullbacks from the notation. For real parameters $b,c,e$, the backgrounds in the real span of $\kappa,aw,w$ are
\[
 \theta_{b,c,e}=b\kappa+Lw,\qquad L=ca+e.
\]

\begin{theorem}\label{thm:warped-gauduchon-volume}\index{Gauduchon variational volume!warped product cones}
This span contains no nonzero smooth $\ddc$ direction, and its intersection with $\mathscr P$ is the cone
\[
 \mathscr C=\{b>0,\ e>0,\ e+cM>0\}.
\]
For $\theta=\theta_{b,c,e}\in\mathscr C$, let $H(c,e)=\int_B\sigma/L$. Then
\begin{equation}\label{eq:warped-gauduchon-volume}
 \mathcal V_\mu(\theta)=\frac{nb^m}{H(c,e)}.
\end{equation}
The minimizing Gauduchon metrics are the constant multiples of the balanced metric
\[
 \gamma=L^{-1/m}(A+Lw),\qquad
 A=b\kappa+\ddc\phi>0,\qquad
 A^m=\frac{b^m}{H}\frac{\sigma}{L},
\]
where $\phi$ is a smooth real function on $B$. Backgrounds with $c\ne0$ are nonclosed.
\end{theorem}
\begin{proof}
Pairing a smooth positive representative $\theta+\ddc u$ with the closed semipositive form $\kappa^{m-1}\wedge w$ gives $b>0$, and its integral over the compact fiber $\{y\}\times C$ is $L(y)>0$, by Stokes' theorem. These are exactly the stated inequalities; conversely, they make $\theta$ itself positive. The same tests applied to a smooth $\ddc$ direction give $b=0$ and $ca+e=0$ everywhere, and since $a$ is nonconstant, $c=e=0$. Also $\mathrm{d}\theta=c\,\mathrm{d}a\wedge w$, which is not identically zero when $c\ne0$.

Next we apply \cref{thm:tw-smooth} on $B$ to the closed positive background $b\kappa$ and the density $(b^m/H)\sigma/L$. Since this density has integral $b^m=\int_B(b\kappa)^m$, Stokes' theorem forces the constant in the equation to be one, and we obtain $A$. Put $\Omega=A+Lw$. Since $A^{m+1}=0$ and $(Lw)^2=0$, the binomial expansion gives
\begin{equation}\label{eq:warped-balanced-cofactor}
 \Omega^n=\frac{nb^m}{H}\mu,\qquad
 \gamma^m=L^{-1}A^m+mA^{m-1}\wedge w.
\end{equation}
The first summand is the pullback of a top-degree form on $B$, and the second is closed. Hence $\gamma$ is balanced, in particular Gauduchon, and it is conformal to the solution form $\Omega=\theta+\ddc\phi$ of $\theta$. Since $n\geq2$, \cref{thm:gaud-volume-dual-formula} proves \eqref{eq:warped-gauduchon-volume} and the description of the minimizers.
\end{proof}

The minimizing metric gives the probability identity
\begin{equation}\label{eq:warped-minimizer-trace}
 \frac{Lw\wedge\gamma^m}{b^m}
       =\frac{\sigma\wedge w}{HL}.
\end{equation}
Its marginal $\sigma/(HL)$ on $B$ is the probability measure that governs the Hessian metric. To compute that metric we use the ratio $v=c/e$ and set
\begin{align*}
 v&=c/e\in(-1/M,\infty),
 &Q(v)&=\int_B\frac{\sigma}{1+va},\\
 \mathrm{d}\nu_v&=\frac{\sigma}{Q(v)(1+va)},
 &\ell_v&=\frac{a}{1+va}.
\end{align*}

\begin{proposition}\label{prop:warped-flat-unit-metric}\index{Hessian metric!flat warped cones}
On the unit-volume hypersurface $\mathscr S=\{\theta\in\mathscr C:\mathcal V_\mu(\theta)=1\}$, the functions $U=\log b$ and $v$ are global coordinates, in which
\begin{equation}\label{eq:warped-flat-unit-metric}
 \begin{aligned}
 g_{\mathscr S}&=m\,\mathrm{d}U^2+\rho(v)\,\mathrm{d}v^2,\\
 \rho(v)&=\frac{Q''/Q-2(Q'/Q)^2}{n}
       =\frac2n\operatorname{Var}_{\nu_v}(\ell_v)>0.
 \end{aligned}
\end{equation}
For any $v_*\in(-1/M,\infty)$, the coordinates
\[
 Y=\sqrt m\,U,\qquad
 R(v)=\int_{v_*}^v\sqrt{\rho(t)}\,\mathrm{d}t
\]
define a global isometry onto $\mathbb R\times I$, where $I$ is an open interval with the Euclidean metric. The intrinsic distance is the Euclidean distance in $(Y,R)$.
\end{proposition}
\begin{proof}
Consider an affine variation $\delta=(\delta b,\delta c,\delta e)$ of the coefficients, write $F=\delta L/L$, and note that $\sigma/(HL)=\mathrm{d}\nu_v$ because $HL=Q(v)(1+va)$. Differentiating under the integral sign gives
\[
 D^2\log H(\delta,\delta)
   =2\mathbb E_{\nu_v} F^2-(\mathbb E_{\nu_v} F)^2,
\]
where $\mathbb E_{\nu_v}$ denotes expectation with respect to $\nu_v$. This is positive unless $\delta L=0$ everywhere, which forces $\delta c=\delta e=0$. Since $\log\lambda_\mu=(\log n+m\log b-\log H)/n$,
\[
 g=-D^2\log\lambda_\mu
   =\frac mn\left(\frac{\mathrm{d}b}{b}\right)^2
                  +\frac1nD^2\log H
\]
is positive definite. On the unit-volume hypersurface, the relation $H(c,e)=Q(v)/e=nb^m$ gives
\[
 e=\frac{Q(v)}{n\exp(mU)},\qquad c=ve.
\]
To pull back the Hessian, write $t=\mathrm{d}e/e$ and $z=Q'/Q$. Along an affine path of coefficients one has $v''=-2tv'$, so differentiating $\log H=-\log e+\log Q(v)$ twice gives
\[
 D^2\log H=t^2-2zt\,\mathrm{d}v
                   +(Q''/Q-z^2)\,\mathrm{d}v^2.
\]
On $\mathscr S$ we have $t=z\,\mathrm{d}v-m\,\mathrm{d}U$; substituting this and using $n=m+1$ gives \eqref{eq:warped-flat-unit-metric}. Since $Q'/Q=-\mathbb E_{\nu_v}\ell_v$ and $Q''/Q=2\mathbb E_{\nu_v}\ell_v^2$, the variance formula follows, and nonconstancy of $a$ makes this variance positive. In the coordinates $(Y,R)$ the metric is Euclidean on the convex set $\mathbb R\times I$, so straight segments prove the distance assertion.
\end{proof}

\subsubsection{Lengths of the two ends}

The interval $I$ has two ends, corresponding to $v\to\infty$ and $v\to-1/M$. Whether they lie at finite distance depends on how $a$ attains its extreme values. Assume that each of the extreme sets $\{a=0\}$ and $\{a=M\}$ is a finite disjoint union of smooth compact real submanifolds $Z_i$ of positive codimension, and that near each component, respectively,
\[
 a\asymp\operatorname{dist}(Z_i)^{k_i},\qquad
 M-a\asymp\operatorname{dist}(Z_i)^{k_i},\qquad k_i>0,
\]
where $\asymp$ means comparability up to positive constant factors and $\operatorname{dist}(Z_i)$ is the distance to $Z_i$ for a fixed Riemannian metric on $B$. The function $a$, respectively $M-a$, is then bounded below by a positive constant outside small disjoint tubes around these components. Define, with minima over the components $Z_i$ of the respective extreme set,
\[
 r_0=\min_{\{a=0\}}
       \frac{\operatorname{codim}_{\mathbb R}Z_i}{k_i},\qquad
 r_M=\min_{\{a=M\}}
       \frac{\operatorname{codim}_{\mathbb R}Z_i}{k_i}.
\]

\Needspace{15\baselineskip}
\begin{lemma}\label{lem:warped-endpoint-length}\index{Gauduchon metric!trace variance}
The end $v\to\infty$ of $I$ has finite length exactly when $r_0>1$, and the end $v\to-1/M$ exactly when $r_M>1$. More precisely, at the end $v\to\infty$ put $f=a$, $\epsilon=1/v$ and $r=r_0$; at the end $v\to-1/M$ put $f=M-a$, $\epsilon=-1/v-M$ and $r=r_M$. Let
\[
 H_f=\int_B\frac{\sigma}{f+\epsilon},\quad
 J_f=\int_B\frac{\sigma}{(f+\epsilon)^2},\quad
 K_f=\int_B\frac{\sigma}{(f+\epsilon)^3}.
\]
Then $\epsilon\downarrow0$ at the end, and
\begin{equation}\label{eq:warped-endpoint-variance}
 \rho(v)\,\mathrm{d}v^2=G_f(\epsilon)\,\mathrm{d}\epsilon^2,
 \quad
 G_f=\frac2n\left[\frac{K_f}{H_f}
                    -\left(\frac{J_f}{H_f}\right)^2\right].
\end{equation}
Moreover, $G_f\asymp\epsilon^{-2}$ for $r<1$ and $G_f\asymp\epsilon^{-2}/\log(1/\epsilon)$ for $r=1$, while for $r>1$ the function $\sqrt{G_f}$ is integrable near $\epsilon=0$.
\end{lemma}
\begin{proof}
At the end $v\to\infty$, where $\epsilon=1/v$, we have
\[
 \ell_v=\epsilon-\frac{\epsilon^2}{a+\epsilon},\qquad
 \mathrm{d}v=-\epsilon^{-2}\mathrm{d}\epsilon.
\]
At the end $v\to-1/M$, where $\epsilon=-1/v-M$, we obtain
\[
 \ell_v=\frac{(M+\epsilon)^2}{M-a+\epsilon}-(M+\epsilon),
 \qquad \mathrm{d}v=(M+\epsilon)^{-2}\mathrm{d}\epsilon.
\]
In both cases $\ell_v$ is an affine function of $1/(f+\epsilon)$ with coefficients depending only on $\epsilon$, and \eqref{eq:warped-endpoint-variance} expresses $\rho(v)\,\mathrm{d}v^2$ through the variance of $1/(f+\epsilon)$ under the probability $\sigma/[H_f(f+\epsilon)]$. This probability is $\nu_v$, the marginal in \eqref{eq:warped-minimizer-trace}, because $L$ is a positive multiple of $f+\epsilon$ and positive scaling of $L$ does not change that marginal.

Integration over tubes around the components $Z_i$, splitting the normal radius at $\epsilon^{1/k_i}$, gives
\[
 H_f\asymp
 \begin{cases}\epsilon^{r-1}&r<1,\\
 \log(1/\epsilon)&r=1,\\1&r>1,
 \end{cases}
 \qquad
 K_f\asymp
 \begin{cases}\epsilon^{r-3}&r<3,\\
 \log(1/\epsilon)&r=3,\\1&r>3.
 \end{cases}
\]
The analogous threshold for $J_f$ is two. For a finite sum of tube integrals, the smallest ratio governs.

If $r<1$, choose a component attaining $r$, and two normal annuli at fixed multiples of $\epsilon^{1/k_i}$, so far apart that the two-sided bounds separate their ranges of $f/\epsilon$. Each annulus has probability bounded below, since its weighted mass and $H_f$ both have order $\epsilon^{r-1}$, and the values of $1/(f+\epsilon)$ on the two annuli differ by at least $c\epsilon^{-1}$. The identity
\[
 \operatorname{Var}_\nu T
   =\frac12\iint(T(x)-T(y))^2\,\mathrm{d}\nu(x)\mathrm{d}\nu(y)
\]
applied to $T=1/(f+\epsilon)$ gives the lower bound, and $T\leq\epsilon^{-1}$ gives the upper bound.

If $r=1$, the near annulus has probability at least $c/\log(1/\epsilon)$. In the same tube, the normal radii between $\epsilon^{1/(2k_i)}$ and a fixed small radius carry probability bounded below, by the logarithmic integral. There $1/(f+\epsilon)\leq C\epsilon^{-1/2}$, while $1/(f+\epsilon)\geq c\epsilon^{-1}$ on the near annulus. The variance identity and the bound $G_f\leq2K_f/(nH_f)$ prove the claimed logarithmic estimate. In both cases $r<1$ and $r=1$, the integral of $\sqrt{G_f}$ diverges at $\epsilon=0$.

For $r>1$, $H_f$ tends to a finite positive integral and $G_f\leq C K_f$. For $1<r<3$, $\sqrt{G_f}$ is therefore bounded by $C\epsilon^{(r-3)/2}$, whose exponent exceeds $-1$. At $r=3$ the bound is $C\sqrt{\log(1/\epsilon)}$, and for $r>3$ it is a constant. All these bounds are integrable.
\end{proof}

\begin{theorem}\label{thm:warped-unit-volume-completion}\index{Metric completion!warped unit-volume hypersurface}
Under the hypotheses of \cref{lem:warped-endpoint-length} on the extreme sets, the metric completion of $\mathscr S$ is a closed Euclidean strip if $r_0,r_M>1$, a closed Euclidean half-plane if exactly one of them exceeds one, and the whole Euclidean plane otherwise. Each boundary line of the completion is parametrized by $b>0$, and its points are represented by the forms
\begin{align*}
 \theta_+(b)&=b\kappa+\frac{H_a}{nb^m}a w,
 &H_a&=\int_B\frac{\sigma}{a}\quad(r_0>1),\\
 \theta_-(b)&=b\kappa+\frac{H_-}{nb^m}(1-a/M)w,
 &H_-&=\int_B\frac{\sigma}{1-a/M}\quad(r_M>1).
\end{align*}
These forms are smooth, semipositive, nef, nonclosed and not big, and their variational volume is one. The completion has no other boundary points.
\end{theorem}
\begin{proof}
The completion of the convex product $\mathbb R\times I$ of \cref{prop:warped-flat-unit-metric} is its Euclidean closure. Along a Cauchy sequence, $Y$ converges, so $b$ has a finite positive limit, and $R$ converges to a point of $I$ or to a finite endpoint of $I$. \Cref{lem:warped-endpoint-length} decides which endpoints are finite and gives the three alternatives.

On $\mathscr S$ we have $e=Q(v)/(nb^m)$ and $c=ve$. At the end $v\to\infty$, $vQ(v)\to H_a$ by monotone convergence, which gives the coefficient limit $\theta_+(b)$. At the end $v\to-1/M$, write $1+va=1-a/M+(v+1/M)a$; dominated convergence gives $Q(v)\to H_-$, and hence the limit $\theta_-(b)$. These are the coefficient limits at all finite ends. The forms are semipositive, hence nef, and they are nonclosed. Their fiber coefficients $L_\pm$ vanish somewhere, so they are not big by \cref{prop:warped-current-cones} below. By the definitions of $H_a$ and $H_-$, $\int_B\sigma/L_\pm=nb^m$. For $\epsilon>0$, the form $\theta_\pm(b)+\epsilon(\kappa+w)$ lies in $\mathscr C$, and by \eqref{eq:warped-gauduchon-volume} and monotone convergence its volume
\[
 \frac{n(b+\epsilon)^m}{\int_B\sigma/(L_\pm+\epsilon)}
\]
tends to one as $\epsilon\downarrow0$. \Cref{thm:gaud-volume-nef-limit}, with the Hermitian form $\kappa+w$, shows that the variational volume of $\theta_\pm(b)$ is one.
\end{proof}

\subsubsection{Current cones and endpoint equations}
\label{sec:warped-current-endpoint-equations}

Retain $X=B\times C$, $\theta=b\kappa+(ca+e)w$ with real parameters $b,c,e$, and the probability $\mu=\sigma\wedge w$ of \cref{sec:warped-gauduchon-completion}. Localized closed tests determine its current cone as well as its smooth cone. At a finite unit-volume end, the reciprocal coefficient is integrable with an exponent greater than one, so the base equation still has a continuous solution. Write $h=\kappa+w$.

\Needspace{14\baselineskip}
\begin{proposition}\label{prop:warped-current-cones}\index{Cone duality!warped product currents}
In this real span, the pseudoeffective and nef cones both equal
\[
 \{b\geq0,\ e\geq0,\ e+cM\geq0\}.
\]
A background admits a strict positive current representative if and only if $b>0$, $e>0$ and $e+cM>0$. Thus its strict-current cone equals the smooth-positive cone $\mathscr C$ in \cref{thm:warped-gauduchon-volume}.
\end{proposition}
\begin{proof}
Put $L=ca+e$. Suppose an arbitrary quasi-plurisubharmonic $u$ satisfies $\theta+\ddc u\geq\delta h$, where $\delta\geq0$. The closed semipositive test $\kappa^{m-1}\wedge w$ and distributional Stokes give $b\geq\delta$. For every smooth nonnegative top-degree form $\eta$ on $B$, its pullback to $X$ is closed by dimension, and
\[
 \int_X(\theta+\ddc u-\delta h)\wedge\eta
   =\int_B(L-\delta)\eta\geq0.
\]
If $L<\delta$ at a point, a nonzero test supported in that open set contradicts this inequality. Hence $L\geq\delta$ everywhere. For $\delta=0$ these are precisely $b\geq0$, $e\geq0$, $e+cM\geq0$; conversely they make $\theta$ semipositive. Such forms are nef because $\theta+\varepsilon h>0$ for every $\varepsilon>0$. Conversely, applying the same tests to smooth positive representatives of $\theta+\varepsilon h$ and sending $\varepsilon$ to zero gives those inequalities. For $\delta>0$ they force the three strict inequalities, and the converse follows because $\theta$ itself is then strictly positive. No restriction of $u$ to an extreme fiber is needed.
\end{proof}

\Needspace{14\baselineskip}
\begin{proposition}\label{prop:warped-continuous-endpoint-equation}\index{Monge--Ampere equation@Monge--Amp\`ere equation!finite warped boundary}\index{Regularity!continuous solution without C11@continuous solution without $C^{1,1}$}
Under the extreme-set hypotheses of \cref{lem:warped-endpoint-length}, let $\theta=b\kappa+L_0w$ be either finite face $\theta_+(b)$ or $\theta_-(b)$ of \cref{thm:warped-unit-volume-completion}, with $b>0$. There is a continuous $\theta$-plurisubharmonic potential $u$, pulled back from $B$, such that
\begin{equation}\label{eq:warped-continuous-endpoint-equation}
 (\theta+\ddc u)^n=\mu.
\end{equation}
The product is the ordinary bounded-potential product. No solution of this equation is globally $C^{1,1}$.
\end{proposition}
\begin{proof}
Let $f=a$ at the first face and $f=1-a/M$ at the second, and let $r>1$ be the corresponding codimension-to-order ratio. The coefficient $L_0$ is a positive constant multiple of $f$. In each extreme tube the integral of $f^{-p}$ has the same convergence condition as
\[
 \int_0^{s_0}s^{\operatorname{codim}_{\mathbb R}Z_i-1-k_ip}\,\mathrm{d}s,
\]
where $s$ is normal radius and $s_0>0$ is small. Thus $1/L_0\in L^p(B,\sigma)$ for any fixed $1<p<r$. The zero set has smooth volume zero, and the density can be assigned any finite value there. The face formulas give
\[
 H_0:=\int_B\frac{\sigma}{L_0}=nb^m.
\]
Apply \cref{thm:kn-continuous} on the positive closed background $b\kappa$ to obtain a continuous $\phi$ with $A=b\kappa+\ddc\phi\geq0$ and $A^m=\lambda\sigma/L_0$ for a scalar $\lambda>0$. The bounded closed mass identity in \cref{thm:mass-invariance-characterization} gives $\int_BA^m=b^m$, hence $\lambda=b^m/H_0$. Set $u=\phi\circ p_B$, where $p_B:X\to B$ is projection. Its solution current is $A+L_0w\geq0$.

The pullback identity of \cref{lem:bounded-products-submersion} gives $A^{m+1}=0$ on $X$, and $(L_0w)^2=0$. The multilinear expansion in \cref{prop:bounded-products} therefore gives
\[
 (A+L_0w)^n=nL_0 A^m\wedge w
       =\frac{nb^m}{H_0}\sigma\wedge w=\mu.
\]
The cancellation $L_0(\sigma/L_0)=\sigma$ holds as measures, since the zero set has smooth volume zero and $A^m$ has the displayed absolutely continuous density.

Suppose a solution $q$ were globally $C^{1,1}$. Since $\mu/h^n$ is smooth and strictly positive on compact $X$, \cref{prop:c11-ellipticity} would give $\theta+\ddc q\geq\delta h$ for some $\delta>0$. This contradicts \cref{prop:warped-current-cones}, because $L_0$ vanishes. The argument applies to every solution, not only to the pulled-back one.
\end{proof}

\input{author/section_warped_endpoint_potential_limits}


\subsubsection{Traces of minimizing metrics}
\label{sec:warped-minimizer-trace-limits}

Uniform convergence of the base potentials does not control the minimizing balanced metrics themselves: by the cofactor formula \eqref{eq:warped-balanced-cofactor}, their $(n-1)$st powers contain the square of the reciprocal fiber coefficient. Depending on the ratio at the end, the traces of these powers stay bounded or blow up, and in the second case the trace-normalized powers concentrate on fibers over the extreme set.

Keep the sequence $\theta_j\to\theta_0$ of \cref{prop:warped-endpoint-potential-limits}, and let $\phi_0$ be the uniform limit of its base potentials $\phi_j$. Put $A_j=b_j\kappa+\ddc\phi_j$, and write $G_j=\gamma_j^m$, where $\gamma_j$ is the balanced minimizing metric of \cref{thm:warped-gauduchon-volume} for $\theta_j$. Traces are taken with respect to $h=\kappa+w$.

\Needspace{14\baselineskip}
\begin{proposition}\label{prop:warped-minimizer-trace}\index{Gauduchon metric!finite trace boundary}
Let $r>1$ be the ratio $r_0$ or $r_M$ of the end. After discarding finitely many terms, write $L_j=s_j(f+\epsilon_j)$ with $s_j>0$, $s_j\to s_0>0$ and $\epsilon_j\to0^+$. Here $f=a$, $s_j=c_j$ at the first face, and $f=M-a$, $s_j=-c_j$ at the second. Set $J_j=\int_B\sigma/L_j^2$. Then
\begin{equation}\label{eq:warped-minimizer-trace-mass}
 G_j=\frac{\sigma}{nL_j^2}+mA_j^{m-1}\wedge w,
 \qquad t_j:=\int_Xh\wedge G_j=\frac{J_j}{n}+mb_j^{m-1}.
\end{equation}
For $1<r<2$, $t_j\asymp\epsilon_j^{r-2}$; for $r=2$, $t_j\asymp\log(1/\epsilon_j)$. For $r>2$, the entire sequence converges weakly to the finite positive closed current
\[
 G_0=\frac{\sigma}{nL_0^2}+mA_0^{m-1}\wedge w,
 \qquad A_0=b\kappa+\ddc\phi_0,
\]
and $t_j\to t_0=\int_Xh\wedge G_0>0$.
\end{proposition}
\begin{proof}
On $\mathscr S$, $H(c_j,e_j)=nb_j^m$, so $A_j^m=\sigma/(nL_j)$. Substituting this in \eqref{eq:warped-balanced-cofactor} gives the formula for $G_j$. Stokes' theorem on $B$, for the closed form $\kappa$, gives $\int_B\kappa\wedge A_j^{m-1}=b_j^{m-1}$, which proves the trace identity. The factorization of $L_j$ uses $\epsilon_j=e_j/c_j$ at the first end and $\epsilon_j=e_j/(-c_j)-M$ at the second, and the coefficient limits give $s_j\to s_0>0$.

On each tube around an extreme component, $f\asymp\operatorname{dist}(Z_i)^{k_i}$, so the integral of $f^{-2}$ is comparable to
\[
 \int_0^{\varrho_0}\frac{\varrho^{\operatorname{codim}_{\mathbb R}Z_i-1}}{(\varrho^{k_i}+\epsilon_j)^2}\,\mathrm{d}\varrho,
\]
where $\varrho$ is normal radius, $\varrho_0>0$ is small and $r_i=\operatorname{codim}_{\mathbb R}Z_i/k_i$. Splitting the integral at $\varrho=\epsilon_j^{1/k_i}$ gives order $\epsilon_j^{r_i-2}$ for $r_i<2$, order $\log(1/\epsilon_j)$ for $r_i=2$, and a finite limit for $r_i>2$. The finite sum over components is governed by $r=\min_i r_i$. Since $b_j$ and $s_j$ have positive finite limits, these estimates prove the three alternatives for $t_j$.

If $r>2$, then $f^{-2}$ is integrable and dominates $L_j^{-2}$ up to a fixed constant factor, so $\sigma/L_j^2\to\sigma/L_0^2$ in total variation. Since $w\wedge w=0$, we have $A_j^{m-1}\wedge w=(\theta_j+\ddc u_j)^{m-1}\wedge w$, which converges weakly by \cref{prop:warped-endpoint-potential-limits}. This proves $G_j\rightharpoonup G_0$, and $G_0$ is positive and closed as a weak limit of positive closed forms. Pairing with $h$ gives the limit of the traces, which is positive because of the term $\sigma/(nL_0^2)$.
\end{proof}

\Needspace{15\baselineskip}
\begin{proposition}\label{prop:warped-normalized-test-limits}\index{Gauduchon metric!concentrating trace}\index{Gauduchon current!warped fiber limit}
Suppose $1<r\leq2$, and put $T_j=G_j/t_j$. Every weak cluster limit of $(T_j)$ has the form $p_B^*\nu$, where $\nu$ is a probability measure on $B$ supported on the extreme components with ratio $r$. Such a limit is positive and closed, has trace one relative to $h$, and satisfies $\theta_0\wedge p_B^*\nu=0$. If those components reduce to a single point $p$, the whole sequence converges to the current of integration $[\{p\}\times C]$.

The trace-normalized smooth metrics $\eta_j=t_j^{-1/m}\gamma_j$ satisfy
\[
 \int_X\theta_j\wedge\eta_j^m=D_\mu(\eta_j)
                  =\frac{nb_j^m}{t_j}\longrightarrow0.
\]
Both quantities have order $\epsilon_j^{2-r}$ when $1<r<2$, and order $1/\log(1/\epsilon_j)$ when $r=2$. Their quotient remains one.
\end{proposition}
\begin{proof}
Let $\nu_j=\sigma/(J_jL_j^2)$. These are probability measures on $B$, and
\[
 T_j=\frac{J_j}{nt_j}p_B^*\nu_j
                   +\frac m{t_j}A_j^{m-1}\wedge w.
\]
The first coefficient tends to one. The second current is positive with trace $mb_j^{m-1}/t_j\to0$, hence converges weakly to zero. Compactness of probability measures on $B$ and weak continuity of pullback under the submersion $p_B$ therefore show that every cluster limit is $p_B^*\nu$ for a cluster limit $\nu$ of $(\nu_j)$. Away from small tubes around the extreme set, the density $1/L_j^2$ is uniformly bounded, while $J_j\to\infty$; that region receives no mass in the limit. The tube integrals for components with ratio greater than $r$ are $o(J_j)$ by the preceding estimates, so only components with ratio $r$ keep mass.

For a limit $T=p_B^*\nu$, $\kappa\wedge T=0$ for degree reasons on $B$, while $L_0w\wedge T=0$ because $L_0$ vanishes on the support of $\nu$. Its trace is one because $\int_Cw=1$. A top-degree measure on $B$ is closed, so its pullback is closed as well. If the support is the single point $p$, every cluster limit of $(\nu_j)$ is $\delta_p$; compactness then gives convergence of the whole sequence, and $p_B^*\delta_p$ is the current of integration over the fiber $\{p\}\times C$.

By Stokes' theorem and \eqref{eq:warped-minimizer-trace-mass}, $\int_X\theta_j\wedge G_j=\int_X(A_j+L_jw)\wedge G_j=n\int_BA_j^m=nb_j^m$. Since $\gamma_j$ minimizes and $\lambda_\mu(\theta_j)=1$ on $\mathscr S$, also $D_\mu(\gamma_j)=nb_j^m$. Both are multiplied by $t_j^{-1}$ under the trace normalization, and the trace estimates give their limits and rates.
\end{proof}

In the projective strip of \cref{ex:warped-projective-strip}, both ratios are two and each extreme set is a single point. At each of the two ends, the normalized minimizing tests therefore converge to the current of integration over the corresponding fiber $\{p\}\times\mathbb P^1$, and the numerator and the denominator of the quotient both decay like $1/\log(1/\epsilon_j)$. This is convergence of currents; it does not give a pointwise limit of the metrics.


\subsubsection{Projective and torus examples}

\begin{example}\label{ex:warped-projective-strip}
On $B=\mathbb P^2$, take normalized Fubini--Study $\kappa$, $\sigma=\kappa^2$, $C=\mathbb P^1$, and
\[
 a([Z])=\frac{|Z_1|^2/2+|Z_2|^2}
                  {|Z_0|^2+|Z_1|^2+|Z_2|^2}.
\]
Then $M=1$, $r_0=r_M=2$, and the completion is a closed strip with finite faces
\[
 b\kappa+\frac{4\log2}{3b^2}a w,\qquad
 b\kappa+\frac{4\log2}{3b^2}(1-a)w.
\]
For $v\ne0$,
\[
 Q(v)=\frac4{v^2}\bigl[(v+1)\log(1+v)
                      -(v+2)\log(1+v/2)\bigr],
 \qquad Q(0)=1.
\]
The width of the strip is the finite positive integral
\[
 \int_{-1}^{\infty}
       \sqrt{\frac{Q''/Q-2(Q'/Q)^2}{3}}\,\mathrm{d}v.
\]
\end{example}
\begin{proof}
The extreme sets are the points $[1:0:0]$ and $[0:0:1]$. In affine coordinates, $a$ and $1-a$ vanish there to exactly second order, so both real codimensions are four and both orders are two. To compute the law of $a$ under $\sigma$, polar integration in the chart $Z_0=1$ gives $2\,\mathrm{d}s_1\mathrm{d}s_2/(1+s_1+s_2)^3$, where $s_j=|z_j|^2$. The transformation $u_j=s_j/(1+s_1+s_2)$ gives the uniform density two on the triangle $u_1,u_2\geq0$, $u_1+u_2\leq1$, and $a=u_1/2+u_2$. For $0\leq s\leq1/2$, the triangle $u_1/2+u_2\leq s$ has area $s^2$, so the law of $a$ has density $4s$ on $[0,1/2]$. With $u_0=1-u_1-u_2$, exchanging $u_0$ and $u_2$ preserves the uniform measure and replaces $a$ by $1-a$; hence the law has density $4(1-s)$ on $[1/2,1]$. Integrating against these two densities gives $H_a=H_-=4\log2$ and the formula for $Q$, and the integral definition of $Q$ shows that $Q(0)=1$ is its continuous extension. \Cref{thm:warped-unit-volume-completion} gives the boundary lines and the finite width, with Euclidean coordinates $Y=\sqrt2\log b$ and $R(v)$.
\end{proof}

\begin{example}\label{ex:warped-torus-plane}
Let $E=\mathbb C/(\mathbb Z+i\mathbb Z)$, $B=E^m$, and write $z_j=u_j+i\tau_j$. Put
\[
 \eta_j=\frac i2\mathrm{d}z_j\wedge\mathrm{d}\bar z_j,
 \qquad s=(m!)^{-1/m},\qquad
 \kappa=s\sum_j\eta_j,\qquad\sigma=\kappa^m.
\]
For $a=\sin^2(\pi u_1)$, the unit-volume hypersurface is itself a complete Euclidean plane, with Euclidean coordinates
\[
 Y=\sqrt m\log b,\qquad
 R=\frac{\log(1+v)}{2\sqrt n},\qquad v>-1.
\]
In this example the minimizing balanced metrics of \cref{thm:warped-gauduchon-volume} can be written down by integrating a periodic function twice.
\end{example}
\begin{proof}
The forms descend to $B$, and $\sigma=\eta_1\wedge\cdots\wedge\eta_m$ is a probability form. Both extreme sets are real hypersurfaces along which $a$, respectively $1-a$, vanishes to second order, so $r_0=r_M=1/2$. Splitting the period and substituting $t=\tan(\pi u)$ gives
\begin{align*}
 Q(v)&=\int_0^1\frac{\mathrm{d}u}{1+v\sin^2(\pi u)}\\
     &=\frac2\pi\int_0^\infty\frac{\mathrm{d}t}{1+(1+v)t^2}
      =\frac1{\sqrt{1+v}}.
\end{align*}
Let $F_v(u)=[Q(v)(1+v\sin^2(\pi u))]^{-1}$, which has mean one over a period. Integrating the equation
\[
 \phi_v''=4\pi bs(F_v-1)
\]
twice, whose right side has mean zero, gives a smooth periodic potential $\phi_v$ of $u=u_1$; the first integration constant is chosen so that $\phi_v'$ has mean zero. In our normalization $\ddc\phi_v=\phi_v''\eta_1/(4\pi)$, hence
\[
 A=bs\left(F_v\eta_1+\sum_{j=2}^m\eta_j\right)>0,
 \qquad A^m=b^mF_v\sigma=\frac{b^m}{H}\frac{\sigma}{L}.
\]
Thus the minimizer of \cref{thm:warped-gauduchon-volume}, with its cofactor formula \eqref{eq:warped-balanced-cofactor}, is explicit in this example, without an appeal to \cref{thm:tw-smooth}. Differentiating $Q$ gives
\[
 \rho(v)=\frac1{4n(1+v)^2},\qquad
 g_{\mathscr S}=m\,\mathrm{d}U^2+
             \frac1{4n}\left(\frac{\mathrm{d}v}{1+v}\right)^2.
\]
As $U$ and $v$ range over $\mathbb R$ and $(-1,\infty)$, the coordinates $(Y,R)$ range over all of $\mathbb R^2$. Both boundary faces of the coefficient cone are therefore at infinite distance on the unit-volume hypersurface.
\end{proof}

These distances are those of the fixed three-parameter cone with the product probability $\mu=\sigma\wedge w$. The endpoint criterion of \cref{lem:warped-endpoint-length} uses the hypothesis that the extreme sets are disjoint unions of submanifolds along which $a$ vanishes to finite power order; intersecting extreme strata or vanishing of infinite order would require further estimates. The variances are computed from the minimizing metrics, but we do not assert that these metrics have a pointwise limit at a finite end.



\Needspace{13\baselineskip}
\section{A compact dual problem over Gauduchon currents}
\label{sec:weak-gauduchon-duality}

We now fix a Hermitian background and vary the density. For a smooth positive density, \cref{thm:gaud-volume-dual-formula} provides a minimizing Gauduchon metric; for a density that is only continuous or integrable, there is no smooth equation to produce one. The extremal problem becomes compact when positive $\ddc$-closed $(n-1,n-1)$-currents are allowed as tests, provided the denominator $D_\mu$ is extended to currents. For this we take the determinant of the matrix density of a current relative to its trace measure. The integral of its $n$th root is an infimum of continuous linear functionals, hence upper semicontinuous under weak convergence, and its maximum over normalized currents is attained. The smooth equation identifies this maximum with the supremum over smooth tests.

From here until \cref{sec:gaud-volume-nef-limit}, we assume $n\geq2$, and we fix a Hermitian form $\omega$ and a smooth positive probability volume form $\mathrm{d}V$ on $X$; integrability and convergence in measure refer to $\mathrm{d}V$. As in \cref{chap:products}, $\mathrm{MA}_\omega(u)=(\omega+\ddc u)^n$ for bounded $u\in\PSH(X,\omega)$, and $\mathrm{MA}^{\mathrm{ext}}_\omega(u)$ is the extended non-pluripolar Monge--Amp\`ere measure of \cref{thm:arbitrary-extended-ma}. Let $\mathscr G_\omega$ be the compact convex set of normalized Gauduchon currents of \cref{sec:gauduchon-linear-duality}, that is, of positive $(n-1,n-1)$-currents $T$ with $\ddc T=0$ and $\int_X\omega\wedge T=1$. Finally, let $F\in C(X,\mathbb R)$.

For $T\in\mathscr G_\omega$, the trace measure $\sigma_T=\omega\wedge T$ is a probability measure, and all coefficient measures of $T$ are absolutely continuous with respect to it. In a unitary frame for $\omega$, let $B$ be the matrix of a smooth real $(1,1)$-form $\beta$ relative to $\omega$. The coefficient densities of $T$ form a measurable positive semidefinite Hermitian matrix $P_T$, determined $\sigma_T$-almost everywhere by
\[
 \beta\wedge T=\frac1n\operatorname{tr}(B P_T)\,\mathrm{d}\sigma_T,
 \qquad \operatorname{tr}P_T=n\quad\sigma_T\text{-almost everywhere}.
\]
The matrix $P_T$ transforms by unitary conjugation when the frame changes, so its determinant is well defined. Define
\begin{equation}\label{eq:weak-gaud-functional}
 f_F=\left(\frac{\mathrm{e}^F\mathrm{d}V}{\omega^n}\right)^{1/n},\qquad
 \Gamma_F(T)=f_F(\det P_T)^{1/n}\sigma_T,\qquad
 \mathfrak D_F(T)=\Gamma_F(T)(X).
\end{equation}
These quantities are finite, since the arithmetic--geometric mean inequality gives $(\det P_T)^{1/n}\leq1$; the measure $\Gamma_F(T)$ may vanish for some currents. The same definitions apply to every positive $(n-1,n-1)$-current $T$ of finite mass, with the finite measure $\sigma_T=\omega\wedge T$ and $\operatorname{tr}P_T=n$, and $\Gamma_F$ is positively homogeneous of degree one in $T$. If the coefficient measures of $T$ are written instead as a matrix density $P'$ relative to another positive measure $\sigma'$ with respect to which they are absolutely continuous, then $(\det P_T)^{1/n}\sigma_T=(\det P')^{1/n}\sigma'$, by homogeneity of the determinant.

For a smooth positive $T=\gamma^{n-1}$, we may therefore use the reference measure $\omega^n$. The coefficient matrix of $T$ relative to $\omega^n$ is the cofactor matrix of the matrix of $\gamma$, and its determinant is $(\gamma^n/\omega^n)^{n-1}$. The definition therefore gives
\[
 \mathfrak D_F(\gamma^{n-1})
 =\int_X\left(\frac{\mathrm{e}^F\mathrm{d}V}{\gamma^n}\right)^{1/n}\gamma^n.
\]
Thus $\mathfrak D_F$ extends the denominator $D_\mu$ of \cref{def:gaud-variational-volume} to currents, with $\mu$ replaced by the possibly unnormalized density $\mathrm{e}^F\mathrm{d}V$.

\begin{lemma}\label{lem:weak-gaud-determinant-dual}
For every positive $(n-1,n-1)$-current $T$ of finite mass on $X$,
\[
 \mathfrak D_F(T)
   =\inf_{\substack{\beta>0\ \mathrm{continuous}\\\beta^n=\mathrm{e}^F\mathrm{d}V}}
                         \int_X\beta\wedge T.
\]
In particular, $\mathfrak D_F$ is concave, positively homogeneous and weakly upper semicontinuous on $\mathscr G_\omega$, and its value does not depend on the auxiliary frame or on the metric used to express the matrix density.
\end{lemma}
\begin{proof}
Let $I$ be the identity matrix of size $n$. For any positive semidefinite Hermitian matrix $P$ of that size,
\[
 (\det P)^{1/n}
      =\inf_{Q>0,\ \det Q=1}\frac1n\operatorname{tr}(QP).
\]
The inequality $\leq$ is the arithmetic--geometric mean inequality. For the reverse inequality, take, for real $\varepsilon>0$,
\[
 Q_\varepsilon=
       \det(P+\varepsilon I)^{1/n}(P+\varepsilon I)^{-1}.
\]
Its determinant is one. If $P$ is positive definite, the trace expression $\frac1n\operatorname{tr}(Q_\varepsilon P)$ converges to $(\det P)^{1/n}$ as $\varepsilon\downarrow0$. If $P$ is singular, it converges to zero: each ratio of an eigenvalue of $P$ to that eigenvalue plus $\varepsilon$ is bounded, while the prefactor tends to zero. When $\operatorname{tr}P=n$, the trace expression is at most $1+\varepsilon$.

We apply this formula to $P_T$ and multiply by $f_F$. Every continuous positive $\beta$ in the asserted infimum has matrix $f_FQ$ with $\det Q=1$, so it gives an upper bound for $\mathfrak D_F(T)$. Conversely, for fixed $\varepsilon>0$ the measurable field $Q_\varepsilon$ formed with $P=P_T$ is bounded above and below by positive constants. It can be approximated in $L^1(\sigma_T)$ by continuous positive definite Hermitian fields with bounds of the same type. One way to construct them is to apply Lusin's theorem in finitely many bundle charts, extend the continuous fields from the resulting compact sets, and glue with a partition of unity; spectral truncation to a fixed positive interval preserves the bounds. Finally we divide each approximating field by the $n$th root of its determinant. This normalization is uniformly continuous on fields with spectrum in a fixed compact positive interval, so the normalized fields still converge in $L^1(\sigma_T)$ and have determinant one. Their trace integrals converge to that of $Q_\varepsilon$.

Now let $\varepsilon\downarrow0$. The bound $1+\varepsilon$ permits dominated convergence, and we obtain equality of the infimum and the integral. For a current of arbitrary positive finite trace mass we use homogeneity, and the zero current is immediate. The infimum is taken over functionals $T\mapsto\int_X\beta\wedge T$ that are linear and weakly continuous, which proves concavity and upper semicontinuity. It also expresses $\mathfrak D_F(T)$ without a matrix frame or an auxiliary metric, which proves the final assertion.
\end{proof}

Let $B_0(F)$ be the normalization scalar of \cref{sec:normalization-scalar} for the background $\omega$ and the reference potential zero. Thus the equation
$\mathrm{MA}_\omega(u)=\mathrm{e}^{B_0(F)+F}\mathrm{d}V$
has a bounded solution $u\in\PSH(X,\omega)$, and $B_0(F)$ is the only constant with this property.

\begin{theorem}\label{thm:weak-gaud-compact-duality}\index{Gauduchon duality!weak compact formulation}
For every $F\in C(X,\mathbb R)$,
\begin{equation}\label{eq:weak-gaud-max-formula}
 M(F):=\max_{T\in\mathscr G_\omega}\mathfrak D_F(T)
                      =\mathrm{e}^{-B_0(F)/n}>0.
\end{equation}
The maximum equals the supremum over the smooth strictly positive members of $\mathscr G_\omega$. The set $\mathscr M(F)$ of maximizing currents is nonempty, compact and convex.

If $F$ is smooth, there is exactly one maximizing current, namely the normalized $(n-1)$st power of the Gauduchon metric conformal to the smooth solution form $\omega+\ddc u$ of $\mathrm{MA}_\omega(u)=\mathrm{e}^{B_0(F)+F}\mathrm{d}V$.
\end{theorem}
\begin{proof}
Upper semicontinuity of $\mathfrak D_F$ on the compact set $\mathscr G_\omega$ gives attainment, and testing a smooth positive member gives a positive value. Concavity makes the set of maximizers convex, and upper semicontinuity makes it closed.

Suppose first that $F$ is smooth, and let $\Omega=\omega+\ddc u>0$ be the smooth solution form with $\Omega^n=\mathrm{e}^{B_0(F)+F}\mathrm{d}V$. The form $\beta=\mathrm{e}^{-B_0(F)/n}\Omega$ is an admissible test in the preceding lemma, so Stokes' theorem for currents gives, for every $T\in\mathscr G_\omega$,
\[
 \mathfrak D_F(T)\leq\int_X\beta\wedge T
  =\mathrm{e}^{-B_0(F)/n}\int_X\omega\wedge T=\mathrm{e}^{-B_0(F)/n}.
\]
If $\gamma$ is the Gauduchon metric conformal to $\Omega$, equality holds in the pointwise matrix inequality for $T$ proportional to $\gamma^{n-1}$; normalizing $T$ so that $\int_X\omega\wedge T=1$, we obtain the formula and its attainment by a smooth test.

For uniqueness, let $T$ be a maximizer. Equality in the integrated inequalities forces equality in the matrix arithmetic--geometric mean inequality $\sigma_T$-almost everywhere. Since the matrix of $\beta$ is positive definite and $\operatorname{tr}P_T=n$, this means that the coefficient measures of $T$ are a nonnegative scalar measure times the smooth coefficients of $\gamma^{n-1}$. Thus $T=a\gamma^{n-1}$, where $a$ is a positive measure, regarded as a distribution. The equation $\ddc(a\gamma^{n-1})=0$ is the adjoint of a scalar elliptic equation, so elliptic regularity makes $a$ smooth. Since $\gamma$ is Gauduchon, \cref{prop:adjoint-gauduchon} identifies the kernel of this adjoint operator with the constants: the trace operator $v\mapsto n\ddc v\wedge\gamma^{n-1}/\gamma^n$ has the constants as kernel and index zero, and integration against $\gamma^n$ annihilates its range. Thus $a$ is constant, and the normalization fixes the constant. This proves uniqueness.

For continuous $F,G$, the pointwise definition gives
\[
 \mathrm{e}^{-\|F-G\|_\infty/n}\mathfrak D_G(T)
 \leq\mathfrak D_F(T)
 \leq \mathrm{e}^{\|F-G\|_\infty/n}\mathfrak D_G(T).
\]
The same bounds hold for the maxima, and for the suprema over smooth positive members. We approximate $F$ uniformly by smooth functions. Since $B_0$ is $1$-Lipschitz for the uniform norm by \cref{prop:relative-normalization-order}, the formula just proved for smooth data gives \eqref{eq:weak-gaud-max-formula} in the limit, including the equality with the smooth supremum. Note that only $F$ is approximated here: some normalized Gauduchon currents are not limits of smooth positive ones, as \cref{ex:gauduchon-current-smoothing-obstruction} shows.
\end{proof}

\section{Supporting probabilities represented by maximizing currents}
\label{sec:gaud-current-adjoints}

By \cref{thm:weak-gaud-compact-duality}, $B_0=-n\log M$, and $M(F)$ is the maximum of the functionals $T\mapsto\mathfrak D_F(T)=\int_X\mathrm{e}^{F/n}\,\mathrm{d}\Gamma_0(T)$, each of which is linear in $\mathrm{e}^{F/n}$. The supporting probabilities of $B_0$ should therefore be the normalized determinant measures of maximizing currents, and we show that they are exactly these. The point that needs proof is the continuity of the determinant measures: on the set of maximizers, upper semicontinuity of the total determinant integral becomes equality. Recall from \cref{sec:normalization-supports} that the set $\mathcal A_0(F)$ of \emph{supporting probabilities} of $B_0$ at $F\in C(X)$ consists of the Borel probability measures $\nu$ with $B_0(G)\leq B_0(F)-\int_X(G-F)\,\mathrm{d}\nu$ for all $G\in C(X)$; by \cref{thm:weak-adjoint-probabilities} it is nonempty, convex and weakly compact.

For $T\in\mathscr M(F)$ put
\[
 \nu_{T,F}=\Gamma_F(T)/M(F).
\]
This is a probability measure.

\begin{lemma}\label{lem:gaud-maximizer-measure-continuity}
Suppose $T_j\to T$ weakly in $\mathscr G_\omega$ and $\mathfrak D_F(T_j)\to M(F)$. Then $T\in\mathscr M(F)$ and $\Gamma_F(T_j)\rightharpoonup\Gamma_F(T)$. On $\mathscr M(F)$ the map $T\mapsto\Gamma_F(T)$ is affine and continuous.
\end{lemma}
\begin{proof}
Upper semicontinuity gives $\mathfrak D_F(T)\geq M(F)$, hence equality. The measures $\Gamma_F(T_j)$ have bounded mass; let $\Gamma'$ be a weak cluster point, whose mass is $M(F)$. For a continuous real function $h$, choose $\varepsilon>0$ with $1\pm\varepsilon h>0$. Since $\Gamma_{F+n\log(1\pm\varepsilon h)}=(1\pm\varepsilon h)\Gamma_F$, upper semicontinuity of $\mathfrak D$ for these two continuous logarithmic densities gives
\[
 \int_X(1\pm\varepsilon h)\,\mathrm{d}\Gamma'
       \leq\int_X(1\pm\varepsilon h)\,\mathrm{d}\Gamma_F(T).
\]
Since both measures have mass $M(F)$, the two signs give equality of the integrals of $h$. Thus every cluster point equals $\Gamma_F(T)$, which proves the convergence.

For affineness, take $T_1,T_2\in\mathscr G_\omega$ and $t\in[0,1]$, and express the coefficient measures of both currents through densities with respect to the sum of their trace measures. Concavity and homogeneity of the $n$th root of the determinant, applied to these matrix densities, give
\[
 \Gamma_F(tT_1+(1-t)T_2)
    \geq t\Gamma_F(T_1)+(1-t)\Gamma_F(T_2),\qquad 0\leq t\leq1.
\]
If $T_1$ and $T_2$ maximize, so does the current on the left, by concavity of $\mathfrak D_F$. Both measures in the inequality then have mass $M(F)$, so equality holds. This proves affineness, and continuity is the special case of the first assertion in which every $T_j$ is a maximizer.
\end{proof}

\Needspace{9\baselineskip}
\begin{theorem}\label{thm:gaud-current-all-adjoints}
For every $F\in C(X,\mathbb R)$,
\[
 \mathcal A_0(F)=\{\nu_{T,F}:T\in\mathscr M(F)\}.
\]
Consequently the determinant measures $\Gamma_F(T)$ of all maximizing currents are absolutely continuous with respect to $\mathrm{d}V$. (This concerns the measures $\Gamma_F(T)$, not the currents $T$ themselves.)
\end{theorem}
\begin{proof}
For a maximizing $T$ and any $G\in C(X)$,
\[
 M(G)\geq\mathfrak D_G(T)
       =M(F)\int_X\mathrm{e}^{(G-F)/n}\,\mathrm{d}\nu_{T,F}.
\]
Since $B_0=-n\log M$, taking $-n$ times the logarithm and applying Jensen's inequality gives
\[
 B_0(G)\leq B_0(F)-n\log\int_X\mathrm{e}^{(G-F)/n}\,\mathrm{d}\nu_{T,F}
          \leq B_0(F)-\int_X(G-F)\,\mathrm{d}\nu_{T,F}.
\]
Hence every probability $\nu_{T,F}$ with $T\in\mathscr M(F)$ belongs to $\mathcal A_0(F)$.

Denote the set of these probabilities by $\mathscr N(F)$; by the preceding lemma it is compact and convex. To prove the reverse inclusion we compute the right derivative $(B_0)'_+(F;h)=\lim_{t\downarrow0}(B_0(F+th)-B_0(F))/t$ for $h\in C(X)$. Since $\mathfrak D_{F+th}(T)=\int_X\mathrm{e}^{th/n}\,\mathrm{d}\Gamma_F(T)$ and the masses of the measures $\Gamma_F(T)$ are bounded on $\mathscr G_\omega$, a Taylor expansion of the exponential gives, as $t\downarrow0$,
\[
 \mathfrak D_{F+th}(T)
 =\mathfrak D_F(T)+\frac tn\int_Xh\,\mathrm{d}\Gamma_F(T)+O(t^2),
\]
with the error bounded uniformly in $T$. Testing maximizers at $F$ gives
\[
 \liminf_{t\downarrow0}\frac{M(F+th)-M(F)}t
    \geq\frac{M(F)}n\max_{\nu\in\mathscr N(F)}\int_Xh\,\mathrm{d}\nu.
\]
For the reverse inequality, choose a maximizer $T_t$ at $F+th$. The multiplicative bounds in the previous proof imply $\mathfrak D_F(T_t)\to M(F)$. Every sequence $t\downarrow0$ has a subsequence along which $T_t$ converges, and the preceding lemma identifies the limit of $\Gamma_F(T_t)$ with the determinant measure of a maximizer at $F$. Since $\mathfrak D_F(T_t)\leq M(F)$, the Taylor expansion gives the matching upper bound for the difference quotients of $M$. Since $B_0=-n\log M$, therefore
\[
 (B_0)'_+(F;h)
       =-\max_{\nu\in\mathscr N(F)}\int_Xh\,\mathrm{d}\nu.
\]
If $\nu\in\mathcal A_0(F)$, its supporting inequality in direction $th$ implies
\[
 \int_Xh\,\mathrm{d}\nu\leq\max_{\widetilde\nu\in\mathscr N(F)}
                                    \int_Xh\,\mathrm{d}\widetilde\nu
 \quad\text{for every }h\in C(X).
\]
If $\nu$ were not in the weakly compact convex set $\mathscr N(F)$, the Hahn--Banach separation theorem in the dual of $C(X)$ would give some $h\in C(X)$ violating this inequality. Hence $\nu\in\mathscr N(F)$, and the two sets are equal. The absolute continuity of their members was proved in \cref{thm:weak-adjoint-probabilities}.
\end{proof}

For continuous $F$, the supporting probability is therefore unique exactly when all maximizing currents have the same normalized determinant measure, which is weaker than uniqueness of the maximizing current. Whether a maximizing current can be written as a conformal multiple of a power $(\omega+\ddc u)^{n-1}$ of a weak solution form is a further question, discussed in \cref{sec:history-gaud-volume}; the construction above works directly with the coefficient measures of the currents.


\section{Integrability of Gauduchon determinant measures}
\label{sec:gaud-determinant-integrability}

The lower bound for equation normalizations gives a uniform $L^q$ estimate for the determinant measures of all normalized Gauduchon currents. Every exponent below $n/(n-1)$ is available. At the endpoint, the optimal constant is exactly $v_-(\omega)^{-1/n}$, linking integrability of the dual tests to the lower mass of the background.

Write $\Gamma_0(T)$ for the determinant measure in \eqref{eq:weak-gaud-functional} with $F=0$.

\begin{theorem}\label{thm:gaud-determinant-subcritical}
For every $1<q<n/(n-1)$ there is $C_q>0$ such that every $T\in\mathscr G_\omega$ satisfies
\[
 \Gamma_0(T)=k_T\,\mathrm{d}V,\qquad k_T\geq0,\qquad
 \|k_T\|_{L^q(\mathrm{d}V)}\leq C_q.
\]
No bounded mass or positive lower-volume hypothesis is required.
\end{theorem}
\begin{proof}
Fix $p>1$ and take any smooth strictly positive function $h$. The smooth equation gives $\Omega_h=\omega+\ddc u_h>0$ and $c_h>0$ such that
\[
 \Omega_h^n=c_hh^n\,\mathrm{d}V.
\]
By \cref{thm:lp-potential-estimate}, the density of this bounded solution satisfies
\[
 c_h\|h^n\|_{L^p(\mathrm{d}V)}\geq c_p>0.
\]
Pointwise matrix arithmetic--geometric mean, applied to $\Omega_h$ and the coefficient density of $T$, gives the measure inequality
\[
 c_h^{1/n}h\,\Gamma_0(T)\leq\Omega_h\wedge T.
\]
Integrating and using $\ddc T=0$ yields
\[
 \int_Xh\,\mathrm{d}\Gamma_0(T)
 \leq c_h^{-1/n}\int_X\Omega_h\wedge T
 =c_h^{-1/n}
 \leq c_p^{-1/n}\|h\|_{L^{np}(\mathrm{d}V)}.
\]
Approximation by smooth strictly positive functions extends this bound to every continuous nonnegative $h$. Applied to $|h|$, it bounds the positive measure functional on all continuous functions by the $L^{np}$ norm. Density of continuous functions in $L^{np}$ and the usual duality of these spaces show that the measure has a density in $L^{np/(np-1)}$, with norm at most $c_p^{-1/n}$. As $p$ ranges over $(1,+\infty)$, the exponent $np/(np-1)$ ranges over $(1,n/(n-1))$. This proves the result.
\end{proof}

In particular, the determinant of the matrix density vanishes on the singular part of the trace measure: if $\sigma_T=\sigma_{T,\mathrm{ac}}+\sigma_{T,\mathrm{s}}$ is its Lebesgue decomposition, then
\[
 \det P_T=0\quad\sigma_{T,\mathrm{s}}\text{-almost everywhere}.
\]
Indeed $f_0=(\mathrm{d}V/\omega^n)^{1/n}$ is everywhere positive and the determinant-weighted measure is absolutely continuous. This assertion concerns the determinant; it does not remove a possible singular part of the current itself.

\Needspace{12\baselineskip}
\begin{theorem}\label{thm:gaud-determinant-endpoint}\index{Gauduchon metric!determinant integrability}
With the extended value $+\infty$ allowed, one has
\[
 \sup_{T\in\mathscr G_\omega}
       \|k_T\|_{L^{n/(n-1)}(\mathrm{d}V)}=v_-(\omega)^{-1/n},
\]
where the right side is interpreted as $+\infty$ when $v_-(\omega)=0$. The same supremum is obtained by restricting to smooth strictly positive normalized Gauduchon forms. Thus a uniform estimate at the endpoint exponent is equivalent to $v_-(\omega)>0$.
\end{theorem}
\begin{proof}
Suppose first that $v_-(\omega)>0$. In the smooth equation of the preceding proof, the definition of $v_-$ gives
\[
 c_h\int_Xh^n\,\mathrm{d}V=\int_X\Omega_h^n\geq v_-(\omega).
\]
The same pointwise inequality and Stokes therefore give
\[
 \int_Xh\,\mathrm{d}\Gamma_0(T)
      \leq v_-(\omega)^{-1/n}\|h\|_{L^n(\mathrm{d}V)}.
\]
Duality proves the asserted upper bound for every current in $\mathscr G_\omega$.

\Needspace{8\baselineskip}
For the reverse bound, let $g$ be any smooth positive probability density relative to $\mathrm{d}V$, and solve
\[
 \Omega_g^n=c(g)g\,\mathrm{d}V.
\]
Let $T_g$ be the trace-normalized $(n-1)$st power of the Gauduchon metric conformal to $\Omega_g$. Equality in the pointwise matrix inequality gives
\[
 \int_Xg^{1/n}\,\mathrm{d}\Gamma_0(T_g)=c(g)^{-1/n}.
\]
Since $\|g^{1/n}\|_{L^n}=1$, its endpoint dual norm is at least $c(g)^{-1/n}$. By \cref{thm:gaud-volume-mass-extrema},
\[
 \inf_{g>0\ \mathrm{smooth},\ \int_Xg\,\mathrm{d}V=1}c(g)=v_-(\omega).
\]
Taking the supremum over these smooth $T_g$ proves the matching lower bound. If $v_-(\omega)=0$, the same argument makes this supremum infinite. This also proves that allowing nonsmooth currents does not change it.
\end{proof}

The endpoint norm is independent of the smooth probability measure used to express it, as the formula shows. Subcritical constants may depend on that choice. Neither assertion identifies $v_-$ with the upper mass $v_+$.

\begin{corollary}\label{cor:gaud-adjoint-lq}
Let $\mathcal A_0(F)$ be the supporting probabilities in \cref{thm:gaud-current-all-adjoints}. For every $R>0$ and $1<q<n/(n-1)$ there is $C_{R,q}$ with the following property. If $F\in C(X)$ and $\|F\|_\infty\leq R$, every $\nu\in\mathcal A_0(F)$ has a density $a$ satisfying
\[
 \nu=a\,\mathrm{d}V,\qquad\|a\|_{L^q(\mathrm{d}V)}\leq C_{R,q}.
\]
If $v_-(\omega)>0$, this holds also for $q=n/(n-1)$.
\end{corollary}
\begin{proof}
Write $M(F)=\exp(-B_0(F)/n)$ and represent $\nu$ by a maximizing current $T\in\mathscr M(F)$, using \cref{thm:gaud-current-all-adjoints}. Then
\[
 \nu=\frac{\mathrm{e}^{F/n}k_T}{M(F)}\,\mathrm{d}V,
 \qquad M(F)\geq \mathrm{e}^{-R/n}M(0)>0.
\]
The two preceding theorems give the bound, for example with constant
$\mathrm{e}^{2R/n}M(0)^{-1}C_q$ in the subcritical case. Use the endpoint constant when it is finite.
\end{proof}

The endpoint fails for the full normalized Gauduchon family on the threefold of \cref{thm:vanishing-lower-mass}, while every subcritical estimate above remains valid. For a fixed smooth right-hand side, \cref{thm:weak-gaud-compact-duality} gives a unique smooth maximizing current, whose determinant density is bounded. This does not give uniform higher moments along sequences approaching its value; \cref{prop:gaud-determinant-concentration} provides a counterexample.

\Needspace{13\baselineskip}
The endpoint can be available on a manifold with infinite upper mass. On the Hopf threefold of \cref{chap:hopf}, take its Gauduchon metric $\omega$, for which $v_-(\omega)=6$ and $v_+(\omega)=+\infty$, and set $\mathrm{d}V=\omega^3/6$. The normalized Gauduchon form $T_0=\omega^2/6$ has $\sigma_{T_0}=\mathrm{d}V$ and $P_{T_0}=I$. Hence $k_{T_0}=6^{-1/3}$, and
\[
 \sup_{T\in\mathscr G_\omega}\|k_T\|_{L^{3/2}(\mathrm{d}V)}=6^{-1/3}.
\]
This gives an attained endpoint determinant estimate alongside failure of bounded mass. The estimate controls these finite Gauduchon measures, not the total mass of arbitrary nonlinear products.


\Needspace{16\baselineskip}
\section{Integrable densities in the compact Gauduchon problem}
\label{sec:integrable-gauduchon-currents}

We keep the Hermitian background $\omega$ and replace the continuous logarithmic density by an integrable density. For an unbounded density, weak compactness of the test currents alone does not pass the determinant integral to the limit; uniform control of the weighted tails provides the missing upper semicontinuity.

We use the following notation from \cref{sec:individual-density-subsolutions,sec:integrable-gauduchon-tests}, with $\omega$ in place of $h$. For a nonzero nonnegative $f\in L^1(\mathrm{d}V)$, the subsolution threshold $b_\omega(f)\geq0$ is, by \cref{thm:individual-density-subsolution}, the largest $b\geq0$ for which some $u\in\PSH(X,\omega)$ satisfies $\mathrm{MA}^{\mathrm{ext}}_\omega(u)\geq bf\,\mathrm{d}V$. For $p\in L^n$, the root functional is $\rho_\omega(p)=b_\omega(|p|^n)^{-1/n}$ if $p\ne0$, and $\rho_\omega(0)=0$. The space $\mathcal D_\omega$ of \cref{prop:individual-density-banach} consists of the $p\in L^n$ with $\rho_\omega(p)<+\infty$, and $\mathcal D_\omega^0$ is the closure of the bounded functions in $\mathcal D_\omega$, characterized by vanishing tails in \cref{prop:root-uniform-tails}. For a Gauduchon metric $\gamma$, $L_{\omega,\gamma}(p)$ denotes the test quotient of \cref{sec:integrable-gauduchon-tests}, and \cref{thm:integrable-gauduchon-formula} gives $\rho_\omega(p)=\sup_\gamma L_{\omega,\gamma}(|p|)$. Finally, $k_T$ is the determinant density of $T\in\mathscr G_\omega$ from \cref{thm:gaud-determinant-subcritical}.

If $p\geq0$ is measurable, define
\[
 D_p(T)=\int_Xp k_T\,\mathrm{d}V\qquad(T\in\mathscr G_\omega),
\]
allowing the value $+\infty$. If $T=\gamma^{n-1}/\int_X\omega\wedge\gamma^{n-1}$ for a Gauduchon metric $\gamma$, then $D_p(T)=L_{\omega,\gamma}(p)$.

\Needspace{11\baselineskip}
\begin{theorem}\label{thm:integrable-current-attainment}\index{Gauduchon duality!current attainment}
For every nonnegative $p\in L^n$,
\begin{equation}\label{eq:integrable-current-supremum}
 \sup_{T\in\mathscr G_\omega}D_p(T)=\rho_\omega(p),
\end{equation}
with the value $+\infty$ allowed. If $p\in\mathcal D_\omega^0$, the functional $D_p$ is finite, concave and weakly upper semicontinuous on $\mathscr G_\omega$, its maximum is attained, and the densities $p k_T$, for $T\in\mathscr G_\omega$, are uniformly integrable.
\end{theorem}
\begin{proof}
For a smooth strictly positive $p$, we have $D_p=\mathfrak D_{n\log p}$ and $\mathrm{e}^{-B_0(n\log p)/n}=\rho_\omega(p)$, so \cref{thm:weak-gaud-compact-duality}, applied to $F=n\log p$, proves \eqref{eq:integrable-current-supremum}. By H\"older's inequality and \cref{thm:gaud-determinant-subcritical}, for every real $r>n$ there is a constant $C_r$ such that, for every $q\in L^r$,
\begin{equation}\label{eq:current-root-test-bound}
 \sup_{T\in\mathscr G_\omega}\int_X|q| k_T\,\mathrm{d}V
       \leq C_r\|q\|_{L^r}.
\end{equation}
Next, approximate a bounded nonnegative $p$ in $L^{2n}$ by smooth strictly positive functions. This estimate and \eqref{eq:root-subcritical-bound} give uniform convergence on $\mathscr G_\omega$ of the corresponding functionals $D$ and convergence of the root functionals, so the supremum formula holds for bounded $p$. For such $p$, the functional $D_p$ is also upper semicontinuous and concave, by \cref{lem:weak-gaud-determinant-dual} and uniform convergence. For general $p\in L^n$, we use the truncations $\min(p,K)$ at positive integer levels $K$. Monotone convergence and the definition of $b_\omega$ through truncations give
\[
 \sup_T D_p(T)=\sup_T\sup_K D_{\min(p,K)}(T)
       =\sup_K\rho_\omega(\min(p,K))=\rho_\omega(p).
\]
This step only interchanges two suprema and gives no attainment.

Now suppose $p\in\mathcal D_\omega^0$. The supremum formula applied to the tail gives
\[
 0\leq D_p(T)-D_{\min(p,K)}(T)
 \leq \rho_\omega(p\mathbf1_{\{p>K\}})\longrightarrow0
\]
uniformly in $T$, by \cref{prop:root-uniform-tails}. Hence $D_p$ is finite and is a uniform limit of upper semicontinuous concave functionals, and compactness of $\mathscr G_\omega$ gives a maximizer. Finally, fix $1<q<n/(n-1)$. By H\"older's inequality and \cref{thm:gaud-determinant-subcritical}, for every measurable set $E$,
\[
 \sup_T\int_Ep k_T\,\mathrm{d}V
 \leq\rho_\omega(p\mathbf1_{\{p>K\}})
       +K C_q\,\mathrm{d}V(E)^{1-1/q}.
\]
Choosing first $K$ large and then the volume of $E$ small proves uniform absolute continuity. Since the total masses are bounded by $\rho_\omega(p)$, the densities are uniformly integrable.
\end{proof}

For a nonzero density, a maximizing current gives a supporting probability for the threshold, as we show below. Whether the threshold is attained by an exact equation depends on the regularity of the maximizing current, which enters through the trace inequalities of \cref{subsec:threshold-trace-inequalities}.

\subsection{The tail condition and perturbations of the density}

Vanishing tails are equivalent to continuity of the threshold under uniformly bounded multiplicative perturbations that converge in measure. This characterizes the precise role of the tail condition: if it fails, a perturbation supported on a shrinking set can lower the threshold by a fixed factor.

\begin{proposition}\label{prop:root-tail-continuity}
Let $0\leq p\in\mathcal D_\omega$ be nonzero. The following are equivalent:
\begin{enumerate}
\item $p\in\mathcal D_\omega^0$;
\item the densities $p k_T$, for $T\in\mathscr G_\omega$, are uniformly integrable;
\item whenever $F_j$ are smooth real functions, uniformly bounded and tending to zero in measure, one has
\[
 b_\omega(\mathrm{e}^{F_j}p^n)\longrightarrow b_\omega(p^n).
\]
\end{enumerate}
If the first condition holds, the convergence in the third also holds for bounded measurable $F_j$ with the same uniform bound and convergence in measure. If it fails, there are smooth $F_j\geq0$, uniformly bounded and tending to zero in measure, such that
\[
 \mathrm{e}^{F_j}p^n\longrightarrow p^n\quad\text{in }L^1,
 \qquad
 b_\omega(\mathrm{e}^{F_j}p^n)\leq2^{-n}b_\omega(p^n).
\]
\end{proposition}
\begin{proof}
The first condition implies the second by \cref{thm:integrable-current-attainment}. Conversely, $\mathrm{d}V(\{p>K\})\to0$, so uniform integrability and \eqref{eq:integrable-current-supremum} give
\[
 \rho_\omega(p\mathbf1_{\{p>K\}})
 =\sup_{T\in\mathscr G_\omega}\int_{\{p>K\}}p k_T\,\mathrm{d}V
 \longrightarrow0.
\]
This is the first condition by \cref{prop:root-uniform-tails}.

If $p\in\mathcal D_\omega^0$ and $|F_j|\leq A$, then $p(\mathrm{e}^{F_j/n}-1)$ tends to zero in measure and its absolute value is at most $(\mathrm{e}^{A/n}+1)p$. The dominated-convergence criterion of \cref{prop:root-uniform-tails} makes its root norm tend to zero. The triangle inequality gives $\rho_\omega(p\mathrm{e}^{F_j/n})\to\rho_\omega(p)>0$. Taking negative $n$th powers proves the third condition, including measurable perturbations.

Suppose now that the tails do not tend to zero, and put $R=\rho_\omega(p)>0$. The smooth Gauduchon formula yields $\delta>0$, numbers $K_j\to\infty$ and smooth strictly positive $T_j\in\mathscr G_\omega$ such that
\[
 \int_{\{p>K_j\}}p k_{T_j}\,\mathrm{d}V>\delta.
\]
The finite measure $p k_{T_j}\,\mathrm{d}V$ is regular. Choose a compact set $C_j\subset\{p>K_j\}$ whose mass exceeds $\delta/2$, and an open neighborhood $O_j$ with
$\mathrm{d}V(O_j)\leq\mathrm{d}V(\{p>K_j\})+1/j$. Take a smooth cutoff $0\leq\chi_j\leq1$ that equals one on $C_j$ and is supported in $O_j$. Choose a fixed $A>0$ with $\mathrm{e}^{A/n}\delta/2>2R$, and set $F_j=A\chi_j$. Then $F_j\to0$ in measure, while
\[
 \rho_\omega(p\mathrm{e}^{F_j/n})
 \geq\int_Xp\mathrm{e}^{F_j/n}k_{T_j}\,\mathrm{d}V
 >2R.
\]
The threshold formula gives the asserted uniform gap. Finally,
$|\mathrm{e}^{F_j}-1|p^n\leq(\mathrm{e}^A-1)p^n\mathbf1_{O_j}$, whose integral tends to zero by absolute continuity of the integral. This proves the last assertion and the converse to the third condition.
\end{proof}

For $f=p^n$, these perturbations satisfy $\mathrm{e}^{-A}f\leq\mathrm{e}^{F_j}f\leq\mathrm{e}^Af$ with a fixed $A$. The continuity statement is restricted to such an order interval. The unrestricted $L^1$ instability in \cref{sec:individual-density-stability} allows approximations without these bounds.

For attainment at one density, uniform integrability of a single maximizing sequence is enough. This weaker condition also identifies the limit of its weighted measures.

\begin{proposition}\label{prop:selected-gauduchon-attainment}
Let $0\leq p\in\mathcal D_\omega$ be nonzero and put $R=\rho_\omega(p)$. Suppose $T_j\in\mathscr G_\omega$ converge weakly to $T$ and satisfy $D_p(T_j)\to R$. If the densities $p k_{T_j}$ are uniformly integrable, then
\[
 D_p(T)=R,\qquad
 p k_{T_j}\,\mathrm{d}V\rightharpoonup p k_T\,\mathrm{d}V.
\]
Consequently, a uniformly integrable maximizing sequence has a maximizing current among its weak cluster limits.
\end{proposition}
\begin{proof}
For positive integers $K$, put $p_K=\min(p,K)$. Since $p$ is finite almost everywhere, uniform integrability gives
\[
 \varepsilon_K:=\sup_j\int_{\{p>K\}}p k_{T_j}\,\mathrm{d}V
       \longrightarrow0.
\]
For every nonnegative continuous $h$, the weight $hp_K$ is bounded and measurable. The bounded-weight upper semicontinuity proved in \cref{thm:integrable-current-attainment} therefore gives
\[
 \limsup_j\int_Xhp k_{T_j}\,\mathrm{d}V
 \leq \int_Xhp_K k_T\,\mathrm{d}V
                   +\|h\|_\infty\varepsilon_K.
\]
Let $K\to\infty$. For $h=1$ this yields $R\leq D_p(T)$; the supremum formula \eqref{eq:integrable-current-supremum} gives the reverse inequality. If $\mu$ is a weak cluster limit of $p k_{T_j}\,\mathrm{d}V$, the same estimate for all nonnegative continuous $h$ gives $\mu\leq p k_T\,\mathrm{d}V$. Both measures have mass $R$, so they agree. Their bounded masses give weak compactness, and every cluster limit has now been identified. The final assertion follows from compactness of $\mathscr G_\omega$.
\end{proof}


\subsection{Supporting probabilities for an integrable density}

Under the vanishing-tail condition, we can also perturb the density by continuous multiplicative factors $\mathrm{e}^F$. The logarithm of the threshold is then a concave function of $F$, and its supporting probabilities are given by the maximizing currents, just as for the continuous densities in \cref{sec:gaud-current-adjoints}.

Let $f\in L^1(\mathrm{d}V)$ be nonzero and nonnegative, and assume $p=f^{1/n}\in\mathcal D_\omega^0$. For a continuous real $F$, set
\[
 B_f(F)=\log b_\omega(\mathrm{e}^Ff)
       =-n\log\rho_\omega(p \mathrm{e}^{F/n}),\qquad
 M_f(F)=\rho_\omega(p \mathrm{e}^{F/n}).
\]
Both quantities are finite and $M_f(F)>0$, since $\mathcal D_\omega^0$ is an order ideal. Let $\mathscr M_f(F)$ be the set of currents in $\mathscr G_\omega$ maximizing $D_{p \mathrm{e}^{F/n}}$, which is nonempty by \cref{thm:integrable-current-attainment}, and put
\[
 \nu_{T,F}=\frac{p \mathrm{e}^{F/n} k_T}{M_f(F)}\,\mathrm{d}V
 \quad(T\in\mathscr M_f(F)).
\]
These are probability measures. Let $\mathcal A_f(F)$, the set of \emph{supporting probabilities} of $B_f$ at $F$, consist of all Borel probability measures $\nu$ satisfying
\begin{equation}\label{eq:integrable-threshold-support}
 B_f(G)\leq B_f(F)-\int_X(G-F)\,\mathrm{d}\nu
 \qquad(G\in C(X,\mathbb R)).
\end{equation}

\Needspace{20\baselineskip}
\begin{proposition}\label{prop:integrable-threshold-adjoints}
The functional $B_f$ is concave, decreasing and $1$-Lipschitz for the uniform norm, and $B_f(F+a)=B_f(F)-a$ for constants $a$. Moreover,
\[
 \mathcal A_f(F)=\{\nu_{T,F}:T\in\mathscr M_f(F)\}.
\]
This set is nonempty, convex and weakly compact, and its measures are absolutely continuous. For each $R>0$, the densities of all these measures, for all $F$ with $\|F\|_\infty\leq R$, are uniformly integrable. For $H\in C(X,\mathbb R)$, the one-sided directional derivatives of $B_f$ at $F$ are
\[
 (B_f)'_+(F;H)=-\max_{\nu\in\mathcal A_f(F)}\int_XH\,\mathrm{d}\nu,
 \qquad
 (B_f)'_-(F;H)=-\min_{\nu\in\mathcal A_f(F)}\int_XH\,\mathrm{d}\nu.
\]
\end{proposition}
\begin{proof}
For each $T$, H\"older's inequality makes $F\mapsto\int_Xp \mathrm{e}^{F/n}k_T\,\mathrm{d}V$ logarithmically convex, and taking suprema preserves the inequality defining logarithmic convexity. Thus $M_f$ is logarithmically convex and $B_f=-n\log M_f$ is concave. Monotonicity and homogeneity of $\rho_\omega$ give the translation rule and, by comparison with constant translates, the Lipschitz estimate.

For fixed $F$, the weight $p\mathrm{e}^{F/n}$ belongs to $\mathcal D_\omega^0$, so the densities $p\mathrm{e}^{F/n}k_T$, $T\in\mathscr G_\omega$, are uniformly integrable by \cref{thm:integrable-current-attainment}. If $T_j\to T$ and $D_{p\mathrm{e}^{F/n}}(T_j)\to M_f(F)$, \cref{prop:selected-gauduchon-attainment} makes $T$ a maximizer and gives convergence of the weighted measures. For $T_1,T_2\in\mathscr G_\omega$ and real $0\leq t\leq1$, matrix determinant concavity gives
\[
 p \mathrm{e}^{F/n}k_{tT_1+(1-t)T_2}\,\mathrm{d}V
 \geq t p \mathrm{e}^{F/n}k_{T_1}\,\mathrm{d}V
       +(1-t)p \mathrm{e}^{F/n}k_{T_2}\,\mathrm{d}V.
\]
For two maximizing currents, the current on the left also maximizes, and equal total masses turn this inequality into an equality. The set $\mathscr N_f(F)=\{\nu_{T,F}:T\in\mathscr M_f(F)\}$ is therefore nonempty, compact and convex.

\Needspace{6\baselineskip}
A maximizing $T$ gives
\[
 M_f(G)\geq M_f(F)\int_X\mathrm{e}^{(G-F)/n}\,\mathrm{d}\nu_{T,F}.
\]
Taking $-n$ times the logarithm and applying Jensen's inequality proves \eqref{eq:integrable-threshold-support}, so $\mathscr N_f(F)\subset\mathcal A_f(F)$. For the reverse inclusion, a Taylor expansion of the exponential gives
\[
 D_{p \mathrm{e}^{(F+tH)/n}}(T)
 =D_{p \mathrm{e}^{F/n}}(T)+\frac tn\int_X H p \mathrm{e}^{F/n} k_T\,\mathrm{d}V+O(t^2),
\]
where the error is uniform in $T$ since $D_{p \mathrm{e}^{F/n}}(T)\leq M_f(F)$. Maximizers at $F+tH$ are near-maximizers at $F$, by the multiplicative bounds $\mathrm{e}^{-t\|H\|_\infty/n}\leq M_f(F+tH)/M_f(F)\leq\mathrm{e}^{t\|H\|_\infty/n}$. The convergence of their weighted measures proved above and this uniform expansion are exactly the two inputs of the one-sided differentiation argument in the proof of \cref{thm:gaud-current-all-adjoints}, which gives
\[
 (B_f)'_+(F;H)=-\max_{\nu\in\mathscr N_f(F)}\int_XH\,\mathrm{d}\nu.
\]
The separation argument of that proof then identifies all supporting probabilities with $\mathscr N_f(F)$, and replacing $H$ by $-H$ gives the left derivative.

Finally, $M_f(F)\geq \mathrm{e}^{-R/n}\rho_\omega(p)$ when $\|F\|_\infty\leq R$, so the densities of the probabilities $\nu_{T,F}$ are at most $\mathrm{e}^{2R/n}\rho_\omega(p)^{-1}p k_T$. The uniform integrability of the densities $pk_T$ in \cref{thm:integrable-current-attainment} proves the remaining assertions.
\end{proof}

When $f=1$, we have $B_f=B_0$, and the proposition reduces to \cref{thm:gaud-current-all-adjoints}. For general $f$, the functional $B_f$ is defined through the subsolution threshold and does not presuppose that the corresponding equation has a solution. If $v_-(\omega)>0$, every nonzero nonnegative $f\in L^1$ satisfies $f^{1/n}\in\mathcal D_\omega^0$ by \cref{prop:root-uniform-tails}, and \cref{cor:positive-lower-mass-exact-threshold} identifies $b_\omega(f)$ as the unique constant for which the equation $\mathrm{MA}^{\mathrm{ext}}_\omega(v)=b\,f\,\mathrm{d}V$ has a solution $v$ in the envelope class $\mathcal E_{\mathrm{env}}(X,\omega,0)$ of that corollary. Without positive lower mass, the uniform tail condition is an explicit sufficient hypothesis for this compact dual description.

\Needspace{16\baselineskip}
\subsection{Trace inequalities at the subsolution threshold}
\label{subsec:threshold-trace-inequalities}

For a nonzero $0\leq f\in L^1(\mathrm{d}V)$ with $b=b_\omega(f)>0$, recall from \cref{sec:individual-density-subsolutions} that a threshold subsolution is a function $u\in\PSH(X,\omega)$ with $\mathrm{MA}^{\mathrm{ext}}_\omega(u)\geq bf\,\mathrm{d}V$, and a threshold solution is one for which equality holds; threshold subsolutions exist by \cref{thm:individual-density-subsolution}. A smooth Gauduchon test can be paired with the full current $\omega+\ddc u$ of a threshold subsolution, and by \cref{lem:unbounded-gauduchon-trace-density} the resulting trace measure is bounded below by the determinant-weighted density. The nonnegative difference has a mass that depends only on the test. Along a maximizing sequence this mass tends to zero, uniformly over all threshold subsolutions. We write $\|\cdot\|_{\mathrm{TV}}$ for the total variation of a finite signed measure.

\Needspace{14\baselineskip}
\begin{proposition}\label{prop:threshold-trace-saturation}\index{Gauduchon duality!trace saturation}
Let $0\leq f\in L^1$ be nonzero with $b=b_\omega(f)>0$, write $p=f^{1/n}$, and let $u$ be a threshold subsolution. For every smooth strictly positive $T\in\mathscr G_\omega$, the measure
\[
 R_{u,T}=(\omega+\ddc u)\wedge T-b^{1/n}p k_T\,\mathrm{d}V
\]
is positive and has mass
\[
 R_{u,T}(X)=1-b^{1/n}D_p(T).
\]
In particular, if $(T_j)$ is a sequence of smooth strictly positive members of $\mathscr G_\omega$ with $D_p(T_j)\to\rho_\omega(p)$, then $\|R_{u,T_j}\|_{\mathrm{TV}}\to0$, uniformly over all threshold subsolutions $u$.

If also $p\in\mathcal D_\omega^0$ and $T_j\to T$ as currents, then $T$ maximizes $D_p$ and
\[
 (\omega+\ddc u)\wedge T_j
       \rightharpoonup b^{1/n}p k_T\,\mathrm{d}V.
\]
\end{proposition}
\begin{proof}
A smooth strictly positive $T\in\mathscr G_\omega$ is $\gamma^{n-1}$ for a Gauduchon metric $\gamma$ normalized by $\int_X\omega\wedge\gamma^{n-1}=1$. \Cref{lem:unbounded-gauduchon-trace-density} proves positivity of $R_{u,T}$, and Stokes' theorem for currents gives $((\omega+\ddc u)\wedge T)(X)=\int_X\omega\wedge T=1$. This proves the mass formula. Since $\rho_\omega(p)=b^{-1/n}$, the nonnegative mass tends to zero along a maximizing sequence; it equals the total variation norm and does not depend on $u$.

Under the tail hypothesis, \cref{thm:integrable-current-attainment} gives uniform integrability of the densities $p k_{T_j}$. Thus \cref{prop:selected-gauduchon-attainment} makes $T$ a maximizer and gives
$p k_{T_j}\,\mathrm{d}V\rightharpoonup p k_T\,\mathrm{d}V$. Adding the error measures $R_{u,T_j}$, which tend to zero in total variation, proves the limit formula.
\end{proof}

The limit formula concerns smooth maximizing sequences. It does not define the product $(\omega+\ddc u)\wedge T$ for a singular current $T$, nor does it show that the limit is the same for different maximizing currents.

If a maximizing test is itself a smooth strictly positive form, the trace inequality becomes an equality of currents, and since $n\geq2$ this forces an exact smooth equation. The local regularity needed for this comes from \cref{lem:smooth-conformal-current}.

\Needspace{12\baselineskip}
\begin{corollary}\label{cor:smooth-dual-exact-equation}\index{Gauduchon duality!smooth maximizing test}\label{cor:smooth-dual-threshold-rigidity}
Let $0\leq f\in L^1(\mathrm{d}V)$ be nonzero with $b=b_\omega(f)>0$. If $D_{f^{1/n}}$ has a smooth strictly positive maximizer $T\in\mathscr G_\omega$, then $f^{1/n}$ agrees almost everywhere with a smooth nonnegative function, and every threshold subsolution $u$ is smooth and satisfies
\[
 \mathrm{MA}_\omega(u)=bf\,\mathrm{d}V.
\]
In particular a smooth threshold solution exists, and any two threshold subsolutions differ by a constant.
\end{corollary}
\Needspace{8\baselineskip}
\begin{proof}
Let $u$ be any threshold subsolution; one exists by \cref{thm:individual-density-subsolution}. Write $T=\gamma^{n-1}$, where $\gamma$ is the smooth Gauduchon metric given by the positive $(n-1)$st root of $T$, and set
\[
 a=\left(\frac{bf\,\mathrm{d}V}{\gamma^n}\right)^{1/n}.
\]
Put $A=\omega+\ddc u\geq0$. Since $T$ maximizes, the positive measure $R_{u,T}$ of \cref{prop:threshold-trace-saturation} has mass zero, which gives the equality
\[
 A\wedge\gamma^{n-1}=a\gamma^n.
\]
The coefficient measures of the positive current $A$ are controlled by this trace measure, which is absolutely continuous. Thus $A$ has a matrix density that is positive semidefinite almost everywhere, with coefficients in $L^n$ because $a\in L^n$.

\Needspace{5\baselineskip}
For every smooth Hermitian form $\beta$, \cref{lem:unbounded-gauduchon-trace-density} also gives
\[
 A\wedge\beta^{n-1}
 \geq\left(\frac{bf\,\mathrm{d}V}{\beta^n}\right)^{1/n}\beta^n.
\]
The resulting pointwise matrix inequalities hold almost everywhere, simultaneously for all positive definite matrices: on each coordinate chart, first use a countable dense family of constant positive matrices, realized on smaller charts by global smooth Hermitian forms, and then use continuity in the matrix. The variational formula for $(\det P)^{1/n}$ in the proof of \cref{lem:weak-gaud-determinant-dual} therefore yields
\[
 \left(\frac{A^n}{\gamma^n}\right)^{1/n}\geq a
 \quad\text{almost everywhere}.
\]
Here $A^n$ denotes only the determinant of the measurable coefficient matrix of $A$; no nonlinear product of currents is used. The trace equality gives the reverse inequality, by the arithmetic--geometric mean inequality. Equality then forces all eigenvalues of $A$ relative to $\gamma$ to equal $a$, also where $a=0$. Consequently $A=a\gamma$ as currents.

By \cref{lem:smooth-conformal-current}, applied with $\theta=\omega$, both $a$ and $u$ are smooth, and $f^{1/n}=b^{-1/n}a(\gamma^n/\mathrm{d}V)^{1/n}$ has a smooth nonnegative representative. Since $u$ is smooth, we may now take the top power: $(\omega+\ddc u)^n=a^n\gamma^n=bf\,\mathrm{d}V$. Finally, the function $a$ was defined from $f$, $b$ and $\gamma$ alone, so two threshold subsolutions $u_1,u_2$ satisfy $\omega+\ddc u_1=a\gamma=\omega+\ddc u_2$. Their smooth difference is therefore pluriharmonic, hence constant on the connected compact manifold $X$.
\end{proof}

If $f$ is smooth and strictly positive, \cref{thm:gaud-volume-dual-formula} provides a smooth strictly positive maximizing Gauduchon form, obtained from the smooth equation. The corollary therefore applies, and shows in particular that all threshold subsolutions, bounded or not, differ from the smooth solution by constants.


\begin{corollary}\label{cor:smooth-dual-exponential-limit}
Under the hypotheses of \cref{cor:smooth-dual-exact-equation}, let $u$ be the threshold solution with $\sup_Xu=0$. If $u_\lambda$ is the unique bounded solution of
\[
 \mathrm{MA}_\omega(u_\lambda)=\mathrm{e}^{\lambda u_\lambda}f\,\mathrm{d}V,
 \qquad\lambda>0,
\]
then $u_\lambda-\sup_Xu_\lambda\to u$ uniformly as $\lambda\downarrow0$.
\end{corollary}
\begin{proof}
By \cref{cor:smooth-dual-exact-equation}, $f$ agrees almost everywhere with a smooth function, so it belongs to $L^p$ for every finite $p$. By \cref{prop:individual-normalizer-lp}, the constant $c$ for which $\mathrm{MA}_\omega(v)=cf\,\mathrm{d}V$ has a bounded solution is $b_\omega(f)$. Thus \cref{prop:symmetric-equation-clusters}, with the trivial symmetry group, shows that the normalized family $u_\lambda-\sup_Xu_\lambda$ is uniformly precompact as $\lambda\downarrow0$ and that every cluster limit is a normalized solution with this constant. Each cluster limit is therefore a normalized threshold subsolution and equals $u$ by \cref{cor:smooth-dual-exact-equation}. If the whole family did not converge uniformly, a sequence staying at a fixed positive uniform distance from $u$ would have a uniformly convergent subsequence, a contradiction.
\end{proof}

\begin{corollary}\label{cor:kahler-smooth-dual-attainment}
Suppose $\omega$ is K\"ahler. For a nonzero nonnegative $f\in L^1(\mathrm{d}V)$, the functional $D_{f^{1/n}}$ has a smooth strictly positive maximizer in $\mathscr G_\omega$ if and only if $f$ agrees almost everywhere with a smooth strictly positive function.
\end{corollary}
\begin{proof}
If a smooth strictly positive maximizer exists, the proof of \cref{cor:smooth-dual-exact-equation} gives
\[
 \omega+\ddc u=a\gamma,
 \qquad a\geq0,
\]
with $a,u$ smooth and $a\not\equiv0$. Since $\mathrm{d}\omega=0$, the first-order equation in \cref{lem:smooth-conformal-current} becomes $\mathrm{d}a=-aC$ for a smooth one-form $C$. Along any smooth path $\sigma$, this is the linear scalar equation
\[
 (a\circ\sigma)'=-(a\circ\sigma)C(\dot\sigma),
\]
whose solution is the initial value multiplied by a positive exponential. Since $X$ is connected and $a\not\equiv0$, it follows that $a>0$ everywhere, and $f=b_\omega(f)^{-1}a^n\gamma^n/\mathrm{d}V$ is smooth and strictly positive. Conversely, if $f$ is smooth and strictly positive, \cref{thm:weak-gaud-compact-duality}, applied to $F=\log f$, for which $\mathfrak D_F=D_{f^{1/n}}$, provides a smooth strictly positive maximizer.
\end{proof}

For nonclosed backgrounds, the strict positivity of $f$ in \cref{cor:kahler-smooth-dual-attainment} can fail on every compact complex manifold of dimension at least two, including K\"ahler manifolds.

\begin{example}\label{ex:nonclosed-dual-zero-density}
On every connected compact complex manifold $X$ of dimension $n\geq2$, with a fixed smooth strictly positive probability volume $\mathrm{d}V$, there is a nonclosed Hermitian form $\omega$ and a nonzero smooth nonnegative density $f$ vanishing on an open set such that $b_\omega(f)=1$ and $D_{f^{1/n}}$ has a smooth strictly positive maximizer.
\end{example}
\begin{proof}
Fix a Gauduchon metric $\gamma$ and a point $x$. Choose a smooth real function $\psi$ such that $\ddc\psi\leq-\gamma$ on a coordinate neighborhood $U$ of $x$; in holomorphic coordinates $z=(z_1,\ldots,z_n)$ centered at $x$, such a function is obtained by multiplying $-M|z|^2$, with $M$ large enough, by a cutoff that equals one near $x$, and then shrinking $U$.

Choose a smooth function $\chi$ with $0\leq\chi\leq1$ that vanishes on a smaller neighborhood of $x$ and equals one outside $U$. Put
\[
 a=L\chi,\qquad \omega=a\gamma-\ddc\psi.
\]
Since $X$ is compact, $\omega$ is strictly positive for sufficiently large $L>0$: on $U$ the term $-\ddc\psi$ is at least $\gamma$, and outside $U$ the term $L\gamma$ dominates $\ddc\psi$. Set $f=a^n\gamma^n/\mathrm{d}V$. The smooth function $\psi$ satisfies $\omega+\ddc\psi=a\gamma\geq0$ and $\mathrm{MA}_\omega(\psi)=f\,\mathrm{d}V$, so $b_\omega(f)\geq1$.

Stokes' theorem and $\ddc\gamma^{n-1}=0$ give
\[
 \int_X\omega\wedge\gamma^{n-1}
 =\int_Xa\gamma^n
 =\int_X(f\,\mathrm{d}V/\gamma^n)^{1/n}\gamma^n.
\]
Testing \eqref{eq:integrable-gauduchon-threshold} with $\gamma$ therefore gives the opposite inequality $b_\omega(f)\leq1$. Hence $b_\omega(f)=1$, the metric $\gamma$ attains the infimum, and the normalized current $\gamma^{n-1}/\int_X\omega\wedge\gamma^{n-1}$ is a smooth strictly positive maximizer of $D_{f^{1/n}}$. If $\omega$ were closed, \cref{cor:kahler-smooth-dual-attainment} would force $f$ to be strictly positive, contrary to its vanishing on an open set. Thus $\omega$ is nonclosed.
\end{proof}

In \cref{ex:nonclosed-dual-zero-density}, \cref{cor:smooth-dual-exact-equation,cor:smooth-dual-exponential-limit} show that every threshold subsolution is a constant translate of $\psi$ and that the normalized solutions of the exponential equations converge uniformly to $\psi-\sup_X\psi$. Thus vanishing of the density on an open set does not by itself force nonuniqueness or the failure of this approximation.

The restriction $n\geq2$ in \cref{cor:smooth-dual-exact-equation} is necessary. On a compact connected Riemann surface $X$ with a Hermitian form $\omega$, a normalized Gauduchon current has bidegree $(0,0)$, and the equation $\ddc T=0$ makes it a harmonic function, hence the constant $(\int_X\omega)^{-1}$. This single smooth strictly positive current maximizes for every nonzero nonnegative $L^1$ density, which need not have a smooth representative.


\subsection{A smooth maximizing form with an analytic degeneracy set}

Strict positivity of a smooth maximizing form can be weakened if its degeneracy is confined to a proper analytic subset. The proof is local away from that subset. The possible ambiguity of the potential is then divisorial.

\begin{theorem}\label{thm:semipositive-dual-exact-equation}
Let $\omega$ be Hermitian on a connected compact complex manifold of dimension $n\geq2$, let $0\leq f\in L^1(\mathrm{d}V)$ be nonzero, and assume $b=b_\omega(f)>0$. Suppose $D_{f^{1/n}}$ has a smooth semipositive maximizer $T\in\mathscr G_\omega$ which is strictly positive on $X\setminus Z$, where $Z$ is a proper analytic subset. Then every threshold subsolution is smooth on $X\setminus Z$ and satisfies
\[
 \mathrm{MA}^{\mathrm{ext}}_\omega(u)=bf\,\mathrm{d}V
\]
on $X$. Such a solution exists and has finite mass. The function $f^{1/n}$ has a smooth nonnegative representative on $X\setminus Z$.

If $u,v$ are threshold subsolutions and $Z_1,\ldots,Z_m$ are the prime divisorial components of $Z$, there are real numbers $c_j$ such that
\begin{equation}\label{eq:semipositive-dual-divisorial-difference}
 \ddc(u-v)=\sum_{j=1}^m c_j[Z_j],\qquad
 \sum_{j=1}^m c_j\{Z_j\}=0
       \quad\text{in }H^{1,1}_{\mathrm{BC}}(X,\mathbb R).
\end{equation}
Consequently threshold solutions are unique modulo constants within each multiplier-ideal equivalence class, hence within each ordinary singularity type. They are unique modulo constants without a prescribed type if the classes $\{Z_j\}$ are linearly independent. In particular this holds when $Z$ has codimension at least two.
\end{theorem}
\begin{proof}
Put $p=f^{1/n}$ and $A=\omega+\ddc u$, where $u$ is any threshold subsolution. Existence of at least one such $u$ is \cref{thm:individual-density-subsolution}. Fix a smooth strictly positive $T_0\in\mathscr G_\omega$ and set $T_\varepsilon=(T+\varepsilon T_0)/(1+\varepsilon)$. Each $T_\varepsilon$ belongs to $\mathscr G_\omega$ and has a smooth Hermitian root, so \cref{prop:threshold-trace-saturation} gives
\[
 A\wedge T_\varepsilon\geq b^{1/n}p k_{T_\varepsilon}\,\mathrm{d}V.
\]
For smooth semipositive forms, $k_T$ is the determinant root of the coefficient matrix, in the normalization of \cref{sec:weak-gauduchon-duality}. This is a continuous function of that matrix. Thus $k_{T_\varepsilon}\to k_T$ uniformly on compact $X$. Since $p\in L^1$, passage to the limit gives
\[
 A\wedge T\geq b^{1/n}p k_T\,\mathrm{d}V.
\]
The left side has mass one by Stokes and $\ddc T=0$. The right side has mass $b^{1/n}D_p(T)=1$, because $T$ maximizes. Their difference is positive and has zero mass, hence
\begin{equation}\label{eq:semipositive-dual-trace-equality}
 A\wedge T=b^{1/n}p k_T\,\mathrm{d}V.
\end{equation}

On $U=X\setminus Z$ write $T=\gamma^{n-1}$, where $\gamma$ is smooth and strictly positive, and set $a=(bf\,\mathrm{d}V/\gamma^n)^{1/n}$. The trace equality becomes $A\wedge\gamma^{n-1}=a\gamma^n$. On relatively compact subsets of $U$ this controls every coefficient measure of $A$ by an $L^n$ density. The all-Hermitian trace inequalities and matrix arithmetic--geometric mean, exactly as in the proof of \cref{cor:smooth-dual-exact-equation}, force $A=a\gamma$ there. This step uses only the local positivity of $\gamma$; the other Hermitian tests extend smoothly from smaller coordinate charts to $X$.

The local regularity lemma \cref{lem:smooth-conformal-current} makes $u,a$ smooth on $U$. Taking the smooth top power proves the equation there, and $f^{1/n}=b^{-1/n}a(\gamma^n/\mathrm{d}V)^{1/n}$ is smooth and nonnegative there. Both $\mathrm{MA}^{\mathrm{ext}}_\omega(u)$ and $f\,\mathrm{d}V$ vanish on the pluripolar set $Z$. Hence the equality holds on $X$, and its total mass is $b\int_Xf\,\mathrm{d}V<+\infty$.

For two threshold subsolutions the formula $\omega+\ddc u=a\gamma$ has the same right side on $U$. Thus $S=\ddc(u-v)$ is a real closed current supported on $Z$. It has order zero, being the difference of the positive currents $\omega+\ddc u$ and $\omega+\ddc v$. The classical support theorem gives $S=\sum_jc_j[Z_j]$; components of higher codimension carry no such current. This is \cite[Chapter III, Corollaries 2.11 and 2.14]{DemCADG}. To check the class relation with smooth representatives, write $[Z_j]=\beta_j+\ddc\ell_j$, using a smooth Hermitian metric on the divisor line bundle. Then
\[
 \sum_jc_j\beta_j
   =\ddc\bigl(u-v-\sum_jc_j\ell_j\bigr).
\]
The left side is smooth. Taking its local trace and applying elliptic regularity makes the global $L^1$ potential on the right smooth. Hence the displayed form is $\ddc$-exact in smooth Bott--Chern cohomology, proving \eqref{eq:semipositive-dual-divisorial-difference}.

Near a regular point of $Z_j$ lying on no other component, choose a local defining function $z_1$. In our normalization $\ddc\log|z_1|^2=[Z_j]$, so
\[
 \ddc\bigl(u-v-c_j\log|z_1|^2\bigr)=0.
\]
The expression in parentheses is a smooth pluriharmonic function by distributional elliptic regularity. Taking generic Lelong numbers therefore gives $c_j=\nu_{Z_j}(u)-\nu_{Z_j}(v)$. If $u\sim_{\mathcal I}v$, these numbers agree by \cref{thm:local-i-inputs}, so all coefficients vanish. Ordinary equivalence implies multiplier-ideal equivalence. Alternatively, independence of the divisor classes forces them all to vanish directly from \eqref{eq:semipositive-dual-divisorial-difference}. In either case $\ddc(u-v)=0$ globally, and a real pluriharmonic distribution on connected compact $X$ is constant. When $Z$ has codimension at least two, the sum is empty.
\end{proof}

The theorem assumes a smooth form. It neither defines a wedge product with an arbitrary singular maximizing current nor proves that such a current has an analytic degeneracy set. The divisor relation is necessary for two threshold solutions to differ; it does not assert that every relation is realized. No bounded mass, positive lower mass, root-tail condition, or convergence of an exponential approximation is used.


\Needspace{22\baselineskip}
\subsection{The limit of the exponential equations}
\label{subsec:semipositive-exponential-setting}

Use the data of \cref{thm:semipositive-dual-exact-equation}: the manifold $X$ of dimension $n\geq2$, Hermitian form $\omega$, probability volume $\mathrm{d}V$, nonzero nonnegative density $f$, positive threshold $b=b_\omega(f)$, and smooth semipositive maximizing form $T$ with analytic degeneracy set $Z$. Put $p=f^{1/n}$ and let $k_T$ be its determinant density. All $L^1$ convergence below uses $\mathrm{d}V$.

A smooth semipositive maximizing form also identifies the normalized exponential limit. The argument passes the linear trace inequalities to the limit before using the local regularity proof.

\begin{proposition}\label{prop:semipositive-dual-exponential-limit}
Assume the hypotheses of \cref{thm:semipositive-dual-exact-equation}, and let $u_0$ be its unique threshold solution with $\sup_Xu_0=0$. For real $\lambda>0$, let $u_\lambda$ be the greatest exponential solution of \cref{cor:individual-positive-threshold-exponential}, and put
\[
 v_\lambda=u_\lambda-\sup_Xu_\lambda,\qquad
 c_\lambda=\mathrm{e}^{\lambda\sup_Xu_\lambda}.
\]
There is a constant $C$ independent of $\lambda>0$ such that
\[
 b\leq c_\lambda\leq b\,\mathrm{e}^{C\lambda},\qquad
 u_0-C\leq v_\lambda\leq0.
\]
Moreover $v_\lambda\to u_0$ in $L^1(X)$ as $\lambda\downarrow0$.
\end{proposition}
\begin{proof}
Write $p=f^{1/n}$ and introduce the probability measure
\[
 \mathrm{d}\nu=b^{1/n}p k_T\,\mathrm{d}V.
\]
It has mass one because $T$ maximizes. Its density belongs to $L^n$, since $p\in L^n$ and the smooth semipositive form $T$ has bounded determinant density $k_T$. By \cref{lem:uniform-exponential}, normalized $\omega$-plurisubharmonic functions have uniformly bounded moments of every finite order. H\"older's inequality with exponents $n$ and $n/(n-1)$ therefore gives
\[
 \int_X(-v_\lambda)\,\mathrm{d}\nu\leq C.
\]
Here $n\geq2$ is used.

Apply the trace inequality to $u_\lambda$, or equivalently to $v_\lambda$, and pass from positive regularizations of $T$ to $T$ as in the preceding proof. After integration,
\[
 1\geq(c_\lambda/b)^{1/n}
       \int_X\mathrm{e}^{\lambda v_\lambda/n}\,\mathrm{d}\nu
 \geq(c_\lambda/b)^{1/n}\mathrm{e}^{-C\lambda/n},
\]
where the last step is Jensen's inequality. This proves the upper scalar bound. The lower bound and the comparison
\[
 v_\lambda\geq u_0-\lambda^{-1}\log(c_\lambda/b)
\]
are \cref{cor:individual-positive-threshold-exponential}, giving the asserted common lower barrier.

Take any sequence $\lambda_j\downarrow0$ and an $L^1$-convergent subsequence $v_{\lambda_j}\to v$, supplied by \cref{thm:normalized-compactness}. Its limit has supremum zero. The normalized $L^1$ bound implies $\lambda v_\lambda\to0$ in measure. Since $v_\lambda\leq0$ and $p\in L^1$, it follows that
\[
 c_\lambda^{1/n}\mathrm{e}^{\lambda v_\lambda/n}p
       \longrightarrow b^{1/n}p
       \quad\text{in }L^1(\mathrm{d}V).
\]
Put $A_\lambda=\omega+\ddc v_\lambda$. The positive trace errors
\[
 A_\lambda\wedge T
   -c_\lambda^{1/n}\mathrm{e}^{\lambda v_\lambda/n}p k_T\,\mathrm{d}V
\]
have masses tending to zero, since the two terms have total masses tending to one. The $L^1$ convergence of the potentials gives weak convergence of their linear currents. Consequently
\[
 (\omega+\ddc v)\wedge T=b^{1/n}p k_T\,\mathrm{d}V.
\]
For every smooth Hermitian form $\beta$, the all-Hermitian trace inequality passes to the same limit and gives
\[
 (\omega+\ddc v)\wedge\beta^{n-1}
   \geq(bf\,\mathrm{d}V/\beta^n)^{1/n}\beta^n.
\]
The local matrix and regularity argument in the proof of \cref{thm:semipositive-dual-exact-equation} now applies on $X\setminus Z$: with $\gamma$ the positive root of $T$ there and $a=(bf\,\mathrm{d}V/\gamma^n)^{1/n}$, it gives $\omega+\ddc v=a\gamma$, then smoothness, and then the exact equation there. Both extended measures ignore $Z$, so $v$ is a global threshold solution. Uniqueness and normalization give $v=u_0$.

Every sequence tending to zero has a subsequence with this same $L^1$ limit. Normalized compactness therefore proves convergence of the whole family.
\end{proof}

The lower barrier gives capacity convergence under the additional envelope condition below. If the maximizing form is strictly positive everywhere, \cref{cor:smooth-dual-exponential-limit} already gives uniform convergence.


\Needspace{15\baselineskip}
\begin{corollary}\label{cor:semipositive-dual-capacity-limit}
Under the hypotheses of \cref{prop:semipositive-dual-exponential-limit}, suppose in addition that $P_{\delta\omega}(u_0)$ is genuine for some $0<\delta<1$. Then $v_\lambda\to u_0$ in capacity as $\lambda\downarrow0$. In particular this holds if the threshold solution $u_0$ is bounded.
\end{corollary}
\begin{proof}
Fix a sequence $\lambda_j\downarrow0$. \Cref{prop:semipositive-dual-exponential-limit} gives $u_0-C\leq v_{\lambda_j}\leq0$ and $v_{\lambda_j}\to u_0$ in $L^1$. For all large $j$,
\[
 \mathrm{MA}^{\mathrm{ext}}_\omega(v_{\lambda_j})
   =c_{\lambda_j}\mathrm{e}^{\lambda_jv_{\lambda_j}}f\,\mathrm{d}V
   \leq b\mathrm{e}^{C}f\,\mathrm{d}V.
\]
The bounded functions $|\mathrm{e}^{v_{\lambda_j}}-\mathrm{e}^{u_0}|$ tend to zero in smooth-volume measure. Since $f\in L^1$, their integrals against $f\,\mathrm{d}V$ tend to zero. Thus the self-measure error in \eqref{eq:common-barrier-self-error} vanishes. Apply \cref{prop:lower-barrier-criterion} with $\theta=\omega$ and $\varphi=u_0$; its additional limit-type hypothesis is automatic because the limit is $u_0$ itself. This proves capacity convergence along every sequence and hence for the family. If $u_0$ is bounded, a constant below it is a competitor for $P_{\delta\omega}(u_0)$, so that envelope is genuine.
\end{proof}

The smaller-background condition is an additional hypothesis. Its validity for an arbitrary threshold solution in this semipositive dual setting is not established here.


\subsection{Envelope criteria and extension across the degeneracy set}
\label{subsec:semipositive-envelope-extension}

The capacity criterion above asks for a genuine smaller-background obstacle envelope. Two sufficient conditions come from different sources: invariant bounded mass controls the contact measures of stretched truncations, while a sublinear singularity bound supplies explicit competitors. Both conditions give zero Lelong numbers. Under dual attainment, a globally smooth root of rank at least two then permits smooth extension across the analytic degeneracy set.

We retain the connected compact manifold $X$ of dimension $n\geq2$, the Hermitian background $\omega$ and the smooth positive probability volume $\mathrm{d}V$. All obstacle inequalities concern canonical potentials. We will use one consequence of \cref{lem:gauduchon-uniform-lelong}: if, for every $0<\delta<1$, a genuine $v_\delta\in\PSH(X,\delta\omega)$ satisfies $v_\delta\leq u$, then
\[
 0\leq\nu(u,x)\leq\nu(v_\delta,x)
       =\delta\nu(v_\delta/\delta,x)\leq\delta C_\omega.
\]
Thus $u$ has zero Lelong numbers. The bound is uniform in $\delta$ because $v_\delta/\delta$ has the fixed background $\omega$.

\begin{theorem}\label{thm:invariant-mass-threshold-envelopes}\index{Envelope!invariant mass criterion}
Suppose $\ddc(\omega^j)=0$ for $1\leq j<n$. Put $V=\int_X\omega^n$. For every nonzero $0\leq f\in L^1(\mathrm{d}V)$, with $M=\int_Xf\,\mathrm{d}V$, one has
\[
 b_\omega(f)=\frac VM.
\]
Every threshold subsolution satisfies $\mathrm{MA}^{\mathrm{ext}}_\omega(u)=(V/M)f\,\mathrm{d}V$ and has mass $V$. Its obstacle envelope $P_{\delta\omega}(u)$ is genuine for every $0<\delta<1$, and $u$ has zero Lelong numbers. More generally, these last two conclusions hold for every $\omega$-plurisubharmonic potential of extended mass $V$.
\end{theorem}
\begin{proof}
By \cref{thm:mass-invariance-characterization}, every bounded $\omega$-plurisubharmonic potential has mass $V$. Truncation therefore bounds every extended mass above by $V$. For $f_k=\min(f,k)$, integration of the bounded equation gives $c_\omega(f_k)=V/\int_Xf_k\,\mathrm{d}V$. The definition \eqref{eq:individual-truncated-normalizer} gives $b_\omega(f)=V/M$. A threshold subsolution has mass at least $b_\omega(f)M=V$. Its mass is consequently $V$, and the nonnegative difference from $b_\omega(f)f\,\mathrm{d}V$ has zero mass.

Let now $u$ have extended mass $V$. Normalize $u\leq0$, put $\mu=\mathrm{MA}^{\mathrm{ext}}_\omega(u)$, and fix $a>1$. For $u_j=\max(u,-j)$, let $w_j=P_\omega(au_j)$. Constants give $-aj\leq w_j\leq0$. The bounded contact theorem \cref{thm:bounded-contact-support} concentrates its measure on $C_j=\{w_j=au_j\}$. Canonical almost-everywhere order gives $w_j\leq au_j$ everywhere. Both potentials are $a\omega$-plurisubharmonic, so positivity and ordered contact give
\[
 \mathrm{MA}_\omega(w_j)
 \leq\mathbf1_{C_j}\mathrm{MA}_{a\omega}(w_j)
 \leq a^n\mathbf1_{C_j}\mathrm{MA}_\omega(u_j).
\]
For integers $j\geq N$, locality and the bounded mass identity yield
\[
 \int_{\{u_j\leq-N\}}\mathrm{MA}_\omega(u_j)
       =V-\mu(\{u>-N\})=\mu(\{u\leq-N\}).
\]
Choose $N$ with $a^n\mu(\{u\leq-N\})<V$. This is possible because $\mu$ is finite and gives no mass to the pole set. If $\sup_Xw_j<-aN$, its contact set lies in $\{u_j<-N\}$, contradicting the preceding inequalities and $\int_X\mathrm{MA}_\omega(w_j)=V$. Hence $\sup_Xw_j\geq-aN$ for $j\geq N$; for smaller $j$ the bound follows from $w_j\geq-aj$.

The envelopes decrease. Their uniform lower supremum prevents collapse to $-\infty$, by normalized quasi-plurisubharmonic compactness. Their genuine limit $w$ satisfies $w\leq au$, including at poles. With $\delta=1/a$, the function $\delta w$ is a genuine $\delta\omega$-plurisubharmonic competitor below $u$. Thus $P_{\delta\omega}(u)$ is genuine at every scale. The Lelong estimate preceding the theorem completes the proof. An additive normalization does not affect any conclusion.
\end{proof}

\begin{proposition}\label{prop:sublinear-threshold-envelope}\index{Envelope!sublinear singularity criterion}
Suppose a genuine quasi-plurisubharmonic $\psi\leq0$ satisfies $\ddc\psi\geq-C\omega$ for some $C>0$, and a genuine $\omega$-plurisubharmonic $u$ satisfies the canonical inequality
\[
 u\geq-A(1-\psi)^\alpha-B,
 \qquad A>0,\quad B\in\mathbb R,\quad0<\alpha<1.
\]
Then every $P_{\delta\omega}(u)$ is genuine, for $0<\delta<1$, and $u$ has zero Lelong numbers.
\end{proposition}
\begin{proof}
Choose $\alpha<\beta<1$ and $\varepsilon>0$ with $\varepsilon\beta C\leq\delta$. The number
\[
 K=B+\sup_{t\geq1}(At^\alpha-\varepsilon t^\beta)
\]
is finite, and $v_\delta=-\varepsilon(1-\psi)^\beta-K\leq u$. This composition is increasing and convex, with first derivative at most $\varepsilon\beta$ and nonnegative second derivative. Apply the quasi-plurisubharmonic chain inequality to $\max(\psi,-j)$ and pass to the decreasing limit. It gives $\ddc v_\delta\geq-\varepsilon\beta C\omega\geq-\delta\omega$. The limit is genuine because $\psi$ is finite almost everywhere. Its canonical pole values are $-\infty$, where the comparison still holds. Thus $v_\delta$ is an actual competitor at each scale. The preceding uniform Lelong estimate proves the last assertion.
\end{proof}

\subsubsection{A globally smooth root of rank at least two}

\begin{theorem}\label{thm:rank-two-dual-smooth-extension}\index{Gauduchon duality!smooth extension across degeneration}
Assume the semipositive attained setting of \cref{thm:semipositive-dual-exact-equation}, with greatest normalized threshold solution $\Phi_f$ and analytic degeneracy set $Z$. Suppose the maximizing form has the representation $T=\gamma^{n-1}$ for a globally smooth semipositive $(1,1)$-form $\gamma$ whose complex rank is at least two everywhere. If $\Phi_f$ has zero Lelong numbers, then it is smooth on $X$. This applies under either sufficient condition above. Every threshold subsolution is then smooth, since these subsolutions differ by constants.
\end{theorem}
\begin{proof}
Put $U=X\setminus Z$ and $A=\omega+\ddc\Phi_f$. The existing attained theorem makes $A$ smooth on $U$ and gives its exact global equation. The root $\gamma$ is positive on $U$. Stokes gives $\int_XA\wedge T=1$, while attainment gives
\[
 \int_U(A^n)^{1/n}(\gamma^n)^{(n-1)/n}
       =b_\omega(f)^{1/n}D_{f^{1/n}}(T)=1.
\]
The arithmetic--geometric mean gap is nonnegative and continuous on $U$, and its integral must vanish. Equality yields $A=r\gamma$ there for smooth $r\geq0$. Where $A^n=0$, equality forces $A=0$, so the same formula holds with $r=\operatorname{tr}_\gamma A/n=0$.

Since $\ddc\Phi_f$ is closed,
\[
 \mathrm dr\wedge\gamma+r\,\mathrm d\gamma=\mathrm d\omega
       \quad\hbox{on }U.
\]
For a semipositive form of rank at least two, the real-linear map $W_\gamma:\xi\mapsto\xi\wedge\gamma$ from one-forms to three-forms is injective. In a diagonal real unitary frame, for each coordinate of $\xi$ choose a positive coordinate two-plane not containing it. Its triple coefficient is that coordinate times the positive eigenvalue; no other disjoint diagonal pair contributes. Vanishing of the wedge therefore forces every coordinate to vanish.

Choose an auxiliary metric. The left inverse $L_\gamma=(W_\gamma^*W_\gamma)^{-1}W_\gamma^*$ is smooth and uniformly bounded on compact $X$. Thus
\begin{equation}\label{eq:rank-two-extension-first-order}
 \mathrm dr=\alpha-r\beta,\qquad
 \alpha=L_\gamma(\mathrm d\omega),\quad
 \beta=L_\gamma(\mathrm d\gamma),
\end{equation}
with smooth one-forms on $X$. In particular $|\mathrm d\log(1+r)|\leq C$ on $U$.

The analytic set permits short avoiding paths. In a coordinate ball, stratify it into countably many smooth strata of real dimension at most $2n-2$. For a fixed outside point, endpoints of segments meeting one stratum lie in a cone of dimension at most $2n-1$, a Lebesgue-null set. For two nearby outside points choose a generic intermediate point in a ball of radius comparable to their distance. Both connecting segments avoid $Z$ and their combined length is bounded by a fixed multiple of that distance. A finite connected coordinate cover, with outside points chosen in its overlaps, also gives a global uniform bound for connecting lengths in $U$.

The logarithmic gradient bound now bounds $r$ on $U$, and \eqref{eq:rank-two-extension-first-order} bounds its derivative. The local short paths give a unique locally Lipschitz extension across $Z$. Since $Z$ is Lebesgue-null, its weak derivative satisfies the same equation on $X$. The right side is continuous, so $r$ is $C^1$; iteration makes it smooth.

Set $S=r\gamma$. It is smooth semipositive and $\mathrm dS=\mathrm d\omega$ everywhere by density. The current $R=A-S$ is closed and supported on $Z$. It is positive: it equals the positive coefficient-measure part of $A$ on $Z$, since $A=S$ off $Z$ and $S$ has no mass there. The analytic support theorem recalled in \cref{sec:psh-divisorial-estimates} expresses $R$ as a nonnegative sum of divisorial integration currents; components of codimension at least two carry no such current. Locally $S-\omega$ is smooth and closed and has a smooth $\ddc$-potential. The divisorial coefficients are therefore the Lelong numbers of $\Phi_f$, and all vanish. Consequently $\ddc\Phi_f=S-\omega$ is smooth. Taking a Euclidean trace and applying elliptic regularity gives a smooth potential; its canonical representative agrees with it everywhere.
\end{proof}

The smooth root and its rank bound are explicit hypotheses. A smooth semipositive maximizing $(n-1,n-1)$-form need not have such a root across its degeneracy set. On a surface the rank bound already forces $\gamma>0$ everywhere.

\subsubsection{Two attained instances}

\begin{example}\label{ex:hopf-surface-invariant-attainment}\index{Gauduchon duality!Hopf surface attainment}
On the Hopf surface $X=(\mathbb C^2\setminus\{0\})/\langle z\mapsto2z\rangle$, the background $\omega=|z|^{-2}\ddc|z|^2$ has invariant bounded mass. It admits nonconstant smooth threshold solutions with actual smooth positive maximizing tests.
\end{example}
\begin{proof}
The form descends under $z\mapsto2z$ and is positive. For $s=|z|^2$ and $F(s)=s^{-1}$, the identity $sF''+2F'=0$ gives $\ddc\omega=\ddc(s^{-1})\wedge\ddc s=0$. It is nonclosed. Choose a nonconstant smooth real $q$ and small nonzero $\varepsilon$ such that $A=\omega+\ddc(\varepsilon q)>0$. Put $V=\int_X\omega^2$, $\mathrm{d}V=\omega^2/V$, $f=A^2/\mathrm{d}V$ and $u=\varepsilon q-\sup_X\varepsilon q$. Invariant mass gives $\int_Xf\,\mathrm{d}V=V$, hence $b_\omega(f)=1$ by \cref{thm:invariant-mass-threshold-envelopes}.

The test $T=A/V$ satisfies $\ddc T=0$, $\int_X\omega\wedge T=1$ and $D_{\sqrt f}(T)=1$. For every normalized smooth positive Gauduchon test $Q$, arithmetic--geometric mean and Stokes give $D_{\sqrt f}(Q)\leq\int_XA\wedge Q=1$. Thus $T$ attains the supremum. The attained uniqueness theorem identifies $\Phi_f=u$. This is a bounded nonconstant instance.
\end{proof}

\begin{example}\label{ex:quadric-rank-two-dual-attainment}\index{Gauduchon duality!quadric ramification}
For $n\geq3$, a smooth quadric double cover of $\mathbb P^n$ has an actual semipositive maximizing test whose smooth root has rank at least two everywhere and degenerates along the ramification divisor.
\end{example}
\begin{proof}
Let
\[
 X=\{[w:z_0:\cdots:z_n]:w^2=\textstyle\sum_{i=0}^nz_i^2\}
       \subset\mathbb P^{n+1},\qquad F([w:z])=[z].
\]
The quadratic is nonsingular and the map is finite of degree two. Its smooth ramification divisor is $Z=\{w=0\}$. Local branch coordinates show that $\mathrm dF$ has rank $n-1$ on $Z$ and $n$ elsewhere. Let $\omega$ be the restricted normalized Fubini--Study form and $A=F^*\omega_{\mathrm{FS}}$. Then
\[
 u=\log\frac{\sum_i|z_i|^2}{|w|^2+\sum_i|z_i|^2}
\]
is global and smooth, $-\log2\leq u\leq0$, $\sup_Xu=0$, and $A=\omega+\ddc u$. Branched change of variables gives $\int_XA^n=2$, and Stokes gives $V=\int_X\omega^n=2$.

Set $\mathrm{d}V=\omega^n/V$, $f=A^n/\mathrm{d}V$, $T=A^{n-1}/V$ and $\gamma=A/V^{1/(n-1)}$. The closed background gives $b_\omega(f)=1$. Stokes gives $\ddc T=0$ and $\int_X\omega\wedge T=1$; the determinant weights give $D_{f^{1/n}}(T)=1$. Arithmetic--geometric mean, valid also for semipositive $A$, bounds every competing normalized positive test by $\int_XA\wedge Q=1$. Thus $T$ is an actual maximizer and $\Phi_f=u$ by attained uniqueness. Its root has rank $n-1\geq2$ on $Z$ and is positive elsewhere, as required.
\end{proof}

The invariant-mass and sublinear conditions answer restricted instances of the envelope question. General moment or entropy bounds have not been shown to supply either condition or the smooth-root rank hypothesis.


\Needspace{14\baselineskip}
\begin{corollary}\label{cor:semipositive-dual-lp-limit}
Under the hypotheses of \cref{thm:semipositive-dual-exact-equation}, suppose in addition that $f\in L^p(\mathrm{d}V)$ for some $p>1$. Then the normalized threshold solution $u_0$ is continuous, and the normalized greatest exponential solutions converge uniformly to $u_0$ as the exponent tends to zero.
\end{corollary}
\begin{proof}
By \cref{prop:individual-normalizer-lp}, the constant $c$ for which $\mathrm{MA}_\omega(v)=cf\,\mathrm{d}V$ has a continuous solution is $b_\omega(f)$. The continuous normalized solution is therefore a threshold solution, and it equals $u_0$ by the uniqueness in \cref{thm:semipositive-dual-exact-equation}. For $L^p$ data the greatest exponential solution coincides with the bounded solution, as explained after \cref{thm:individual-exponential-subsolution}. By \cref{prop:symmetric-equation-clusters}, applied with the trivial group, the normalized family is uniformly precompact, and every cluster limit solves the threshold equation. Uniqueness identifies every cluster limit with $u_0$, which proves uniform convergence of the whole family.
\end{proof}


\subsection{A full-mass point pole without a smooth maximizing test}
\label{subsec:point-pole-dual-nonattainment}

The smaller-background criteria do not require $L^p$ integrability. The following closed-background example has genuine obstacle envelopes at every scale and an unbounded threshold solution. Its dual supremum is approached by smooth tests but is not attained by any smooth semipositive test.

\begin{proposition}\label{prop:point-pole-full-mass-nonattainment}\index{Gauduchon duality!point-pole nonattainment}\index{Monge--Amp\`ere density!point-pole integrability}
Let $X=\mathbb P^n$, $n\geq2$, and let $\omega$ be the normalized Fubini--Study form, with $\int_X\omega^n=1$ and $\mathrm{d}V=\omega^n$. For $0<\alpha<1$ put
\[
 \psi=\log\frac{\sum_{i=1}^n|z_i|^2}{\sum_{i=0}^n|z_i|^2},
 \qquad u=1-(1-\psi)^\alpha,
 \qquad o=[1:0:\cdots:0].
\]
The function $u$ is a genuine normalized $\omega$-plurisubharmonic potential with sole pole $o$ and extended mass one. Every $P_{\delta\omega}(u)$ is genuine, for $0<\delta<1$. Its measure is $f\,\mathrm{d}V$ with $f\in L^1(\mathrm{d}V)$ and $f\notin L^p(\mathrm{d}V)$ for every $p>1$.

One has $b_\omega(f)=1$ and $\rho_\omega(f^{1/n})=1$. Thus $u$ is an actual threshold solution. No smooth semipositive $T\in\mathscr G_\omega$ attains $D_{f^{1/n}}(T)=1$.
\end{proposition}
\begin{proof}
The numerator vanishes only at $o$. Locally $\omega+\ddc\psi=\ddc\log\sum_{i=1}^n|z_i|^2\geq0$. The function $\chi(t)=1-(1-t)^\alpha$ is increasing and convex on $t\leq0$, with $0\leq\chi'\leq\alpha$. Apply its chain inequality to bounded truncations of $\psi$ and take the decreasing limit. It gives a genuine $u\in\PSH(X,\omega)$, with
\[
 A=\omega+\ddc u\geq(1-\alpha)\omega.
\]
The function is smooth off $o$ and has supremum zero on $\{z_0=0\}$. Since $u=-(1-\psi)^\alpha+1$, \cref{prop:sublinear-threshold-envelope} gives all the smaller-background envelopes.

To compute its mass, use the chart $z_0=1$ and write $s=\sum_{i=1}^n|z_i|^2$. A local potential for $A$ and its radial flux are
\begin{equation}\label{eq:point-pole-radial-flux}
 \begin{aligned}
 F(s)&=\log(1+s)+1-\left(1-\log\frac{s}{1+s}\right)^\alpha,\\
 q(s)&=sF'(s)=\frac{s+\alpha(1-\log(s/(1+s)))^{\alpha-1}}{1+s}.
 \end{aligned}
\end{equation}
The tangential Hessian eigenvalue is $q/s$ and the radial one is $q'$. Polar integration gives the radial top-degree measure $nq^{n-1}q'\,\mathrm{d}s$. Its angular normalization follows from the Fubini--Study case $q=s/(1+s)$, whose radial density is $ns^{n-1}(1+s)^{-n-1}\,\mathrm{d}s$ and whose total mass is one. Consequently
\[
 \int_{\{a<|z|^2<b\}}A^n=q(b)^n-q(a)^n.
\]
Here $q(s)\to0$ as $s\downarrow0$ and $q(s)\to1$ as $s\to\infty$. On each compact annulus the extended measure equals the smooth top power. It gives no mass to $o$ or to the infinity hyperplane, since both are pluripolar. Exhausting the affine chart by annuli therefore gives extended mass exactly one. This is a mass calculation, independent of the vanishing of Lelong numbers.

Near $o$ write $L=1-\log s$. Direct differentiation of \eqref{eq:point-pole-radial-flux} gives
\[
 \begin{gathered}
 \frac qs\sim\alpha s^{-1}L^{\alpha-1},\qquad
 q'\sim\alpha(1-\alpha)s^{-1}L^{\alpha-2},\\
 f(s)\asymp s^{-n}L^{-1-n(1-\alpha)}.
 \end{gathered}
\]
The smooth radial volume is comparable to $s^{n-1}\,\mathrm{d}s$. Its $L^1$ tail is thus comparable to $\int^\infty L^{-1-n(1-\alpha)}\,\mathrm{d}L$, which is finite. For every $p>1$, the corresponding tail has an exponentially growing factor $\exp(n(p-1)L)$ and diverges. The measure is smooth off $o$, so it is globally $f\,\mathrm{d}V$. Since its total mass is one, \cref{thm:invariant-mass-threshold-envelopes} gives $b_\omega(f)=1$, and \cref{thm:integrable-gauduchon-formula} gives $\rho_\omega(f^{1/n})=1$.

Suppose a smooth semipositive normalized test $T$ attained this value. The positive current $A$ satisfies $\int_XA\wedge T=1$ by distributional Stokes. On $X\setminus\{o\}$, matrix arithmetic--geometric mean gives
\[
 A\wedge T\geq f^{1/n}k_T\,\mathrm{d}V.
\]
The right side has total mass one by attainment. Thus the nonnegative gap vanishes off $o$. There $A$ is smooth strictly positive, so equality gives $T=tA^{n-1}$, including at points where $T=0$. This equality case follows by applying arithmetic--geometric mean to the positive semidefinite coefficient matrix conjugated by $A^{1/2}$: all its eigenvalues agree, possibly at zero. The function $t=(T\wedge A)/A^n$ is smooth and nonnegative on the punctured manifold.

Since $A$ is closed, $\ddc T=0$ becomes $\ddc t\wedge A^{n-1}=0$, a locally strictly elliptic scalar equation there. All eigenvalues of $A$ tend to infinity at $o$, so, as $n\geq2$, all complementary-product eigenvalues of $A^{n-1}$ do also. Boundedness of the smooth coefficient matrix of $T$ forces $t\to0$ at $o$.

Extend $t(o)=0$. This continuous nonnegative function on compact $X$ is nonzero by the normalization of $T$. It has a positive maximum away from $o$. The local strong maximum principle on connected $X\setminus\{o\}$ forces it to be constant there, contrary to its zero limit at $o$. This excludes every smooth semipositive maximizing test; no assumption about an analytic degeneracy set was needed.
\end{proof}

\begin{corollary}\label{cor:point-pole-entropy-energy-threshold}\index{Entropy!point-pole threshold}
For the density and potential of \cref{prop:point-pole-full-mass-nonattainment},
\[
 \int_X f\log f\,\mathrm{d}V<+\infty
 \quad\Longleftrightarrow\quad \alpha<1-\frac1n.
\]
For every real $r>0$,
\[
 \int_X(-u)^r f\,\mathrm{d}V<+\infty
 \quad\Longleftrightarrow\quad \alpha r<n(1-\alpha).
\]
In particular, finite entropy of this density can coexist with an unbounded full-mass threshold solution and absence of a smooth semipositive maximizing test.
\end{corollary}
\begin{proof}
All possible divergence is near $o$. The preceding estimates give $f\,\mathrm{d}V\asymp L^{-1-n(1-\alpha)}\,\mathrm{d}L$, $\log f\sim nL$, and $-u\asymp L^\alpha$. The entropy tail converges exactly when $n(1-\alpha)>1$; the weighted tail converges exactly when $\alpha r<n(1-\alpha)$. At equality each tail is comparable to $\int^\infty L^{-1}\,\mathrm{d}L$ and diverges. The negative part of $f\log f$ is boundedly integrable on the probability volume, so the entropy criterion is unambiguous.
\end{proof}

These are entropy and weighted integrals of the single equation density. The original S16 question concerns uniform moment or entropy conditions for determinant-weighted probabilities of a maximizing sequence under smooth semipositive attainment. The present example lies outside that attained setting. The proposition concerns the explicit threshold solution $u$; its identification with the greatest threshold envelope would require a separate uniqueness argument.


\Needspace{19\baselineskip}
\subsection{Quantitative convergence of the currents}

The bound for the constants $c_\lambda$ in \cref{prop:semipositive-dual-exponential-limit} also bounds the trace error linearly in $\lambda$. Away from the degeneracy set, the deficit in the arithmetic--geometric mean inequality converts this bound into a rate of convergence for all coefficients of the currents. We keep the notation $u_0$, $v_\lambda$, $c_\lambda$ of that proposition. The total variation of a current of order zero is measured with respect to a fixed smooth Hermitian metric; on a relatively compact chart it is equivalent to the sum of the total variations of its coefficients.

\begin{proposition}\label{prop:semipositive-dual-current-rate}
Under the hypotheses and notation of \cref{prop:semipositive-dual-exponential-limit}, put
\[
 A_\lambda=\omega+\ddc v_\lambda,\qquad
 A_0=\omega+\ddc u_0.
\]
There is a constant $C$ such that
\[
 \|A_\lambda\wedge T-A_0\wedge T\|_{\mathrm{TV}}
       \leq C\lambda\qquad(\lambda>0).
\]
For each compact set $K\Subset X\setminus Z$, there is a constant $C_K$ such that
\[
 \|A_\lambda-A_0\|_{\mathrm{TV}(K)}
       \leq C_K\sqrt\lambda\qquad(0<\lambda\leq1).
\]
\end{proposition}
\begin{proof}
Put $p=f^{1/n}$, use the probability measure $\nu=A_0\wedge T=b^{1/n}p k_T\,\mathrm{d}V$ from the proof of \cref{prop:semipositive-dual-exponential-limit}, where the equality is \eqref{eq:semipositive-dual-trace-equality}, and set
\[
 q_\lambda=(c_\lambda/b)^{1/n}\mathrm{e}^{\lambda v_\lambda/n},
 \qquad R_\lambda=A_\lambda\wedge T-q_\lambda\nu\geq0.
\]
Here $R_\lambda\geq0$ by the trace inequality used in that proof, which also gives $c_\lambda\geq b$ and $\int_X(-v_\lambda)\,\mathrm{d}\nu\leq C_0$ for a constant $C_0$. Write $s_+=\max(s,0)$ for the positive part of a real number $s$. Since $q_\lambda\geq\mathrm{e}^{\lambda v_\lambda/n}$ and $1-\mathrm{e}^{x}\leq-x$, we obtain
\[
 (1-q_\lambda)_+\leq1-\mathrm{e}^{\lambda v_\lambda/n},
 \qquad
 \int_X(1-q_\lambda)_+\,\mathrm{d}\nu\leq C_0\lambda/n.
\]
Since $A_\lambda\wedge T$ and $\nu$ both have mass one,
\[
 \|A_\lambda\wedge T-\nu\|_{\mathrm{TV}}
 \leq R_\lambda(X)+\int_X|q_\lambda-1|\,\mathrm{d}\nu
 =2\int_X(1-q_\lambda)_+\,\mathrm{d}\nu.
\]
This proves the first estimate, with $C=2C_0/n$, and also gives
$R_\lambda(X)\leq C_0\lambda/n$ and
$\int_X|q_\lambda-1|\,\mathrm{d}\nu\leq2C_0\lambda/n$.

On $U=X\setminus Z$, write $T=\gamma^{n-1}$ and $A_0=a\gamma$ as in the proof of \cref{thm:semipositive-dual-exact-equation}, and put $\alpha_\lambda=q_\lambda a$. Decompose the positive current $A_\lambda$ into its absolutely continuous part, with matrix density $B_\lambda$ relative to $\gamma$, and its positive singular part. Since $\mathrm{MA}^{\mathrm{ext}}_\omega(v_\lambda)=c_\lambda\mathrm{e}^{\lambda v_\lambda}f\,\mathrm{d}V$, the trace inequalities for all Hermitian forms imply, as in the proof of \cref{cor:smooth-dual-exact-equation},
\[
 (\det B_\lambda)^{1/n}\geq\alpha_\lambda
 \quad\text{almost everywhere on }U.
\]
\Needspace{9\baselineskip}
This is a statement about determinants of matrices; no nonlinear product of coefficient measures is involved. Put $s_\lambda=\operatorname{tr}B_\lambda/n$ and $\delta_\lambda=s_\lambda-\alpha_\lambda\geq0$. On $U$, the absolutely continuous part of the trace error $R_\lambda$ is $\delta_\lambda\gamma^n$, and its singular part is the trace measure of the singular part of $A_\lambda$. Thus
\[
 \int_U s_\lambda\gamma^n\leq1,\qquad
 \int_U\delta_\lambda\gamma^n\leq C_0\lambda/n,
\]
and the singular trace mass on $U$ has the same upper bound.

For nonnegative numbers $\ell_1,\ldots,\ell_n$ with geometric mean at least $\alpha\geq0$ and arithmetic mean $s$, the arithmetic--geometric mean inequality and the Cauchy--Schwarz inequality give
\[
 \sum_i(\sqrt{\ell_i}-\sqrt\alpha)^2\leq n(s-\alpha),
 \qquad
 \sum_i|\ell_i-\alpha|\leq2n\sqrt{s(s-\alpha)}.
\]
Indeed $\sum_i\sqrt{\ell_i}\geq n\sqrt\alpha$, while
$\sum_i(\sqrt{\ell_i}+\sqrt\alpha)^2\leq4ns$.
Applying this to the eigenvalues of $B_\lambda$ relative to $\gamma$ and integrating, with the Cauchy--Schwarz inequality once more, we see that the absolutely continuous part of $A_\lambda-\alpha_\lambda\gamma$ has integrated trace norm at most $2n\sqrt{C_0\lambda/n}$ on $U$. Its singular part has norm at most $C_0\lambda$ in the same metric. Finally
\[
 \int_U|\alpha_\lambda-a|\gamma^n
   =\int_U|q_\lambda-1|\,\mathrm{d}\nu
   \leq2C_0\lambda/n.
\]
Since $\gamma$ is comparable to the fixed metric on the compact set $K\subset U$, these estimates and $\lambda\leq\sqrt\lambda$ for $0<\lambda\leq1$ give the second conclusion.
\end{proof}

The coefficient estimate is local on $X\setminus Z$. Convergence in capacity across $Z$ requires the additional envelope condition of \cref{cor:semipositive-dual-capacity-limit}.


\Needspace{16\baselineskip}
\subsection{The classes of the divisorial components}

Uniqueness constrains the divisors along which the threshold solution can have positive Lelong numbers. A linear relation between their Bott--Chern classes would produce a global potential that changes the coefficients of these divisors while preserving the equation off the degeneracy set, contradicting uniqueness. We write $\{D\}\in H^{1,1}_{\mathrm{BC}}(X,\mathbb R)$ for the Bott--Chern class of the current of integration $[D]$ along a prime divisor $D$, and $\nu(u,D)$ for the generic Lelong number of $u$ along $D$.

\begin{corollary}\label{cor:semipositive-dual-divisor-independence}
Under the hypotheses of \cref{thm:semipositive-dual-exact-equation}, let $u$ be a threshold solution and let $D_1,\ldots,D_k$ be the prime divisors with $\nu(u,D_j)>0$. Then the classes $\{D_j\}$ are linearly independent in $H^{1,1}_{\mathrm{BC}}(X,\mathbb R)$. In particular their number is at most $\dim_{\mathbb R}H^{1,1}_{\mathrm{BC}}(X,\mathbb R)$.
\end{corollary}
\begin{proof}
Since $u$ is smooth off $Z$, these divisors are among the finitely many components of $Z$ of codimension one. Put $a_j=\nu(u,D_j)>0$. By \cref{lem:nonclosed-divisorial-removal},
\[
 R=\omega+\ddc u-\sum_j a_j[D_j]\geq0.
\]
Suppose that real numbers $c_1,\ldots,c_k$ satisfy $\sum_j c_j\{D_j\}=0$. Fix a Gauduchon metric $\eta$ and apply \cref{lem:constant-trace-bc-representative} to the currents of integration $[D_j]$, with the integrable potentials $v_\eta([D_j])$ given there. Since the smooth representative in that lemma depends linearly on the Bott--Chern class, we obtain
\[
 h=\sum_j c_j v_\eta([D_j])\in L^1(X),\qquad
 \ddc h=\sum_j c_j[D_j].
\]
Choose $t\ne0$ sufficiently small that $a_j+tc_j\geq0$ for every $j$. The distribution $u_t=u+th$ satisfies
\[
 \omega+\ddc u_t=R+\sum_j(a_j+tc_j)[D_j]\geq0,
\]
so it is represented by a function $u_t\in\PSH(X,\omega)$. On $X\setminus Z$, the function $h$ is smooth and pluriharmonic, so $u_t$ has the same Monge--Amp\`ere measure as $u$ there. Since neither extended measure charges the pluripolar set $Z$, $u_t$ is another threshold solution. Uniqueness makes $th=u_t-u$ constant. Hence $\sum_jc_j[D_j]=0$, and testing near a generic point of each of the distinct divisors $D_j$ gives $c_j=0$.
\end{proof}


\Needspace{20\baselineskip}
\subsection{A branched projective example}

In this subsection $X=\mathbb P^n$, with integer $n\geq2$ and homogeneous coordinates $[Z_0:\cdots:Z_n]$. Let $\omega$ be the Fubini--Study form normalized by $\int_X\omega^n=1$, and put $\mathrm{d}V=\omega^n$. The current set $\mathscr G_\omega$, determinant density $k_T$, functionals $D_p$, $\rho_\omega$ and threshold $b_\omega$ have the meanings of \cref{sec:integrable-gauduchon-currents}.

A smooth semipositive maximizer need not be strictly positive, even when the threshold solution is smooth. The following example also shows why smoothness of the density root was asserted only off the degeneracy set.

\begin{example}\label{ex:semipositive-projective-maximizer}
For every $n\geq2$, there is a smooth nonnegative density $f$ on $\mathbb P^n$, vanishing precisely on the coordinate hyperplanes, such that $b_\omega(f)=1$ and $D_{f^{1/n}}$ has a smooth semipositive maximizer strictly positive off those hyperplanes. The threshold solution is smooth, but $f^{1/n}$ is not differentiable at a general point of any coordinate hyperplane. There is no smooth strictly positive maximizing form. The normalized greatest exponential solutions converge uniformly to the normalized threshold solution.
\end{example}
\begin{proof}
Define the holomorphic map $F:\mathbb P^n\to\mathbb P^n$ and the semipositive form $A$ by
\[
 F[Z_0:\cdots:Z_n]=[Z_0^2:\cdots:Z_n^2],\qquad
 A=\tfrac12F^*\omega.
\]
The homogeneous formula
\[
 u=\tfrac12\log\!\left(\sum_{i=0}^n|Z_i|^4\right)
       -\log\!\left(\sum_{i=0}^n|Z_i|^2\right)
\]
defines a smooth function on projective space, and $A=\omega+\ddc u$. This form is semipositive and is strictly positive exactly where every homogeneous coordinate is nonzero. Define
\[
 f=A^n/\omega^n,\qquad T=A^{n-1}.
\]
Closed Stokes gives $\int f\omega^n=\int\omega\wedge T=1$. Thus $T\in\mathscr G_\omega$. Every genuine $\omega$-plurisubharmonic potential has extended mass at most one: each bounded truncation has mass one, and restriction away from its truncation level can only decrease that mass. Consequently an inequality $\mathrm{MA}^{\mathrm{ext}}_\omega(v)\geq c f\omega^n$ forces $c\leq1$. The smooth solution $u$ gives the reverse bound, so $b_\omega(f)=1$.

The determinant formula for smooth semipositive forms gives
\[
 k_T=f^{(n-1)/n},\qquad
 D_{f^{1/n}}(T)=\int_{\mathbb P^n}f\omega^n=1.
\]
By \eqref{eq:integrable-current-supremum} and $\rho_\omega(f^{1/n})=b_\omega(f)^{-1/n}=1$, this is a maximizer. Its strict positivity away from the coordinate hyperplanes follows from that of $A$; on a hyperplane $A$ has deficient rank, and so does its cofactor matrix.

For the density root, use the chart $Z_0=1$ with coordinates $z_1,\ldots,z_n$. Pulling back the Fubini--Study volume by $(z_1,\ldots,z_n)\mapsto(z_1^2,\ldots,z_n^2)$ and including the factor $2^{-n}$ gives
\[
 f(z)=2^n\left(\prod_{i=1}^n|z_i|^2\right)
       \frac{(1+\sum_i|z_i|^2)^{n+1}}
            {(1+\sum_i|z_i|^4)^{n+1}}.
\]
Near a point of $\{z_1=0\}$ with all the other coordinates nonzero, its $n$th root is $|z_1|^{2/n}$ times a smooth strictly positive function. Restriction to a real line transverse to that hyperplane proves nondifferentiability, including the absolute-value case $n=2$. Coordinate symmetry gives the same conclusion on every component. The absence of a smooth strictly positive maximizer follows from \cref{cor:kahler-smooth-dual-attainment}, since $f$ has zeros. Finally $f$ is bounded, so \cref{cor:semipositive-dual-lp-limit} gives the uniform exponential limit.
\end{proof}


\subsection{A moment condition on maximizing tests}

For a general Hermitian background, a uniform bound on the moments of normalized potentials with respect to the probabilities of one smooth maximizing sequence can replace the regularity of the maximizing current. Together with the envelope condition in a smaller background, it shows that the cluster limits of the solutions of the exponential equations are exact threshold solutions.

For a nonzero $0\leq f\in L^1(\mathrm{d}V)$ with $b_\omega(f)>0$, let $\mathcal S_f$ be the class of threshold subsolutions with supremum zero and $\Phi_f$ its greatest member, from \cref{prop:greatest-normalized-threshold}. As in \cref{prop:semipositive-dual-exponential-limit}, $u_\lambda$ denotes the greatest solution of the exponential equation $\mathrm{MA}^{\mathrm{ext}}_\omega(u)=\mathrm{e}^{\lambda u}f\,\mathrm{d}V$, and $v_\lambda=u_\lambda-\sup_Xu_\lambda$, $c_\lambda=\mathrm{e}^{\lambda\sup_Xu_\lambda}$.

\Needspace{12\baselineskip}
\begin{theorem}\label{thm:gauduchon-moment-threshold}\index{Moment constraint}
Let $0\leq f\in L^1(\mathrm{d}V)$ be nonzero with $b=b_\omega(f)>0$, and put $p=f^{1/n}$ and $R=\rho_\omega(p)=b^{-1/n}$. Suppose that there are smooth strictly positive $T_j\in\mathscr G_\omega$, with associated probability measures $\nu_j$, such that
\[
 d_j=D_p(T_j)\longrightarrow R,\qquad
 \mathrm{d}\nu_j=d_j^{-1}p k_{T_j}\,\mathrm{d}V,
\]
and a constant $C$ satisfying
\begin{equation}\label{eq:maximizing-gauduchon-moments}
 \int_X(-v)\,\mathrm{d}\nu_j\leq C
 \quad\text{for all }j\text{ and all }v\in\PSH(X,\omega)
                  \text{ with }\sup_Xv=0.
\end{equation}
Then, for every $\lambda>0$,
\[
 b\leq c_\lambda\leq b\mathrm{e}^{C\lambda},\qquad
 \Phi_f-C\leq v_\lambda\leq0.
\]
Suppose in addition that $P_{\delta\omega}(\Phi_f)$ is not identically $-\infty$ for some $0<\delta<1$. Whenever $\lambda_k\downarrow0$ and $v_{\lambda_k}\to v$ in $L^1(X)$, the convergence holds in capacity and
\[
 \mathrm{MA}^{\mathrm{ext}}_\omega(v)=bf\,\mathrm{d}V,
 \qquad \Phi_f-C\leq v\leq\Phi_f.
\]
In particular an exact threshold solution exists.
\end{theorem}
\begin{proof}
Since $\mathrm{MA}^{\mathrm{ext}}_\omega(v_\lambda)=c_\lambda\mathrm{e}^{\lambda v_\lambda}f\,\mathrm{d}V$, the trace inequality of \cref{lem:unbounded-gauduchon-trace-density} and Stokes' theorem give, for each $j$ and $\lambda>0$,
\[
 1\geq c_\lambda^{1/n}d_j
              \int_X\mathrm{e}^{\lambda v_\lambda/n}\,\mathrm{d}\nu_j
   \geq c_\lambda^{1/n}d_j\mathrm{e}^{-C\lambda/n}.
\]
\Needspace{9\baselineskip}
The second inequality follows from Jensen's inequality and \eqref{eq:maximizing-gauduchon-moments}. Letting $j\to\infty$ and using $R=b^{-1/n}$, we obtain the upper bound for $c_\lambda$. The lower bound and the comparison with every normalized threshold subsolution were proved in \cref{cor:individual-positive-threshold-exponential}; applying that comparison to $\Phi_f$, we obtain
\[
 v_\lambda\geq\Phi_f-\lambda^{-1}\log(c_\lambda/b)
                \geq\Phi_f-C.
\]

Now let $v_{\lambda_k}\to v$ in $L^1$. By \cref{prop:greatest-normalized-threshold}, the limit $v$ lies in $\mathcal S_f$, hence below $\Phi_f$. The lower bound passes to the limit, so $\Phi_f-C\leq v\leq\Phi_f$; in particular $v\sim\Phi_f$ in the sense of \cref{def:singularity-order}. For all $k$ with $\lambda_k\leq1$,
\[
 \mathrm{MA}^{\mathrm{ext}}_\omega(v_{\lambda_k})
   =c_{\lambda_k}\mathrm{e}^{\lambda_k v_{\lambda_k}}f\,\mathrm{d}V
   \leq b\mathrm{e}^{C}f\,\mathrm{d}V.
\]
The functions $|\mathrm{e}^{v_{\lambda_k}}-\mathrm{e}^v|$ are uniformly bounded and tend to zero in $\mathrm{d}V$-measure, so their integrals against the fixed $L^1$ density $f$ tend to zero. Hence the error term in \eqref{eq:common-barrier-self-error} tends to zero, as required by \cref{prop:lower-barrier-criterion} with $\varphi=\Phi_f$. The additional hypothesis $v\preceq\Phi_f$ of that proposition holds since $v\leq\Phi_f$, and the condition on $P_{\delta\omega}(\Phi_f)$ is assumed. That proposition therefore gives convergence in capacity.

The continuity of damped measures under convergence in capacity, \cref{thm:damped-continuity}, now identifies the equation without a mass-conservation argument. Write $g_k=c_{\lambda_k}\mathrm{e}^{\lambda_k v_{\lambda_k}}f$. We know $g_k\to bf$ in $L^1$ and $0\leq g_k\leq b\mathrm{e}^Cf$ for large $k$. Hence
\[
 \mathrm{e}^{v_{\lambda_k}}\mathrm{MA}^{\mathrm{ext}}_\omega(v_{\lambda_k})
   =\mathrm{e}^{v_{\lambda_k}}g_k\,\mathrm{d}V
   \longrightarrow\mathrm{e}^v bf\,\mathrm{d}V
\]
in total variation, while \cref{thm:damped-continuity} gives the weak limit
$\mathrm{e}^v\mathrm{MA}^{\mathrm{ext}}_\omega(v)$. Dividing by $\mathrm{e}^v$ off the pole set of $v$, which neither measure charges, gives the equation and the finiteness of the mass. Finally, compactness of normalized $\omega$-plurisubharmonic functions provides an $L^1$ cluster limit along any sequence $\lambda_k\downarrow0$, which proves existence.
\end{proof}

The moment condition holds if the densities of the $\nu_j$ are uniformly bounded in $L^q(\mathrm{d}V)$ for some $q>1$, by H\"older's inequality and the uniform bounds for the moments of normalized $\omega$-plurisubharmonic functions. More generally, a uniform entropy bound suffices. Write $a_j=\mathrm{d}\nu_j/\mathrm{d}V$ and suppose $\int_Xa_j\log a_j\,\mathrm{d}V\leq H$ for all $j$. (By \cref{prop:hermitian-dual-entropy-bounds}, a uniform entropy bound is equivalent to a uniform bound for the dual entropy $H_0(\nu_j)$ with zero reference defined there.) Choose $\alpha>0$ and $A$ from \cref{lem:uniform-exponential}, so that $\int_X\mathrm{e}^{-\alpha v}\,\mathrm{d}V\leq A$ for every $v\in\PSH(X,\omega)$ with $\sup_Xv=0$. The elementary entropy inequality, applied first to the truncations $v_T=\max(v,-T)$ at levels $T>0$, gives
\[
 \alpha\int_X(-v_T)a_j\,\mathrm{d}V
 \leq \int_Xa_j\log a_j\,\mathrm{d}V
             +\log\int_X\mathrm{e}^{-\alpha v_T}\,\mathrm{d}V
 \leq H+\log A.
\]
To prove this inequality, apply $t\log t-t+1\geq0$ to $t=a_j/q$, integrate against $q\,\mathrm{d}V$, where $q=\mathrm{e}^{-\alpha v_T}/\int_X\mathrm{e}^{-\alpha v_T}\,\mathrm{d}V$ is a probability density, and rearrange. Monotone convergence as $T\to\infty$ proves \eqref{eq:maximizing-gauduchon-moments} with $C=(H+\log A)/\alpha$.

The same entropy bound also gives attainment of the dual supremum along this sequence. Indeed, $t\log t\geq-\mathrm{e}^{-1}$ and, for $K>1$,
\[
 \int_{\{a_j>K\}}a_j\,\mathrm{d}V
       \leq\frac{H+\mathrm{e}^{-1}}{\log K}.
\]
This proves uniform integrability of the $a_j$, and since $d_j\to R>0$, also of the densities $p k_{T_j}=d_ja_j$. By \cref{prop:selected-gauduchon-attainment}, every weak cluster limit $T$ of $(T_j)$ maximizes $D_p$, and along a subsequence converging to $T$ the probabilities $\nu_j$ converge to $R^{-1}p k_T\,\mathrm{d}V$. For each normalized $v$, the function $-v$ is nonnegative and lower semicontinuous, so the lower semicontinuity of its integral under weak convergence of measures passes the moment bound to this limiting probability. These conclusions give no regularity of $T$.

The envelope condition holds if $\Phi_f$ is bounded. In general, it and the moment bound are additional hypotheses; the tail condition on $f^{1/n}$ alone does not provide them in this argument. The theorem identifies the equation satisfied by each cluster limit. Uniqueness, which gave convergence of the whole family in \cref{prop:semipositive-dual-exponential-limit}, is not available here.




\Needspace{11\baselineskip}
\section{Nef limits and nonattainment}
\label{sec:gaud-volume-nef-limit}

We return to arbitrary dimension $n\geq1$. Along a nef perturbation, the infimum over Gauduchon metrics commutes with the decreasing limit of the background. Thus the boundary value is determined by smooth equations even when no metric attains it. For a nef $\theta$ and any Hermitian $h$, the backgrounds $\theta+\epsilon h$ lie in $\mathscr P$ for every $\epsilon>0$, by \cref{sec:affine-model-space}.

\Needspace{12\baselineskip}
\begin{theorem}\label{thm:gaud-volume-nef-limit}\index{Gauduchon variational volume!nef limit}
For every nef background $\theta$ and every Hermitian form $h$,
\[
 \lambda_\mu(\theta)
     =\lim_{\epsilon\downarrow0}\lambda_\mu(\theta+\epsilon h),
 \qquad
 \mathcal V_\mu(\theta)
     =\lim_{\epsilon\downarrow0}c_{\theta+\epsilon h}(\mu).
\]
These limits are decreasing and independent of $h$.
\end{theorem}
\begin{proof}
The quantities decrease by monotonicity. The infima over $\epsilon>0$ and over Gauduchon metrics commute. For a fixed metric the quotient is linear in $\epsilon$, so
\[
 \begin{split}
 \inf_{\epsilon>0}\lambda_\mu(\theta+\epsilon h)
 &=\inf_\gamma\inf_{\epsilon>0}
   \frac{\int_X\theta\wedge\gamma^{n-1}
                +\epsilon\int_Xh\wedge\gamma^{n-1}}
        {D_\mu(\gamma)}\\
 &=\lambda_\mu(\theta).
 \end{split}
\]
The coefficient of $\epsilon$ is positive, so the inner infimum is its limit at zero. Apply the exact formula on $\mathscr P$ and raise to the $n$th power. The resulting expression is the original definition at $\theta$, proving independence of $h$.
\end{proof}

If $X$ is K\"ahler and $\theta$ is closed and nef, choose $h$ K\"ahler. Closed Stokes in every perturbed equation gives
\[
 c_{\theta+\epsilon h}(\mu)=\int_X(\theta+\epsilon h)^n.
\]
It follows that $\mathcal V_\mu(\theta)=\int_X\theta^n$ in this regime.

\begin{proposition}\label{prop:zero-gaud-volume-attainment}
Suppose $\theta$ is pseudoeffective and $\lambda_\mu(\theta)=0$. The infimum is attained if and only if $\theta=\ddc f$ for a smooth real function $f$.
\end{proposition}
\begin{proof}
The denominator $D_\mu(\gamma)$ is strictly positive. Thus a minimizing
metric at value zero has $\int_X\theta\wedge\gamma^{n-1}=0$.
Apply \cref{lem:zero-gauduchon-degree}, including its dimension-one case.
Conversely, Stokes' theorem makes every quotient zero when $\theta$ is
smoothly $\ddc$-exact.
\end{proof}

For the ruled surface $X$ and section $D$ of \cref{lem:ruled-explicit-nef}, the nef class $c_1(\mathcal O_X(D))$ has zero square and a nonzero positive representative $[D]$. Its root volume is zero, but no smooth Gauduchon metric minimizes the defining quotient.

\Needspace{11\baselineskip}
\section{Subtraction with a closed nef core}
\label{sec:nef-boundary-subtraction}

Consider a path of backgrounds that starts in $\mathscr P$ and reaches its boundary at a nef point. If the path contains a closed nef class of positive top intersection, this closed core keeps the root volume positive up to and including the boundary, and adding a semipositive form, closed or not, preserves this lower bound. By \cref{prop:gaud-volume-first-variation}, the total decrease of the root volume along the path is the integral of an explicit ratio formed with the minimizing Gauduchon metrics; we call it the subtraction cost. Two test metrics chosen at the initial point bound this cost from above, with explicit dependence on the density and on the derivatives of the added form.

Throughout this section, $X$ is a compact K\"ahler manifold of dimension $n\geq2$ and $\Omega$ is a fixed K\"ahler form on $X$. Let $\alpha$ be a smooth closed nef real $(1,1)$-form, let $B\geq0$ be a smooth real $(1,1)$-form, and let $t>0$. For real $s$, put
\[
 a_s=\alpha+(t-s)\Omega,\qquad
 \theta_s=a_s+B,\qquad V_s=\int_Xa_s^n.
\]
The class of $a_0$ is K\"ahler. Choose a K\"ahler representative $\kappa_0$ of $[a_0]$, and let $u_0$ be the smooth solution of the Calabi--Yau equation
$\gamma_0=\kappa_0+\ddc u_0>0$, $\gamma_0^n=V_0\mu$.
Define the fixed quantities
\begin{equation}\label{eq:closed-core-fixed-tests}
 \begin{gathered}
 I_B=\int_XB\wedge\kappa_0^{n-1},\qquad
 U_\mu=\frac{V_0+I_B}{D_\mu(\kappa_0)},\\
 \Delta_\mu(B)=V_0^{-(n-1)/n}\int_XB\wedge\gamma_0^{n-1}.
 \end{gathered}
\end{equation}

\begin{theorem}\label{thm:closed-core-subtraction}
Suppose that $V:=V_t>0$ and that there is a smooth compact complex curve $C\subset X$ with $\int_C\alpha=\int_CB=0$. Then $\theta_s$ has a smooth positive representative for $s<t$, $\theta_t$ is nef but has no smooth positive representative, and $\theta_s$ is not nef for $s>t$. Moreover,
\[
 \lambda_\mu(a_s)=V_s^{1/n}\quad(0\leq s\leq t),
 \qquad \lambda_\mu(\theta_t)\geq V^{1/n}>0.
\]
For every choice of minimizing Gauduchon metrics $g_s$ for $\theta_s$, $0\leq s<t$, the ratio $\int_X\Omega\wedge g_s^{n-1}/D_\mu(g_s)$ is a continuous function of $s$, and the subtraction cost satisfies
\begin{equation}\label{eq:closed-core-cost}
 \begin{split}
 0\leq\mathcal C
 &:=\int_0^t\frac{\int_X\Omega\wedge g_s^{n-1}}{D_\mu(g_s)}\,\mathrm{d}s
   =\lambda_\mu(\theta_0)-\lambda_\mu(\theta_t)\\
 &\leq\min\{U_\mu,V_0^{1/n}+\Delta_\mu(B)\}-V^{1/n}<\infty,
 \end{split}
\end{equation}
where the integral is improper at $t$. If $B\leq M\kappa_0$ for a constant $M\geq0$, then
\[
 \Delta_\mu(B)\leq MV_0^{1/n},\qquad
 \mathcal C\leq(1+M)V_0^{1/n}-V^{1/n}.
\]
Finally, there are a quasi-plurisubharmonic function $\rho$ and a constant $\varepsilon>0$ with
$\alpha+\ddc\rho\geq\varepsilon\Omega$ and
$\theta_t+\ddc\rho\geq\varepsilon\Omega$.
\end{theorem}
\begin{proof}
We first locate the boundary of $\mathscr P$ along the path. Let $s<t$. Since $\alpha$ is nef, for every $\delta>0$ there is a smooth function $v_\delta$ with $\alpha+\ddc v_\delta\geq-\delta\Omega$. Taking $\delta<t-s$ makes $a_s+\ddc v_\delta$ K\"ahler, and adding $B\geq0$ gives a positive representative of $\theta_s$. The same approximations show that $\theta_t$ is nef. On the other hand, the hypotheses on $C$ and Stokes' theorem on the compact curve give
\[
 \int_C(\theta_s+\ddc v)=(t-s)\int_C\Omega
\]
for every smooth $v$. At $s=t$ the right side vanishes, which excludes a positive representative of $\theta_t$. For $s>t$, a representative with $\theta_s+\ddc v\geq-\delta\Omega$ and $\delta<s-t$ would contradict this identity, so $\theta_s$ is not nef.

Since every $a_s$ with $s\leq t$ is closed and nef, the closed nef formula following \cref{thm:gaud-volume-nef-limit} gives $\lambda_\mu(a_s)=V_s^{1/n}$, including at $s=t$. The difference $\theta_t-a_t=B$ is semipositive, so monotonicity gives $\lambda_\mu(\theta_t)\geq V^{1/n}$, while any fixed Gauduchon test bounds $\lambda_\mu(\theta_t)$ from above. Since $\theta_s=\theta_t+(t-s)\Omega$, \cref{thm:gaud-volume-nef-limit} with the Hermitian form $\Omega$ identifies $\lambda_\mu(\theta_t)$ with the limit of $\lambda_\mu(\theta_s)$ as $s\uparrow t$.

Write $f(s)=\lambda_\mu(\theta_s)$ for $s<t$. The path has constant velocity $-\Omega$, so \cref{prop:gaud-volume-first-variation} shows that $f$ is smooth and
\[
 H(s):=\frac{\int_X\Omega\wedge g_s^{n-1}}{D_\mu(g_s)}=-f'(s).
\]
In particular, the ratio $H$ is continuous and independent of the choice of minimizing metrics. For $r<t$, the fundamental theorem of calculus gives
\[
 \int_0^rH(s)\,\mathrm{d}s=f(0)-f(r).
\]
Since $H\geq0$, letting $r\uparrow t$ and using the limit just established yields the identity in \eqref{eq:closed-core-cost}.

It remains to bound $f(0)=\lambda_\mu(\theta_0)$ from above by two test metrics. The K\"ahler form $\kappa_0$ is Gauduchon, and since the closed forms $\kappa_0$ and $a_0$ are cohomologous, Stokes' theorem gives $\int_X\theta_0\wedge\kappa_0^{n-1}=V_0+I_B$; hence $\lambda_\mu(\theta_0)\leq U_\mu$. For the second test $\gamma_0$, the equation $\gamma_0^n=V_0\mu$ and the cohomology class of $\gamma_0$ give
\[
 D_\mu(\gamma_0)=V_0^{(n-1)/n},\qquad
 \int_Xa_0\wedge\gamma_0^{n-1}=V_0,
\]
so its quotient for $\theta_0$ is $V_0^{1/n}+\Delta_\mu(B)$. Combined with the lower bound $\lambda_\mu(\theta_t)\geq V^{1/n}$, these two tests prove the upper bound in \eqref{eq:closed-core-cost}. If $B\leq M\kappa_0$, positivity and Stokes' theorem yield
$\int_XB\wedge\gamma_0^{n-1}\leq M\int_X\kappa_0\wedge\gamma_0^{n-1}=MV_0$,
which gives the remaining cost bounds.

Finally, $V=\int_X\alpha^n>0$, so the closed nef class of $\alpha$ is big by Demailly and P{\u a}un's criterion for nef classes on compact K\"ahler manifolds. Hence it contains a K\"ahler current $\alpha+\ddc\rho\geq\varepsilon\Omega$. Adding $B\geq0$ gives the last assertion. Unlike the first current, the second is not closed when $\mathrm dB\neq0$, since its exterior derivative is $\mathrm{d}B$.
\end{proof}

Such data arise from a bimeromorphic holomorphic map $p\colon X\to Y$ onto a compact K\"ahler manifold $Y$ that contracts a smooth curve $C$ to a point. Take $\alpha=p^*\omega_Y$, with $\omega_Y$ K\"ahler, and any smooth $B\geq0$ vanishing on $TC$. Since $p$ has degree one, $\int_X\alpha^n=\int_Y\omega_Y^n>0$, while both curve integrals vanish.

The second initial test separates the closed intersection number $I_B$ from a correction that involves only $\ddc B$.

\begin{proposition}\label{prop:closed-core-differential-cost}
Set $S=\sum_{j=0}^{n-2}\gamma_0^j\wedge\kappa_0^{n-2-j}$. Then
\begin{equation}\label{eq:closed-core-differential-cost}
 \Delta_\mu(B)=V_0^{-(n-1)/n}
       \left(I_B+\int_Xu_0\,\ddc B\wedge S\right).
\end{equation}
If $-C_B\kappa_0^2\leq\ddc B\leq C_B\kappa_0^2$ in the sense of weak positivity of $(2,2)$-forms, then
\[
 \Delta_\mu(B)\leq V_0^{-(n-1)/n}I_B
       +(n-1)C_B\operatorname{osc}_X(u_0)V_0^{1/n}.
\]
Such a finite constant $C_B$ exists for every fixed smooth $B$.
\end{proposition}
\Needspace{8\baselineskip}
\begin{proof}
The factors of $S$ are closed, and
$\gamma_0^{n-1}-\kappa_0^{n-1}=\ddc u_0\wedge S$.
Integration by parts therefore gives \eqref{eq:closed-core-differential-cost}. Since $\int_X\ddc B\wedge S=0$, the right side does not change when a constant is added to $u_0$, so we may assume $\sup_Xu_0=0$. Then $-\operatorname{osc}_X(u_0)\leq u_0\leq0$, and the lower bound for $\ddc B$ together with the positivity of the summands of $S$ bounds the correction term above by
\[
 C_B\operatorname{osc}_X(u_0)\int_X\kappa_0^2\wedge S
       =(n-1)C_B\operatorname{osc}_X(u_0)V_0.
\]
The equality holds because the cohomologous K\"ahler forms $\gamma_0,\kappa_0$ have all mixed top intersections equal to $V_0$. Finally, restrict $\ddc B$ to the complex two-planes of $TX$ and normalize by $\kappa_0$; compactness of the bundle of these planes bounds the restrictions and gives $C_B$.
\end{proof}

The potential $u_0$ solves a single complex Monge--Amp\`ere equation, with density $V_0\mu/\kappa_0^n$ relative to $\kappa_0^n$. Its oscillation therefore carries the dependence on $\mu$, whereas the derivatives $\mathrm{d}\theta_s=\mathrm{d}B$ and $\ddc\theta_s=\ddc B$ do not change along the path.

For the closed core alone the cost is explicit. Put $P_s=\int_X\Omega\wedge a_s^{n-1}$. Since $V_s'=-nP_s$,
\[
 -\frac{\mathrm{d}}{\mathrm{d}s}\lambda_\mu(a_s)
       =\frac{P_s}{V_s^{(n-1)/n}},\qquad
 \int_0^t\frac{P_s}{V_s^{(n-1)/n}}\,\mathrm{d}s=V_0^{1/n}-V^{1/n}.
\]
Moreover $P_s'=-(n-1)\int_X\Omega^2\wedge a_s^{n-2}\leq0$, since all classes on the segment are nef. Integrating $V_s'=-nP_s$ and using $P_s\leq P_0$ gives
\begin{equation}\label{eq:closed-core-nef-morse}
 V\geq V_0-nt\int_X\Omega\wedge a_0^{n-1}.
\end{equation}
This argument does not use $V>0$, so a positive right side guarantees the hypothesis $V>0$ of \cref{thm:closed-core-subtraction}. The following example shows that this sufficient condition is not necessary.

\begin{example}\label{ex:nonclosed-nef-subtraction}
On the blowup of $\mathbb P^2$ at a point there are data as above with $t=1$, $\mathrm{d}B\neq0$,
\[
 \lambda_\mu(\theta_1)\geq\frac32,\qquad
 0\leq\mathcal C\leq\frac{11}{\sqrt8}-\frac32,
\]
and an explicit current $\theta_1+\ddc\rho\geq\frac12\Omega$ at the nef endpoint.
\end{example}
\begin{proof}
Realize the blowup as
$X=\{(x,[u:v])\in\mathbb P^2\times\mathbb P^1:x_1v=x_2u\}$,
with projections $p,q$ and exceptional curve $E=p^{-1}([1:0:0])$.
Let $L=p^*\omega_{\mathrm{FS},\mathbb P^2}$, $Q=q^*\omega_{\mathrm{FS},\mathbb P^1}$ and $\Omega=L+Q$, with the Fubini--Study forms normalized to have integral one on projective lines. The incidence embedding makes $\Omega$ K\"ahler. The degree-one blowdown, a fiber of $q$, and the isomorphism $q|_E\colon E\to\mathbb P^1$ give
\[
 \int_XL^2=\int_XL\wedge Q=1,\quad Q^2=0,\quad L|_{TE}=0,\quad\int_EQ=1.
\]
Take $t=1$, $\alpha=L$, $\kappa_0=2L+Q$ and $\mu=\kappa_0^2/8$. With $r=1-s$, the core volumes are $V_s=1+4r+3r^2$; in particular $V_0=8$ and $V=V_1=1$.
Choose a nonconstant smooth function $0\leq\chi\leq1$ supported away from $E$, and put
\[
 b=\frac13\int_X\chi L\wedge\kappa_0,\qquad
 f=\chi-b,\qquad B=(1+f/2)L.
\]
Since $\int_XL\wedge\kappa_0=3$, we have $|f|\leq1$ and
$\int_XfL\wedge\kappa_0=0$. Thus $B\geq L/2$ and $B|_{TE}=0$, so the curve $C=E$ satisfies the hypotheses of \cref{thm:closed-core-subtraction}.
Where $\mathrm{d}\chi\ne0$, the form $L$ is symplectic, and wedging with $L$ is injective on real one-forms in real dimension four. Hence $\mathrm{d}B=\tfrac12\mathrm{d}f\wedge L\ne0$.

At the endpoint $\theta_1=(2+f/2)L\geq3L/2$, and $\lambda_\mu(L)=1$ by the closed nef formula; this gives the lower bound for $\lambda_\mu(\theta_1)$. Since $\kappa_0^2=8\mu$, the initial solution is $\gamma_0=\kappa_0$ with $u_0=0$. Therefore
$D_\mu(\kappa_0)=\sqrt8$, $I_B=3$, $\Delta_\mu(B)=3/\sqrt8$ and $U_\mu=11/\sqrt8$, so $\lambda_\mu(\theta_0)\leq11/\sqrt8$. The cost identity and the lower bound at the endpoint give the asserted estimate for $\mathcal C$.

The pullback of $\{x_1=0\}$ is $E$ plus a fiber of $q$, so
$\mathcal O(E)=p^*\mathcal O_{\mathbb P^2}(1)\otimes q^*\mathcal O_{\mathbb P^1}(-1)$.
Give this bundle its product metric $h_E$, with curvature $L-Q$, and let $s_E$ be its canonical section. The Poincar\'e--Lelong formula gives, for $\rho=\tfrac12\log|s_E|_{h_E}^2$,
\[
 \theta_1+\ddc\rho=B+\tfrac12\Omega+\tfrac12[E]\geq\tfrac12\Omega.
\]
Finally, $\int_X\Omega\wedge a_0=5$, so the right side of \eqref{eq:closed-core-nef-morse} is $8-10=-2$, although $V_1=1$.
\end{proof}


\Needspace{18\baselineskip}
\section{Varying the probability measure}

Changing the probability measure changes the denominator of each Gauduchon quotient by a uniformly bounded factor. The value of the volume may change, but its vanishing and positivity do not. On the positive cone, complete independence of the measure is equivalent to invariance of total mass under changes of potential.

\begin{proposition}\label{prop:gaud-volume-density-comparison}
For smooth positive probability measures $\mu,\nu$, choose constants $m,M>0$ such that $m\mu\leq\nu\leq M\mu$. Then for every pseudoeffective $\theta$,
\[
 M^{-1/n}\lambda_\mu(\theta)
 \leq\lambda_\nu(\theta)
 \leq m^{-1/n}\lambda_\mu(\theta).
\]
In particular the zero set and the positivity set are independent of the chosen smooth probability measure.
\end{proposition}
\begin{proof}
The density inequalities give
$m^{1/n}D_\mu(\gamma)\leq D_\nu(\gamma)
 \leq M^{1/n}D_\mu(\gamma)$.
The numerators are nonnegative. Divide and take infima.
\end{proof}

\begin{proposition}\label{prop:gaud-volume-density-independence}
Fix a Hermitian representative $\omega$. The value $\mathcal V_\mu(\omega)$ is the same for all smooth positive probability measures $\mu$ if and only if
\[
 \int_X(\omega+\ddc v)^n
\]
is independent of the smooth potential $v$ with $\omega+\ddc v>0$.
\end{proposition}
\begin{proof}
For such a potential, put
\[
 M_v=\int_X(\omega+\ddc v)^n,
 \qquad \mu_v=\frac{(\omega+\ddc v)^n}{M_v}.
\]
It is a smooth positive probability measure, and $v$ already solves its normalized equation with scalar $M_v$. Scalar uniqueness and \cref{thm:gaud-volume-dual-formula} give
$\mathcal V_{\mu_v}(\omega)=M_v$.
Thus measure independence implies mass invariance. Conversely, if all these masses equal one constant, every smooth normalized equation has that constant as its scalar after integration. The volume is then independent of its smooth probability measure.
\end{proof}

The non-closed torus example of \cref{ex:torus-mass} therefore produces different values for different smooth probability measures. On a Gauduchon surface, smooth mass is conserved, so the value equals $\int_X\omega^2$ for every such measure. On the Hopf threefold of \cref{chap:hopf}, the value is finite and positive for each fixed $\mu$, even though $v_+(\omega)=+\infty$.

\Needspace{25\baselineskip}
\section{The lower-mass quotient}
\label{sec:lower-mass-quotient}

From now on $X$ is again an arbitrary connected compact complex manifold of dimension $n\geq1$, and $\omega$ is a Hermitian form on $X$. By \cref{thm:gaud-volume-mass-extrema}, the lower mass $v_-(\omega)$ is the infimum of the volumes $\mathcal V_\mu(\omega)$ over the set $\mathscr M$ of smooth strictly positive probability measures, and \eqref{eq:gaud-lower-mass-formula} expresses it through Gauduchon metrics alone. We relate attainment of this infimum to smooth minimizers of the mass, and then extend the quotient to nef backgrounds.

\begin{proposition}\label{prop:lower-mass-smooth-attainment}\index{Monge--Ampere mass@Monge--Amp\`ere mass!smooth minimum}
The infimum in \eqref{eq:gaud-lower-mass-formula} is attained if and only if the infimum of the masses $\int_X(\omega+\ddc v)^n$ over smooth functions $v$ with $\omega+\ddc v>0$ is attained. For every such minimizing potential $v$, the metric $\omega+\ddc v$ is Gauduchon.
\end{proposition}
\begin{proof}
Suppose that a Gauduchon metric $\gamma$ attains the infimum in \eqref{eq:gaud-lower-mass-formula}, and put $m=v_-(\omega)^{1/n}$. For the measure $\mu_\gamma=\gamma^n/\int_X\gamma^n$, equality holds in the H\"older inequality of the proof of \cref{thm:gaud-volume-mass-extrema}, so the test $\gamma$ gives
$\lambda_{\mu_\gamma}(\omega)\leq m$.
\Cref{thm:gaud-volume-mass-extrema} gives the opposite inequality. Hence the solution of the smooth equation for $\mu_\gamma$ has mass $m^n=v_-(\omega)$, which gives a minimizing smooth potential. For $n\geq2$, $\gamma$ also minimizes the quotient for $\mu_\gamma$, so by \cref{thm:gaud-volume-dual-formula} the solution form is conformal to $\gamma$. Since both have volume forms proportional to $\mu_\gamma$, the conformal factor is constant, and the solution metric is Gauduchon.

Conversely, let $v$ be a smooth potential minimizing the mass. For every smooth real $f$, the potentials $v\pm tf$ are admissible for sufficiently small $t>0$. The vanishing of the first variation of the mass therefore gives
\[
 \int_X f\,\ddc\bigl((\omega+\ddc v)^{n-1}\bigr)=0
 \quad\hbox{for every smooth real }f.
\]
Thus $\Omega=\omega+\ddc v$ is Gauduchon. Taking $\gamma=\Omega$ in the quotient and using Stokes' theorem for the Gauduchon metric $\Omega$ gives
\[
 \frac{\int_X\omega\wedge\Omega^{n-1}}
      {(\int_X\Omega^n)^{(n-1)/n}}
             =(\int_X\Omega^n)^{1/n}=v_-(\omega)^{1/n},
\]
so the infimum in \eqref{eq:gaud-lower-mass-formula} is attained. In dimension one all metrics are Gauduchon and all masses equal $\int_X\omega$.
\end{proof}

\Needspace{9\baselineskip}
The infimum over densities also makes sense beyond the positive cone. For a pseudoeffective smooth real $(1,1)$-form $\theta$ on $X$, put
\[
 \lambda_*(\theta)=
 \inf_{\gamma\text{ Gauduchon}}
 \frac{\int_X\theta\wedge\gamma^{n-1}}
      {(\int_X\gamma^n)^{(n-1)/n}}.
\]
The H\"older argument in the proof of \cref{thm:gaud-volume-mass-extrema} proves
\[
 \lambda_*(\theta)=\inf_{\mu\in\mathscr M}\lambda_\mu(\theta).
\]
Like each $\lambda_\mu$, the function $\lambda_*$ is finite, nonnegative, concave, positively homogeneous, and invariant under smooth $\ddc$ changes of $\theta$. Suppose $\theta$ is nef, fix a Hermitian form $h$, and for real $\epsilon>0$ choose a Hermitian representative $\omega_\epsilon$ of $\theta+\epsilon h$. Interchanging the infima over the test metric and over $\epsilon>0$, as in \cref{thm:gaud-volume-nef-limit}, gives
\[
 \lambda_*(\theta)
 =\lim_{\epsilon\downarrow0}\lambda_*(\theta+\epsilon h)
 =\lim_{\epsilon\downarrow0}v_-(\omega_\epsilon)^{1/n}.
\]
By invariance under smooth $\ddc$ changes, the last expression does not depend on the choice of the representatives $\omega_\epsilon$. This extends the lower-mass formula to the nef boundary. We do not compare $\lambda_*$ with other notions of lower volume for singular classes.


\section{Curvature conditions for mass extrema}
\label{sec:mass-extremum-geometry}

At a smooth local extremum of the mass, the first two variations give a Gauduchon condition and a sign for one lower background power. On a threefold these necessary conditions are sufficient because the mass polynomial stops at its quadratic term. Fix a Hermitian form $\omega$ and use $\HH_\omega$ and $\Mass_\omega$ from \cref{chap:mass}.

For $n\geq2$, \cref{thm:mass-invariance-characterization,prop:gaud-volume-density-independence} identify the differential conditions for $\mathcal V_\mu(\omega)$ to be independent of the smooth positive probability measure $\mu$. By \cref{prop:entropy-mass-extrema}, these are also the conditions for $H_0-\operatorname{Ent}_{\mathrm{d}V}$ to be constant on finite-entropy probabilities, for a fixed smooth positive probability volume $\mathrm{d}V$; testing smooth positive probabilities suffices. In dimension one the mass is always $\int_X\omega$, so both conclusions hold automatically.

For the next argument, a smooth real $(n-1,n-1)$-form $T$ is nonnegative when $\mathrm i\alpha\wedge\bar\alpha\wedge T\geq0$ for every $(1,0)$-covector $\alpha$. This is the usual positivity in this bidegree. We need an elementary local test for its sign.

\begin{lemma}\label{lem:gradient-test-positive-form}
Suppose a smooth real $(n-1,n-1)$-form $T$ satisfies
\[
 \int_X \mathrm{d}f\wedge\dc f\wedge T\geq0
 \quad\hbox{for every smooth real }f.
\]
Then $T\geq0$.
\end{lemma}
\begin{proof}
If the conclusion fails at a point, choose a real linear function $\ell$ in local real coordinates, whose differential at that point has the offending complex direction. Shrinking the coordinate ball makes
$\mathrm{d}\ell\wedge\dc\ell\wedge T$ strictly negative there. Choose a nonzero smooth cutoff $\chi$ supported in this ball and, for positive integers $k$, set
\[
 f_k=k^{-1}\chi\sin(k\ell),
\]
extended by zero to $X$. Since $\mathrm{d}\ell$ is a nonzero constant covector in these coordinates,
\[
 \mathrm{d}f_k\wedge\dc f_k
 =\chi^2\cos^2(k\ell)\,\mathrm{d}\ell\wedge\dc\ell+O(k^{-1}),
\]
where the remainder has uniformly bounded coefficients after multiplication by $k$. The oscillatory part of $\cos^2(k\ell)=(1+\cos(2k\ell))/2$ has integral tending to zero, by integration by parts in a coordinate direction where $\ell$ has nonzero derivative. Hence
\[
 \lim_{k\to\infty}\int_X\mathrm{d}f_k\wedge\dc f_k\wedge T
 =\frac12\int_X\chi^2\mathrm{d}\ell\wedge\dc\ell\wedge T<0,
\]
a contradiction. Every complex covector is the $(1,0)$ part of a real covector, so this tests all required directions.
\end{proof}

\Needspace{12\baselineskip}
\begin{proposition}\label{prop:local-mass-extremum-necessary}
Let $n\geq3$, and let $v\in\HH_\omega$ be a local minimum of $\Mass_\omega$ in the $C^2$ topology. Then $\Omega=\omega+\ddc v$ satisfies
\[
 \ddc(\Omega^{n-1})=0,
 \qquad \ddc(\Omega^{n-2})\leq0.
\]
At a local maximum the first condition is unchanged and the second inequality is reversed.
\end{proposition}
\begin{proof}
Every smooth direction is admissible for a sufficiently short two-sided interval. The first variation gives the Gauduchon condition. Differentiating a second time gives
\[
 \begin{aligned}
 \left.\frac{\mathrm{d}^2}{\mathrm{d}t^2}\right|_{t=0}\Mass_\omega(v+tf)
 &=n(n-1)\int_X(\ddc f)^2\wedge\Omega^{n-2}\\
 &=-n(n-1)\int_X\mathrm{d}f\wedge\dc f\wedge\ddc(\Omega^{n-2}).
 \end{aligned}
\]
For the last equality, integrate once against $f\ddc f$, then use that $\ddc(\Omega^{n-2})$ is $\mathrm{d}$- and $\dc$-closed. At a minimum this second variation is nonnegative for every $f$. Apply \cref{lem:gradient-test-positive-form} to $-\ddc(\Omega^{n-2})$. The maximum case is identical with the opposite sign.
\end{proof}

In dimension three these necessary conditions are also sufficient, and a local extremum is automatically global.

\begin{theorem}\label{thm:threefold-mass-minimizer}
Assume $n=3$ and $v\in\HH_\omega$. The following are equivalent:
\begin{enumerate}
\item $v$ is a local minimum of $\Mass_\omega$;
\item $v$ is a global minimum of $\Mass_\omega$ on $\HH_\omega$;
\item $\Omega=\omega+\ddc v$ is Gauduchon and $\ddc\omega\leq0$.
\end{enumerate}
In this case $v_-(\omega)=\int_X\Omega^3$, and the Gauduchon infimum in \eqref{eq:gaud-lower-mass-formula} is attained. The corresponding statements for maxima hold with $\ddc\omega\geq0$, and then $v_+(\omega)=\int_X\Omega^3<+\infty$.
\end{theorem}
\begin{proof}
The preceding proposition gives the necessary conditions, since $\ddc\Omega=\ddc\omega$. Conversely, if $\Omega$ is Gauduchon, the dimension-three mass polynomial becomes, for every smooth real $f$,
\begin{equation}\label{eq:threefold-mass-dirichlet}
 \int_X(\Omega+\ddc f)^3-\int_X\Omega^3
       =-3\int_X\mathrm{d}f\wedge\dc f\wedge\ddc\omega.
\end{equation}
The form $\mathrm{d}f\wedge\dc f$ is nonnegative. The displayed difference is therefore nonnegative if $\ddc\omega\leq0$, proving a global minimum. Global minimality implies local minimality. Exact-background regularization identifies the infimum over smooth strictly positive potentials with $v_-(\omega)$, and \cref{prop:lower-mass-smooth-attainment} proves attainment in the Gauduchon formula. Reversing the sign in the same identity proves the maximum statements, including the equality with the supremum over bounded potentials.
\end{proof}

The local test supplies the necessity for smooth attainment in \eqref{eq:threefold-mass-dirichlet}. Neither that test nor the sufficient lower-bound argument produces such a representative in an arbitrary background.

\Needspace{8\baselineskip}
For the Hopf metric $\omega$ and the pulled-back Fubini--Study form $\alpha$ of \cref{chap:hopf}, $\ddc\omega=-\alpha^2\leq0$ and $\omega$ is Gauduchon. Thus the minimum is attained by the constant potential and equals $6$, consistently with its infinite upper mass. On a surface, by contrast, the mass is affine in the potential; a local extremum exists precisely when $\omega$ is Gauduchon, in which case all masses are equal. In higher dimensions the necessary sign condition of \cref{prop:local-mass-extremum-necessary} does not control all higher terms of the mass polynomial, and sufficiency is not claimed.


\section{Geometric bounds and determinant concentration}
\label{sec:gaud-determinant-concentration}

By \cref{thm:gaud-determinant-endpoint}, the determinant densities of all normalized Gauduchon currents are bounded in the critical space $L^{n/(n-1)}$ exactly when the lower mass is positive. Together with the mass-invariance criterion of \cref{thm:mass-invariance-characterization}, this gives geometric sufficient conditions for the critical bound. The critical exponent cannot be raised, even for currents that are powers of K\"ahler forms in one fixed class. The same concentrating family shows that currents approaching the maximum of the determinant functional for a fixed smooth density need not satisfy any uniform higher integrability, although the maximizer itself is smooth. In the subsections that follow, we compute the complete determinant profile of a background that satisfies the conditions of \cref{thm:mass-invariance-characterization}, show that this profile does not control the ellipticity of the background, and construct exact maximizers for varying densities whose weak limit has zero determinant measure.

Throughout this section $n\geq2$ and $\mathrm{d}V$ is a smooth positive probability volume form. For a Hermitian form $h$, $k_T$ denotes the determinant density relative to $\mathrm{d}V$ of a current $T\in\mathscr G_h$, as in \cref{sec:gaud-determinant-integrability}. Put $q=n/(n-1)$.

\Needspace{14\baselineskip}
\begin{corollary}\label{cor:geometric-determinant-bound}
Suppose $g$ is a Hermitian form with $\ddc g=\ddc(g^2)=0$, and put $V_g=\int_Xg^n$. For every Hermitian form $h\geq cg$, $c>0$,
\[
 v_-(h)\geq c^nV_g,\qquad
 \sup_{T\in\mathscr G_h}\|k_T\|_{L^q(\mathrm{d}V)}
       \leq c^{-1}V_g^{-1/n}.
\]
Equality holds in the determinant bound when $h=g$ and $c=1$. In particular, on a compact K\"ahler manifold or a compact complex surface, the critical norm is uniformly bounded over the normalized Gauduchon currents of each fixed Hermitian form.
\end{corollary}
\begin{proof}
The identity $\ddc(g^2)=2g\wedge\ddc g+2\mathrm{d}g\wedge\dc g$ shows that $\mathrm{d}g\wedge\dc g=0$, so \cref{thm:mass-invariance-characterization} gives $v_-(g)=V_g$. \Cref{cor:gaud-determinant-background-comparison} then gives both bounds, and \cref{thm:gaud-determinant-endpoint} gives equality in the determinant bound when $h=g$.

A K\"ahler form satisfies both differential identities. On a surface, \cref{thm:gauduchon} gives a metric $g$ with $\ddc g=0$, and $\ddc(g^2)=0$ for degree reasons. By compactness, each fixed $h$ satisfies $h\geq cg$ for some $c>0$.
\end{proof}

For a fixed smooth $F$, \cref{thm:weak-gaud-compact-duality} gives a unique maximizing current for the determinant functional $\mathfrak D_F$ of \eqref{eq:weak-gaud-functional}, and this current is smooth and strictly positive, so its determinant density is bounded. The following construction shows that currents approaching this maximum satisfy no uniform bound on higher moments of their determinant densities.

\begin{proposition}\label{prop:gaud-determinant-concentration}
Let $X=E^n$, where $E=\mathbb C/(\mathbb Z+\mathrm i\mathbb Z)$ carries its flat area form of total area one, and write $\omega_i$ for the pullback of this form from the $i$th factor. Set
\[
 h=\sum_{i=1}^n\omega_i,\qquad
 \mathrm{d}V=\omega_1\wedge\cdots\wedge\omega_n,\qquad V=n!,\qquad
 \alpha=\frac{n-1}{n}.
\]
For $F=0$, the unique maximizing current is
$T_*=h^{n-1}/V$, with $k_{T_*}=V^{-1/n}$ and
$\mathfrak D_0(T_*)=V^{-1/n}$.
There are smooth strictly positive closed forms $S_\varepsilon\in\mathscr G_h$, $0<\varepsilon<1/4$, such that
\[
 \|k_{S_\varepsilon}\|_{L^q}=V^{-1/n},\qquad
 \mathfrak D_0(S_\varepsilon)\longrightarrow V^{-1/n},
\]
while for every fixed finite $p>q$,
\begin{equation}\label{eq:gaud-concentration-moments}
 \int_X k_{S_\varepsilon}^{p}\,\mathrm{d}V
 \geq C_{n,p}[\log(1/\varepsilon)]^{-\alpha p}
                    \varepsilon^{2-2\alpha p}\longrightarrow+\infty.
\end{equation}
Their $L^\infty$ norms also diverge. Each $S_\varepsilon$ is the $(n-1)$st power of a K\"ahler form in the fixed class $V^{-1/(n-1)}[h]$, and has the same Bott--Chern, de Rham and Aeppli classes as $T_*$. No $S_\varepsilon$ attains the maximum of $\mathfrak D_0$.
\end{proposition}
\begin{proof}
Put $H_i=\bigwedge_{j\ne i}\omega_j$. Expanding the wedge powers gives
$h^{n-1}=(n-1)!\sum_iH_i$ and $h\wedge H_i=\mathrm{d}V$.
Thus $T_*=(1/n)\sum_iH_i$ is normalized, and the determinant formula gives $k_{T_*}=V^{-1/n}$. The smooth equation for $F=0$ is solved by the zero potential, since $h^n=V\,\mathrm{d}V$; \cref{thm:weak-gaud-compact-duality} therefore gives the maximal value and the uniqueness of the maximizer.

To construct $S_\varepsilon$, choose a nonnegative smooth function $\beta$ compactly supported in the unit disk in $\mathbb C$, of integral one, with $\beta\geq\beta_0>0$ on a disk of radius $r_0\in(0,1)$. On an embedded disk in the first elliptic factor put
$b_\varepsilon(z)=\varepsilon^{-2}\beta(z/\varepsilon)$, extended by zero; it has integral one. For $0<t<1$, set
\[
 \begin{gathered}
 f_{\varepsilon,t}=1-t+tb_\varepsilon,\qquad a=V^{-1/(n-1)},\\
 \gamma_{\varepsilon,t}=a(f_{\varepsilon,t}\omega_1+\omega_2+\cdots+\omega_n),
 \qquad T_{\varepsilon,t}=\gamma_{\varepsilon,t}^{n-1}.
 \end{gathered}
\]
The form $\gamma_{\varepsilon,t}$ is strictly positive, and it is closed because $f_{\varepsilon,t}$ depends only on the first factor, which has real dimension two. Direct expansion gives
\begin{equation}\label{eq:gaud-concentration-density}
 T_{\varepsilon,t}=\frac1n\left(H_1+f_{\varepsilon,t}\sum_{i=2}^nH_i\right),
 \qquad k_{T_{\varepsilon,t}}=V^{-1/n}f_{\varepsilon,t}^{\alpha}.
\end{equation}
Since $\int_Xf_{\varepsilon,t}\,\mathrm{d}V=1$, the trace integral is one, so $T_{\varepsilon,t}\in\mathscr G_h$, and
$\|k_{T_{\varepsilon,t}}\|_{L^q}=V^{-1/n}$ because $\alpha q=1$.

The mean-zero function $f_{\varepsilon,t}-1$ has a smooth real potential $\psi$ on $E$ with
$\ddc\psi=(f_{\varepsilon,t}-1)\omega_1$;
for example, the mean-zero Fourier series solving
$\Delta\psi=4\pi(f_{\varepsilon,t}-1)$ in the area-one coordinates.
Consequently $\gamma_{\varepsilon,t}-ah=\ddc(a\psi)$. Since $T_*=(ah)^{n-1}$, the closed form
$R=\sum_{j=0}^{n-2}\gamma_{\varepsilon,t}^{n-2-j}\wedge(ah)^j$ satisfies
\[
 T_{\varepsilon,t}-T_*=\ddc(a\psi R).
\]
This proves all the assertions about classes, including the case $n=2$, where $R=1$.

Jensen's inequality and $f_{\varepsilon,t}\geq1-t$ give
\[
 V^{-1/n}(1-t)^\alpha
 \leq\mathfrak D_0(T_{\varepsilon,t})
 =V^{-1/n}\int_Xf_{\varepsilon,t}^{\alpha}\,\mathrm{d}V
 <V^{-1/n},
\]
where the inequality is strict because $f_{\varepsilon,t}$ is nonconstant and $0<\alpha<1$.
On the product of the disk of radius $r_0\varepsilon$ with the other factors, we have
$f_{\varepsilon,t}\geq t\beta_0\varepsilon^{-2}$, and this set has volume $\pi r_0^2\varepsilon^2$. Hence
\[
 \begin{split}
 \int_Xk_{T_{\varepsilon,t}}^p\,\mathrm{d}V
   &\geq V^{-p/n}\pi r_0^2\beta_0^{\alpha p}t^{\alpha p}
                    \varepsilon^{2-2\alpha p},\\
 \|k_{T_{\varepsilon,t}}\|_\infty
   &\geq V^{-1/n}\beta_0^\alpha t^\alpha\varepsilon^{-2\alpha}.
 \end{split}
\]
Now take $t=1/\log(1/\varepsilon)$, which lies in $(0,1)$ because $\varepsilon<1/4<\mathrm e^{-1}$, and $S_\varepsilon=T_{\varepsilon,t}$. Then $\mathfrak D_0(S_\varepsilon)$ converges to its maximum, while both lower bounds diverge for $p>q$ because $\alpha p>1$. All constants are independent of $\varepsilon$.
\end{proof}

If instead $t\in(0,1)$ is kept fixed in the same construction, then
$\int_Xk_{T_{\varepsilon,t}}^p\,\mathrm{d}V\geq C_{n,p,t}\varepsilon^{2-2\alpha p}$ and
$\|k_{T_{\varepsilon,t}}\|_\infty\geq C_{n,t}\varepsilon^{-2\alpha}$.
Thus there are no uniform bounds on higher moments, either over the full normalized cone or along a maximizing sequence for one smooth density. Every member of these families attains the supremum $V^{-1/n}$ of the critical norm in \cref{cor:geometric-determinant-bound}, whereas the maximum of $\mathfrak D_0$ is attained only by $T_*$.


\section{Vanishing lower mass}
\label{sec:vanishing-lower-mass}

By \eqref{eq:gaud-lower-mass-formula}, the lower mass $v_-(h)$ vanishes as soon as there are Gauduchon metrics $\eta$ whose degrees $\int_Xh\wedge\eta^{n-1}$ stay bounded while their volumes $\int_X\eta^n$ tend to infinity. We construct such metrics on a compact threefold. Writing $\Gamma=\eta^2$, the volume of $\eta$ is a determinant integral of the coefficient matrix of the positive $\ddc$-closed form $\Gamma$. The construction first concentrates the trace of a normalized form $\Gamma$ in one direction near a hypersurface, and then adds an average of compact rational curves that fills the two complementary directions. The degree stays bounded, while the determinant, and hence the volume of the positive root, becomes large. The smooth equation then converts these metrics into potentials whose masses tend to zero.

\begin{theorem}\label{thm:vanishing-lower-mass}\index{Positive lower mass!failure}
There is a connected compact complex threefold $X$ such that, for every Hermitian form $h$ on $X$, there are smooth real functions $u_j$ satisfying
\[
 h+\ddc u_j>0,\qquad \sup_Xu_j=0,\qquad
 \int_X(h+\ddc u_j)^3\longrightarrow0.
\]
In particular, $v_-(h)=0$ for every Hermitian form $h$ on this manifold.
\end{theorem}

\subsection{The elliptic quotient}\label{subsec:lower-mass-elliptic-quotient}

Write points of $\mathbb P^3$ as $[z:w]$, with $z,w\in\C^2$, and set
\[
 \Omega=\{[z:w]:z\ne0,\ w\ne0\},\qquad
 q=\mathrm{e}^{2\pi},\qquad g[z:w]=[z:qw].
\]
Thus $g$ is an automorphism of $\Omega$; let $\langle g\rangle$ be the cyclic group it generates. This group acts freely. Indeed, if $g^k$ fixed a point, the same scalar would have to equal $1$ on the nonzero block $z$ and $q^k$ on the nonzero block $w$, which forces $k=0$.

\Needspace{8\baselineskip}
The function
\[
 t=\log\frac{\|w\|^2}{\|z\|^2}
\]
is well defined on $\Omega$ and satisfies $t\circ g=t+4\pi$. Consequently the action is properly discontinuous: $t$ is bounded on compact subsets, so only finitely many translates of one compact set can meet another. Moreover, the strip $\{0\leq t\leq4\pi\}$ is compact. In homogeneous coordinates of total norm one, its defining inequalities bound both blocks away from zero, so it is a closed subset of $\mathbb P^3$ contained in $\Omega$. Since the strip meets every orbit,
\[
 X=\Omega/\langle g\rangle
\]
is a connected compact complex threefold.

The map $[z:w]\mapsto([z],[w])$ induces a holomorphic submersion
$X\to\mathbb P^1\times\mathbb P^1$. Using a logarithm of the scalar relating two points with the same image, one sees that its fibers are the elliptic curves
$E=\C/(2\pi\mathbb Z+2\pi\mathrm{i}\mathbb Z)$. Let $\alpha,\beta$ be the pullbacks of the Fubini--Study forms of the two factors, each normalized to have integral one. Although $t$ does not descend to $X$, its differentials $\mathrm dt$ and $\dc t$ do. The form
\begin{equation}\label{eq:lower-mass-quotient-metric}
 \omega=\alpha+\beta+\mathrm{d}t\wedge\dc t
\end{equation}
is Hermitian: $\alpha+\beta$ is semipositive with kernel the vertical line, while $\mathrm{d}t\wedge\dc t$ is semipositive and strictly positive on the vertical line.

\Needspace{7\baselineskip}
On the open set $\{z_0w_0\ne0\}$ we use the coordinates
\[
 x=\frac{z_1}{z_0},\qquad y=\frac{w_1}{w_0},\qquad
 \zeta=\log\frac{w_0}{z_0}\in E.
\]
Near $0\in E$, the last coordinate is represented by an ordinary complex number. For any one of these coordinates, denoted by $v$, put
$\sigma_v=\frac{\mathrm{i}}{2\pi}\mathrm{d}v\wedge\mathrm{d}\bar v$.
Choose a smooth probability density $p$ relative to $\sigma_v$ that is positive on $\{|v|<1\}$ and vanishes on $\{|v|\geq1\}$. Choose also a smooth function $\chi_0$ with values in $[0,1]$, supported in $\{|v|<1\}$ and equal to one on $\{|v|\leq1/2\}$, and write $\chi(x,y)=\chi_0(x)\chi_0(y)$. For $0<\delta<1/100$ define
\[
 B_\delta=\chi(x,y)\delta^{-2}p(\zeta/\delta)\sigma_\zeta,
\]
extended by zero. This is a smooth semipositive $(1,1)$-form on $X$, because the chosen functions are smooth with compact support in their coordinate domains. As $\delta\downarrow0$ it concentrates on the hypersurface $\{\zeta=0\}$.

The first step shows that normalized positive $\ddc$-closed $(2,2)$-forms can pair arbitrarily strongly with $B_\delta$.

\begin{lemma}\label{lem:lower-mass-trace-concentration}
For $\delta>0$ put
\[
 C_\delta=\sup\left\{\int_XB_\delta\wedge\Theta:
 \Theta>0,\ \ddc\Theta=0,\ \int_X\omega\wedge\Theta=1\right\},
\]
the supremum being taken over smooth real $(2,2)$-forms $\Theta$. These numbers are finite and tend to infinity as $\delta\downarrow0$.
\end{lemma}
\Needspace{12\baselineskip}
\begin{proof}
For fixed $\delta$, choose a constant $M_\delta>0$ with $B_\delta\leq M_\delta\omega$; this shows that $C_\delta\leq M_\delta$. Suppose instead that $C_{\delta_l}\leq C$ along a sequence $\delta_l\to0$. Then $\int_X(C\omega-B_{\delta_l})\wedge\Theta\geq0$ for every admissible $\Theta$, and in particular for the squares of Gauduchon metrics. By \cref{thm:lamari}, there are quasi-plurisubharmonic functions $v_l$ such that
\[
 C\omega+\ddc v_l\geq B_{\delta_l},\qquad \sup_Xv_l=0.
\]
By \cref{thm:normalized-compactness}, a subsequence converges in $L^1$ to a $C\omega$-plurisubharmonic function $v$, not identically $-\infty$. Since $\delta^{-2}p(\zeta/\delta)\sigma_\zeta$ is an approximate identity in the variable $\zeta$,
\[
 B_{\delta_l}\rightharpoonup\chi[\zeta=0],\qquad
 C\omega+\ddc v\geq\chi[\zeta=0].
\]
On a coordinate ball where $\chi=1$, add to $v$ a smooth strictly plurisubharmonic function whose complex Hessian dominates $C\omega$. The resulting plurisubharmonic function has a current that dominates $[\zeta=0]$, so its Lelong number is at least one at every point of an open subset of this hypersurface. Smooth corrections do not change Lelong numbers, so the same holds for $v$. Siu's analyticity theorem, recalled in \cref{subsec:psh-lelong-levels}, shows that
$Z=\{\nu(v,\cdot)\geq1\}$ is an analytic subset of $X$. It is a proper subset because $v$ is not identically $-\infty$.

Let $\widetilde Z$ be the inverse image of $Z$ in $\Omega$. It contains an open subset of
$H_0\cap\Omega$, where $H_0=\{w_0=z_0\}$. This surface is a projective plane with two points removed, so it is connected and irreducible, and the identity principle gives $H_0\cap\Omega\subset\widetilde Z$. Since $\widetilde Z$ is invariant under $g$, also
\[
 H_k\cap\Omega\subset\widetilde Z,\qquad H_k=\{w_0=q^kz_0\},
 \qquad k\in\mathbb Z.
\]
Now use the chart $z_1=1$ near the point $[0:1:1:0]\in\Omega$. There $H_k$ is the graph $z_0=q^{-k}w_0$. For each fixed $(w_0,w_1)$ near $(1,0)$, every local holomorphic defining function of $\widetilde Z$ vanishes at the distinct values $z_0=q^{-k}w_0$, which tend to zero as $k\to+\infty$. The one-variable identity theorem makes each defining function vanish identically on a polydisk. Thus $\widetilde Z$ contains a nonempty open subset of the connected manifold $\Omega$, contradicting properness. Hence no subsequence of $C_\delta$ stays bounded as $\delta\downarrow0$.
\end{proof}

By the lemma we can choose $\delta_j\downarrow0$, each smaller than half its predecessor, and then, by the definition of the supremum, smooth strictly positive $\ddc$-closed $(2,2)$-forms $\Theta_j$ with
\begin{equation}\label{eq:lower-mass-large-trace}
 \int_X\omega\wedge\Theta_j=1,\qquad
 J_j=\int_XB_{\delta_j}\wedge\Theta_j>j^4.
\end{equation}

\subsection{Averaging the trace and compact curves}

A large pairing with $B_{\delta_j}$ means that $\Theta_j$ has a large coefficient in one direction, but possibly only on a small set. Averaging by translations spreads this coefficient over a whole product region of size comparable to $\delta_j$ in the vertical direction. A second average, this time of compact rational curves, fills the two complementary directions, which is needed for a lower bound on a determinant. Fix $j$ and write $\delta=\delta_j$.

For complex parameters $a,b,\tau$, the automorphisms
\[
 T_{a,b,\tau}[z_0:z_1:w_0:w_1]
 =[z_0:z_1+az_0:\mathrm{e}^{\tau}w_0:
                  \mathrm{e}^{\tau}(w_1+bw_0)]
\]
of $\Omega$ commute with $g$, and in the product chart they act by
$(x,y,\zeta)\mapsto(x+a,y+b,\zeta+\tau)$.
\Needspace{13\baselineskip}
Average $T_{a,b,\tau}^*\Theta_j$ against the probability density
\[
 3^{-4}p(a/3)p(b/3)(3\delta)^{-2}p(\tau/(3\delta))
       \sigma_a\wedge\sigma_b\wedge\sigma_\tau,
\]
and denote the result by $\widehat\Theta_j$. It is smooth, strictly positive and $\ddc$-closed. The parameters range in a fixed compact set of automorphisms, so $(T_{a,b,\tau}^{-1})^*\omega\leq C\omega$ with one constant $C$, and a change of variables gives
\begin{equation}\label{eq:lower-mass-averaged-degree}
 \int_X\omega\wedge\widehat\Theta_j\leq C.
\end{equation}

Let $F_j$ be the coefficient of $\sigma_x\wedge\sigma_y$ in the restriction of $\Theta_j$ to the slices $\{\zeta=\text{constant}\}$, and let $\widehat F_j$ be the corresponding coefficient of $\widehat\Theta_j$. Positivity gives $F_j\geq0$, and since the translations preserve $\sigma_x\wedge\sigma_y$, $\widehat F_j$ is an average of translates of $F_j$. At a point of
\[
 V_j=\{|x|,|y|<1/4,\ |\zeta|<\delta/4\},
\]
every translation carrying it into $\operatorname{supp}B_\delta$ has
$|a|,|b|\leq5/4$ and $|\tau|\leq5\delta/4$. On this range the first two averaging factors are bounded below by a positive constant, and the third by a positive constant times $\delta^{-2}$. On the other hand, $\chi\leq1$ and $p$ is bounded, so the number $J_j$ of \eqref{eq:lower-mass-large-trace} is at most a constant times $\delta^{-2}$ times the integral of $F_j$ over $\operatorname{supp}B_\delta$. Comparing the two integrals gives
\begin{equation}\label{eq:lower-mass-spread-trace}
 \widehat F_j\geq cJ_j\quad\hbox{on }V_j,
\end{equation}
with $c>0$ independent of $j$.

The second average fills the two horizontal directions. Put
\[
 f(a,b)=p(4(a-1))p(b/2),\qquad
 \rho_j(\tau)=4\delta^{-2}p(2\tau/\delta).
\]
For parameters in the supports of these functions and $|u|\leq1/16$, consider the holomorphic maps $\iota_{\tau,u,a,b}\colon\mathbb P^1\to X$ defined by
\begin{equation}\label{eq:lower-mass-rational-curves}
 \iota_{\tau,u,a,b}[z_0:z_1]
 =[z_0:z_1:\mathrm{e}^{\tau}(z_0+uz_1):
                \mathrm{e}^{\tau}(az_1+bz_0)]\in X.
\end{equation}
For $0<\varepsilon\leq1/16$, average their integration currents against
\[
 \rho_j(\tau)f(a,b)\varepsilon^{-2}p(u/\varepsilon)
       \sigma_\tau\wedge\sigma_a\wedge\sigma_b\wedge\sigma_u,
\]
and call the result $T_{j,\varepsilon}$. Define $T_{j,0}$ in the same way, but with $u=0$ and without the integral over $u$.

\Needspace{22\baselineskip}
\begin{lemma}\label{lem:lower-mass-curve-average}
For $\varepsilon>0$, the form $T_{j,\varepsilon}$ is smooth on $X$, semipositive and closed, and
\[
 \int_X\omega\wedge T_{j,\varepsilon}\leq C'
\]
with $C'$ independent of $j$ and $\varepsilon$. For each fixed $j$, $T_{j,\varepsilon}$ converges smoothly to $T_{j,0}$ on compact subsets of the product chart as $\varepsilon\downarrow0$. In this chart
\begin{equation}\label{eq:lower-mass-curve-local}
 T_{j,0}=\rho_j(\zeta)\sigma_\zeta\wedge K_f,\qquad
 K_f=\int_\C f(a,y-ax)\frac{\mathrm{i}}{2\pi}
       (\mathrm{d}y-a\,\mathrm{d}x)\wedge
       (\mathrm{d}\bar y-\bar a\,\mathrm{d}\bar x)\,\sigma_a,
\end{equation}
and on $\{|x|,|y|\leq1/4\}$ the form $K_f$ is bounded below by a positive multiple of $\sigma_x+\sigma_y$.
\end{lemma}
\Needspace{8\baselineskip}
\begin{proof}
On the parameter supports, $|a-bu|\geq3/4-2/16=5/8$. Hence the matrix
$\left(\begin{smallmatrix}1&u\\ b&a\end{smallmatrix}\right)$ is uniformly invertible, and the projective line \eqref{eq:lower-mass-rational-curves} lies in $\Omega$. Distinct deck translates of this line are disjoint: equality at a point fixes $[z_0:z_1]$, and invertibility of the matrix then forces the deck scalar to be one. Thus each map is an embedding onto a compact rational curve in $X$.

The evaluation map from $\mathbb P^1$ times the open parameter domain to $X$ is a holomorphic submersion. Indeed, after fixing a projective gauge, variations of $u$ and $a$ span the two $w$-coordinates at points with $z_1\ne0$, variations of $\tau$ and $b$ span them at points with $z_1=0$, and the coordinate on $\mathbb P^1$ supplies the remaining direction. For $\varepsilon>0$, the average is the integral along the fibers of this submersion of a smooth compactly supported parameter density, which proves smoothness. As an average of closed positive integration currents, $T_{j,\varepsilon}$ is closed and semipositive.

All curve parameters lie in one compact subset of the domain where the matrix is invertible, so the $\omega$-areas of the curves are uniformly bounded. The weights in $\tau$ and $u$ have mass one, and the weight in $(a,b)$ has a fixed finite mass. This proves the degree bound.

When $u=0$, the curve is given in the product chart by $\zeta=\tau$ and $y=ax+b$. Integrating in $b$ and $\tau$ yields \eqref{eq:lower-mass-curve-local}; the $(1,1)$-factor there is the integration current of the graph of $y=ax+b$. On $\{|x|,|y|\leq1/4\}$, the density $f(a,y-ax)$ is bounded below by a positive constant for $a$ in a fixed small disk about $1$. The covectors $\mathrm{d}y-a\,\mathrm{d}x$, for $a$ in that disk, have no common nonzero annihilator, so integrating their squared absolute values and using compactness of the unit sphere in the $(x,y)$-directions gives the asserted lower bound.

For general $u$, the curve is given by
\[
 \zeta=\tau+\log(1+ux),\qquad y=\frac{ax+b}{1+ux}.
\]
On a fixed compact subset of the chart, $1+ux$ stays away from zero for all sufficiently small $u$. Solving for the parameters gives
$\tau=\zeta-\log(1+ux)$ and $b=(1+ux)y-ax$. Regarding $\rho_j$ as a smooth function on $E$, this change of variables also applies near the boundary of its support. The Jacobians, the density coefficients and any fixed number of their derivatives depend smoothly on $u$. Averaging against the approximate identity in $u$ therefore gives local $C^\infty$ convergence. This convergence is needed only for each fixed $j$; it is not uniform in $j$.
\end{proof}

\subsection{The volume of the positive root}

The two averages contribute complementary blocks of the coefficient matrix of a $(2,2)$-form: the translated forms give a large vertical entry, and the curves give a positive horizontal block. A Schur-complement estimate combines them into a lower bound for the volume of the positive root. In the coframe $(x,y,\zeta)$ put
$\mathrm{d}V_0=\sigma_x\wedge\sigma_y\wedge\sigma_\zeta$. For a real $(1,1)$-form $\xi$, let $H_\xi$ be its coefficient matrix in this normalization. A real $(2,2)$-form $\Theta$ corresponds to a Hermitian matrix $Q_\Theta$, which we call its cofactor matrix, by
\[
 \xi\wedge\Theta=\operatorname{tr}(H_\xi Q_\Theta)\,\mathrm{d}V_0.
\]
Strict positivity of $\Theta$ is equivalent to $Q_\Theta>0$.

\Needspace{8\baselineskip}
If $\eta^2=\Theta$ and $H=H_\eta$, expanding the wedge products gives
\[
 Q_\Theta=2(\det H)H^{-1},\qquad
 \eta^3=\frac{3}{\sqrt2}(\det Q_\Theta)^{1/2}\,\mathrm{d}V_0.
\]
Conversely, these formulas uniquely recover a smooth positive $H$ from $Q_\Theta>0$; this is the three-dimensional case of the root construction preceding the proof of \cref{thm:lamari}. We write $\Theta^{3/2}=\eta^3$.

\begin{lemma}\label{lem:lower-mass-unbounded-root}
There are smooth strictly positive $\ddc$-closed $(2,2)$-forms $\Gamma_j$ such that
\[
 \int_X\omega\wedge\Gamma_j\leq C_2,\qquad
 R_j=\int_X\Gamma_j^{3/2}\longrightarrow+\infty.
\]
\end{lemma}
\begin{proof}
Write the cofactor matrix of $\widehat\Theta_j$ in block form, with a $2\times2$ Hermitian block $A_j$ in the $(x,y)$-directions and a column $\kappa_j\in\C^2$:
\[
 Q_j=\begin{pmatrix}A_j&\kappa_j\\\kappa_j^*&\widehat F_j\end{pmatrix}.
\]
The cofactor matrix of $T_{j,0}$ has the form $\operatorname{diag}(B_j,0)$. Let $I_2$ denote the $2\times2$ identity matrix. The lower bound for $K_f$ and the definition of $\rho_j$ give a constant $c_0>0$ such that $B_j\geq c_0\delta_j^{-2}I_2$ on $V_j$. Since $Q_j>0$, the Schur complement $A_j-\kappa_j\kappa_j^*/\widehat F_j$ is positive definite, and
\[
 \det\bigl(Q_j+\operatorname{diag}(B_j,0)\bigr)
 =\widehat F_j\det\left(B_j+A_j-
             \frac{\kappa_j\kappa_j^*}{\widehat F_j}\right)
 \geq\widehat F_j\det B_j,
\]
by monotonicity of the determinant on positive Hermitian matrices. Combining the root formula with \eqref{eq:lower-mass-spread-trace}, we obtain
\[
 (\widehat\Theta_j+T_{j,0})^{3/2}
       \geq c\delta_j^{-2}\sqrt{J_j}\,\mathrm{d}V_0
       \quad\hbox{on }V_j.
\]
\Needspace{8\baselineskip}
The smaller region $W_j=\{|x|,|y|<1/8,\ |\zeta|<\delta_j/8\}$ has $\mathrm{d}V_0$-volume equal to a fixed positive constant times $\delta_j^2$. Consequently
\[
 I_j=\int_{W_j}(\widehat\Theta_j+T_{j,0})^{3/2}
       \geq c_1\sqrt{J_j}>c_1j^2.
\]
For this fixed $j$, \cref{lem:lower-mass-curve-average} gives smooth convergence of $T_{j,\varepsilon}$ to $T_{j,0}$ on the closure of $W_j$. The cofactor matrix of $\widehat\Theta_j+T_{j,0}$ is positive definite there, so the integrand of the root volume depends continuously on these smooth coefficients. We can therefore choose $0<\varepsilon_j\leq1/16$ so small that
\[
 \int_{W_j}(\widehat\Theta_j+T_{j,\varepsilon_j})^{3/2}
       \geq I_j/2.
\]
Set $\Gamma_j=\widehat\Theta_j+T_{j,\varepsilon_j}$. It is smooth, strictly positive and $\ddc$-closed. Its degree is at most $C_2=C+C'$, by \eqref{eq:lower-mass-averaged-degree} and \cref{lem:lower-mass-curve-average}, whereas its root volume is at least $c_1j^2/2$.
\end{proof}

\Needspace{15\baselineskip}
\begin{proof}[of \cref{thm:vanishing-lower-mass}]
Fix a Hermitian form $h$ on $X$, and choose $A>0$ with $h\leq A\omega$. Let $\eta_j$ be the positive root of $\Gamma_j$, so that $\eta_j^2=\Gamma_j$. These metrics are Gauduchon, and
\[
 P_j=\int_Xh\wedge\eta_j^2\leq AC_2,\qquad
 R_j=\int_X\eta_j^3\longrightarrow+\infty.
\]
Set $\mu_j=\eta_j^3/R_j$ and solve the smooth equation in the background $h$:
\[
 h_j=h+\ddc u_j>0,\qquad h_j^3=c_h(\mu_j)\mu_j,
 \qquad \sup_Xu_j=0.
\]
Since $D_{\mu_j}(\eta_j)=R_j^{2/3}$, using $\eta_j$ as a test in \cref{thm:gaud-volume-dual-formula} gives
\[
 0<\int_Xh_j^3=c_h(\mu_j)=\mathcal V_{\mu_j}(h)
       \leq\frac{P_j^3}{R_j^2}\longrightarrow0.
\]
The potentials $u_j$ are smooth, hence bounded, so $v_-(h)=0$.
\end{proof}

\subsection{Consequences for dimension and equations}

The example occurs in the smallest possible dimension: on a surface, comparison with a Gauduchon metric gives positive lower mass. Products with projective spaces carry the threefold degeneration to every higher dimension, and a product with the Hopf threefold produces, in addition, unbounded upper mass.

\begin{corollary}\label{cor:lower-mass-dimension}
Every Hermitian form on a compact complex manifold of dimension at most two has positive lower mass. In every dimension $n\geq3$ there is a connected compact complex manifold on which every Hermitian form has zero lower mass.
\end{corollary}
\Needspace{10\baselineskip}
\begin{proof}
In dimension one the mass is $\int_Xh>0$, and the surface case is \cref{cor:geometric-determinant-bound}. In higher dimensions we use the monotonicity and the homogeneity of degree $n$ of $v_-$, which follow directly from \eqref{eq:gaud-lower-mass-formula}, to pass from a split metric to every Hermitian metric.

\Needspace{11\baselineskip}
For $n=3$ this is \cref{thm:vanishing-lower-mass}. For $n>3$, take the product of the threefold $X$ with $\mathbb P^{n-3}$, let $\kappa$ be a Fubini--Study form on the projective factor, and consider the split metric $h_0=\omega+\kappa$, with pullbacks understood. For the potentials $u_j$ given by the theorem for the background $\omega$,
\[
 \int_{X\times\mathbb P^{n-3}}(h_0+\ddc u_j)^n
 =\binom n3\left(\int_X(\omega+\ddc u_j)^3\right)
                   \left(\int_{\mathbb P^{n-3}}\kappa^{n-3}\right)
 \longrightarrow0.
\]
For another Hermitian form $h$, choose $A$ with $h\leq Ah_0$. Monotonicity and homogeneity give $v_-(h)\leq A^n v_-(h_0)=0$.
\end{proof}

\Needspace{16\baselineskip}
\begin{corollary}\label{cor:both-mass-properties-fail}
In every dimension $n\geq6$ there is a connected compact complex manifold with
\[
 v_-(h)=0,\qquad v_+(h)=+\infty
\]
for every Hermitian form $h$.
\end{corollary}
\begin{proof}
Take the product of the threefold above with the Hopf threefold of \cref{chap:hopf}, with a split metric. Potentials pulled back from the first factor make the masses tend to zero, while the smooth potentials of \cref{thm:hopf-unbounded}, pulled back from the second factor, make them tend to infinity. In each case the product integral is $\binom63$ times the product of the two factor masses. The metric comparison in the proof of \cref{cor:lower-mass-dimension} handles the lower mass for every Hermitian form, and the metric independence in \cref{prop:bmp-metric} handles the upper mass. Products with projective spaces give every higher dimension by the same expansion of wedge powers.
\end{proof}

Finally, fix a smooth positive probability measure $\mu$ on the threefold $X$. The smooth equation still gives $\mathcal V_\mu(h)>0$ for each Hermitian form $h$. With $\mathscr M$ and $\lambda_*$ as in \cref{sec:gaud-volume-extrema}, \cref{thm:gaud-volume-mass-extrema} gives
\[
 \inf_{\nu\in\mathscr M}\mathcal V_\nu(h)=0,
 \qquad \lambda_*(h)=0.
\]
Thus even a smooth strictly positive background can have vanishing infimum over densities. For the metrics $\eta_j$ and numbers $R_j,P_j$ in the proof of \cref{thm:vanishing-lower-mass}, the measures $\mu_j=\eta_j^3/R_j$ satisfy
$\mathcal V_{\mu_j}(h)\leq P_j^3/R_j^2\to0$.
Hence the densities $\mu_j/\mu$ cannot be uniformly bounded above: a bound $\mu_j\leq M\mu$ would give $\mathcal V_{\mu_j}(h)\geq M^{-1}\mathcal V_\mu(h)>0$, by the comparison of denominators in the proof of \cref{prop:gaud-volume-density-comparison}.

The example also shows that the endpoint estimate in \cref{thm:gaud-determinant-endpoint} fails for the full family of normalized Gauduchon currents: the supremum of the $L^{3/2}$ norms of their determinant densities is infinite, although every subcritical $L^q$ norm, $1<q<3/2$, remains uniformly bounded by \cref{thm:gaud-determinant-subcritical}. For the same fixed background, \cref{thm:gaud-local-adjoint-endpoint,cor:gaud-local-density-family} give uniform $L^\infty$ bounds for the exact maximizers and their adjoint probabilities in a small $C^{2,\alpha}$ family of smooth densities. A larger sufficient family is given by \cref{thm:gaud-hessian-adjoint-endpoint}: uniformly comparable smooth solution forms yield uniform bounds for both densities, without derivative bounds on the varying Hessian. Endpoint estimates for arbitrary bounded families of continuous densities remain open.


\Needspace{13\baselineskip}
\subsection{A balanced metric and a finite upper mass}
\label{sec:lower-mass-metric-geometry}

Although every Hermitian form on the threefold of \cref{thm:vanishing-lower-mass} has zero lower mass, the metric $\omega$ of \eqref{eq:lower-mass-quotient-metric} is balanced and has finite upper mass. The threefold is a primitive elliptic bundle, so the loss of mass under a potential is a vertical Dirichlet energy along the elliptic fibers, and the degeneration cannot be realized by potentials that are the sum of a periodic function of $t$ and a function on the base: their masses stay above $8$. Put $D=\mathrm{d}t\wedge\dc t$, so that $\omega=\alpha+\beta+D$.

Let $B=\mathbb P^1\times\mathbb P^1$, let $\alpha_B,\beta_B$ be the Fubini--Study forms of the two factors, normalized to have integral one, so that $\alpha,\beta$ are their pullbacks, and put $\eta=\alpha_B+\beta_B$. On $B$ consider the line bundle $\mathcal O(1)\boxtimes\mathcal O(-1)$, whose fiber at $([z],[w])$ is the space of linear maps from the line $\mathbb Cz$ to the line $\mathbb Cw$, with the metric induced by the standard norms. Sending $[z:w]\in\Omega$ to the linear map $z\mapsto w$ identifies $\Omega$ with the complement of the zero section, the function $\log(\|w\|^2/\|z\|^2)$ with the logarithm of the squared norm, and $g$ with multiplication by $q=\mathrm e^{2\pi}$. The curvature representative of minus the first Chern class is $c=\beta_B-\alpha_B$, and $\eta\wedge c=0$. Hence $X$ is the primitive elliptic bundle of \cref{sec:primitive-elliptic-bundles} for these data, with the same function $t$, and $\omega$ is its metric $\pi^*\eta+D$. The constants of that subsection are
\[
 T=4\pi,\qquad A=2,\qquad c^2=-2\alpha_B\wedge\beta_B,\qquad \kappa=1,
\]
and the eigenvalues of $c$ relative to $\eta$ are $\pm1$, so that $M_0=1$.

\begin{corollary}\label{prop:lower-mass-balanced-metric}\index{Balanced metric!finite upper mass}\index{Elliptic fibration!vanishing lower mass}
On the threefold $X$ of \cref{thm:vanishing-lower-mass}, the Hermitian form $\omega=\alpha+\beta+D$ satisfies
\[
 \mathrm{d}(\omega^2)=0,\qquad \ddc\omega=2\alpha\wedge\beta\geq0,
 \qquad \int_X\omega^3=24\pi.
\]
For every smooth real function $u$,
\begin{equation}\label{eq:lower-mass-vertical-energy}
 \int_X(\omega+\ddc u)^3
 =24\pi-6\int_X\mathrm{d}u\wedge\dc u\wedge\alpha\wedge\beta.
\end{equation}
Consequently $v_+(\omega)=24\pi$. In particular, $X$ has the bounded mass property, while $v_-(h)=0$ for every Hermitian form $h$ on $X$.
\end{corollary}
\begin{proof}
\Cref{prop:primitive-elliptic-metric}, with $T=4\pi$, $A=2$ and $-\pi^*(c^2)=2\alpha\wedge\beta$, gives the balanced condition, the formula for $\ddc\omega$, the volume $3TA=24\pi$, the identity \eqref{eq:lower-mass-vertical-energy} and $v_+(\omega)=24\pi$. The metric independence in \cref{prop:bmp-metric} then gives the bounded mass property of $X$. The assertion about the lower mass is \cref{thm:vanishing-lower-mass}.
\end{proof}

For an admissible smooth $u$, equality in the upper-mass bound holds precisely when $u$ is constant along every elliptic fiber. Indeed, $\alpha\wedge\beta$ is a positive volume form in the base directions, so the nonnegative integral in \eqref{eq:lower-mass-vertical-energy} vanishes exactly when the derivative of $u$ along the fibers vanishes. The fibers are connected, and smooth local sections then show that $u$ is the pullback of a smooth function on $\mathbb P^1\times\mathbb P^1$. For the potentials with masses tending to zero in \cref{thm:vanishing-lower-mass}, the same integral instead tends to $4\pi$.

A function $f(t)$ descends to $X$ when $f$ is $4\pi$-periodic, since $t\circ g=t+4\pi$. Even after adding an arbitrary potential pulled back from the base, such functions leave a positive mass gap. The sharp gap follows from \cref{lem:periodic-variance-bound}, through the bounds on the first two derivatives of $f$ imposed by positivity.

\begin{corollary}\label{prop:elliptic-radial-mass}\index{Monge--Ampere mass@Monge--Amp\`ere mass!radial gap}
Let $\pi\colon X\to B=\mathbb P^1\times\mathbb P^1$ be the elliptic fibration above, and put $L=4\pi$. Consider the smooth potentials
\[
 u=f(t)+\pi^*v,\qquad \omega+\ddc u>0,
\]
where $f\colon\mathbb R\to\mathbb R$ is smooth and $L$-periodic and $v\in C^\infty(B,\R)$. The infimum of $\int_X(\omega+\ddc u)^3$ over these potentials is $8$. Every such mass is strictly greater than $8$, and the infimum is approached already with $v=0$.
\end{corollary}
\begin{proof}
Here $L=T$ and $\kappa^{-1/2}=M_0=1$. Since $T\geq2$, the definition of $F$ gives $F(1,4\pi)=1-1/(3\pi)$, so both bounds in \cref{prop:primitive-elliptic-separated} equal $3TA\bigl(1-F(1,4\pi)\bigr)=8$. That proposition therefore shows that every mass in the statement is strictly greater than $8$, and that potentials $f(t)$ with $v=0$ have masses tending to $8$. In terms of $p=f'$, the mass formula in its proof reads
\begin{equation}\label{eq:elliptic-radial-mass}
 \int_X(\omega+\ddc u)^3
 =6\left(L-\int_0^L p(s)^2\,\mathrm{d}s\right).
\end{equation}
\end{proof}

\Needspace{12\baselineskip}
\begin{corollary}\label{prop:lower-mass-no-pluriclosed}
The threefold $X$ admits neither a pluriclosed Hermitian metric nor a closed K\"ahler current.
\end{corollary}
\begin{proof}
The first Chern class of $\mathcal O(1)\boxtimes\mathcal O(-1)$ is represented by $\alpha_B-\beta_B$, which is nonzero in real cohomology, so this is part of \cref{prop:primitive-elliptic-metric}.
\end{proof}

Since a compact complex manifold carries a closed K\"ahler current exactly when it belongs to Fujiki's class $\mathcal C$, as recalled in \cref{sec:non-kahler-geometric-problems}, the threefold $X$ lies outside this class. The example shows that a balanced metric does not guarantee positive lower mass, even together with bounded upper mass. The Hopf threefold, in contrast, has positive lower mass and infinite upper mass, so the two properties are logically independent, and the products in \cref{cor:both-mass-properties-fail} show that both can fail at once. In particular, in a theorem asserting solvability for every nonzero finite non-pluripolar measure, bounded mass alone cannot replace the hypothesis of positive lower mass.

The exact equations of \cref{sec:zero-lower-mass-product-equations} on products $X_0\times Y$ with $v_-(h_0)=0$ apply to this threefold. We use the notation of that section with $Y=\mathbb P^2$: $\kappa$ is the Fubini--Study form normalized by $\int_Y\kappa^2=1$, $\mathrm{d}V$ is a smooth positive probability volume form on $X_0\times Y$, the functions $R$, $q_\alpha$ and $u_\alpha$, for $0<\alpha<1$, are defined there, and $b_H$ is the threshold of \cref{sec:individual-density-subsolutions}.

\Needspace{6\baselineskip}
\begin{corollary}\label{cor:elliptic-fivefold-exact-threshold}\index{Elliptic fibration!exact product equation}
Let $X_0$ be the elliptic threefold of \cref{thm:vanishing-lower-mass}, and let $h_0$ be the balanced metric of \cref{prop:lower-mass-balanced-metric}. On $X_0\times\mathbb P^2$, with background $H=h_0+\kappa$, the densities $f=Rq_\alpha$ of \cref{thm:zero-lower-mass-product-equation} satisfy
\[
 v_-(H)=0,\qquad b_H(f)=240\pi,
 \qquad\mathrm{MA}^{\mathrm{ext}}_H(u_\alpha)=240\pi f\,\mathrm{d}V.
\]
\end{corollary}
\begin{proof}
\Cref{thm:vanishing-lower-mass} gives $v_-(h_0)=0$, and \cref{prop:lower-mass-balanced-metric} gives $V_0=24\pi$. \Cref{thm:zero-lower-mass-product-equation} applies with $m=3$, $k=2$ and $C=\binom53=10$.
\end{proof}


\section{Orders of degeneration at the nef boundary}
\label{sec:nef-volume-degeneration}

If $\theta$ is nef and $h$ is Hermitian, the volumes $\mathcal V_\mu(\theta+\epsilon h)$ are positive and, by \cref{thm:gaud-volume-nef-limit}, decrease to $\mathcal V_\mu(\theta)$ as $\epsilon\downarrow0$. When this limit is zero, the rate of decay measures how degenerate $\theta$ is at the nef boundary. For a closed nef class on a compact K\"ahler manifold, the volume of the perturbation is the intersection polynomial $\int_X(\theta+\epsilon h)^n$, so it decays like an integral power of $\epsilon$ determined by the numerical dimension. For a nonclosed nef background the logarithmic rate of decay can be fractional, and it need not converge. We therefore define lower and upper orders of degeneration as lower and upper limits of this rate. Changing the probability measure or the perturbing metric only introduces bounded factors, so the orders are invariants of the background.

Let $X$ again be a connected compact complex manifold of dimension $n\geq1$, and $\mu$ a smooth strictly positive probability form on $X$. Let $\theta$ be nef and let $h$ be Hermitian. For $0<\epsilon<1$, put
\[
 a_{\mu,h}(\theta;\epsilon)
   =\frac{\log\mathcal V_\mu(\theta+\epsilon h)}{\log\epsilon}.
\]
The logarithm is finite because $\theta+\epsilon h$ belongs to $\mathscr P$, so that its volume is positive.

\Needspace{16\baselineskip}
\begin{proposition}\label{prop:nef-degeneration-orders}\index{Nef degeneration!orders}
The limits
\[
 \underline a(\theta)=\liminf_{\epsilon\downarrow0}
                         a_{\mu,h}(\theta;\epsilon),
 \qquad
 \overline a(\theta)=\limsup_{\epsilon\downarrow0}
                         a_{\mu,h}(\theta;\epsilon)
\]
are independent of $\mu$ and $h$, and satisfy
\[
 0\leq\underline a(\theta)\leq\overline a(\theta)\leq n.
\]
They are unchanged when $\theta$ is replaced by $\theta+\ddc f$, with $f$ smooth, or by $r\theta$, with $r>0$. If $\mathcal V_\mu(\theta)>0$, both limits are zero.
\end{proposition}
\begin{proof}
Since $\theta$ is pseudoeffective, monotonicity and homogeneity give
\[
 \epsilon\lambda_\mu(h)
       \leq\lambda_\mu(\theta+\epsilon h)
       \leq\lambda_\mu(\theta+h),
 \qquad 0<\epsilon\leq1.
\]
Thus there are constants $c,C>0$ such that
\[
 c\epsilon^n\leq\mathcal V_\mu(\theta+\epsilon h)\leq C,
 \qquad 0<\epsilon\leq1.
\]
Taking logarithms and dividing by the negative number $\log\epsilon$ shows that the quotients $a_{\mu,h}(\theta;\epsilon)$ stay bounded as $\epsilon\downarrow0$, and that their lower and upper limits lie in $[0,n]$.

By \cref{prop:gaud-volume-density-comparison}, changing the smooth probability measure multiplies the volume by a factor between two fixed positive constants. The resulting error in the logarithm is bounded, and after division by $\log\epsilon$ it tends to zero.

\Needspace{8\baselineskip}
For another Hermitian metric $h'$, choose constants $c_2\geq c_1>0$ with $c_1h\leq h'\leq c_2h$. Then
\[
 \mathcal V_\mu(\theta+c_1\epsilon h)
 \leq\mathcal V_\mu(\theta+\epsilon h')
 \leq\mathcal V_\mu(\theta+c_2\epsilon h).
\]
For any fixed $t>0$ and small $\epsilon$,
\[
 \frac{\log\mathcal V_\mu(\theta+t\epsilon h)}{\log\epsilon}
 =a_{\mu,h}(\theta;t\epsilon)
       \frac{\log(t\epsilon)}{\log\epsilon}.
\]
The last factor tends to one and the first is bounded, while the rescaling $\epsilon\mapsto t\epsilon$ does not change lower and upper limits as $\epsilon\downarrow0$. Hence the left side has the same lower and upper limits as $a_{\mu,h}(\theta;\epsilon)$. Dividing the preceding inequalities by $\log\epsilon<0$ reverses them, and taking lower and upper limits proves independence of the metric.

Invariance under smooth $\ddc$ changes follows from that of $\mathcal V_\mu$. For scaling, homogeneity gives
\[
 \mathcal V_\mu(r\theta+\epsilon h)
       =r^n\mathcal V_\mu(\theta+(\epsilon/r)h).
\]
The bounded error $n\log r$ and the same rescaling argument prove scale invariance. Finally, if $\mathcal V_\mu(\theta)>0$, then by \cref{thm:gaud-volume-nef-limit} the perturbed volumes lie between $\mathcal V_\mu(\theta)$ and $\mathcal V_\mu(\theta+h)$ for $0<\epsilon\leq1$, so both limits are zero.
\end{proof}

We call $\underline a(\theta)$ and $\overline a(\theta)$ the \emph{lower} and \emph{upper orders of degeneration} of the nef background $\theta$. When they agree, their common value is the \emph{common order}, and we then also say that $\theta$ has this order. The two orders can differ, and a common order can be irrational. Even when both are zero, the volume can tend to zero, at a logarithmic rate. The elliptic examples in \cref{sec:nef-surface-orders} exhibit each of these behaviors. Closed nef classes on K\"ahler manifolds, in contrast, degenerate polynomially.

\Needspace{15\baselineskip}
\begin{proposition}\label{prop:nef-degeneration-kahler}
Suppose $X$ is K\"ahler, $\theta$ is closed and nef, and $h$ is K\"ahler. Set
\[
 r=\max\left\{0\leq k\leq n:
          \int_X\theta^k\wedge h^{n-k}>0\right\}.
\]
Then
\[
 \underline a(\theta)=\overline a(\theta)=n-r.
\]
Thus $n$ minus the common order is the numerical dimension of the closed nef class of $\theta$.
\end{proposition}
\begin{proof}
All the intersection numbers in the definition of $r$ are nonnegative. Indeed, the mixed integrals against $h$ of smooth positive representatives of $\theta+\delta h$ are positive, agree with the cohomological integrals by closed Stokes, and converge to the stated numbers as $\delta\downarrow0$. The number for $k=0$ is $\int_Xh^n>0$, so $r$ is well defined.

By the closed nef formula following \cref{thm:gaud-volume-nef-limit},
\[
 \mathcal V_\mu(\theta+\epsilon h)
       =\int_X(\theta+\epsilon h)^n
       =\sum_{k=0}^n\binom nk\epsilon^{n-k}
                         \int_X\theta^k\wedge h^{n-k}.
\]
As $\epsilon\downarrow0$, the leading nonzero term is a positive constant times $\epsilon^{n-r}$, and taking the logarithmic quotient proves the assertion. By \cref{prop:nef-degeneration-orders}, the result holds for every Hermitian perturbation, not only for K\"ahler ones.
\end{proof}

\Needspace{17\baselineskip}
The orders can also be computed on the non-K\"ahler Hopf threefold $X$ of \cref{chap:hopf}. Keep its forms $\omega=\alpha+\beta$, where $\alpha\geq0$ is closed of rank two, $\beta\geq0$ has rank one, $\alpha^3=\beta^2=0$, and $\int_X\omega^3=6$. Take $\theta=\alpha$, $h=\omega$, and $\mu_0=\omega^3/6$. The metric
\[
 \Omega_\epsilon=\alpha+\epsilon\omega
                    =(1+\epsilon)\alpha+\epsilon\beta
\]
is strictly positive. Its volume is
\[
 \Omega_\epsilon^3=\epsilon(1+\epsilon)^2\omega^3
                  =6\epsilon(1+\epsilon)^2\mu_0.
\]
Hence the zero potential solves the smooth normalized equation, and
\[
 \mathcal V_{\mu_0}(\alpha+\epsilon\omega)
             =6\epsilon(1+\epsilon)^2,
 \qquad
 \underline a(\alpha)=\overline a(\alpha)=1.
\]
By \cref{prop:nef-degeneration-orders}, the same orders are obtained for every smooth positive probability measure and every Hermitian perturbation.

Each perturbed metric nevertheless has infinite upper mass. If $V_j$ are the smooth potentials of \cref{thm:hopf-unbounded}, then
\[
 \Omega_\epsilon+\ddc(\epsilon V_j)
       =\alpha+\epsilon(\omega+\ddc V_j)>0,
\]
and monotonicity of the determinant of positive matrices gives
\[
 \int_X\bigl(\Omega_\epsilon+\ddc(\epsilon V_j)\bigr)^3
       \geq\epsilon^3\int_X(\omega+\ddc V_j)^3\longrightarrow+\infty.
\]
Thus a finite order of degeneration at the boundary coexists with infinite upper mass at every positive perturbation.

\subsection{Elliptic surface volumes and order spectra}\label{sec:nef-surface-orders}

A coefficient that vanishes along one elliptic factor determines the harmonic mean in the volume formula. Finite-order zeros give fractional or logarithmic decay; flat coefficients allow the logarithmic slope to oscillate.

Let $E=\mathbb C/(\mathbb Z+\mathrm{i}\mathbb Z)$, write $X=E_z\times E_w$, and fix
\[
 \omega_1=\frac{\mathrm{i}}{2}\mathrm{d} z\wedge\mathrm{d}\bar z,\qquad
 \omega_2=\frac{\mathrm{i}}{2}\mathrm{d} w\wedge\mathrm{d}\bar w,\qquad
 \mu=\omega_1\wedge\omega_2,\qquad h=\omega_1+\omega_2.
\]
Both area forms have integral one. Let $a,q\geq0$ and $r$ be smooth functions of $z$ alone such that $|r|^2\leq aq$, and set
\[
 \theta=\frac{\mathrm{i}}{2}\bigl(a\,\mathrm{d} z\wedge\mathrm{d}\bar z+
 r\,\mathrm{d} z\wedge\mathrm{d}\bar w+\bar r\,\mathrm{d} w\wedge\mathrm{d}\bar z+
 q\,\mathrm{d} w\wedge\mathrm{d}\bar w\bigr).
\]
This is a smooth semipositive, hence nef, form.

\begin{proposition}\label{prop:nef-elliptic-exact-volume}
For $\epsilon>0$, put $g=q+\epsilon$ and $I=\int_Eg^{-1}\omega_1$. Then
\begin{equation}\label{eq:nef-elliptic-volume}
 \mathcal V_\mu(\theta+\epsilon h)
 =\frac{2}{I}\left(\int_E(a+\epsilon)\omega_1-
 \int_E\frac{|r|^2}{g}\omega_1\right).
\end{equation}
The infimum defining this volume is attained by a smooth positive Gauduchon metric. In particular, if $a=b>0$ and $r=0$, then $\mathcal V_\mu(\theta+\epsilon h)=2(b+\epsilon)/I$. If $|r|^2=aq$, then
\begin{equation}\label{eq:nef-rank-one-elliptic-volume}
 \mathcal V_\mu(\theta+\epsilon h)=\frac{2\epsilon(1+K)}{I},
 \qquad K=\int_E\frac{a}{q+\epsilon}\omega_1.
\end{equation}
\end{proposition}

\begin{proof}
Define
\[
 c=I^{-1}\left(\int_E(a+\epsilon)\omega_1-
 \int_E\frac{|r|^2}{g}\omega_1\right),\qquad
 A=\frac{c+|r|^2}{g}.
\]
The inequality $(a+\epsilon)g-|r|^2>0$ gives $c>0$, while the definition of $c$ makes $A-a-\epsilon$ have zero $\omega_1$-mean. Division of its nonconstant Fourier coefficients by the nonzero flat Laplace eigenvalues produces a smooth real function $u$ on $E$ with $\ddc u=(A-a-\epsilon)\omega_1$. Pull it back to $X$. The form $\Omega=\theta+\epsilon h+\ddc u$ has coefficient matrix $\left(\begin{smallmatrix}A&r\\\bar r&g\end{smallmatrix}\right)$, determinant $c$, and $\Omega^2=2c\mu$.

The metric $\gamma_*=\Omega/g$ has constant fiber coefficient $1$; its remaining coefficients depend only on $z$. Applying $\ddc$ to its base or mixed terms repeats $\mathrm{d} z$ or $\mathrm{d}\bar z$, so $\ddc\gamma_*=0$. It is smooth, positive, and pointwise conformal to $\Omega$. The arithmetic--geometric mean argument in \cref{thm:gaud-volume-dual-formula} and equality for $\gamma_*$ prove \eqref{eq:nef-elliptic-volume}. The two specializations follow by direct substitution.
\end{proof}

\begin{proposition}\label{prop:nef-diagonal-fractional-and-zero}\index{Nef degeneration!fractional and zero orders}
There are smooth semipositive nonclosed forms on the fixed surface $X$ with common order $1-1/m$ for each integer $m>1$. There is also one with
\[
 \mathcal V_\mu(\theta+\epsilon h)\sim\frac{2b}{\pi\log(1/\epsilon)}
 \quad(\epsilon\downarrow0),
\]
for any chosen $b>0$; it has boundary volume zero and both orders zero.
\end{proposition}

\begin{proof}
In \cref{prop:nef-elliptic-exact-volume} take $a=b>0$ and $r=0$. Choose $q$ smooth, positive outside $z_0\in E$, and equal to $|\zeta|^{2m}$ in a small flat coordinate disk about $z_0$. A positive smooth extension exists away from the disk. For $m=1$, polar coordinates give
\[
 I=\int_E\frac{\omega_1}{q+\epsilon}
   =\pi\log(1/\epsilon)+O(1).
\]
Equation \eqref{eq:nef-elliptic-volume} gives the displayed asymptotic. If $m>1$, the change $s=|\zeta|^2=\epsilon^{1/m}t$ gives
\[
 I\sim C_m\epsilon^{1/m-1},\qquad
 C_m=\pi\int_0^\infty\frac{\mathrm{d} t}{1+t^m}>0,
\]
and the common order is $1-1/m$. In the logarithmic case $\lambda_\mu(\theta)\geq0$ by semipositivity, while $\lambda_\mu(\theta)\leq\lambda_\mu(\theta+\epsilon h)\to0$ by monotonicity. Hence $\mathcal V_\mu(\theta)=0$. The quotient of the displayed volume by $\epsilon^\delta$ diverges for every $\delta>0$. Since $\mathrm{d}\theta=\mathrm{d} q\wedge\omega_2\ne0$, these backgrounds are nonclosed.
\end{proof}

\begin{proposition}\label{prop:nef-oscillatory-orders}\index{Nef degeneration!unequal orders}
For every $0<\alpha<\beta<1$ there is a smooth semipositive nonclosed background on the same $(X,\mu,h)$ with
\[
 \bigl(\underline a(\theta),\overline a(\theta)\bigr)
 =(1+\alpha,1+\beta).
\]
For every $0<\sigma<1$ there is such a background with common order $1+\sigma$.
\end{proposition}

\begin{proof}
For all sufficiently large $t$, define
\[
 p(t)=\frac{\alpha+\beta}{2}
 +\frac{\beta-\alpha}{2}\sin(\log\log t),\qquad P(t)=tp(t).
\]
Choose $T$ large enough that $0<c_0\leq P'(t)\leq d_0<1$ for $t\geq T$; this follows from
\[
 P'(t)=p(t)+\frac{\beta-\alpha}{2\log t}\cos(\log\log t).
\]
On $E\setminus\{z_0\}$ choose a smooth $t(z)\geq T$ equal to $|\zeta|^{-2}$ near $z_0$ and bounded outside every neighborhood of $z_0$. Put
\[
 q=\mathrm{e}^{-t},\qquad a=\mathrm{e}^{-P(t)},\qquad
 r=\mathrm{e}^{-(P(t)+t)/2},
\]
and set all three equal to zero at $z_0$. These functions are smooth and flat there: $P(t)\geq\alpha t$ and every fixed derivative of $P$ grows at most polynomially in $t$, so the chain rule bounds each coordinate derivative by a polynomial in $|\zeta|^{-1}$ times $\mathrm{e}^{-c/|\zeta|^2}$. We have $|r|^2=aq$, so the resulting $\theta$ is semipositive. It is nonclosed: the component $\mathrm{d} q\wedge\omega_2$ of $\mathrm{d}\theta$ has two fiber differentials, which derivatives of the remaining terms cannot cancel.

Write $\epsilon=\mathrm{e}^{-L}$. Choose a fixed integration endpoint $T_0\geq T$ inside the coordinate disk, with the transition to the bounded global $t(z)$ outside a smaller disk. The complementary contributions to $I$ and $K$ in \eqref{eq:nef-rank-one-elliptic-volume} are $O(1)$, while polar coordinates on the disk give
\[
 I_{\mathrm{loc}}=\pi\int_{T_0}^\infty
 \frac{t^{-2}\,\mathrm{d} t}{\mathrm{e}^{-t}+\mathrm{e}^{-L}},\qquad
 K_{\mathrm{loc}}=\pi\int_{T_0}^\infty
 \frac{\mathrm{e}^{-P(t)}t^{-2}\,\mathrm{d} t}{\mathrm{e}^{-t}+\mathrm{e}^{-L}}.
\]
Integrating $I$ over $t\geq L$ gives a lower bound by a fixed multiple of $\mathrm{e}^L/L$. For the upper bound, use $(\mathrm{e}^{-t}+\mathrm{e}^{-L})^{-1}\leq\mathrm{e}^t$ for $t\leq L$ and at most $\mathrm{e}^L$ for $t\geq L$, then split $[T_0,L]$ at $L/2$. Thus
\begin{equation}\label{eq:nef-elliptic-I-asymptotic}
 \log I=L-\log L+O(1).
\end{equation}
The derivatives of $P$ and $t-P(t)$ are positive and bounded away from zero. Splitting the $K$-integral at $L$ therefore gives $K\leq C\mathrm{e}^{L-P(L)}$. On $[L,L+1]$ its denominator is at most $2\mathrm{e}^{-L}$ and $P(t)\leq P(L)+1$, so
\[
 c_1L^{-2}\mathrm{e}^{L-P(L)}\leq K\leq c_2\mathrm{e}^{L-P(L)}.
\]
The complementary $O(1)$ term is absorbed because $P(L)\leq\beta L$. Consequently $\log K=L-P(L)+O(\log L)$ and $K\to\infty$. Combining this with \eqref{eq:nef-rank-one-elliptic-volume} and \eqref{eq:nef-elliptic-I-asymptotic} yields
\begin{equation}\label{eq:nef-oscillatory-log-volume}
 \log \mathcal V_\mu(\theta+\mathrm{e}^{-L}h)=-L-P(L)+O(\log L),
 \qquad
 \frac{\log \mathcal V_\mu(\theta+\epsilon h)}{\log\epsilon}
 =1+p(L)+o(1).
\end{equation}
The sine in $p(L)$ has limits $-1$ and $1$ along sequences, proving the first assertion. If instead $P(t)=\sigma t$, the same smoothness and integral estimates give common order $1+\sigma$. In particular $\sigma=\sqrt2-1$ gives an irrational common order.
\end{proof}

\subsection{Comparison under an arbitrary modification}\label{sec:nef-modification-comparison}

The orders of degeneration are defined on a fixed manifold. We now show that pulling a nef background back by a modification can only increase them. Recall that a modification $\pi:Y\to X$ is a proper surjective holomorphic map that restricts to a biholomorphism between dense Zariski open subsets. The proof compares Gauduchon tests on $X$ with tests on $Y$. The pullback of a Gauduchon test is degenerate along the exceptional locus, and adding a small positive test on $Y$ repairs this. To control the effect of the added test uniformly, we first bound the pairing of pulled-back tests with a Hermitian form on $Y$ by their pairing with a Hermitian form on $X$, uniformly over the whole cone of tests.

\Needspace{8\baselineskip}
In this subsection $\pi:Y\to X$ is a modification of connected compact complex manifolds of dimension $n\geq2$. We use Hironaka's analytic elimination of indeterminacies in the following form: there are a finite composition $p:Z\to X$ of blowups along smooth centers and a modification $F:Z\to Y$ with $p=\pi\circ F$. It is obtained by resolving the indeterminacies of the inverse meromorphic map $\pi^{-1}\circ p$; the identity $p=\pi\circ F$ holds on the dense set where the maps are biholomorphic, hence everywhere. This classical input is discussed in \cref{sec:history-gaud-volume}. For an effective real divisor $D$, we write $\eta_D$ for the curvature form of a smooth Hermitian metric on the corresponding real line bundle, as in \cref{sec:divisorial-gaud-volume}.

\begin{lemma}\label{lem:nef-blowup-metric}
Let $p:Z\to X$ be a finite composition of blowups along smooth centers of codimension at least two, and let $h_X$ be a Hermitian form on $X$. There are a constant $C_0>0$, an effective real divisor $D$ on $Z$ whose components are exceptional for $p$, and a smooth curvature form $\eta_D$ such that
\begin{equation}\label{eq:nef-blowup-positive-metric}
 h_Z=C_0p^*h_X-\eta_D>0.
\end{equation}
Every component of $D$ is mapped by $p$ into a set of complex codimension at least two in $X$.
\end{lemma}

\begin{proof}
First let $q:M'\to M$ be one blowup along a smooth center, with exceptional divisor $E$, and let $h_M$ be a Hermitian form on $M$. The restriction of $\mathcal O(-E)$ to $E$ is the tautological line bundle of the projectivized normal bundle of the center, which is positive on the projective fibers. A Hermitian metric on that normal bundle therefore induces a metric on $\mathcal O(-E)|_E$ whose curvature is positive on the tangent spaces of the fibers. Extend this metric to $\mathcal O(-E)$ on $M'$ by extending the smooth logarithm of its ratio to a fixed reference metric from the closed submanifold $E$, and let $\beta$ be its curvature. The semipositive form $q^*h_M$ is positive off $E$, and at points of $E$ its kernel is exactly the tangent space of the fiber, on which $\beta$ is positive. Hence $Cq^*h_M+\beta>0$ for every sufficiently large $C$: otherwise, unit tangent vectors violating the inequality for $C\to\infty$ would converge, by compactness, to a vector in the kernel of $q^*h_M$ on which $\beta$ is nonpositive. With $\eta_E=-\beta$, the curvature of the induced metric on $\mathcal O(E)$, this reads $Cq^*h_M-\eta_E>0$.

For a composition of blowups, apply this step successively, with $h_M$ the positive form already obtained on the previous blowup. Each step multiplies the previous divisor by a positive constant, pulls it back and adds the new exceptional divisor, so the divisor stays effective and exceptional. This proves \eqref{eq:nef-blowup-positive-metric}. Each exceptional divisor maps into the image of its center, which has dimension at most $n-2$.
\end{proof}

\begin{lemma}\label{lem:nef-uniform-pairing}
Fix Hermitian forms $h_X,h_Y$ for a modification $\pi:Y\to X$. There is a constant $K>0$ such that
\begin{equation}\label{eq:nef-uniform-pairing}
 \int_Y h_Y\wedge\pi^*\Psi\leq K\int_X h_X\wedge\Psi
\end{equation}
for every smooth strictly positive $(n-1,n-1)$-form $\Psi$ on $X$ with $\ddc\Psi=0$.
\end{lemma}

\begin{proof}
Use the factorization $p=\pi\circ F$ above and \cref{lem:nef-blowup-metric}. By compactness there is $a>0$ with $h_Z\geq aF^*h_Y$. Pulling back the semipositive form $\Psi$ and integrating, we obtain
\[
 a\int_Yh_Y\wedge\pi^*\Psi
 \leq\int_Zh_Z\wedge p^*\Psi
 =C_0\int_Xh_X\wedge\Psi-\int_Dp^*\Psi.
\]
Here the changes of variables use that $F$ and $p$ have degree one, and the Poincar\'e--Lelong formula, together with $\ddc p^*\Psi=0$, replaces the pairing with $\eta_D$ by the integral over $D$. This last integral vanishes: on the regular locus of each component of $D$, the restriction of $p$ has rank at most $n-2$, so the restriction of the $(n-1,n-1)$-form $p^*\Psi$ is zero. Taking $K=C_0/a$ proves the assertion.
\end{proof}

\Needspace{17\baselineskip}
\begin{proposition}\label{prop:nef-modification-volume}\index{Nef degeneration!modification}
Let $\pi:Y\to X$ be a modification of connected compact complex manifolds of dimension $n\geq2$, and let $\theta$ be a smooth nef background on $X$. Then $\pi^*\theta$ is nef. For fixed Hermitian forms $h_X,h_Y$ and smooth positive probability forms $\mu_X,\mu_Y$, there are constants $C,K>0$, independent of $\epsilon$, such that
\begin{equation}\label{eq:nef-main-modification-comparison}
 \mathcal V_{\mu_Y}(\pi^*\theta+\epsilon h_Y)
 \leq C\mathcal V_{\mu_X}(\theta+K\epsilon h_X)
 \qquad(\epsilon>0).
\end{equation}
Consequently
\[
 \underline a(\pi^*\theta)\geq\underline a(\theta),\qquad
 \overline a(\pi^*\theta)\geq\overline a(\theta).
\]
In dimension one a modification is an isomorphism and preserves both orders.
\end{proposition}

\begin{proof}
\emph{Step 1: nefness.} Choose $C'$ with $\pi^*h_X\leq C'h_Y$. For $\epsilon>0$, pull back a smooth representative $\theta+\ddc u_\tau\geq-\tau h_X$ with $\tau<\epsilon/C'$; this gives a positive representative of $\pi^*\theta+\epsilon h_Y$. Thus $\pi^*\theta$ is nef, and both volumes in \eqref{eq:nef-main-modification-comparison} are positive.

\emph{Step 2: regularized tests.} The ratio $\pi^*\mu_X/\mu_Y$ is smooth, nonnegative and bounded on the compact manifold $Y$, so there is $c>0$ with $\mu_Y\geq c\pi^*\mu_X$. Fix a smooth Gauduchon metric $\gamma_0$ on $Y$, given by \cref{thm:gauduchon}. For an arbitrary Gauduchon metric $\gamma$ on $X$, write $\Psi=\gamma^{n-1}$, and for $\delta>0$ define $\Gamma_\delta$ by
\[
 \Gamma_\delta^{n-1}=\pi^*\Psi+\delta\gamma_0^{n-1}.
\]
The right side is a strictly positive $(n-1,n-1)$-form, and such a form has a unique smooth positive $(n-1)$st root, obtained by solving for the Hermitian matrix of the root from the positive cofactor matrix. Since the right side is $\ddc$-closed, $\Gamma_\delta$ is a smooth Gauduchon metric on $Y$.

\emph{Step 3: the numerator.} For $\epsilon>0$, \cref{lem:nef-uniform-pairing} gives
\begin{equation}\label{eq:nef-modification-numerator}
 \lim_{\delta\downarrow0}\int_Y(\pi^*\theta+\epsilon h_Y)
 \wedge\Gamma_\delta^{n-1}
 =\int_X\theta\wedge\Psi+\epsilon\int_Yh_Y\wedge\pi^*\Psi
 \leq\int_X(\theta+K\epsilon h_X)\wedge\Psi.
\end{equation}
The first term in the middle is nonnegative because $\theta$ is nef, and the second is positive because $\pi^*\Psi$ is positive on the set where $\pi$ is biholomorphic, so the limit is positive.

\emph{Step 4: the denominator.} We claim that
\begin{equation}\label{eq:nef-modification-denominator}
 \lim_{\delta\downarrow0}D_{\mu_Y}(\Gamma_\delta)
 =\int_Y(\pi^*\gamma^n)^{(n-1)/n}\mu_Y^{1/n}
 \geq c^{1/n}D_{\mu_X}(\gamma)>0.
\end{equation}
To justify the limit on the exceptional locus, represent $\Gamma_\delta^{n-1}$ by its positive cofactor matrix. The density of $(\Gamma_\delta^n)^{(n-1)/n}$ is a fixed dimensional constant times the $1/n$th power of the determinant of that matrix, and this expression extends continuously to positive semidefinite matrices. Since $Y$ is compact, dominated convergence applies. The inequality follows from $\mu_Y\geq c\pi^*\mu_X$ and the change of variables for the degree-one map $\pi$.

\emph{Step 5: conclusion.} Test the infimum defining $\lambda_{\mu_Y}(\pi^*\theta+\epsilon h_Y)$ with $\Gamma_\delta$, use \eqref{eq:nef-modification-numerator} and \eqref{eq:nef-modification-denominator}, and let $\delta\downarrow0$. For every Gauduchon metric $\gamma$ on $X$ this gives
\[
 \lambda_{\mu_Y}(\pi^*\theta+\epsilon h_Y)
 \leq c^{-1/n}
 \frac{\int_X(\theta+K\epsilon h_X)\wedge\gamma^{n-1}}{D_{\mu_X}(\gamma)}.
\]
Taking the infimum over $\gamma$ and then the $n$th power proves \eqref{eq:nef-main-modification-comparison} with $C=c^{-1}$.

Dividing the logarithm of \eqref{eq:nef-main-modification-comparison} by $\log\epsilon<0$ reverses the inequality. The bounded error $\log C$ and the rescaling $\epsilon\mapsto K\epsilon$ do not affect lower and upper limits, as in the proof of \cref{prop:nef-degeneration-orders}. This gives $\underline a(\pi^*\theta)\geq\underline a(\theta)$ and $\overline a(\pi^*\theta)\geq\overline a(\theta)$. Finally, a proper holomorphic map of degree one between compact Riemann surfaces is an isomorphism, which settles the case $n=1$.
\end{proof}

\subsection{Strict changes under one blowup}\label{sec:nef-strict-blowups}

A blowup changes the amount of smooth volume near the vanishing set of a coefficient. A product equation turns that change into a strict increase of the degeneration order. Let $B=\mathbb P^m$, $m\geq2$, and normalize the Fubini--Study form $H$ by $\int_BH^m=1$. Let $\pi:\widetilde B\to B$ be the blowup of the affine origin. The incidence projection $\rho:\widetilde B\to\mathbb P^{m-1}$ gives $L=\rho^*\omega_{\rm FS}$; we also write $H$ for $\pi^*H$. The incidence embedding in $\mathbb P^m\times\mathbb P^{m-1}$ makes $\kappa=H+L$ K\"ahler.

\begin{proposition}\label{prop:nef-strict-blowup-orders}\index{Nef degeneration!strict blowup jump}
For every choice of integers $m\geq2$ and $p,k\geq1$, there is a single smooth-center blowup of a projective $(m+p)$-fold and a smooth semipositive nonclosed background whose common orders downstairs and after pullback are
\[
 a_{\rm down}=\max\{p-m/k,0\},\qquad a_{\rm up}=p-1/k.
\]
The orders exist on both sides. The increase can be any positive integer in varying dimension. The same family can lose positive limiting volume while both orders remain zero.
\end{proposition}

\begin{proof}
Put $W_0=\mathbb P^p$, with normalized Fubini--Study form $\omega$ satisfying $\int_{W_0}\omega^p=1$. Write $s=|z|^2$ on the affine chart of $B$. Choose a smooth $Q(s)\geq0$ equal to $s^k$ near zero, positive for $s>0$, and constant and positive for large $s$; it defines a smooth function $q$ on $B$. Put $n=m+p$, choose $b>0$, and define
\[
 \theta=bH+q\omega\quad\text{on }B\times W_0,\qquad
 \widetilde\theta=bH+(q\circ\pi)\omega
 \quad\text{on }\widetilde B\times W_0.
\]
The map $\Pi=\pi\times\mathrm{id}$ is the single blowup of the smooth center $\{0\}\times W_0$, and $\widetilde\theta=\Pi^*\theta$. Fix the K\"ahler directions $h=H+\omega$, $\widetilde h=\kappa+\omega$ and smooth positive probability densities $\mu=H^m\wedge\omega^p$, $\widetilde\mu=\kappa^m\wedge\omega^p/(2^m-1)$. Both backgrounds are smooth semipositive and nonclosed.

Write $\nu_{\rm down}=H^m$ and $\nu_{\rm up}=\kappa^m/(2^m-1)$ for the base probability forms, whose normalization is checked below. For either product write its perturbation as $\xi=A+g\omega$. Downstairs $A=(b+\epsilon)H$; upstairs $A=(b+\epsilon)H+\epsilon L$; in both cases $g=q+\epsilon$, with pullback understood upstairs. The form $A$ is K\"ahler on the relevant base. Use $\nu$ for the corresponding base form, and put $M=\int A^m$, $J=\int g^{-p}\nu$. The smooth equation \cref{thm:tw-smooth} and closed Stokes give $u$ with
\[
 \alpha=A+\ddc u>0,\qquad \alpha^m=(M/J)g^{-p}\nu.
\]
\Needspace{7\baselineskip}
The radial cumulative masses of the base forms are
\begin{equation}\label{eq:nef-radial-reference-masses}
 \begin{split}
 \Xi_{\rm down}(s)&=\left(\frac{s}{1+s}\right)^m,\\
 \Xi_{\rm up}(s)&=\frac{\left(1+\frac{s}{1+s}\right)^m-1}{2^m-1}.
 \end{split}
\end{equation}
To check these formulas and the constants $M$, write $t=\log|z|^2$ on the punctured affine chart. For a radial potential $F(t)$, direct polar integration of its Hessian gives
\[
 \int_{s_1<|z|^2<s_2}(\ddc F(\log|z|^2))^m
 =F'(\log s_2)^m-F'(\log s_1)^m.
\]
The potentials of $H$ and $L$ are $\log(1+\mathrm{e}^t)$ and $t$. Thus $xH+yL$ has endpoint slopes $y$ and $x+y$, and its total mass is $(x+y)^m-y^m$. The cases $(x,y)=(1,0),(1,1)$ give \eqref{eq:nef-radial-reference-masses}; the exceptional divisor and hyperplane at infinity have zero smooth volume. The form
\[
 \Omega=\alpha+g\omega=\xi+\ddc u>0,\qquad
\Omega^n=\binom{n}{m}\frac{M}{J}\nu\wedge\omega^p
\]
has constant volume density. The lower bound for all tests in \cref{thm:gaud-volume-dual-formula} is attained by
\[
 \gamma_*=g^{-p/(n-1)}\Omega.
\]
Indeed its $(n-1)$st power is
\begin{equation}\label{eq:nef-balanced-product-power}
 \gamma_*^{n-1}
 =\binom{n-1}{m}g^{-1}\alpha^m\wedge\omega^{p-1}
 +\binom{n-1}{m-1}\alpha^{m-1}\wedge\omega^p.
\end{equation}
The first coefficient multiplies a top-degree form on the base, and every factor in the second term is closed; hence both terms are closed. Thus $\gamma_*$ is balanced and Gauduchon. \Cref{thm:gaud-volume-dual-formula} gives the exact identities
\begin{equation}\label{eq:nef-strict-blowup-volumes}
 \mathcal V_\mu(\theta+\epsilon h)
 =\binom{n}{m}\frac{(b+\epsilon)^m}{J_{\rm down}},
 \qquad
 \mathcal V_{\widetilde\mu}(\widetilde\theta+\epsilon\widetilde h)
 =\binom{n}{m}
 \frac{(b+2\epsilon)^m-\epsilon^m}{J_{\rm up}},
\end{equation}
where $J_{\rm down}=\int_B(q+\epsilon)^{-p}\nu_{\rm down}$ and $J_{\rm up}=\int_{\widetilde B}(q\circ\pi+\epsilon)^{-p}\nu_{\rm up}$.

Near zero, differentiation of \eqref{eq:nef-radial-reference-masses} gives
\begin{equation}\label{eq:nef-radial-measure-asymptotics}
 \mathrm{d}\Xi_{\rm down}=(m s^{m-1}+O(s^m))\,\mathrm{d} s,\qquad
 \mathrm{d}\Xi_{\rm up}=\left(\frac{m}{2^m-1}+O(s)\right)\mathrm{d} s.
\end{equation}
\Needspace{10\baselineskip}
For a density $(a s^{d-1}+O(s^d))\,\mathrm{d} s$ and $q=s^k$, the singular integral has the following leading behavior as $\epsilon\downarrow0$:
\begin{equation}\label{eq:nef-radial-singular-integral}
 \begin{cases}
 \displaystyle\frac{a}{k}B(d/k,p-d/k)\epsilon^{d/k-p}
     &kp>d,\\[4pt]
 \displaystyle\frac{a}{k}\log(1/\epsilon)&kp=d,\\[4pt]
 \text{a finite positive limit}&kp<d.
 \end{cases}
\end{equation}
Here $B(u,v)=\int_0^\infty t^{u-1}(1+t)^{-u-v}\,\mathrm{d} t$. In the power case substitute $s=\epsilon^{1/k}v$, then $t=v^k$. At equality the integrand on $\epsilon^{1/k}\ll s\ll1$ is $a/s$. The $O(s^d)$ error is negligible: on a sufficiently small fixed interval its relative size is arbitrarily small, while the complementary integral is bounded. For $kp<d$, monotone convergence applies to the integrable limiting density. Since the numerators of \eqref{eq:nef-strict-blowup-volumes} tend to $\binom{n}{m}b^m>0$, formulas \eqref{eq:nef-radial-measure-asymptotics}--\eqref{eq:nef-radial-singular-integral}, with $d=m$ downstairs and $d=1$ upstairs, give the stated orders.

For $m=k=2$, $p=1$, the denominators satisfy $J_{\rm down}\sim\log(1/\epsilon)$ and $J_{\rm up}\sim(\pi/3)\epsilon^{-1/2}$; the volumes are respectively asymptotic to $3b^2/\log(1/\epsilon)$ and $(9b^2/\pi)\epsilon^{1/2}$. The orders change from $0$ to $1/2$. For $m=2$, $k=p=1$, the downstairs denominator has a finite positive limit, while $J_{\rm up}\sim(2/3)\log(1/\epsilon)$, so positive limiting volume downstairs becomes $9b^2/[2\log(1/\epsilon)]$ upstairs, with order zero on each side. For an arbitrary integer $j\geq1$, choose $m=j+2$, $p=j+1$, $k=1$. Then $p<m$ gives a positive downstairs limit, while the orders are $0$ and $j$.
\end{proof}

\subsection{Positive-current subtraction and divisorial residuals}\label{sec:nef-divisorial-residuals}

Another way to change the order of degeneration is to subtract a positive current from the background. If $\theta-\beta$ differs from a positive current by a $\ddc$-exact term, then every Gauduchon pairing of $\theta$ dominates the corresponding pairing of $\beta$. When the residual $\beta$ is still nef, its perturbed volumes are therefore smaller, and its orders are at least those of $\theta$. The main example is the subtraction of the curvature form of an effective divisor. Such a subtraction can preserve the order, raise it by an integer or by a fraction, or produce unequal lower and upper orders; and the residual need not be nef.

\Needspace{15\baselineskip}
\begin{proposition}\label{prop:nef-positive-current-subtraction}\index{Nef degeneration!positive-current subtraction}
Let $\theta,\beta$ be smooth nef real $(1,1)$-forms on a compact complex manifold $X$ of dimension $n$, and suppose that
\[
 S=(\theta-\beta)+\ddc\varphi\geq0
\]
is a positive $(1,1)$-current for some integrable function $\varphi$. For every Hermitian form $h$, every smooth strictly positive probability form $\mu$ and every $\epsilon>0$,
\begin{equation}\label{eq:nef-residual-volume-comparison}
 \mathcal V_\mu(\beta+\epsilon h)\leq \mathcal V_\mu(\theta+\epsilon h).
\end{equation}
Consequently $\underline a(\beta)\geq\underline a(\theta)$ and $\overline a(\beta)\geq\overline a(\theta)$.
\end{proposition}

\begin{proof}
For every Gauduchon metric $\gamma$, distributional integration by parts gives
\[
 \int_X(\theta-\beta)\wedge\gamma^{n-1}
 =\langle S,\gamma^{n-1}\rangle\geq0.
\]
Hence, for each test metric, the quotient for $\theta+\epsilon h$ is at least the quotient for $\beta+\epsilon h$, and taking infima gives $\lambda_\mu(\beta+\epsilon h)\leq\lambda_\mu(\theta+\epsilon h)$. Both root volumes are nonnegative, so raising to the $n$th power gives \eqref{eq:nef-residual-volume-comparison}. Dividing the logarithms by the negative number $\log\epsilon$ reverses the inequality and proves both comparisons of orders.
\end{proof}

For an effective real divisor $D$ with curvature form $\eta_D$, the Poincar\'e--Lelong formula gives $[D]=\eta_D+\ddc\varphi$ for a quasi-plurisubharmonic function $\varphi$. Thus \cref{prop:nef-positive-current-subtraction} applies to $\beta=\theta-\eta_D$, with $S=[D]$, whenever $\beta$ is nef. If $T=\theta+\ddc u$ is a positive current with $T-S\geq0$, then $T-S=\beta+\ddc(u-\varphi)$, so $\beta$ is the background of the positive current that remains after $S$ is removed from $T$.

\begin{proposition}\label{prop:nef-specified-residual-examples}
Subtracting the curvature form of an effective divisor from a smooth nef background, with a smooth nef residual, can preserve the common order, raise it by any integer between $1$ and the dimension or by a fraction, or produce a residual whose lower and upper orders differ. The residual can also fail to be nef.
\end{proposition}

\begin{proof}
\emph{Integer jumps.} On $E^n$, where $E$ is an elliptic curve, let $\omega_i$ be the pullback of a flat area form of integral one from the $i$th factor, and set $h=\sum_i\omega_i$ and $\mu=\omega_1\wedge\cdots\wedge\omega_n$. The divisor $F_i$ obtained by fixing a point in the $i$th factor has $\omega_i$ as a smooth curvature form. Indeed, on that elliptic curve any curvature form of the corresponding line bundle differs from $\omega_i$ by $\ddc$ of a solution of a zero-mean Poisson equation, so the metric can be adjusted on the curve and then pulled back. For $\theta=h$ and $D=\sum_{i=1}^kF_i$, the residual $\beta=\sum_{i=k+1}^n\omega_i$ is smooth and nef. The perturbed forms are closed, positive and of constant density with respect to $\mu$, so the zero potential solves their equations, and \cref{thm:gaud-volume-dual-formula} gives
\begin{equation}\label{eq:nef-integer-residual-jump}
 \mathcal V_\mu(\theta+\epsilon h)=n!(1+\epsilon)^n,\qquad
 \mathcal V_\mu(\beta+\epsilon h)=n!\epsilon^k(1+\epsilon)^{n-k}.
\end{equation}
These examples realize each integer jump $1,\ldots,n$. A nonzero divisor with jump zero is given by $\theta=2h$ and $D=F_1$, since then the residual is still strictly positive.

\emph{A fractional jump.} Use the background $\beta=b\omega_1+q\omega_2$ of \cref{prop:nef-diagonal-fractional-and-zero}, where $q=|\zeta|^{2m}$ near its unique zero and $m>1$. Choose $c>0$ and put $\theta=\beta+c\omega_2>0$ and $D=c(E\times\{w_0\})$, whose smooth curvature form is $c\omega_2$. The original background is strictly positive, so both its orders are zero, whereas the residual has common order $1-1/m$.

\emph{Unequal orders.} Let $\beta_{\mathrm{res}}$ be the oscillatory rank-one background of \cref{prop:nef-oscillatory-orders}, and put
\[
 \theta=\beta_{\mathrm{res}}+ch>0,\qquad
 D=c(\{z_0\}\times E)+c(E\times\{w_0\}),\qquad
 \eta_D=ch.
\]
The background $\theta$ has both orders zero: since $\theta\geq ch$, every test quotient is bounded below by a positive constant, while one fixed test gives a finite upper bound. Its smooth nef residual is exactly $\beta_{\mathrm{res}}$, whose lower and upper orders differ. In each of these examples, the residual plus $[D]$ is a positive current in the smooth $\ddc$-class of $\theta$, and its divisorial part is exactly $[D]$ because the residual is smooth. We do not need this current to have minimal singularities; that case is treated in \cref{cor:nef-boundary-minimal-removal}.

\emph{A residual that is not nef.} Blow up $\mathbb P^2$ at two distinct points. Let $H$ be the pullback of the class of the normalized Fubini--Study form, let $E_1,E_2$ be the exceptional curves, and let $L$ be the strict transform of the line through the two points. The pullback $\theta$ of the Fubini--Study form is smooth, semipositive and nef. Subtract curvature forms of $D=E_1+E_2$. The residual $\beta$ has class $H-E_1-E_2$, which is the class of $L$, and the Poincar\'e--Lelong formula for the total transform of the line shows that it contains the positive current $[L]$. The intersection numbers $H^2=1$, $H\cdot E_i=0$, $E_i^2=-1$ and $E_1\cdot E_2=0$ give
\[
 \int_L\beta=(H-E_1-E_2)^2=-1.
\]
If $\beta$ were nef, then, for a Hermitian form $h$, integrating smooth positive representatives $\beta+\epsilon h+\ddc u_\epsilon$ over $L$ would give $-1+\epsilon\int_Lh\geq0$ for all $\epsilon>0$, which is impossible.
\end{proof}

For a prime divisor $D$ and $\delta>0$, consider the infimum of the generic Lelong numbers along $D$ of the positive currents $\theta+\delta h+\ddc u$; its limit as $\delta\downarrow0$ is the minimal multiplicity of $\theta$ along $D$. For a nef background all minimal multiplicities vanish, since every $\theta+\delta h$ has a smooth positive representative. Removing the divisor defined by these multiplicities therefore leaves a nef background unchanged. Removing the divisorial part of a particular positive current in the class of $\theta$ is a different operation, even when that current has minimal singularities.

\begin{corollary}\label{cor:nef-boundary-minimal-removal}
On the ruled surface $X$ of \cref{lem:ruled-explicit-nef}, keep its nef curvature form $\theta$, its section $D$ at infinity and its K\"ahler form $\kappa=\theta+\pi^*\sigma_0$. For every smooth positive probability form $\mu$,
\[
 \mathcal V_\mu(\theta+\epsilon\kappa)=2\epsilon(1+\epsilon),
 \qquad \mathcal V_\mu(\epsilon\kappa)=2\epsilon^2.
\]
The only positive current of the form $\theta+\ddc u$ is $[D]$. It has minimal singularities, and removing its whole divisor leaves the zero background and changes the common order from one to two.
\end{corollary}
\begin{proof}
\Cref{lem:ruled-explicit-nef} gives the uniqueness of $[D]$ and the nefness of $\theta$, and the calculation in the proof of \cref{prop:nonrecovery-zero-residual-volume} gives $\int_X\theta^2=0$. Put $B=\pi^*\sigma_0$. Then $B^2=0$, and
\[
 \int_X\theta\wedge B=\int_D B=1,
\]
by the Poincar\'e--Lelong formula, Stokes' theorem and the fact that the section $D$ maps isomorphically onto the elliptic base. The closed nef formula following \cref{thm:gaud-volume-nef-limit} now computes both volumes by expanding the squares, for every $\mu$. Every $\theta$-plurisubharmonic function $u$ satisfies $\theta+\ddc u=[D]$, so any two such functions differ by a pluriharmonic function, hence by a constant; in particular each has minimal singularities. Since $\theta$ is itself a curvature form of $\mathcal O_X(D)$, removing the whole divisor leaves the zero background. The two volume polynomials give the orders one and two.
\end{proof}

\subsection{Integer orders in fixed projective cones}
\label{sec:nef-projective-integer-cones}
\index{Nef degeneration!nonclosed projective cones}

The variable-rank examples show that nefness alone does not force an integral order. Strictly positive coefficients on fixed projective generators give a different regime. A logarithmic divisor potential compares their full Gauduchon quotient with a closed form, even when the background is nonclosed and a pullback loses rank along an exceptional divisor.

\Needspace{14\baselineskip}
\begin{lemma}\label{lem:nef-logarithmic-rank-comparison}
Let $X$ be a compact K\"ahler $n$-fold, $\Omega$ a K\"ahler form, and $B$ a smooth closed semipositive form. Suppose a smooth nef background $\theta$ satisfies
\[
 \theta\leq CB,\qquad \theta+\ddc\varphi\geq cB
\]
as currents, for $c,C>0$ and a global real $\varphi\in L^1(X)$. For every Hermitian $h$ and smooth strictly positive probability $\mu$, there are $k,K>0$ such that, for $0<\epsilon\leq1$,
\begin{equation}\label{eq:nef-logarithmic-quotient-comparison}
 k\lambda_\mu(B+\epsilon\Omega)
 \leq\lambda_\mu(\theta+\epsilon h)
 \leq K\lambda_\mu(B+\epsilon\Omega).
\end{equation}
If $B$ has constant rank $R$, then for $L,U>0$,
\[
 L\epsilon^{n-R}\leq\mathcal V_\mu(\theta+\epsilon h)
                       \leq U\epsilon^{n-R}.
\]
In particular the common order is $n-R$.
\end{lemma}
\begin{proof}
Choose $c_h,C_h>0$ with $c_h\Omega\leq h\leq C_h\Omega$. For every smooth positive Gauduchon test $\gamma$, distributional Stokes gives
\[
 \int_X\ddc\varphi\wedge\gamma^{n-1}
       =\langle\varphi,\ddc\gamma^{n-1}\rangle=0.
\]
The lower current bound therefore gives the lower numerator inequality with $k=\min(c,c_h)$; the upper bound gives the upper one with $K=\max(C,C_h)$. Divide by the same strictly positive $D_\mu(\gamma)$ and take infima over all Gauduchon tests. This proves \eqref{eq:nef-logarithmic-quotient-comparison} without a rank hypothesis.

When $B$ has constant rank, its $R$ positive eigenvalues relative to $\Omega$ are uniformly bounded above and away from zero. With $A_\epsilon=B+\epsilon\Omega$, compactness gives
\begin{equation}\label{eq:nef-projective-pointwise-volume}
 l\epsilon^{n-R}\mu\leq A_\epsilon^n
                         \leq u\epsilon^{n-R}\mu
\end{equation}
for $l,u>0$. Matrix arithmetic--geometric mean gives, for every test,
\[
 A_\epsilon\wedge\gamma^{n-1}
 \geq\left(\frac{A_\epsilon^n}{\gamma^n}\right)^{1/n}\gamma^n.
\]
Thus $\lambda_\mu(A_\epsilon)\geq l^{1/n}\epsilon^{(n-R)/n}$. The positive closed form $\gamma=A_\epsilon$ is itself an admissible test. Writing $A_\epsilon^n=f_\epsilon\mu$ gives
\[
 \begin{aligned}
 D_\mu(A_\epsilon)&=\int_X f_\epsilon^{(n-1)/n}\mu\\
 &\geq l^{(n-1)/n}\epsilon^{(n-R)(n-1)/n},\\
 \int_X A_\epsilon^n&\leq u\epsilon^{n-R}.
 \end{aligned}
\]
Its quotient is at most $ul^{-(n-1)/n}\epsilon^{(n-R)/n}$. Raise both bounds to the $n$th power and apply the comparison. The calculation includes $R=0$ and $n=1$, where $D_\mu=1$. Fixed positive factors vanish after division of their logarithms by $\log\epsilon$.
\end{proof}

\subsubsection{Linear blowups and the fixed generators}
For a plain factor $Z_i=\mathbb P^{m_i}$, $m_i\geq1$, let $H_i$ be normalized Fubini--Study and set $\kappa_i=H_i$. A blown factor has the form
\[
 Z_i=\operatorname{Bl}_{\mathbb P^{\ell_i}}\mathbb P^{m_i},
 \qquad m_i\geq2,\quad 0\leq\ell_i\leq m_i-2,
 \qquad r_i=m_i-\ell_i-1.
\]
Suppress the index temporarily. In coordinates $[z_0:\cdots:z_\ell:w_0:\cdots:w_r]$, its incidence model is
\[
 Z=\{(x,[t])\in\mathbb P^m\times\mathbb P^r:
                         w_at_b=w_bt_a\text{ for all }a,b\}.
\]
Let $p:Z\to\mathbb P^m$ and $q:Z\to\mathbb P^r$ be the projections. In a chart $t_j=1$, the equations $w_a=t_aw_j$ are the charts of the blowup of $(w_0,\ldots,w_r)$. Writing $w=\lambda t$ identifies $q$ locally as a $\mathbb P^{\ell+1}$-bundle, with fiber coordinates $[z:\lambda]$. Thus $Z$ is smooth, $q$ is a submersion, and the incidence inclusion is an embedding. Set
\[
 H=p^*\omega_m,\qquad Q=q^*\omega_r,\qquad \kappa=H+Q.
\]
The form $\kappa$ is K\"ahler as the restricted product form, whereas $Q$ has constant rank $r$.

Let $E=\{w=0\}$ be the exceptional divisor. The coordinate section $w_j=\lambda t_j$ has divisor $E+q^{-1}\{t_j=0\}$, so $\mathcal O(E)=p^*\mathcal O(1)\otimes q^*\mathcal O(-1)$. Give it the product metric of curvature $H-Q$. For its canonical section, $\psi_E=\log|s_E|^2$ is locally integrable and the curvature identity of \cref{sec:divisorial-gaud-volume} gives
\begin{equation}\label{eq:nef-projective-exceptional-potential}
 \ddc\psi_E=[E]-H+Q,\qquad
 H+\ddc(\psi_E/2)=\tfrac12(\kappa+[E])\geq\kappa/2.
\end{equation}

\Needspace{9\baselineskip}
On $X=\prod_i Z_i$, put $n=\sum_i m_i$ and pull back all factor forms. Fix smooth strictly positive functions $f_i$ and, on blown factors, $g_i$ on $X$. These functions may depend on every product coordinate. The cone consists of
\begin{equation}\label{eq:nef-fixed-projective-cone}
 \theta=\sum_i\bigl((a_i+b_if_i)H_i+(c_i+d_ig_i)Q_i\bigr),
 \qquad a_i,b_i,c_i,d_i\geq0,
\end{equation}
with the $Q_i$ terms and their parameters omitted on plain factors. All coefficient functions and generators are fixed. Every member is smooth semipositive, hence nef for every Hermitian perturbation direction. Call $H_i$ active if $a_i+b_i>0$, and $Q_i$ active if $c_i+d_i>0$. Define
\begin{equation}\label{eq:nef-projective-effective-rank}
 R_i(\theta)=
 \begin{cases}
 m_i,&H_i\text{ active},\\
 r_i,&H_i\text{ inactive and }Q_i\text{ active},\\
 0,&\text{both inactive},
 \end{cases}
 \qquad R(\theta)=\sum_i R_i(\theta).
\end{equation}
The middle case occurs only on blown factors.

\Needspace{9\baselineskip}
\begin{theorem}\label{thm:nef-projective-integer-orders}
Every fixed-cone member in \eqref{eq:nef-fixed-projective-cone} has common order $n-R(\theta)$. For arbitrary smooth positive $\mu$ and $h$, its variational volume has two-sided bounds for $0<\epsilon\leq1$:
\begin{equation}\label{eq:nef-projective-power-bounds}
 L\epsilon^{n-R(\theta)}
 \leq\mathcal V_\mu(\theta+\epsilon h)
 \leq U\epsilon^{n-R(\theta)},
\end{equation}
with $L,U>0$ depending on the fixed background and data.
\end{theorem}
\begin{proof}
Let $B_\theta$ be the sum of $\kappa_i$ on $H_i$-active blocks and $Q_i$ on $Q_i$-only blocks, with zero on inactive blocks. It is closed semipositive of constant rank $R(\theta)$. Bounded coefficients and $H_i,Q_i\leq\kappa_i$ give $\theta\leq CB_\theta$.

For every active blown $H_i$, let $u_i=\min_X(a_i+b_if_i)>0$ and put
\[
 \varphi=\sum_{\substack{i\text{ blown}\\H_i\text{ active}}}
                          (u_i/2)\psi_{E_i}.
\]
By \eqref{eq:nef-projective-exceptional-potential}, $u_iH_i$ becomes at least $(u_i/2)\kappa_i$, with a nonnegative divisor current. The remaining $(a_i+b_if_i-u_i)H_i$ and all $Q_i$ terms are nonnegative. On plain active blocks and $Q_i$-only blocks, use the positive minimum of their active coefficient. Taking the minimum over the active blocks gives $\theta+\ddc\varphi\geq cB_\theta$. For the zero member take $\varphi=0$ and any $c>0$. Apply \cref{lem:nef-logarithmic-rank-comparison} with $\Omega=\sum_i\kappa_i$.
\end{proof}

Nonclosed members occur whenever a fixed generator derivative is nonzero, since
\[
 \mathrm d\theta=\sum_i b_i\,\mathrm df_i\wedge H_i
                  +\sum_{i\text{ blown}}d_i\,\mathrm dg_i\wedge Q_i.
\]
Activating that generator alone prevents cancellation. On a blown factor, a small positive bump perturbation of one varying in a $q$-fiber has $\mathrm dg\wedge Q\ne0$, since $r\geq1$. On $\mathbb P^p\times\mathbb P^q$, $p,q\geq1$, a positive $f$ varying on the second factor gives nonclosed $fH_1$ of order $q$. Adding a smooth $\ddc$ potential preserves the exterior derivative.

\subsubsection{Modification equality}
\Needspace{11\baselineskip}
\begin{proposition}\label{prop:nef-projective-modification-equality}
Every product modification obtained by replacing plain factors with the indicated linear blowups preserves the order of a fixed-cone member. More generally, for every modification $\pi:Y\to X$ with $Y$ a compact K\"ahler manifold, the pullback has common order $n-R(\theta)$ and two-sided power bounds with that exponent, for arbitrary smooth positive probability and Hermitian perturbation on $Y$.
\end{proposition}
\begin{proof}
For product linear blowups, pull back the old coefficient functions. The new $Q_i$ coefficients are zero. Active plain $H_i$ become active blown $H_i$ of the same effective rank $m_i$, and inactive blocks stay inactive. Apply \cref{thm:nef-projective-integer-orders} on both products.

For general K\"ahler $Y$, pull back the explicit nonnegative smooth remainders and divisor-section identities above. Each nonzero divisor section remains nonzero under the surjective modification. Its pulled-back squared logarithm is locally integrable and gives the curvature of its effective zero divisor, so
\[
 \pi^*\theta\leq C\pi^*B_\theta,\qquad
 \pi^*\theta+\ddc(\varphi\circ\pi)\geq c\pi^*B_\theta.
\]
Use the rank-free part of \cref{lem:nef-logarithmic-rank-comparison} with a K\"ahler form $\Omega_Y$. The pulled-back background is semipositive and hence nef. Put $R=R(\theta)$. Pullback cannot increase rank, so $(\pi^*B_\theta)^{R+1}=0$, while
\[
 \int_Y(\pi^*B_\theta)^R\wedge\Omega_Y^{n-R}>0.
\]
For $R>0$ this is positive on the nonempty isomorphism locus; for $R=0$ it is the K\"ahler volume. The closed-background formula following \cref{thm:gaud-volume-nef-limit} gives, for the chosen probability $\mu_Y$,
\[
 \mathcal V_{\mu_Y}(\pi^*B_\theta+\epsilon\Omega_Y)
 =\sum_{k=0}^{R}\binom nk\epsilon^{n-k}
       \int_Y(\pi^*B_\theta)^k\wedge\Omega_Y^{n-k}.
\]
Every coefficient is nonnegative and the coefficient at $k=R$ is positive. Thus this polynomial has two-sided bounds by $\epsilon^{n-R}$ for $0<\epsilon\leq1$. The quotient comparison gives the same bounds for $\pi^*\theta+\epsilon h_Y$. This also proves equality with the downstairs order. No pullback of a positive probability is assumed to remain strictly positive on the exceptional locus.
\end{proof}

\Needspace{13\baselineskip}
\subsubsection{Specified effective divisors}
Fix the pullback Fubini--Study metrics on $\mathcal O(H_i)$ and $\mathcal O(Q_i)$. Allow tensor products of their integer powers, including negative powers. For finite nonnegative real sums of zero divisors of nonzero holomorphic sections of these bundles, write
\[
 \begin{gathered}
 D=\sum_\nu t_\nu\operatorname{div}(s_\nu),\quad t_\nu\geq0,\\
 \eta_D=\sum_i(\delta_iH_i+\rho_iQ_i),\qquad
 \Phi_D=\sum_\nu t_\nu\log|s_\nu|^2.
 \end{gathered}
\]
The metrics fix the representatives, and $[D]=\eta_D+\ddc\Phi_D$. Omit $Q_i$ terms on plain factors.

\Needspace{8\baselineskip}
\begin{lemma}\label{lem:nef-projective-effective-divisor-range}
The exact coefficient range of this real divisor family is
\[
 \delta_i\geq0\text{ on plain factors};\qquad
 \delta_i\geq0,\quad\delta_i+\rho_i\geq0
                         \text{ on blown factors}.
\]
\end{lemma}
\begin{proof}
For one nonzero holomorphic section $s$, choose a point where it is nonzero and outside all exceptional divisors, then fix the other factor coordinates. A projective line through that point in a $q_i$-fiber has $H_i$-degree one and $Q_i$-degree zero. The restriction of $s$ is nonzero, so its $H_i$ exponent is nonnegative. On a plain factor the line argument gives the same inequality.

On a blown factor, write the chosen point as $([z:w],[w])$ with $w\ne0$, and choose a linear map $A$ with $Aw=z$. The section of $q_i$ given by $[t]\mapsto([At:t],[t])$ contains the point. Both $H_i$ and $Q_i$ have degree one on any of its projective lines. Such a line exists because $r_i\geq1$, and the original holomorphic section $s$ restricts nontrivially to it. Its degree gives the nonnegative sum of the two exponents. Nonnegative real sums preserve these inequalities.

Conversely, let $F_i$ be a $Q_i$-hyperplane divisor. A permitted pair is represented by $\delta_iE_i+(\delta_i+\rho_i)F_i$, since $E_i$ has curvature $H_i-Q_i$. Plain coefficients are represented by hyperplanes. These are allowed effective real sums and give every coefficient tuple in the stated range.
\end{proof}

\Needspace{13\baselineskip}
\begin{theorem}\label{thm:nef-projective-residual-equality}
Let $\theta$ have the fixed-cone parameters above, and let $D$ be in the specified divisor family. Assume
\[
 a_i\geq\delta_i\text{ for all }i,\qquad
 c_i\geq\rho_i\text{ for every blown }i.
\]
Then $\beta=\theta-\eta_D$ is in the same cone, and its common order satisfies
\[
 a(\beta)-a(\theta)=\sum_i\bigl(R_i(\theta)-R_i(\beta)\bigr)\geq0.
\]
Equality holds if and only if
\begin{align*}
 a_i-\delta_i+b_i&>0&&\text{whenever }a_i+b_i>0,\\
 c_i-\rho_i+d_i&>0&&\text{whenever }a_i=b_i=0,\quad c_i+d_i>0.
\end{align*}
Here $a$ denotes the common order; the second condition concerns blown factors only.
\end{theorem}
\begin{proof}
The residual parameters $(a_i-\delta_i,b_i,c_i-\rho_i,d_i)$ are nonnegative, so $\beta$ is smooth semipositive and nef. Also
\[
 (\theta-\beta)+\ddc\Phi_D=[D]\geq0,
\]
which is the actual positive-current comparison of \cref{prop:nef-positive-current-subtraction}. Apply \cref{thm:nef-projective-integer-orders} to both members. If an active $H_i$ disappears, its rank drops from $m_i$ to either $r_i<m_i$ or zero. If $H_i$ was absent, then $a_i=b_i=0$; the divisor range and residual conditions force $\delta_i=0$ and $\rho_i\geq0$. Such a block cannot gain $Q_i$. Therefore each rank loss is nonnegative and the losses cannot cancel. Their sum vanishes precisely when every initially leading generator remains active, giving the two conditions.
\end{proof}

The parameter inequalities ensure that this displayed representation stays in the fixed cone. They do not characterize every nef residual or every alternative representation by its generators.

\Needspace{10\baselineskip}
\begin{example}\label{ex:nef-nonclosed-projective-order-jumps}
On $Z=\operatorname{Bl}_{\mathbb P^\ell}\mathbb P^m$, choose $g>0$ with $\mathrm dg\wedge Q\ne0$. For $t>0$, the genuinely nonclosed positive background $\theta_t=tH+gQ$ has order zero. Removing $sE$, with its fixed curvature $s(H-Q)$ and $0\leq s\leq t$, gives
\[
 \beta_s=(t-s)H+(g+s)Q,\qquad
 \theta_t+\ddc(s\psi_E)=s[E]+\beta_s.
\]
Every residual is nef and nonclosed. Equality of orders holds exactly for $s<t$; at $s=t$ the jump is $\ell+1$. These examples realize each integer jump $1,\ldots,m-1$.
\end{example}
\begin{proof}
The lower bound $\theta_t\geq\min(t,\min g)\kappa>0$ verifies positivity. For $s<t$, $H$ is active and the effective rank is $m$; at $s=t$ only $Q$ is active, of rank $r=m-\ell-1$. The order formula gives zero and $m-r=\ell+1$, respectively. The exterior derivative of every residual is $\mathrm dg\wedge Q\ne0$. Vary $\ell$ from zero to $m-2$.
\end{proof}

At the endpoint the chosen current contains the actual divisor $t[E]$ before removal. The perturbative minimal divisorial multiplicities are nevertheless zero, since all residual backgrounds admit smooth positive perturbations. Strict positivity of the fixed coefficient functions, the specified cone-preserving divisor inequalities and a K\"ahler source for the general modification statement are essential hypotheses of this construction. The general problem in \cref{qu:nef-degeneration-exponent} remains open.

\subsection{A maximum-rank comparison criterion}
\label{sec:nef-maximum-rank-comparison}
\index{Nef degeneration!maximum-rank comparison}

The constant-rank hypothesis in \cref{lem:nef-logarithmic-rank-comparison} can be removed. For a closed semipositive form $B$, the intersection polynomial $\int_X(B+\epsilon\Omega)^n$ has a leading term of order $\epsilon^{n-R}$, where $R$ is the largest rank that $B$ attains anywhere on $X$. Combined with the comparison of quotients, this gives the order $n-R$ for every background that is comparable to $B$. The criterion applies to backgrounds that vanish on an open set, and to positive variable multiples of arbitrary closed semipositive generators.

\Needspace{12\baselineskip}
\begin{proposition}\label{prop:nef-maximum-rank-comparison}
Let $X$ be a compact K\"ahler manifold of dimension $n$, let $B$ be a smooth closed semipositive $(1,1)$-form, and put $R=\max_X\operatorname{rank}B$. Suppose a smooth nef background $\theta$ satisfies
\[
 \theta\leq CB,\qquad \theta+\ddc\varphi\geq cB
\]
as currents, for constants $c,C>0$ and a real function $\varphi\in L^1(X)$. For every smooth positive probability form $\mu$ and Hermitian form $h$, there are constants $L,U>0$ such that
\[
 L\epsilon^{n-R}\leq\mathcal V_\mu(\theta+\epsilon h)
                       \leq U\epsilon^{n-R},\qquad 0<\epsilon\leq1.
\]
Thus $\theta$ has common order $n-R$, without any constant-rank hypothesis.
\end{proposition}
\begin{proof}
Fix a K\"ahler form $\Omega$. The first part of \cref{lem:nef-logarithmic-rank-comparison}, which does not use the rank, shows that $\lambda_\mu(\theta+\epsilon h)$ and $\lambda_\mu(B+\epsilon\Omega)$ differ by bounded positive factors. Since $B^{R+1}=0$, the closed nef formula following \cref{thm:gaud-volume-nef-limit} gives
\[
 \mathcal V_\mu(B+\epsilon\Omega)
 =\sum_{k=0}^{R}\binom nk\epsilon^{n-k}
                         \int_X B^k\wedge\Omega^{n-k}.
\]
All coefficients are nonnegative. If $R>0$, then at a point where $B$ has rank $R$, some $R\times R$ minor of its coefficient matrix is positive and stays positive nearby, so $B^R\wedge\Omega^{n-R}$ has positive integral. If $R=0$, this integral is the K\"ahler volume. Hence the coefficient of $\epsilon^{n-R}$ is positive. For $0<\epsilon\leq1$, every other power of $\epsilon$ that occurs is at most $\epsilon^{n-R}$, so the polynomial is bounded below by its leading term and above by the sum of its coefficients times $\epsilon^{n-R}$. Raising the comparison of quotients to the $n$th power gives the asserted bounds, and taking logarithms gives the order.
\end{proof}

\Needspace{12\baselineskip}
\begin{corollary}\label{cor:nef-variable-semipositive-generators}
On a compact K\"ahler manifold $X$ of dimension $n$, fix finitely many smooth closed semipositive $(1,1)$-forms $B_j$ and smooth strictly positive functions $f_j$. For $t_j\geq0$, put
\[
 \theta=\sum_jt_jf_jB_j,\qquad
 B=\sum_{j:t_j>0}B_j,\qquad R=\max_X\operatorname{rank}B.
\]
The background $\theta$ is nef, has common order $n-R$, and satisfies the two-sided power bounds of \cref{prop:nef-maximum-rank-comparison}. For every modification $\pi:Y\to X$ with $Y$ a compact K\"ahler manifold, the pullback $\pi^*\theta$ has the same order and power bounds with the same exponent, for every smooth positive probability form and Hermitian perturbation on $Y$.
\end{corollary}
\begin{proof}
Semipositivity gives nefness. If some $t_j$ is positive, compactness and strict positivity of the functions $f_j$ give
\[
 \begin{gathered}
 cB\leq\theta\leq CB,\qquad
 c=\min_{j:t_j>0}\bigl(t_j\min_X f_j\bigr)>0,\\
 C=\max_{j:t_j>0}\bigl(t_j\max_X f_j\bigr),
 \end{gathered}
\]
and \cref{prop:nef-maximum-rank-comparison} applies with $\varphi=0$. If all parameters vanish, then $B=\theta=0$, $R=0$, any $c,C>0$ may be used, and the order is $n$.

These inequalities between smooth forms are preserved by pullback, and the rank of $\pi^*B$ is at most $R$. If $R>0$, the set where $B$ has rank $R$ is nonempty and open, so it meets the dense open set where $\pi$ is biholomorphic; hence the maximal rank of $\pi^*B$ is $R$. This is immediate if $R=0$. Since $\dim Y=n$, the same proposition applies on $Y$, with any choice of data there. No singular potential needs to be pulled back.
\end{proof}

The constants depend on the parameters $t_j$ and need not stay bounded as a positive parameter tends to zero. The rank of a sum records how the generators overlap, and it need not equal the sum of their ranks. Members of the family are not closed when $\sum_jt_j\,\mathrm df_j\wedge B_j\ne0$.

\Needspace{10\baselineskip}
\begin{example}\label{ex:nef-open-zero-rank-background}
Let $X=\mathbb P^1\times\mathbb P^q$, $q\geq1$, with projections $p_1,p_2$. Choose a smooth nonnegative function $b\not\equiv0$ on $\mathbb P^1$ that vanishes on a nonempty open set, and a smooth strictly positive nonconstant function $f$ on $\mathbb P^q$. With $\omega_1$ the normalized Fubini--Study form on $\mathbb P^1$, set
\[
 B=p_1^*(b\omega_1),\qquad \theta=(f\circ p_2)B.
\]
The smooth nef background $\theta$ is not closed, vanishes on a nonempty open set, and has common order $q$, with two-sided bounds by multiples of $\epsilon^q$.
\end{example}
\begin{proof}
The form $b\omega_1$ is closed, because the first factor is a curve, and semipositive. Its pullback has rank one where $b>0$ and rank zero where $b=0$, so its maximal rank is one. Since $n=q+1$, \cref{cor:nef-variable-semipositive-generators} gives the order $q$. A nonconstant smooth function $f$ on the connected second factor has $\mathrm df\ne0$ somewhere. At a point of the product where $b>0$ and $\mathrm df\neq0$, the horizontal two-form $B$ wedges nontrivially with the vertical one-form $p_2^*\mathrm df$. Thus $\mathrm d\theta=p_2^*\mathrm df\wedge B\ne0$.
\end{proof}

These are geometric sufficient conditions for an integral common order, in the direction of \cref{qu:nef-degeneration-exponent}. They do not cover backgrounds without a closed semipositive comparison form, coefficient functions that are allowed to vanish, or the comparison of orders under residuals or non-K\"ahler modifications in general.

\subsection{Flat anticanonical curves}\label{subsec:nef-flat-anticanonical-curves}

A flat neighborhood of a divisor gives a closed semipositive form whose curvature is confined to a thin annulus around the divisor. On a surface this form has rank at most one. The maximum-rank criterion then applies to its positive variable multiples, even when the divisor line has no moving positive multiple. We first construct the flat neighborhood for a square-zero anticanonical curve.

\begin{lemma}\label{lem:flat-elliptic-derivative-inverse}
Let $D$ be an elliptic curve with a flat Hermitian metric, let $\vartheta$ be a parallel unit $(1,0)$-form, and let $v$ be its dual vector field. For every unitary flat line bundle $L$ on $D$ and every smooth $L$-valued $(0,1)$-form $u$, there is a unique smooth section $f$ orthogonal to the parallel sections such that
\[
 \bar\partial_L f=\iota_v\partial_Lu.
\]
If $\Pi_L^{0,1}$ is the orthogonal projection onto the parallel $L$-valued $(0,1)$-forms, then, for every flat Sobolev norm $H^k$,
\[
 \|f\|_{H^k}=\|(I-\Pi_L^{0,1})u\|_{H^k}
 \leq\|u\|_{H^k}.
\]
The bound is independent of the holonomy of $L$.
\end{lemma}
\begin{proof}
Put $P=\iota_v\partial_L$. Flatness gives $P=\nabla_v$, $P^*=-\nabla_{\bar v}$, and $\Delta_L=P^*P=PP^*$ on functions and $(0,1)$-forms. Let $G_q$ be the inverse of $\Delta_L$ on the orthogonal complement of its kernel in degree $q$, extended by zero on that kernel. The operators $P$ and $P^*$ commute with $G_q$. Set
\[
 f=\bar\partial_L^*G_1Pu.
\]
The form $Pu$ is orthogonal to the parallel forms, so $\bar\partial_Lf=Pu$; also $f$ is orthogonal to the parallel sections. Uniqueness follows from the kernel of $\bar\partial_L$. Moreover,
\[
 \|f\|_{L^2}^2
 =\langle Pu,G_1Pu\rangle
 =\langle u,(I-\Pi_L^{0,1})u\rangle.
\]
Covariant differentiation along the global parallel real frame commutes with these operators. Applying the identity to every derivative proves the $H^k$ assertion. In particular, this argument never divides by a lower bound for a nonzero eigenvalue; the derivative in $Pu$ cancels that small eigenvalue.
\end{proof}

For the neighborhood argument, let $S$ be a complex surface, possibly noncompact, and let $D\subset S$ be a smooth compact connected anticanonical divisor with $D^2=0$. Adjunction makes $D$ elliptic. Its normal line $N$ has degree zero, hence admits a unitary flat metric. Let $\pi:N\to D$ be its projection, $N^\times$ its punctured total space, and $J_N$ its complex structure. The reciprocal anticanonical section gives a meromorphic two-form $\Omega_S$, nowhere zero on $S\setminus D$, with simple pole along $D$. Denote its nonzero residue by $\vartheta_D$.

Choose parallel unitary local frames of $N$ with fiber coordinate $w$. On $N^\times$ put
\[
 \xi=w\partial_w,\qquad \Theta=\frac{\mathrm dw}{w},\qquad
 \vartheta=\pi^*\vartheta_D,\qquad
 z=-\log|w|^2,\qquad \Omega_0=\Theta\wedge\vartheta.
\]
These objects are global. Give $D$ the flat metric making $\vartheta_D$ unit, and let $v$ be the parallel vector field dual to $\vartheta_D$. Equip $N^\times$ with the flat cylindrical metric $g$ determined by the parallel coframe $\vartheta,\Theta$, and write $\nabla$ for its connection. The end $z\to+\infty$ approaches the zero section.

\begin{lemma}\label{lem:anticanonical-cylindrical-identification}
There is a smooth tubular identification $\Upsilon$ from a neighborhood of the zero section in $N$ to a neighborhood of $D$ in $S$, fixing $D$, such that each fiber disc is holomorphic for $\Upsilon^*J_S$ and
\[
 |\nabla^k(\Upsilon^*\Omega_S-\Omega_0)|_g
 =O(\mathrm e^{-z/2})\qquad(k\geq0).
\]
\end{lemma}
\begin{proof}
Start with a normal exponential map for a smooth Hermitian metric near $D$. Its differential is complex linear along $D$ and induces the identity on the normal quotient. In a parallel local frame and a coordinate $z_0$ with $\vartheta_D=\mathrm dz_0$, choose holomorphic surface coordinates $Z,W$ with $D=\{W=0\}$. The pullbacks under this first map satisfy
\[
 Z=z_0+b(z_0,\bar z_0)w+O(|w|^2),\qquad
 W=w+O(|w|^2).
\]
The residue normalization writes $\Omega_S=(1+Wa(Z,W))\mathrm dW/W\wedge\mathrm dZ$ with $a$ holomorphic. Substitution gives $\Omega_S-\Omega_0=O(|w|)$ in the logarithmic coframe. The possible term containing $\Theta\wedge\mathrm dw$ vanishes. Every derivative by $\partial_{z_0},\partial_{\bar z_0},w\partial_w,\bar w\partial_{\bar w}$ preserves this bound, proving the stated cylindrical estimate for the first map.

It remains to make the fiber discs holomorphic without changing its first jet. On a fiber disc $\Delta_{p,\rho}\subset N_p$, vary the map by $\Phi_Y(q)=\exp_q(Y_q)$. Identify its target tangent spaces with a fixed complex vector space over $p$ by Chern parallel transport along the radial paths and the short exponential paths. In these identifications its antiholomorphic differential is
\[
 \mathcal F_{p,\rho}(Y)
 =\tfrac12\bigl(\mathrm d\Phi_Y+J_S\mathrm d\Phi_YJ_N\bigr)
       |_{T\Delta_{p,\rho}}.
\]
For a fixed Hölder exponent $0<\nu<1$, use rescaled $C^{k,\nu}$ norms on the discs. The Cauchy--Green operator $T_\rho$ gives the normalized right inverse
\[
 C_\rho a=T_\rho a-T_\rho a(0)
                  -w\partial_wT_\rho a(0).
\]
Restrict its domain to $(0,1)$-forms vanishing at the center. Its image consists of fields with zero value and first jet there; on this image $\bar\partial$ and $C_\rho$ are inverse. The differential of $\mathcal F_{p,\rho}$ at zero has the form $\bar\partial+K_{p,\rho}$, with
\[
 \|K_{p,\rho}C_\rho\|\leq C\rho,\qquad
 \|\mathcal F_{p,\rho}(0)\|\leq C\rho^2.
\]
Indeed, the complex structures agree along the zero section, while pulling a one-form back by $w\mapsto\rho w$ supplies the additional factor $\rho$. The inverse is $C_\rho(I+K_{p,\rho}C_\rho)^{-1}$ for small $\rho$. The quantitative implicit-function theorem therefore gives a correction with zero first jet and $Y=O(|w|^2)$.

Compactness of $D$ makes the choice uniform in $p$. The normalized Cauchy operator is equivariant under the unitary changes of fiber frame; the intrinsic transports and uniqueness of the small solution make the corrections agree on overlaps. The parametric implicit-function theorem gives smooth dependence on $p$, and interior elliptic regularity on slightly smaller fiber discs gives smoothness in the fiber. The fixed first jet then gives the quadratic ordinary error, with its cylindrical derivative bounds. Composing the first map with this correction preserves the estimate and proves the lemma.
\end{proof}

\begin{theorem}\label{thm:flat-anticanonical-neighborhood}
\index{Anticanonical bundle!flat elliptic neighborhood}\index{Divisor!flat neighborhood}
Let $S$ be a complex surface and $D\subset S$ a smooth compact connected anticanonical divisor with $D^2=0$. The line bundle $\mathcal O_S(D)$ is unitary flat on a neighborhood of $D$.
\end{theorem}
\begin{proof}
Use the cylindrical identification of \cref{lem:anticanonical-cylindrical-identification}, and omit its pullback notation. A form is horizontal if contraction with $\xi$ and $\bar\xi$ vanishes. For a horizontal form define
\[
 \mathrm d_H\eta=\mathrm d\eta
       -\Theta\wedge\mathcal L_\xi\eta
       -\bar\Theta\wedge\mathcal L_{\bar\xi}\eta.
\]
The degree-$m$ horizontal lift of a form $\eta_m$ valued in $N^{-m}$ is $w^m\eta_m$ in a parallel frame. Direct differentiation identifies $\mathrm d_H$ with its flat covariant derivative. Since the fibers are holomorphic,
\[
 \iota_{\bar\xi}\Omega_S=0,\qquad
 \Omega_S=\Theta\wedge\alpha+\beta,
 \qquad\alpha=\iota_\xi\Omega_S,
\]
with $\alpha,\beta$ horizontal. Closedness gives $\mathcal L_{\bar\xi}\alpha=\mathcal L_{\bar\xi}\beta=0$ and $\mathcal L_\xi\beta=\mathrm d_H\alpha$.

Fix $0<\delta<1/2$, and let $U_T=\{z\geq T\}$. For tensors define the weighted norm
\[
 \|\eta\|_{k,\delta,T}^2
 =\sum_{j=0}^k\int_{U_T}\mathrm e^{2\delta(z-T)}
                    |\nabla^j\eta|_g^2\,\mathrm dV_g.
\]
Let $\mathcal H^{k+1}_{\delta,T}$ be the functions of finite such norm with $\mathcal L_{\bar\xi}f=0$, and let $\mathcal D^k_{\delta,T}$ be the complex two-forms of finite norm with $\mathrm d\gamma=0$ and $\iota_{\bar\xi}\gamma=0$. We construct an operator
\[
 Q_T:\mathcal D^k_{\delta,T}\longrightarrow\mathcal H^{k+1}_{\delta,T},
 \qquad \|Q_T\gamma\|_{k+1,\delta,T}\leq C_{k,\delta}\|\gamma\|_{k,\delta,T},
\]
with constant independent of $T$, satisfying
\begin{equation}\label{eq:anticanonical-horizontal-inverse}
 \mathrm d_H(Q_T\gamma\,\vartheta)
 =\gamma-\Theta\wedge\iota_\xi\gamma,
 \qquad Q_T(\mathrm df\wedge\vartheta)=f.
\end{equation}

Write $\gamma=\Theta\wedge\alpha+\beta$. Its defining equations give $\mathcal L_{\bar\xi}\gamma=0$ and, in the fiberwise homogeneous expansion, $m\beta_m=\mathrm d_H\alpha_m$. Only $m\geq1$ occur: the finite positive weighted norm excludes the zero and negative Laurent modes. Decompose
\[
 \alpha_m=-u_m\vartheta+v_m,
 \qquad u_m\in\Omega^0(D,N^{-m}),\quad
 v_m\in\Omega^{0,1}(D,N^{-m}).
\]
Apply \cref{lem:flat-elliptic-derivative-inverse} to get $w_m$ with $\bar\partial w_m=\iota_v\partial v_m$ and $\|w_m\|_{H^k}\leq\|v_m\|_{H^k}$. Put $f_m=-(u_m+w_m)/m$. Then
\[
 \mathrm d_H(f_m\vartheta)=\beta_m,\qquad
 \|mf_m\|_{H^k}\leq C\|\alpha_m\|_{H^k}.
\]
Also $\|\bar\partial f_m\|_{H^k}\leq C\|\beta_m\|_{H^k}$. Flatness gives the same bound for $\partial f_m$ by integration by parts, and hence controls the pure horizontal derivatives of order $k+1$. A vertical derivative of a homogeneous mode multiplies it by $m$, so every derivative involving a vertical direction is controlled by the preceding bound for $mf_m$. Circle orthogonality and integration of the weight give the required estimate for $f=\sum_{m\geq1}f_m$: the squared radial factor is $\mathrm e^{-mz}$ and $m-2\delta>0$.

For Sobolev coefficients, regularize each mode by the flat base heat operator, which commutes with the defining differential equations. The uniform estimates and orthogonality make the finite sums Cauchy in the target space. Passing to their limit proves the first identity in \eqref{eq:anticanonical-horizontal-inverse}. If $\gamma=\mathrm df\wedge\vartheta$, then $\alpha_m=mf_m\vartheta$, so its $v_m$ and normalized $w_m$ vanish and the construction returns $f_m$. This proves the second identity.

Put $\eta=\Omega_S-\Omega_0$. Choose an integer $k\geq4$. The cylindrical estimate bounds its weighted $H^{k+2}$ norm on $U_{T-1}$ and its first $k+2$ uniform derivatives by a number $a_T\to0$. For small $f\in\mathcal H^{k+1}_{\delta,T}$, the map $\Phi_f(\zeta)=\mathrm e^{f(\zeta)}\zeta$ is a $C^2$ diffeomorphism onto its image and sends $U_T$ into $U_{T-1}$. Its restriction to every fiber is holomorphic. Thus $\Phi_f^*\Omega_S-\Omega_0$ belongs to the actual domain of $Q_T$. Since
\[
 \Phi_f^*\Omega_0-\Omega_0=\mathrm df\wedge\vartheta,
\]
the desired equation becomes
\[
 Q_T(\Phi_f^*\Omega_S-\Omega_0)=f+E_T(f)=0,
 \qquad E_T(f)=Q_T(\Phi_f^*\eta).
\]
On a fixed sufficiently small ball, weighted Sobolev multiplication and composition on the flat cylinder give
\[
 \|E_T(f)\|_{k+1,\delta,T}\leq Ca_T,\qquad
 \|E_T(f)-E_T(g)\|_{k+1,\delta,T}
 \leq Ca_T\|f-g\|_{k+1,\delta,T}.
\]
The constants are independent of $T$: translation in the cylindrical direction is an isometry, the weight is normalized at $T$, and the shifts of $z$ are uniformly bounded. Differentiating the pullback along the segment between $f$ and $g$ takes one Lie derivative of the fixed form $\eta$ and at most one derivative of $f-g$; the gain of one derivative in $Q_T$ restores the indicated norm. For $T$ sufficiently large, $-E_T$ is a contraction of the closed ball of radius $2Ca_T$ into itself. Let $f$ be its fixed point.

The first identity in \eqref{eq:anticanonical-horizontal-inverse} makes the horizontal component of $\Phi_f^*\Omega_S$ vanish. Consequently $\Theta\wedge\Phi_f^*\Omega_S=0$. The pulled-back volume form is nowhere zero and of type $(2,0)$, so $\Theta$ is of type $(1,0)$ for $\Phi_f^*J_S$. Since $\mathrm dz=-(\Theta+\bar\Theta)$ and $\mathrm d\Theta=0$, the function $z$ is pluriharmonic for that complex structure.

Return to $S$ through $A=\Upsilon\circ\Phi_f$. On each fiber, boundedness extends $f$ holomorphically across zero. Its positive-mode expansion gives $f(0)=0$. The uniform small bound on $f$ lets Rouch\'e's theorem compare $w\mathrm e^{f(w)}-a$ with $w-a$ on a fixed smaller circle, for every sufficiently small $a$. Thus the image of $A$ contains a smaller punctured neighborhood of $D$. The function $r=z\circ A^{-1}$ is $C^2$ and pluriharmonic there, hence smooth by elliptic regularity. If $w_S$ is a holomorphic defining function for $D$, the fixed first normal jet of $\Upsilon$ and the bounded factor $\mathrm e^f$ give
\[
 C^{-1}|w|\leq|w_S\circ A|\leq C|w|.
\]
Thus $r+\log|w_S|^2$ is bounded and pluriharmonic off $D$. Both signs have plurisubharmonic extensions across the divisor; their sum is zero almost everywhere and hence everywhere. The extension is therefore pluriharmonic. It follows that $F=-r$, with its logarithmic pole on $D$, satisfies $\ddc F=[D]$ near $D$. Locally, the norm assigning $\exp(F)$ to the squared norm of the defining section is a smooth flat metric, since $F-\log|w_S|^2$ is smooth pluriharmonic. These norms agree on overlaps and give the asserted unitary flat metric.
\end{proof}

\begin{proposition}\label{prop:flat-divisor-rank-one-curvature}
Let $D$ be a smooth nonempty compact divisor in a complex surface $S$, and suppose $\mathcal O_S(D)$ is unitary flat near $D$. It has a smooth semipositive metric on $S$ whose curvature $B$ has rank at most one and vanishes near $D$ and outside a relatively compact neighborhood of $D$. If $S$ is compact K\"ahler, then $\max_S\operatorname{rank}B=1$.
\end{proposition}
\begin{proof}
Let $s_D$ be the defining section and $h_0$ a flat metric on a neighborhood $U$ of $D$. Put $F=\log|s_D|_{h_0}^2$. Choose a relatively compact tubular neighborhood $U_0\Subset U$. Since $F$ tends uniformly to $-\infty$ along $D$ and is finite on $\partial U_0$, adding a constant makes $F>0$ near $\partial U_0$. Choose a smooth concave function $\chi$ with $\chi(t)=t+c$ for $t\leq-1$ and $\chi(t)=0$ for $t\geq0$. Define $\widetilde F=\chi(F)$ on $U_0$ and extend it by zero away from $U_0$. On $S\setminus D$ prescribe
\[
 |s_D|_h^2=\mathrm e^{\widetilde F}.
\]
Near $D$ this gives $h=\mathrm e^c h_0$, so the frame metric extends smoothly and positively across $D$. It also matches smoothly where the cutoff is constant. Poincar\'e--Lelong gives $B=[D]-\ddc\widetilde F$, and off $D$ one has
\[
 B=-\chi''(F)\,\mathrm dF\wedge\dc F\geq0.
\]
This smooth closed form has rank at most one and the asserted vanishing regions. On a compact K\"ahler surface with K\"ahler form $\kappa$, its mixed integral is
\[
 \int_SB\wedge\kappa=\int_D\kappa>0.
\]
Thus it is not identically zero and its maximum rank is one.
\end{proof}

\Needspace{15\baselineskip}
\begin{corollary}\label{cor:flat-divisor-variational-order}
\index{Nef degeneration!flat divisor curvature}
Let $S$ be a compact K\"ahler surface and $D$ a smooth nonempty divisor whose line bundle is unitary flat near $D$. Let $B$ be the curvature from \cref{prop:flat-divisor-rank-one-curvature}. For every fixed smooth $f>0$, every Hermitian form $h$ and every smooth positive probability $\mu$, there are $L,U>0$ such that
\[
 L\epsilon\leq\mathcal V_\mu(fB+\epsilon h)
              \leq U\epsilon,\qquad0<\epsilon\leq1.
\]
For a K\"ahler perturbation $\kappa$, the closed background itself satisfies
\[
 \mathcal V_\mu(B+\epsilon\kappa)
 =2\epsilon\int_D\kappa+\epsilon^2\int_S\kappa^2.
\]
There are choices of $f$ for which $fB$ is nonclosed.
\end{corollary}
\begin{proof}
The inequalities $(\min f)B\leq fB\leq(\max f)B$ and \cref{prop:nef-maximum-rank-comparison} give the two-sided bounds. The polynomial follows from the closed-volume formula and $B^2=0$. At a point where $B$ has rank one, choose a smooth positive $f$ whose differential has a nonzero component tangential to the kernel of $B$. Then $\mathrm d(fB)=\mathrm df\wedge B\ne0$ there.
\end{proof}

\Needspace{12\baselineskip}
\begin{proposition}\label{prop:nine-point-anticanonical-curvature}
\index{Anticanonical bundle!nine-point blowup}
Let $S$ be obtained from $\mathbb P^2$ by nine successive point blowups, each centered on the strict transform of a fixed smooth cubic. Its final strict transform $D$ is a smooth elliptic anticanonical curve with $D^2=0$. The anticanonical bundle has a smooth semipositive metric with rank-one curvature as above. It is semiample exactly when $N=\mathcal O_D(D)$ is torsion. If $N$ has exact order $m$, then $|mD|$ is basepoint-free and no smaller positive multiple moves. If $N$ is non-torsion, then $|kD|=\{kD\}$ for every $k\geq1$.
\end{proposition}
\begin{proof}
A blowup at a smooth point of the cubic replaces its class by its pullback minus the exceptional curve, exactly as for the anticanonical class. The strict transform remains isomorphic to the cubic, and its self-intersection decreases by one. After nine steps $D^2=9-9=0$. The flat neighborhood theorem and curvature construction apply.

For $k\geq1$, use the restriction sequences
\[
 0\longrightarrow\mathcal O_S((k-1)D)
 \longrightarrow\mathcal O_S(kD)
 \longrightarrow N^k\longrightarrow0.
\]
A nontrivial degree-zero line on an elliptic curve has no section: any nonzero section would have a zero divisor of degree zero and trivialize the line. If $N$ is non-torsion, induction gives $h^0(S,\mathcal O_S(kD))=h^0(S,\mathcal O_S)=1$, proving the singleton systems and failure of semiampleness.

If $N$ has exact order $m$, then $H^0(D,N^k)=H^1(D,N^k)=0$ for $1\leq k<m$, the latter by Serre duality on the elliptic curve. Since $S$ is rational, $H^1(S,\mathcal O_S)=0$. The same sequences give $h^0(S,kD)=1$ and $H^1(S,\mathcal O_S(kD))=0$ for $0\leq k<m$. The constant section of $N^m\simeq\mathcal O_D$ therefore lifts to a section of $\mathcal O_S(mD)$ nonzero everywhere on $D$. Together with $s_D^m$, it generates that line globally. Thus $|mD|$ is basepoint-free, and no smaller multiple moves.

Non-torsion configurations exist with nine distinct centers. Fix the first eight points on the cubic and vary the ninth. Under the unchanged identification with that elliptic curve, the normal line is $\mathcal O_D(3)\otimes\mathcal O_D(-p_1-\cdots-p_9)$. Varying $p_9$ parametrizes its degree-zero Picard group by the elliptic Abel map. The torsion subgroup is countable, while the curve is uncountable; excluding the first eight points still leaves non-torsion choices. Hence \cref{cor:flat-divisor-variational-order} supplies nonclosed backgrounds of order one even when the associated anticanonical line is not semiample.
\end{proof}

\subsection{Flat divisors and Hartogs extension}
\label{subsec:flat-divisor-complements}

All the nine-point blowups of \cref{prop:nine-point-anticanonical-curvature} carry backgrounds of order one, so the order of degeneration does not distinguish the torsion case, in which a multiple of $D$ moves in a basepoint-free pencil, from the non-torsion case, in which no multiple of $D$ moves. The difference is visible in the complex geometry near $D$. On the tube where the curvature form of \cref{prop:flat-divisor-rank-one-curvature} vanishes, the logarithm of the flat norm of the defining section is pluriharmonic. Its regular level sets are therefore Levi-flat, that is, real hypersurfaces foliated by complex hypersurfaces. These complex leaves are governed by the unitary holonomy of the normal bundle: finite holonomy gives compact fibers of a local holomorphic function, whereas infinite holonomy makes every leaf dense and forces every holomorphic function on the punctured tube to be constant. In both cases the complement of $D$ is not Stein.

In this subsection $S$ is a complex manifold of dimension $n\geq2$ and $D\subset S$ is a smooth compact connected divisor. Suppose that $L=\mathcal O_S(D)$ has a unitary flat metric $h$ on a neighborhood of $D$. Write $s_D$ for its defining section and $N=L|_D$ for the normal bundle, with the induced metric and flat connection, and let $\chi\colon\pi_1(D)\to\mathrm U(1)$ be the holonomy character of $N$, which describes parallel transport around closed loops. In parallel unit frames, $s_D$ is represented by holomorphic functions $W_j$ satisfying $W_j=t_{jk}W_k$ with locally constant $t_{jk}\in\mathrm U(1)$. The flat connection defines on each circle bundle $\{v\in N:|v|=r\}$ a foliation whose leaves are given locally, in parallel frames, by $\{w=\text{constant}\}$; we call it the flat horizontal foliation.

\Needspace{16\baselineskip}
\begin{lemma}\label{lem:flat-divisor-foliated-tube}
There are $R>0$, a neighborhood $U$ of $D$ and a smooth diffeomorphism
\[
 \Upsilon\colon\{v\in N:|v|<R\}\longrightarrow U
\]
that is the identity on the zero section and preserves the defining function: in each parallel normal frame, $W_j(\Upsilon(p,w))=w$. In particular, $|s_D|_h\circ\Upsilon=|w|$. For $0<r<R$, the set
\[
 M_r=\{x\in U:|s_D(x)|_h=r\}
\]
is a compact connected Levi-flat hypersurface, and $\Upsilon$ maps the leaves of the flat horizontal foliation of $\{|w|=r\}$ onto the complex leaves of $M_r$.
\end{lemma}
\begin{proof}
Choose a smooth normal exponential map $\Upsilon_0$ whose first normal derivative is complex linear; the derivative of $s_D$ identifies the normal bundle with $L|_D$, and we use this identification to fix that first derivative. Parallel transport along the normal rays identifies $\Upsilon_0^*L$ with the pullback of $N$. In compatible parallel frames one then has
\[
 W_j(\Upsilon_0(p,a))=a+O(|a|^2).
\]
For each $p$, this is a smooth map between real two-dimensional disks whose derivative at zero is the identity. The inverse function theorem with parameters, uniformly over the compact base $D$, gives a unique small $a=a(p,w)$ with $W_j(\Upsilon_0(p,a))=w$. It depends smoothly on $(p,w)$, and its first derivative in $w$ at $w=0$ is the identity. On an overlap of frames, $a$ and $w$ change by the same constant unitary factor, so uniqueness makes these corrections agree. After decreasing $R$, the map $(p,w)\mapsto(p,a(p,w))$ is a diffeomorphism of a tube onto its image, and its composition with $\Upsilon_0$ is the required map $\Upsilon$.

Off $D$, the function $F=\log|s_D|_h^2$ is pluriharmonic, and its radial derivative in the tube is nonzero. Hence every $M_r$ is a smooth hypersurface; it is compact and connected, being a circle bundle over $D$. Locally its complex tangent space is $\ker\mathrm dW_j$, so it is Levi-flat, with complex leaves the hypersurfaces $\{W_j=\text{constant}\}$. The identity $W_j\circ\Upsilon=w$ pulls this distribution back to $\{\mathrm dw=0\}$ in parallel frames, which is the flat horizontal foliation of the circle bundle. The map $\Upsilon$ is only a smooth tubular identification; its base projection need not be holomorphic.
\end{proof}

\begin{proposition}\label{prop:flat-divisor-finite-holonomy}
In this setting, $N$ is torsion as a holomorphic line bundle if and only if its unitary holonomy has finite image. If $N$ has exact order $m$, then a sufficiently small tube $U$ admits a proper surjective holomorphic map
\[
 \phi\colon U\longrightarrow\{z\in\mathbb C:|z|<R^m\}
\]
with connected fibers and $\operatorname{div}(\phi)=mD$. Its fibers over nonzero values are smooth compact complex hypersurfaces.
\end{proposition}
\begin{proof}
If $\chi^m=1$, a parallel nowhere-zero section trivializes $N^m$ holomorphically. Conversely, if a holomorphic section $e$ trivializes $N^m$, then $\log|e|^2$ is pluriharmonic on the compact connected manifold $D$, hence constant by the maximum principle. In parallel holomorphic frames the coefficient of $e$ is a holomorphic function of constant modulus, hence constant. Therefore $e$ is parallel and $\chi^m=1$. Consequently the least such $m$ is both the order of $N$ as a torsion line bundle and the order of the finite holonomy group.

By \cref{lem:flat-divisor-foliated-tube}, the tube $U$ retracts onto $D$. Hence the flat bundle $L^m$ on $U$ has trivial holonomy and a global parallel section $e_U$ of unit length. Define $\phi=s_D^m/e_U$. Then $\phi$ is holomorphic, $\operatorname{div}(\phi)=mD$ and $|\phi|=|s_D|_h^m$. The inverse image of a compact subset of the target lies in a smaller closed normal disk bundle and is closed there, which proves properness, and the tubular model shows that every value in the disk is attained.

Choose local parallel unit frames $e_j$ of $L$ with $e_j^m=e_U$, so that $\phi=W_j^m$. In the tubular model, the fiber over a nonzero value consists, over each point of $D$, of the $m$ roots of that value. Parallel transport acts on these roots through $\chi$, whose image has exact order $m$ and therefore acts transitively. Thus this fiber is a connected covering of $D$, as a real manifold. Since $\mathrm d\phi=mW_j^{m-1}\mathrm dW_j$, it is a smooth complex hypersurface, and it is compact because $\phi$ is proper. The fiber over zero has connected support $D$, which completes the proof that all fibers are connected.
\end{proof}

\begin{theorem}\label{thm:flat-divisor-holomorphic-rigidity}
\index{Hartogs extension!non-torsion flat divisor}
\index{Levi-flat hypersurface!flat divisor tube}
Suppose $N$ is not torsion. Then every complex leaf of $M_r$ is dense in $M_r$, and every holomorphic function on $U\setminus D$ is constant:
\[
 H^0(U\setminus D,\mathcal O_S)=\mathbb C.
\]
Suppose moreover that $S$ is connected. For every compact set $K\subset S\setminus D$ such that $(S\setminus D)\setminus K$ is connected, every holomorphic function on $(S\setminus D)\setminus K$ is constant, and in particular extends holomorphically to $S\setminus D$. Similarly, for every connected open neighborhood $V$ of $D$, every holomorphic function on $V\setminus D$ is constant and extends to $S\setminus D$.
\end{theorem}
\begin{proof}
By \cref{prop:flat-divisor-finite-holonomy}, $\chi(\pi_1(D))$ is an infinite subgroup of $\mathrm U(1)$. Such a subgroup is dense, since a proper closed subgroup of the circle is finite, as one sees by lifting it to a closed subgroup of $\mathbb R$ containing $2\pi\mathbb Z$. The circle bundle with its flat horizontal foliation is the suspension of $\chi$, that is, the quotient of $\widetilde D\times S^1$ by the diagonal action of $\pi_1(D)$, acting on $S^1$ through $\chi$. A leaf projects onto $D$, and its intersection with each fiber contains the holonomy orbit of any one of its points, which is dense in the fiber. In a local trivialization, horizontal transport maps the orbits in nearby fibers to each other by a fixed rotation, so the leaf is dense in the whole circle bundle. By \cref{lem:flat-divisor-foliated-tube}, every complex leaf is dense in $M_r$.

Let $f$ be holomorphic on $U\setminus D$. Its modulus attains a maximum on the compact set $M_r$. On the complex leaf through a maximum point, $|f|$ is bounded by this maximum and attains it, so the maximum modulus principle makes $f$ constant on the connected leaf. By density and continuity, $f$ equals the same constant $c$ on $M_r$. Near any point of $M_r$, choose holomorphic coordinates $(z,W_j)$ with $z\in\mathbb C^{n-1}$. For each fixed $z$, the holomorphic function $W_j\mapsto f(z,W_j)-c$ vanishes on an arc of the circle $|W_j|=r$, so it vanishes on the coordinate neighborhood by the one-variable identity theorem. Analytic continuation on the connected punctured tube gives $f\equiv c$.

A compact set $K\subset S\setminus D$ is disjoint from a sufficiently small tube around the compact divisor $D$. A holomorphic function on $(S\setminus D)\setminus K$ is therefore constant on this punctured tube, and by the identity theorem on the connected set $(S\setminus D)\setminus K$ it is constant everywhere; the constant function is the required extension. Similarly, every open neighborhood $V$ of $D$ contains a full tube, and $V\setminus D$ is connected, because paths in the connected manifold $V$ can be perturbed off the real codimension-two submanifold $D$. The identity theorem gives the assertion for $V$. The same path argument shows that $S\setminus D$ is connected.
\end{proof}

Recall that a complex manifold $Y$ is \emph{holomorphically separable} if its holomorphic functions separate points, and \emph{holomorphically convex} if the holomorphic hull
\[
 \widehat K_Y=\{y\in Y:|f(y)|\leq\max_K|f|\text{ for every }f\in\mathcal O(Y)\}
\]
of every compact set $K\subset Y$ is compact. A Stein manifold has both properties.

\begin{corollary}\label{cor:flat-divisor-nonstein}
\index{Stein manifold!flat divisor obstruction}
For any holonomy of $N$, the punctured tube $U\setminus D$ admits no smooth strictly plurisubharmonic function, and $S\setminus D$ is not Stein. In the non-torsion nine-point blowups of \cref{prop:nine-point-anticanonical-curvature}, every holomorphic function on the complement of $D$ is constant, and the Hartogs extension of \cref{thm:flat-divisor-holomorphic-rigidity} across compact holes holds.
\end{corollary}
\begin{proof}
The restriction of a smooth strictly plurisubharmonic function on $U\setminus D$ to a complex leaf of $M_r$ is strictly plurisubharmonic. At a maximum point on the compact set $M_r$, this restriction has a local maximum on the leaf, which has positive dimension because $n\geq2$; this contradicts the maximum principle.

In the finite-holonomy case, each fiber of the map of \cref{prop:flat-divisor-finite-holonomy} over a nonzero value is a connected compact complex hypersurface of positive dimension. By the maximum modulus principle, every holomorphic function on $S\setminus D$ is constant on it, so two distinct points of the fiber cannot be separated. In the infinite-holonomy case, \cref{thm:flat-divisor-holomorphic-rigidity} makes every holomorphic function constant on the punctured tube, which again contradicts holomorphic separability. In both cases $S\setminus D$ is not Stein. Finally, \cref{thm:flat-anticanonical-neighborhood} gives the flat neighborhood for the nine-point blowups, whose normal bundle is non-torsion by assumption, and \cref{thm:flat-divisor-holomorphic-rigidity}, also with $K=\varnothing$, gives the last assertion.
\end{proof}

The nine-point blowups show that a background with a well-defined order of degeneration can have very different holomorphic geometry near its null locus; \cref{app:cotangent-torsors} gives a second instance of this contrast, for K\"ahler classes on a ruled surface over an elliptic curve.


\begin{problem}\label{qu:nef-degeneration-exponent}
Which geometric conditions on a nonclosed nef background force a common degeneration order, or an integral one? When is equality attained in the modification and residual order comparisons?
\end{problem}


\Needspace{12\baselineskip}
\section{Current tests for smooth positivity}
\label{sec:smooth-cone-current-duality}

Let $X$ again be a connected compact complex manifold of dimension $n\geq1$. Testing against all positive $\ddc$-closed currents detects smooth positivity of a background. Hahn--Banach separation gives the exact threshold, with zero as the nef boundary. The full current set is needed here: the smooth tests in Lamari duality detect pseudoeffectivity.

Fix a Hermitian metric $h$ and use $\mathscr G_h$ from \cref{sec:gauduchon-linear-duality}. That set is compact for $n\geq2$. For $n=1$, elliptic regularity and the maximum principle give $\mathscr G_h=\{1/\int_Xh\}$, so compactness holds in this case as well.
\Needspace{6\baselineskip}
For any smooth real $(1,1)$-form $\theta$, define
\[
 m_h(\theta)=\min_{T\in\mathscr G_h}\int_X\theta\wedge T.
\]
The minimum exists by weak compactness, and it is finite since $-Ch\leq\theta\leq Ch$ for some $C$.

\begin{theorem}\label{thm:smooth-cone-current-duality}\index{Cone duality!smooth forms and currents}
The background $\theta$ has a smooth positive representative if and only if $m_h(\theta)>0$. Moreover $\theta$ is nef if and only if $m_h(\theta)\geq0$, and
\[
 m_h(\theta)=\sup\{a\in\mathbb R:\ \exists u\in C^\infty(X,\mathbb R),
                                  \ \theta+\ddc u>ah\}.
\]
\end{theorem}
\begin{proof}
If $\theta+\ddc u\geq ah$ for $a>0$, then distributional integration by parts gives $\int_X\theta\wedge T\geq a$ for every normalized test current. This proves one implication.

For the converse, work in the real Fr\'echet space of smooth $(1,1)$-forms. Let $\mathcal P$ be its open convex cone of strictly positive forms, and let $L=\ddc C^\infty(X,\mathbb R)$. If the affine space $\theta+L$ is disjoint from $\mathcal P$, the Hahn--Banach separation theorem for an open convex set and a disjoint affine space gives a nonzero continuous real linear functional $T$ with
\[
 T(\theta+\ell)\leq T(\eta)\qquad(\ell\in L,\ \eta\in\mathcal P).
\]
Varying $\ell$ by arbitrary real multiples shows $T(L)=0$. Scaling any $\eta\in\mathcal P$ to zero shows $T(\theta)\leq0$, while scaling it to infinity shows $T(\eta)\geq0$. Thus $T$ is a positive current: every smooth semipositive test form is a limit of strictly positive forms. Positivity also makes it a current of order zero. The equality $T(L)=0$ is precisely $\ddc T=0$. Since $T\ne0$ and $h>0$, its trace $T(h)$ is positive. Normalizing produces an element of $\mathscr G_h$ with nonpositive pairing against $\theta$. This contradicts $m_h(\theta)>0$ and proves the converse.

The identity $m_h(\theta-ah)=m_h(\theta)-a$ and the equivalence just proved show that the set inside the supremum is exactly $(-\infty,m_h(\theta))$. Finally, nefness means that $\theta+\varepsilon h$ has a smooth positive representative for every $\varepsilon>0$. This is equivalent to $m_h(\theta)+\varepsilon>0$ for all such $\varepsilon$, or $m_h(\theta)\geq0$.
\end{proof}

\Needspace{5\baselineskip}
For $n\geq2$, let $\mathscr C_h$ be the weak closure of the smooth positive normalized tests and let $\kappa_h(\theta)$ be their linear positivity threshold, as in \cref{sec:gauduchon-linear-duality}. The inclusion $\mathscr C_h\subseteq\mathscr G_h$ gives
$m_h(\theta)\leq\kappa_h(\theta)$.
The larger test set detects smooth positivity and nefness; the weak
closure of the smooth tests detects positivity and bigness of current
representatives. The inequality can be strict, as
\cref{ex:gauduchon-current-smoothing-obstruction} shows.

\begin{proposition}\label{prop:nef-annihilator-determinant}
Let $\theta$ be nef and suppose $\lambda_\mu(\theta)>0$. If $T\in\mathscr G_h$ satisfies $\int_X\theta\wedge T=0$, then its determinant-weighted measure vanishes.
\end{proposition}
\begin{proof}
In dimension one the hypotheses cannot hold, since
$\int_X\theta\wedge T=\lambda_\mu(\theta)/\int_Xh>0$.
Suppose $n\geq2$, and write $D_\mu(T)$ for the determinant functional of \cref{sec:weak-gauduchon-duality}, with $h$ as reference metric and volume form $\mu$. If $\eta\in\mathscr P$, its solution form satisfies $\Omega^n=\lambda_\mu(\eta)^n\mu$. Test the infimum in \cref{lem:weak-gaud-determinant-dual} with $\Omega/\lambda_\mu(\eta)$. Stokes gives
\[
 \int_X\eta\wedge T=\int_X\Omega\wedge T
              \geq\lambda_\mu(\eta)D_\mu(T).
\]
Apply this to $\eta=\theta+\varepsilon h$. The left side is $\varepsilon$, while \cref{thm:gaud-volume-nef-limit} gives $\lambda_\mu(\eta)\to\lambda_\mu(\theta)>0$. Thus $D_\mu(T)=0$. The determinant-weighted measure is positive and has this total mass, so it vanishes.
\end{proof}

A nef boundary with positive variational volume therefore has annihilating currents whose coefficient matrices are rank deficient almost everywhere. Their trace measures need not be supported on analytic subsets merely because these determinants vanish.

\Needspace{10\baselineskip}
\subsection{Maximal time of the Chern--Ricci flow}
\label{subsec:chern-ricci-maximal-time}

A smooth Chern--Ricci flow gives positive representatives along an affine ray. Conversely, a positive representative at a later time supplies a uniformly positive reference path and prevents an earlier singularity. The current criterion then computes the lifetime. Let $\omega_0$ be the initial Hermitian metric. For a Hermitian form $\omega$ on $X$, use the conventional Chern--Ricci form
\[
 \operatorname{Ric}(\omega)=-\mathrm{i}\partial\bar\partial\log\omega^n
                          =-2\pi\ddc\log\omega^n,
\]
interpreted in local coordinates; coordinate changes do not affect the form. This fixes the time normalization of the flow $\partial_t\omega=-\operatorname{Ric}(\omega)$. Put $\chi=\operatorname{Ric}(\omega_0)$ and $\theta_t=\omega_0-t\chi$.

\begin{theorem}\label{thm:chern-ricci-maximal-time}\index{Chern--Ricci flow!maximal time}
The Chern--Ricci flow starting from $\omega_0$ has a unique maximal smooth solution on $[0,\tau)$, where
\[
 \tau=\sup\{t\geq0:\ \theta_t\text{ has a smooth positive representative}\}.
\]
\end{theorem}
\begin{proof}
The admissible times form an interval containing zero: positive representatives at two times give positive representatives at the intervening times by convex combination. Positivity is also open in time. Fix an admissible $T>0$ and choose a smooth real function $a$ such that $\theta_T+\ddc a>0$. The reference forms
\[
 \gamma_t=\theta_t+\frac tT\ddc a
          =\left(1-\frac tT\right)\omega_0
             +\frac tT(\theta_T+\ddc a),\qquad 0\leq t\leq T,
\]
are uniformly positive, with all their derivatives uniformly bounded. Put
\[
 \Omega=\mathrm{e}^{a/(2\pi T)}\omega_0^n.
\]
Then $\dot\gamma_t=-\operatorname{Ric}(\Omega)$, with the Ricci form of a volume form defined by the same local expression. The scalar equation is
\begin{equation}\label{eq:chern-ricci-scalar-flow}
 \dot u=2\pi\log\frac{(\gamma_t+\ddc u)^n}{\Omega},
 \qquad \gamma_t+\ddc u>0,\qquad u(0)=0.
\end{equation}
Applying $\ddc$ shows that $\omega(t)=\gamma_t+\ddc u(t)$ solves the metric flow. Conversely, integrating $\dot u=2\pi\log(\omega(t)^n/\Omega)$ from zero recovers this identity from any metric solution. The linearization of the scalar equation is $\partial_t-2\pi\operatorname{tr}_{\omega(t)}\ddc$, a strictly parabolic operator. Classical short-time existence for smooth strictly parabolic equations therefore supplies a solution. Two scalar solutions have a difference satisfying a linear parabolic equation, whose positive principal coefficients are the averages of the inverse metrics along the segment between them. The maximum principle proves uniqueness.

Suppose this solution has maximal time $T_*<T$. We establish estimates independent of $t<T_*$ to continue it. For a large fixed $C$, the functions $Ct$ and $-Ct$ are respectively a supersolution and a subsolution of \eqref{eq:chern-ricci-scalar-flow}: at a spatial extremum of $u$ the complex Hessian has the required sign, and $\gamma_t^n/\Omega$ is bounded above and away from zero. The maximum principle gives $|u|\leq CT$.

Set $v=\dot u$ and $\mathcal L=2\pi\operatorname{tr}_{\omega(t)}\ddc$. Differentiation gives
\[
 (\partial_t-\mathcal L)v=2\pi\operatorname{tr}_{\omega(t)}\dot\gamma_t,
 \qquad
 (\partial_t-\mathcal L)u=v-2\pi n+2\pi\operatorname{tr}_{\omega(t)}\gamma_t.
\]
Since $\gamma_t$ is affine in time, these imply
\begin{align*}
 (\partial_t-\mathcal L)(t v-u-2\pi nt)
     &=-2\pi\operatorname{tr}_{\omega(t)}\omega_0<0,\\
 (\partial_t-\mathcal L)((T-t)v+u+2\pi nt)
     &=2\pi\operatorname{tr}_{\omega(t)}\gamma_T>0.
\end{align*}
The first expression is bounded above by its initial value, zero; the second is bounded below by $\min_X T v(0)$. The bound for $u$ and $T-t\geq T-T_*>0$ give a lower bound for $v$. The first inequality gives its upper bound on $[T_*/2,T_*)$, and the initial smooth solution bounds it on $[0,T_*/2]$. In particular, $\omega(t)^n/\Omega=\mathrm{e}^{v/(2\pi)}$ is bounded above and away from zero.

\Needspace{7\baselineskip}
We next bound the trace, retaining the time derivative of the reference metric. In the coordinates of \cref{lem:hermitian-trace-inequality}, write $g(t)$ for the matrix of $\gamma_t$ and $h=g+(u_{i\bar j})$ for that of $\omega(t)$. Put
\[
 S=\operatorname{tr}_g h,\qquad B=\operatorname{tr}_h g,
 \qquad F=\log\frac{\Omega}{\gamma_t^n}.
\]
Thus $\log(\det h/\det g)=v/(2\pi)+F$. The family $g(t)$ has uniformly bounded Chern torsion and curvature. At each fixed time \cref{lem:hermitian-trace-inequality} applies with these uniform constants. Direct differentiation of $S$ gives
\[
 \partial_t\log S
   =\frac{\Delta_g v+\operatorname{tr}_g\dot g}{S}
       -\frac{\operatorname{tr}(g^{-1}\dot g\,g^{-1}h)}{S}.
\]
The $\Delta_g v/S$ term cancels the corresponding term in $\mathcal L\log S$. If $\tau_k=\sum_iT_{ik}^i$ is the torsion trace of $g(t)$, the resulting inequality on $\{S\geq1\}$ is
\begin{equation}\label{eq:chern-ricci-parabolic-trace}
 (\partial_t-\mathcal L)\log S
 \leq-\frac{4\pi}{S^2}\operatorname{Re}(h^{k\bar q}\tau_k S_{\bar q})+CB+C.
\end{equation}
Indeed $\Delta_gF$ and $\operatorname{tr}_g\dot g$ are uniformly bounded, $S^{-1}\leq B/n^2$, and $-Cg\leq\dot g\leq Cg$ bounds the last displayed quotient. These bounds account for every added time-dependent term.

Choose a constant $C_0$ so that $1\leq s=u+C_0\leq D$ on the whole time interval, and maximize
\[
 Q=\log S-Au+s^{-1}
\]
on $X\times[0,t_0]$, where $t_0<T_*$ and $A\geq1$ will be fixed independently of $t_0$. A maximum at time zero gives the desired trace bound from the initial data. At a later maximum, $\partial\log S=(A+s^{-2})\partial u$. Since $|\tau|_h^2\leq CB$, Young's inequality bounds the absolute value of the first term on the right of \eqref{eq:chern-ricci-parabolic-trace} by
\[
 2\pi s^{-3}|\partial u|_h^2
       +\frac{C(A+1)^2s^3}{S^2}B.
\]
If $S^2\leq(A+1)^2D^3$, the trace is already bounded at the maximum. In the remaining case the last coefficient is at most $C$, and
\[
 \begin{split}
 0\leq(\partial_t-\mathcal L)Q
 &\leq CB+C-(A+s^{-2})(v-2\pi n+2\pi B)
                       -2\pi s^{-3}|\partial u|_h^2\\
 &\leq(C-2\pi A)B+C(A+1).
 \end{split}
\]
Choose $A$ so that $2\pi A>C+1$. This bounds $B$ at the maximum. The eigenvalues give $S\leq n(\det h/\det g)B^{n-1}$ for $n\geq2$, and $S=\det h/\det g$ for $n=1$. The determinant bound therefore bounds $S$ at the maximum in every case. The bound for $u$ transfers this estimate from $Q$ to $S$ everywhere, uniformly in $t_0$. The determinant lower bound then bounds the smallest eigenvalue away from zero. Thus $\omega(t)$ and $\gamma_t$ are uniformly equivalent for $t<T_*$.

Uniform ellipticity now gives the higher estimates by classical local parabolic regularity. In a coordinate ball, write the logarithmic determinant as a concave function of the complex part of the real Hessian. On the fixed compact set of positive matrices just obtained, its supporting affine functions have uniformly positive, bounded slopes. Their infimum extends the operator to a concave uniformly elliptic function of every real symmetric matrix. The spatial and time dependence enters through the smooth family $g(t)$ and the fixed density $\Omega$. The parabolic Evans--Krylov estimate gives uniform $C^{2+\alpha,1+\alpha/2}$ bounds away from the initial time. The original smooth solution covers a fixed initial interval. Differentiating the equation and applying parabolic Schauder estimates then gives uniform bounds for all spatial and time derivatives. This is the parabolic counterpart of the local extension used in \cref{cor:classical-higher-estimate}.

The bounds for the time derivatives make $u(t)$ Cauchy in every spatial $C^k$ norm as $t\uparrow T_*$. Its smooth limit still has a strictly positive solution form by uniform ellipticity. Short-time existence with this limit as initial datum extends the solution beyond $T_*$, a contradiction. Hence the flow exists up to every admissible time $T$, and uniqueness patches these solutions on $[0,\tau)$.

\Needspace{5\baselineskip}
Finally, integrating the metric flow gives
\[
 \omega(t)=\theta_t+\ddc\left(2\pi\int_0^t
                   \log\frac{\omega(s)^n}{\omega_0^n}\,\mathrm{d}s\right).
\]
Every time at which a smooth positive solution exists is therefore admissible. This proves maximality.
\end{proof}

Combining the continuation theorem with the current criterion gives a numerical description of the lifetime.

\Needspace{14\baselineskip}
\begin{corollary}\label{cor:chern-ricci-current-lifetime}
Let
\[
 b=\max_{T\in\mathscr G_{\omega_0}}\int_X\chi\wedge T.
\]
Then the maximal smooth existence time is
\[
 \tau=\begin{cases}1/b,&b>0,\\+\infty,&b\leq0.\end{cases}
\]
If $\tau<\infty$, the limiting background $\theta_\tau$ is nef and there is a normalized current $T_*$ with $\int_X\theta_\tau\wedge T_*=0$. If also $\lambda_\mu(\theta_\tau)>0$, every such annihilating current has zero determinant-weighted measure.
\end{corollary}
\Needspace{6\baselineskip}
\begin{proof}
For $t\geq0$, normalization of the tests gives the exact identity
\[
 m_{\omega_0}(\theta_t)=1-tb.
\]
The smooth-positivity theorem and \cref{thm:chern-ricci-maximal-time} prove the formula for $\tau$. When $b>0$, the maximum defining $b$ is attained by weak compactness. Its maximizer $T_*$ annihilates $\theta_\tau$. The same identity gives $m_{\omega_0}(\theta_\tau)=0$, hence nefness. The last assertion is \cref{prop:nef-annihilator-determinant}.
\end{proof}

\Needspace{8\baselineskip}
Along the flow, $\omega(t)=\theta_t+\ddc\varphi_t$ with a smooth potential. For every fixed $T\in\mathscr G_{\omega_0}$ one consequently has
\[
 \int_X\omega(t)\wedge T
       =1-t\int_X\chi\wedge T.
\]
These finite traces describe a boundary obstruction without any total Monge--Amp\`ere mass bound. They do not locate the singularities of the flow or prove smooth convergence away from an analytic set. Those geometric questions require local estimates in addition to the scalar pairings.

\Needspace{11\baselineskip}
\subsection{A positive-volume endpoint that is not big}
\label{subsec:actual-flow-positive-volume-nonbig}

Backgrounds reached at the maximal time of a Chern--Ricci flow are not better behaved, as far as the relation between positive volume and bigness is concerned. The radial background of \cref{sec:radial-positive-volume-nonbig} occurs, up to a smooth $\ddc$ change, as the limiting background of a Chern--Ricci flow at its maximal time. The normalized currents annihilating it can also be described explicitly.

We keep the notation of \cref{sec:radial-positive-volume-nonbig}: $X=\mathbb P^m\times\mathbb P^1$ with $n=m+1\geq3$, the normalized Fubini--Study forms $\omega$ and $w$ pulled back from the two factors, the nonnegative function $a=x/t$ on the first factor, which vanishes only at $p$, the background $\theta=\omega+aw$ and the probability form $\mu=\omega^m\wedge w$. Put $Z=\{p\}\times\mathbb P^1$, and start the flow from the Hermitian metric
\begin{equation}\label{eq:radial-flow-initial-metric}
 \omega_0=(n+1)\omega+(a+2)w.
\end{equation}

\begin{theorem}\label{thm:actual-flow-positive-volume-nonbig}\index{Chern--Ricci flow!nonbig positive-volume endpoint}
With the normalization of the Ricci form in \cref{subsec:chern-ricci-maximal-time}, the maximal smooth existence time of the Chern--Ricci flow starting from $\omega_0$ is $\tau=1/(2\pi)$. The limiting background $\alpha_\tau=\omega_0-\tau\operatorname{Ric}(\omega_0)$ is nef and not big, and
\[
 \alpha_\tau=\theta+\ddc\log(a+2),\qquad
 \mathcal V_\mu(\alpha_\tau)=n.
\]
Its variational volume is positive for every smooth strictly positive probability form on $X$.
\end{theorem}
\begin{proof}
The initial metric is strictly positive, and
\[
 \omega_0^n=n(n+1)^m(a+2)\mu.
\]
With our normalization, the Ricci forms of the normalized Fubini--Study metrics are $2\pi n\omega$ and $4\pi w$. Consequently
\[
 \operatorname{Ric}(\omega_0)
   =2\pi\bigl(n\omega+2w-\ddc\log(a+2)\bigr).
\]
Write $s$ for the flow time and put $r=2\pi s$. The affine background is
\begin{equation}\label{eq:radial-flow-affine-background}
 \alpha_s=(n+1-nr)\omega+(a+2-2r)w
                          +r\ddc\log(a+2).
\end{equation}
For $0\leq r<1$, adding $\ddc$ of the smooth function $-r\log(a+2)$ gives a strictly positive representative: the horizontal coefficient is at least one and the vertical coefficient is at least $2(1-r)>0$. At $r=1$ we obtain the identity $\alpha_\tau=\theta+\ddc\log(a+2)$. Nefness, bigness and the variational volume are unchanged by this smooth $\ddc$ change, by the conventions of \cref{chap:gauduchon-variational-volume} and \cref{prop:gaud-volume-elementary}, so \cref{thm:radial-positive-volume-nonbig} shows that the endpoint background is nef, not big, and of volume $n$.

For $r>1$, choose a nonzero smooth nonnegative function $f$ on the base supported in $\{a+2-2r<0\}$, a neighborhood of $p$. Its pullback $S=f\omega^m$ is a closed semipositive $(n-1,n-1)$-form. For every quasi-plurisubharmonic function $v$, Stokes' theorem for currents gives
\[
 \int_X(\alpha_s+\ddc v)\wedge S
       =\int_{\mathbb P^m}f(a+2-2r)\omega^m<0,
\]
so $\alpha_s$ is not even pseudoeffective for $r>1$. At $r=1$, a smooth positive representative would contradict nonbigness. Thus $r=1$ is the supremum of the admissible times, and \cref{thm:chern-ricci-maximal-time} gives the smooth flow on $[0,\tau)$ with exactly this lifetime; the reference representatives used above need not solve the flow equation. Finally, \cref{prop:gaud-volume-density-comparison} gives positivity of the volume for every smooth positive probability form.
\end{proof}

\begin{proposition}\label{prop:radial-flow-annihilator-support}\index{Gauduchon current!flow endpoint support}
The current $[Z]/2$ belongs to $\mathscr G_{\omega_0}$ and annihilates $\alpha_\tau$. Every $R\in\mathscr G_{\omega_0}$ with $\int_X\alpha_\tau\wedge R=0$ has its trace probability $\sigma_R=\omega_0\wedge R$ supported on $Z$. In the coefficient-matrix convention of \cref{sec:weak-gauduchon-duality}, its matrix is
\[
 P_R=n\Pi_{\mathrm{vert}}\qquad\sigma_R\text{-almost everywhere},
\]
where $\Pi_{\mathrm{vert}}$ is the $\omega_0$-orthogonal projection onto the tangent line of the fiber $\{p\}\times\mathbb P^1$. In particular, $P_R$ has rank one, and the determinant measure of $R$ vanishes for every smooth positive density.
\end{proposition}
\begin{proof}
Only smooth forms need to be restricted to $Z$, and their restrictions satisfy $\omega_0|_Z=2w$ and $\alpha_\tau|_Z=0$. This proves that $[Z]/2$ is normalized and annihilates $\alpha_\tau$.

Let $R$ be a general annihilator. Since $\alpha_\tau$ and $\theta$ differ by $\ddc$ of a smooth function, Stokes' theorem gives $\int_X\theta\wedge R=0$. The measure $\theta\wedge R$ is nonnegative, so it vanishes. On every compact set $K\subset X\setminus Z$, the explicit form $\theta$ satisfies $\theta\geq c_K\omega_0$ for some $c_K>0$. Thus $\sigma_R(K)=0$, and a countable exhaustion by compact sets shows that $\sigma_R$ is supported on $Z$.

Relative to $\omega_0$, the matrix of $R$ is characterized, for smooth test forms $\beta$ with matrix $B_\beta$, by
\[
 \beta\wedge R=\frac1n\operatorname{tr}(B_\beta P_R)\sigma_R,
 \qquad P_R\geq0,\qquad\operatorname{tr}P_R=n.
\]
At points of $Z$, the form $\theta$ is positive on the horizontal $m$-plane and vanishes on the vertical line. The vanishing of $\theta\wedge R$ therefore forces all horizontal diagonal entries of $P_R$ to vanish, and the inequality $|(P_R)_{ij}|^2\leq(P_R)_{ii}(P_R)_{jj}$ for positive matrices removes their rows and columns. The remaining vertical entry is $n$, which proves the formula and the vanishing of the determinant.
\end{proof}

For each such annihilator $R$, the trace identity at the end of \cref{subsec:chern-ricci-maximal-time} becomes
\[
 \int_X\omega(s)\wedge R=1-2\pi s,
 \qquad 0\leq s<1/(2\pi).
\]
The support statement uses the explicit semipositive representative $\theta$ of the endpoint class; the vanishing of the determinant alone would not give it.



\Needspace{13\baselineskip}
\section{Divisorial constraints and residual volume}
\label{sec:divisorial-gaud-volume}

Suppose we require a $\theta$-plurisubharmonic function to have generic Lelong numbers at least $a_j$ along prescribed prime divisors $D_j$. Removing the divisor $D=\sum_ja_jD_j$ leaves the residual background $\theta-\eta_D$, where $\eta_D$ is a curvature form of $D$, and in each Gauduchon quotient this removal subtracts the Gauduchon degree of $D$ from the numerator. Lamari's duality therefore decides when such functions exist, and the volume of the residual background measures what remains. By the correspondence of \cref{thm:ordinary-divisor-normalization}, equations in the ordinary singularity type of the logarithmic potential of $D$ reduce to equations in the residual background. When the residual background has a smooth positive representative, the smooth maximizer of the dual problem forces every subsolution with the prescribed Lelong numbers and the residual volume as scalar to be the logarithmic potential of $D$ plus a smooth function, up to a constant.

Fix a smooth real $(1,1)$-form $\theta$ on $X$. Let $D=\sum_{j=1}^N a_jD_j$ be an effective real divisor, with distinct prime divisors $D_j$ and coefficients $a_j>0$. Choose a smooth Hermitian metric on each line bundle $\mathcal O_X(D_j)$, write $s_j$ for its canonical section, and let $\eta_j$ be its Chern curvature form. Set
\[
 \sigma_D=\sum_j a_j\log|s_j|^2,
 \qquad \eta_D=\sum_j a_j\eta_j,
 \qquad \beta_D=\theta-\eta_D.
\]
The Poincar\'e--Lelong formula reads
\[
 \ddc\sigma_D=[D]-\eta_D.
\]
For a Gauduchon metric $\gamma$ on $X$, define the Gauduchon degree
\[
 \deg_\gamma D=\int_X[D]\wedge\gamma^{n-1}
              =\sum_j a_j\int_{D_j}\gamma^{n-1},
\]
where integration over a singular divisor means integration over its regular locus, or equivalently against its current of integration. This degree is finite, positive for $D\ne0$, and equal to $\int_X\eta_D\wedge\gamma^{n-1}$ by Stokes' theorem for currents. We write $\nu(u,D_j)$ for the generic Lelong number of $u$ along $D_j$, as in \cref{sec:analytic-types}.

\Needspace{12\baselineskip}
\begin{proposition}\label{prop:divisorial-gaud-feasibility}
The following are equivalent:
\begin{enumerate}
\item there is $u\in\PSH(X,\theta)$ with $\nu(u,D_j)\geq a_j$ for every $j$;
\item the residual background $\beta_D$ is pseudoeffective;
\item for every Gauduchon metric $\gamma$,
\[
 \int_X\theta\wedge\gamma^{n-1}\geq\deg_\gamma D.
\]
\end{enumerate}
The first condition does not require $u-\sigma_D$ to be bounded, nor the Lelong numbers of $u$ to equal the prescribed ones.
\end{proposition}
\begin{proof}
Suppose first that $u$ satisfies the first condition. By \cref{lem:nonclosed-divisorial-removal}, these Lelong numbers can be subtracted: $\theta+\ddc u-[D]\geq0$. Consequently
\[
 \beta_D+\ddc(u-\sigma_D)=\theta+\ddc u-[D]\geq0,
\]
where the difference $u-\sigma_D$ is first taken in $L^1$ and then replaced by its quasi-plurisubharmonic representative. This proves the second condition. Conversely, if $v\in\PSH(X,\beta_D)$, then $u=\sigma_D+v$ satisfies
\[
 \theta+\ddc u=[D]+\beta_D+\ddc v\geq0,
\]
so its generic Lelong numbers along the $D_j$ are at least $a_j$. For $n\geq2$, Lamari's duality, \cref{thm:lamari}, says that a smooth background is pseudoeffective exactly when all its Gauduchon integrals are nonnegative; applied to $\beta_D$, together with the degree identity, it gives the equivalence of the second and third conditions. For $n=1$ the same criterion reads $\int_X\beta_D\geq0$: necessity follows from Stokes' theorem, and for sufficiency one solves $\ddc f=(\int_X\beta_D)\mu-\beta_D$ for a smooth real function $f$ with the scalar Green operator.
\end{proof}

For a nonzero effective real divisor $D$ and a pseudoeffective $\theta$, the exact threshold is therefore
\begin{equation}\label{eq:divisorial-gaud-threshold}
 \tau_\theta(D)=
 \inf_{\gamma\text{ Gauduchon}}
       \frac{\int_X\theta\wedge\gamma^{n-1}}{\deg_\gamma D}.
\end{equation}
It is finite and nonnegative, and a function $u\in\PSH(X,\theta)$ with $\nu(u,D_j)\geq ta_j$ for all $j$ exists if and only if $0\leq t\leq\tau_\theta(D)$. The endpoint is included, because the non-strict Gauduchon inequalities still hold there and duality gives a positive current. This is a threshold for the existence of a positive current, not for the existence of a bounded residual potential.

Whenever $\beta_D$ is pseudoeffective, define the residual volume
\[
 \mathcal V_\mu(\theta;D)=\mathcal V_\mu(\beta_D).
\]
The definition and the degree identity give the exact formula
\begin{equation}\label{eq:divisorial-gaud-volume}
 \mathcal V_\mu(\theta;D)^{1/n}
   =\inf_{\gamma\text{ Gauduchon}}
     \frac{\int_X\theta\wedge\gamma^{n-1}-\deg_\gamma D}
          {D_\mu(\gamma)},
\end{equation}
and the residual volume is finite and nonnegative. Changing the auxiliary metrics on the line bundles changes $\eta_D$ by a smooth $\ddc$-exact term, so the residual volume does not depend on these metrics, nor on a smooth $\ddc$ change of $\theta$. It depends on the prescribed divisor and on the smooth probability measure $\mu$.

Here $\mathrm{MA}^{\mathrm{ext}}_\theta$ denotes the extended non-pluripolar Monge--Amp\`ere measure of \cref{thm:arbitrary-extended-ma}.

\begin{theorem}\label{thm:divisorial-residual-equation}\index{Divisor!residual equation}
Suppose that $\beta_D$ has a smooth positive representative, and put $c=\mathcal V_\mu(\theta;D)$. There is a $\theta$-plurisubharmonic function $u_0=\sigma_D+v$, with $v$ smooth, satisfying
\[
 \mathrm{MA}^{\mathrm{ext}}_\theta(u_0)=c\mu.
\]
Every $u\in\PSH(X,\theta)$ such that
\[
 \nu(u,D_j)\geq a_j\quad(1\leq j\leq N),\qquad
 \mathrm{MA}^{\mathrm{ext}}_\theta(u)\geq c\mu
\]
differs from $u_0$ by a constant. In particular, a solution of a normalized equation $\mathrm{MA}^{\mathrm{ext}}_\theta(u)=c'\mu$ in the ordinary singularity type of $\sigma_D$ has scalar $c'=c$, has a smooth remainder $u-\sigma_D$, and is unique up to constants.
\end{theorem}
\begin{proof}
\emph{Existence.} The smooth equation and \cref{thm:gaud-volume-dual-formula} give a smooth $v$ with
\[
 \Omega=\beta_D+\ddc v>0,\qquad \Omega^n=c\mu.
\]
Then $\theta+\ddc(\sigma_D+v)=[D]+\Omega\geq0$. Off $D$, the extended measure of $u_0=\sigma_D+v$ is $\Omega^n$, and both measures give zero mass to $D$. This proves existence.

\emph{Reduction to the residual background.} Let $u$ satisfy the two inequalities. Removing the divisor as in the proof of \cref{prop:divisorial-gaud-feasibility}, we obtain
\[
 w=u-\sigma_D-v\in\PSH(X,\Omega),\qquad
 \Omega+\ddc w=\theta+\ddc u-[D]\geq0,
\]
where the difference is first taken in $L^1$ and then replaced by its canonical representative. On $X\setminus D$, the two currents differ only by smooth changes in the decomposition of their local potentials, so their extended measures agree there, and both measures vanish on $D$. Hence
\[
 \mathrm{MA}^{\mathrm{ext}}_\Omega(w)
       =\mathrm{MA}^{\mathrm{ext}}_\theta(u)\geq c\mu.
\]

\emph{Rigidity.} Take $\mathrm{d}V=\mu$. The zero function solves the smooth equation $\Omega^n=c\mu$ in the background $\Omega$, so its subsolution threshold is $b_\Omega(1)=c$, by \cref{prop:individual-normalizer-lp}. For $n\geq2$, the normalized $(n-1)$st power of the Gauduchon metric conformal to $\Omega$ is a smooth strictly positive maximizer of the dual problem. By \cref{cor:smooth-dual-exact-equation}, any two threshold subsolutions differ by a constant; comparing $w$ with zero shows that $w$ is constant. For $n=1$, the positive current $\Omega+\ddc w$ dominates its extended measure, hence dominates $c\mu=\Omega$, while both have total mass $\int_X\Omega=c$ by Stokes' theorem. Thus $\ddc w=0$, and $w$ is again constant.

\emph{Solutions in the ordinary type.} Finally, let $u$ be in the ordinary singularity type of $\sigma_D$ with $\mathrm{MA}^{\mathrm{ext}}_\theta(u)=c'\mu$ for some $c'\geq0$. Then $\nu(u,D_j)=a_j$, and the same reduction gives a bounded $w\in\PSH(X,\Omega)$ with $\mathrm{MA}_\Omega(w)=c'\mu$. Suppose $c'<c$. For every constant $C$, the bounded functions $w$ and $C$ satisfy $\mathrm{MA}_\Omega(w)=(c'/c)\mathrm{MA}_\Omega(C)$, so \cref{thm:bounded-strict-domination} gives $w\geq C$, which is impossible for a bounded $w$. The case $c'>c$ is symmetric, with the roles of $w$ and the constants exchanged. Therefore $c'=c$, and the rigidity just proved gives $u=u_0+\mathrm{constant}$.
\end{proof}


When $\theta$ is Hermitian and $n\geq2$, write $\mu=f\,\mathrm{d}V$ relative to a smooth strictly positive probability volume $\mathrm{d}V$. For every $u\in\PSH(X,\theta)$ with $\theta+\ddc u\geq[D]$ and $\mathrm{MA}^{\mathrm{ext}}_\theta(u)\geq cf\,\mathrm{d}V$, in particular for the potential $u_0$ of \cref{thm:divisorial-residual-equation} with $c=\mathcal V_\mu(\theta;D)$, \cref{prop:threshold-divisorial-constraint} bounds the coefficient $c$ in terms of the subsolution threshold $b_\theta(f)$ of \cref{thm:individual-density-subsolution} and the linear threshold of the class of the divisor. The residual problem prescribes the divisor, whereas the threshold $b_\theta(f)$ allows all $\theta$-plurisubharmonic functions.

\Needspace{12\baselineskip}
Along a family $D_t=tD$, $t\geq0$, in the feasible interval, the root of \eqref{eq:divisorial-gaud-volume} is concave and nonincreasing in $t$: every test quotient is affine in $t$ with nonpositive slope, and both properties pass to the infimum. On a subinterval where $\theta-t\eta_D$ has a smooth positive representative, choose a minimizing Gauduchon metric $\gamma_t$ for that background. \Cref{prop:gaud-volume-first-variation} then gives the sharper formula
\begin{equation}\label{eq:divisorial-volume-derivative}
 \frac{\mathrm{d}}{\mathrm{d}t}\mathcal V_\mu(\theta;tD)^{1/n}
       =-\frac{\deg_{\gamma_t}D}{D_\mu(\gamma_t)}<0.
\end{equation}
In dimension one the derivative is $-\deg D$, since the denominator is one. If the residual background at an endpoint of the interval is nef, \cref{thm:gaud-volume-nef-limit} also expresses its volume as the limit of the scalars of the smooth normalized equations for the perturbations by $\epsilon h$, with $h$ a fixed Hermitian form and $\epsilon\downarrow0$.

\Needspace{13\baselineskip}
\subsection{Varying divisorial coefficients}

For a fixed support, the residual equation depends smoothly on the coefficients as long as the residual background stays in $\mathscr P$. The smooth remainders then give convergence in capacity of the solutions, and they determine all generic Lelong numbers of the limit on every modification. Keep $D_j$, $s_j$ and $\eta_j$ as above, and put $\sigma_j=\log|s_j|^2$.

For $a=(a_1,\ldots,a_N)\in\mathbb R_{\geq0}^N$ put
\[
 D_a=\sum_j a_jD_j,\qquad
 \sigma_a=\sum_{a_j>0}a_j\sigma_j,\qquad
 \beta_a=\theta-\sum_j a_j\eta_j,
\]
where zero coefficients are omitted from the logarithmic sum. Let $\mathcal A$ be the set of such $a$ for which $\beta_a$ has a smooth positive representative.

\Needspace{14\baselineskip}
\begin{proposition}\label{prop:finite-divisorial-equation-limits}
There are unique smooth functions $v_a$, $a\in\mathcal A$, such that
\[
 \int_Xv_a\,\mathrm{d}\mu=0,\qquad
 \Omega_a=\beta_a+\ddc v_a>0,\qquad
 \Omega_a^n=c(a)\mu,
 \quad c(a)=\mathcal V_\mu(\beta_a)>0.
\]
The maps $a\mapsto v_a$ and $a\mapsto c(a)$ are smooth, and every smooth norm of $v_a$ is bounded on compact subsets of $\mathcal A$. Here smoothness at a parameter with a zero coefficient means that the family is the restriction of a smooth family defined on an open subset of $\mathbb R^N$.

Set $u_a=\sigma_a+v_a$. If $a^{(k)}\in\mathcal A$ and $a^{(k)}\to a\in\mathcal A$, then $u_{a^{(k)}}\to u_a$ in $L^1$ and in capacity, and smoothly on compact subsets of $X\setminus\bigcup_jD_j$. Moreover,
\[
 \bigl\|\mathrm{MA}^{\mathrm{ext}}_\theta(u_{a^{(k)}})
       -\mathrm{MA}^{\mathrm{ext}}_\theta(u_a)\bigr\|_{\mathrm{TV}}
       =|c(a^{(k)})-c(a)|\longrightarrow0,
\]
where $\|\cdot\|_{\mathrm{TV}}$ is the total variation norm. For every modification $\pi:Y\to X$ with $Y$ smooth and every prime divisor $E\subset Y$,
\[
 \nu(u_a\circ\pi,E)=\sum_j a_j\operatorname{ord}_E(\pi^*D_j),
\]
where $\operatorname{ord}_E(\pi^*D_j)$ is the vanishing order along $E$ of the pullback of a local equation of $D_j$. In particular, the limit $u_a$ has exactly the ordinary singularity type of $\sigma_a$, and no further singularity appears in this finite family.
\end{proposition}
\begin{proof}
The set of parameters in $\mathbb R^N$ for which $\beta_a$ has a smooth positive representative is open, and the smooth equation gives $v_a$ and $c(a)$ on it. The argument in the proof of \cref{prop:gaud-volume-first-variation} gives smooth dependence of both quantities on the parameter, with the fixed mean normalization; uniqueness makes the local families agree. Covering a compact set of parameters by finitely many such neighborhoods gives the asserted bounds.

The Poincar\'e--Lelong formula gives
\[
 \theta+\ddc u_a=[D_a]+\Omega_a,
\]
so $u_a\in\PSH(X,\theta)$. Its extended measure equals $\Omega_a^n=c(a)\mu$ off the union of the divisors, and both measures vanish on that union; this also covers the case $D_a=0$. Each fixed $\sigma_j$ is locally integrable, so the convergence of the coefficients and the smooth convergence of $v_a$ give convergence in $L^1$. Off the union of the divisors every $\sigma_j$ is smooth, which gives smooth convergence there. For convergence in capacity, remove a neighborhood of this analytic set of arbitrarily small Hermitian capacity; on its compact complement the convergence is uniform. Since the union of the divisors has capacity zero, this proves convergence in capacity. The identity for the total variation distances follows from the equations and $\mu(X)=1$.

Finally, after pullback the logarithmic potential $\sigma_j\circ\pi$ has generic Lelong number $\operatorname{ord}_E(\pi^*D_j)$ along $E$, whereas $v_a\circ\pi$ is smooth. Additivity of Lelong numbers for the finite sum proves the formula on every modification. Since $v_a$ is smooth also at the limiting parameter, $u_a$ has the ordinary singularity type of $\sigma_a$.
\end{proof}

The smooth bounds above use that the limiting parameter lies in $\mathcal A$. When the sections and weights move, \cref{thm:moving-divisorial-potential-stability} gives $L^1$ convergence of the potentials normalized by $\sup_X=0$; when the residual background is fixed, it also allows uniformly bounded measurable right-hand sides converging in $L^1$. On a modification, the pullback of a positive probability density can vanish along the exceptional divisor, so the smooth equation with positive density no longer applies directly. For a minimal exceptional pole, \cref{prop:exceptional-density-obstruction} gives a vanishing order of the density that is necessary for a smooth residual, but this order alone does not ensure smoothness.

Multiplier ideals need not stabilize in this family. At a smooth point of $D_j$ outside the other divisors, take a local coordinate $z_1$ defining $D_j$. There $u_a=a_j\log|z_1|^2+O(1)$, so
\[
 \mathcal I(u_a)=(z_1^{\lfloor a_j\rfloor})
\]
by the one-variable integral in \cref{ex:ideal-disc} and Fubini's theorem. Thus, if the parameter stays in $\mathcal A$ while $a_j$ increases to an integer from below, this local ideal jumps at the limit. Convergence of the Lelong numbers to their limiting values does not imply that the multiplier ideals are eventually equal to those of the limit.


\subsection{A ruled residual and its subtraction cost}
\label{subsec:ruled-residual-gauduchon-cost}

On the ruled surface of \cref{sec:nonclosed-ruled-trace-masses}, the Gauduchon variational volume has a different dependence on the data from the two volumes of an ordinary type. Keep $Y=\mathbb P^1\times\mathbb P^1$, $\theta=cw_1+aw_2$, $0\leq b<c$ and $\beta=(c-b)w_1+aw_2$. Let $m>0$ be smooth on the first curve, with $\int mw=1$, and put
\[
 \mu=mw_1\wedge w_2,\qquad k=c-b,\qquad H=\int_{\mathbb P^1}\frac{m}{a}w.
\]
In dimension two the Gauduchon tests are all smooth positive $(1,1)$-forms $T$ with $\ddc T=0$, and $D_\mu(T)=\int_Y\sqrt{T^2/\mu}\,\mu$.

\enlargethispage{4\baselineskip}
\begin{proposition}\label{prop:ruled-residual-gauduchon-volume}\index{Gauduchon variational volume!ruled residual}
The all-test infimum is attained and
\begin{equation}\label{eq:ruled-residual-gauduchon-volume}
 \mathcal V_\mu(\theta;bD)=\mathcal V_\mu(\beta)=\frac{2(c-b)}{H},
 \qquad \mathcal V_\mu(\theta)=\frac{2c}{H}.
\end{equation}
Writing $V=2k/H$ and $g=Vm/(2a)$, a minimizing test is the closed positive form
\[
 T=\frac ga w_1+w_2.
\]
\end{proposition}
\begin{proof}
Since $\int gw=k$, solve $\ddc f=(g-k)w$ on the first sphere. Then $\Omega=\beta+\ddc f=gw_1+aw_2>0$ and $\Omega^2=V\mu$. For every Gauduchon test, Stokes and the matrix arithmetic--geometric mean inequality give
\[
 \int_Y\beta\wedge T=\int_Y\Omega\wedge T
       \geq\sqrt V\,D_\mu(T).
\]
The displayed $T$ is closed, since $g/a$ depends only on the first factor, and $T=a^{-1}\Omega$, so equality holds pointwise. Explicitly,
\[
 \int_Y\beta\wedge T=2k,
 \qquad D_\mu(T)=\sqrt V\int\frac ma w=\sqrt V H.
\]
This proves the infimum over all tests. The curvature of $D$ is $w_1$, so the divisorial residual is indeed $\beta$. Replacing $k$ by $c$ gives the formula for $\theta$.
\end{proof}

Fix a real constant $d$ and a smooth real first-factor function $e$. On an interval where $c_s=c-sd>b$ and $a_s=a-se>0$, keep $b,m$ fixed and put
\[
 k_s=c_s-b,\qquad H_s=\int\frac{m}{a_s}w,
 \qquad V_s=\frac{2k_s}{H_s},\qquad
 T_s=\frac{V_sm}{2a_s^2}w_1+w_2.
\]
\begin{corollary}\label{cor:ruled-residual-subtraction-cost}
For the residual background $\beta_s=k_sw_1+a_sw_2$,
\begin{equation}\label{eq:ruled-residual-subtraction-cost}
 \frac{\mathrm d}{\mathrm ds}\sqrt{V_s}
     =-\frac{\int_Y(d w_1+e w_2)\wedge T_s}{D_\mu(T_s)}.
\end{equation}
Consequently, for every compact segment $[s_0,s_1]$ in that interval,
\[
 \sqrt{V_{s_0}}-\sqrt{V_{s_1}}
   =\int_{s_0}^{s_1}
       \frac{\int_Y(d w_1+e w_2)\wedge T_s}{D_\mu(T_s)}\,\mathrm ds.
\]
\end{corollary}
\begin{proof}
We have $H_s'=\int em/a_s^2\,w$, $D_\mu(T_s)=\sqrt{V_s}H_s$ and
\[
 \int_Y(d w_1+e w_2)\wedge T_s=d+\frac{V_s}{2}H_s'.
\]
Differentiating $V_s=2k_s/H_s$ gives \eqref{eq:ruled-residual-subtraction-cost}; integration proves the second identity. This is the explicit first variation of \cref{prop:gaud-volume-first-variation} along the residual family.
\end{proof}

For $c=1$, $b=1/4$, $a=2+\tau$ and $m=1$, the height $\tau$ has distribution $\mathrm d\tau/2$ under $w$. Hence $H=(\log3)/2$, and the residual and original variational volumes are $3/\log3$ and $4/\log3$. The ordinary-type volumes from \eqref{eq:nonclosed-ruled-ordinary-volumes} are instead $3/2$ and $9/2$, independent of $m$. The subtraction identity differentiates the density-dependent variational volume. Pure-pullback invariance of the ordinary-type masses does not identify an arbitrary smooth positive probability density upstairs with the degenerate pullback of $\mu$.


\section{A divisorial threshold on the Hopf threefold}
\label{sec:hopf-divisorial-volume}

On the Hopf threefold, the pullback of a hyperplane has an explicit divisorial threshold. The residual metric solves its normalized equation with zero potential, so the residual volume can be computed directly, also at the nef endpoint.

In this section $X$ is the Hopf threefold of \cref{chap:hopf}, with its projection $\pi\colon X\to\mathbb P^2$, and we write $\zeta$ for its vertical form, so that $\omega=\alpha+\zeta$. These forms satisfy
\[
 \alpha^3=0,\quad\zeta^2=0,\quad
 \ddc\zeta=-\alpha^2,\quad\int_X\alpha^2\wedge\zeta=2.
\]
Let $H$ be the inverse image under $\pi$ of a projective line in $\mathbb P^2$. The pulled-back Fubini--Study metric gives $\eta_H=\alpha$ and a logarithmic potential $\sigma_H$ with $\ddc\sigma_H=[H]-\alpha$.

\begin{proposition}\label{prop:hopf-divisorial-threshold}
The divisorial threshold is $\tau_\omega(H)=1$. For $\mu_0=\omega^3/6$ and $0\leq a\leq1$,
\[
 \mathcal V_{\mu_0}(\omega;aH)=6(1-a)^2.
\]
For $a<1$ this is the scalar of the equation solved, as in \cref{thm:divisorial-residual-equation}, by $a\sigma_H$ plus a smooth function. At $a=1$ the residual background is nef, and the logarithmic potential $\sigma_H$ has zero non-pluripolar Monge--Amp\`ere measure. For $a<1$, every residual background has infinite upper mass.
\end{proposition}
\begin{proof}
The residual background is
\[
 R_a=(1-a)\alpha+\zeta.
\]
It is positive for $0\leq a<1$ and semipositive for $a=1$, so every coefficient $a\leq1$ is feasible. For the reverse inequality, consider the metrics $\gamma_t=\alpha+t\zeta$, $t>0$. They are Gauduchon, since
\[
 \ddc(\gamma_t^2)=-2t\alpha^3=0.
\]
\Needspace{6\baselineskip}
Using $\zeta^2=0$, we obtain
\[
 \frac{\int_X\omega\wedge\gamma_t^2}{\deg_{\gamma_t}H}
       =\frac{1+2t}{2t}\longrightarrow1\qquad(t\to\infty),
\]
and the threshold formula gives $\tau_\omega(H)=1$.

For $a<1$,
\[
 R_a^3=3(1-a)^2\alpha^2\wedge\zeta
       =(1-a)^2\omega^3=6(1-a)^2\mu_0.
\]
Hence the zero potential solves the residual smooth equation, and uniqueness of the scalar gives the volume. At $a=1$, the form $R_1=\zeta$ is nef, because $\zeta+\epsilon\omega>0$ for $\epsilon>0$, and the same computation gives
\[
 \mathcal V_{\mu_0}(\zeta+\epsilon\omega)
       =6\epsilon^2(1+\epsilon)\longrightarrow0.
\]
By \cref{thm:gaud-volume-nef-limit}, the residual volume at $a=1$ is zero. Moreover $\omega+\ddc\sigma_H=[H]+\zeta$, so the non-pluripolar measure of $\sigma_H$ is $\zeta^3=0$.

Finally, fix $a<1$ and choose $0<\delta\leq\min\{1-a,1\}$, so that $R_a-\delta\omega\geq0$. For the smooth functions $V_j$ of \cref{thm:hopf-unbounded},
\[
 R_a+\ddc(\delta V_j)
       =(R_a-\delta\omega)+\delta(\omega+\ddc V_j)>0.
\]
Monotonicity of the determinant on positive matrices bounds the mass of this form below by $\delta^3\int_X(\omega+\ddc V_j)^3$, which tends to infinity. Thus a finite residual volume and an explicit divisorial threshold coexist with unbounded upper mass.
\end{proof}


\par\endgroup
